\documentclass[11pt,a4paper]{article}
\usepackage[margin=24mm]{geometry}
\usepackage{amsmath,amssymb,bm}
\usepackage{fontspec}
\usepackage[hypertexnames=false]{hyperref}
\usepackage{longtable}
\usepackage{xurl}
\setmainfont{Latin Modern Roman}
\title{First-Principles Derivation and Verification of Diffusive Source-Term Release from Multilayer TRISO Fuel}
\author{Ray}
\date{}
\begin{document}
\maketitle
\section*{First-Principles Derivation and Verification of Diffusive Source-Term Release from Multilayer TRISO Fuel}

\textbf{Detailed derivation manuscript. The hard five-layer transform is independently verified with non-blocking findings. Controlled inverse-Laplace verification is authorized but has not yet been executed.}

This is the single continuous derivation path for the Ray TRISO project. It preserves the user's original notebook as a separate evidence stream and incorporates the corrections established during the foundation audit.

The supplied LaTeX notebook contains 33 substantive displayed equations. Every one is mapped in 13\_canonical\_equation\_register.md.

\subsection*{1. Physical problem}

A TRISO particle has five concentric regions: fuel kernel, buffer, IPyC, SiC, and OPyC. A fission-product species is created in the kernel, diffuses through the layers, may decay or undergo other modelled reactions, and may leave the particle.

Let r be distance from the particle centre, t be time, c\_i(r,t) be concentration in layer i, and D\_i be the layer diffusivity.

\textbf{[ASSUMPTION]} One representative spherical particle.
\textbf{[ASSUMPTION]} Spherical symmetry.
\textbf{[ASSUMPTION]} Each layer is homogeneous during one model evaluation.
\textbf{[ASSUMPTION]} Initial ideal interfaces have no interfacial storage or resistance.
\textbf{[QUESTION FOR SUPERVISOR]} The first canonical species, decay/trapping model, partition coefficients, and coolant concentration still need to be fixed.

\subsection*{1.1 Nomenclature}

The notation is grouped here for reference; important symbols are also defined locally at first use.

\begin{longtable}{p{0.12\textwidth} p{0.15\textwidth} p{0.40\textwidth} p{0.16\textwidth}}
Group & Symbol & Definition & SI unit \\
\hline
Geometry & $r$ & radial coordinate from particle centre & m \\
Geometry & $R_i$ & outer radius of TRISO layer $i$ & m \\
Transport & $c_i$ & concentration of tracked species in layer $i$ & mol m$^{-3}$ \\
Transport & $D_i$ & diffusion coefficient in layer $i$ & m$^2$ s$^{-1}$ \\
Transport & $\mathbf J$ & diffusive molar flux & mol m$^{-2}$ s$^{-1}$ \\
Sources & $S_i$ & net volumetric production term & mol m$^{-3}$ s$^{-1}$ \\
FV & $V_P$ & spherical control-volume volume & m$^3$ \\
FV & $A_f$ & spherical face area & m$^2$ \\
FV & $G_f$ & diffusive face conductance & m$^3$ s$^{-1}$ \\
Stochastic & $T$ & first-passage or release time & s \\
Stochastic & $\epsilon$ & interface capture distance & m \\
Stochastic & $\delta=\alpha\epsilon$ & reinsertion displacement & m \\
First passage & $G_a,G_b$ & inner/outer joint exit-time Laplace transforms & dimensionless \\
First passage & $H$ & centred-ball first-exit Laplace transform & dimensionless \\
Renewal & $\Phi_i(s)$ & release-time transform from state $i$ & dimensionless \\
Renewal & $\mathbf K(s)$ & transient renewal-transform matrix & dimensionless \\
Renewal & $\mathbf B(s)$ & direct-absorption transform vector & dimensionless \\
Numerical & $\kappa_2$ & spectral 2-norm condition number & dimensionless \\
\end{longtable}

\subsection*{2. Conservation from a control volume}

\textbf{[EXACT]} Begin with the amount of the conserved species inside an arbitrary fixed control volume V.

\begin{equation}
N_V(t)=\int_V c(\mathbf{x},t)\,dV.
\tag{TRISO-GOV-020}
\label{eq:triso-gov-020}
\end{equation}

The units are

\begin{equation}
[N_V]=\mathrm{mol}.
\tag{TRISO-GOV-021}
\label{eq:triso-gov-021}
\end{equation}

\textbf{[EXACT]} The rate of accumulation is

\begin{equation}
\frac{dN_V}{dt}=\frac{d}{dt}\int_V c\,dV.
\tag{TRISO-GOV-022}
\label{eq:triso-gov-022}
\end{equation}

Let $\mathbf J$ be the diffusive molar flux vector. Its units are

\begin{equation}
[\mathbf J]=\mathrm{mol\,m^{-2}\,s^{-1}}.
\tag{TRISO-GOV-023}
\label{eq:triso-gov-023}
\end{equation}

Let $S$ be the net volumetric production rate. Its units are

\begin{equation}
[S]=\mathrm{mol\,m^{-3}\,s^{-1}}.
\tag{TRISO-GOV-024}
\label{eq:triso-gov-024}
\end{equation}

\textbf{[EXACT]} For outward unit normal $\mathbf n$, the outward amount crossing a boundary element in time $dt$ is proportional to $\mathbf J\cdot\mathbf n$. The outward rate is therefore

\begin{equation}
\dot N_{\mathrm{out}}=\int_{\partial V}\mathbf J\cdot\mathbf n\,dA.
\tag{TRISO-GOV-025}
\label{eq:triso-gov-025}
\end{equation}

\textbf{[EXACT]} The production rate inside the control volume is

\begin{equation}
\dot N_{\mathrm{gen}}=\int_V S\,dV.
\tag{TRISO-GOV-026}
\label{eq:triso-gov-026}
\end{equation}

\textbf{[EXACT]} Accumulation equals production minus outward flux:

\begin{equation}
\frac{dN_V}{dt}=-\dot N_{\mathrm{out}}+\dot N_{\mathrm{gen}}.
\tag{TRISO-GOV-027}
\label{eq:triso-gov-027}
\end{equation}

Substitute the definitions of the two rates:

\begin{equation}
\frac{d}{dt}\int_V c\,dV=-\int_{\partial V}\mathbf J\cdot\mathbf n\,dA+\int_V S\,dV.
\tag{TRISO-GOV-028}
\label{eq:triso-gov-028}
\end{equation}

\textbf{[EXACT]} Because $V$ is fixed in space, the time derivative passes through the volume integral:

\begin{equation}
\frac{d}{dt}\int_V c\,dV=\int_V\frac{\partial c}{\partial t}\,dV.
\tag{TRISO-GOV-029}
\label{eq:triso-gov-029}
\end{equation}

\textbf{[EXACT]} Apply the divergence theorem to the surface term:

\begin{equation}
\int_{\partial V}\mathbf J\cdot\mathbf n\,dA=\int_V\nabla\cdot\mathbf J\,dV.
\tag{TRISO-GOV-030}
\label{eq:triso-gov-030}
\end{equation}

Substitute this result:

\begin{equation}
\int_V\frac{\partial c}{\partial t}\,dV=-\int_V\nabla\cdot\mathbf J\,dV+\int_VS\,dV.
\tag{TRISO-GOV-031}
\label{eq:triso-gov-031}
\end{equation}

Move the flux and source terms into one integrand:

\begin{equation}
\int_V\left(\frac{\partial c}{\partial t}+\nabla\cdot\mathbf J-S\right)dV=0.
\tag{TRISO-GOV-032}
\label{eq:triso-gov-032}
\end{equation}

\textbf{[EXACT]} Since the control volume is arbitrary, the integrand must vanish almost everywhere:

\begin{equation}
\boxed{\frac{\partial c}{\partial t}+\nabla\cdot\mathbf J=S.}
\tag{TRISO-GOV-033}
\label{eq:triso-gov-033}
\end{equation}


\subsection*{3. Constitutive law and general heterogeneous diffusion equation}

\textbf{[CONSTITUTIVE]} Fickian diffusion relates flux to the concentration gradient:

\begin{equation}
\boxed{\mathbf J=-D\nabla c.}
\tag{TRISO-GOV-034}
\label{eq:triso-gov-034}
\end{equation}

The gradient has units

\begin{equation}
[\nabla c]=\mathrm{mol\,m^{-4}}.
\tag{TRISO-GOV-035}
\label{eq:triso-gov-035}
\end{equation}

Multiplying by $D$ gives

\begin{equation}
[D\nabla c]=\mathrm{m^2\,s^{-1}}\times\mathrm{mol\,m^{-4}}.
\tag{TRISO-GOV-036}
\label{eq:triso-gov-036}
\end{equation}

Hence

\begin{equation}
[D\nabla c]=\mathrm{mol\,m^{-2}\,s^{-1}},
\tag{TRISO-GOV-037}
\label{eq:triso-gov-037}
\end{equation}
which matches the flux units.

\textbf{[EXACT]} Substitute Fick's law into conservation:

\begin{equation}
\frac{\partial c}{\partial t}+\nabla\cdot(-D\nabla c)=S.
\tag{TRISO-GOV-038}
\label{eq:triso-gov-038}
\end{equation}

\textbf{[EXACT]} Pull the minus sign through the divergence:

\begin{equation}
\frac{\partial c}{\partial t}-\nabla\cdot(D\nabla c)=S.
\tag{TRISO-GOV-039}
\label{eq:triso-gov-039}
\end{equation}

\textbf{[EXACT]} Rearrange:

\begin{equation}
\boxed{\frac{\partial c}{\partial t}=\nabla\cdot(D\nabla c)+S.}
\tag{TRISO-GOV-040}
\label{eq:triso-gov-040}
\end{equation}

\textbf{[IMPORTANT]} This is the general conservative form. The diffusivity must remain inside the divergence until a later layer-specific assumption establishes that it is constant with respect to the differentiated coordinate.


\subsection*{4. Spherical-coordinate derivation}

\textbf{[EXACT]} In spherical coordinates, the gradient of a scalar field is

\begin{equation}
\nabla c=\mathbf e_r\frac{\partial c}{\partial r}+\mathbf e_\theta\frac1r\frac{\partial c}{\partial\theta}+\mathbf e_\varphi\frac1{r\sin\theta}\frac{\partial c}{\partial\varphi}.
\tag{TRISO-SPH-020}
\label{eq:triso-sph-020}
\end{equation}

\textbf{[EXACT]} Write a general vector flux as $\mathbf J=J_r\mathbf e_r+J_\theta\mathbf e_\theta+J_\varphi\mathbf e_\varphi$.

\textbf{[EXACT]} Its spherical divergence is

\begin{equation}
\nabla\cdot\mathbf J=\frac1{r^2}\frac{\partial}{\partial r}(r^2J_r)+\frac1{r\sin\theta}\frac{\partial}{\partial\theta}(\sin\theta J_\theta)+\frac1{r\sin\theta}\frac{\partial J_\varphi}{\partial\varphi}.
\tag{TRISO-SPH-021}
\label{eq:triso-sph-021}
\end{equation}

\textbf{[ASSUMPTION]} Spherical symmetry means the concentration is independent of both angular coordinates:

\begin{equation}
c=c(r,t).
\tag{TRISO-SPH-022}
\label{eq:triso-sph-022}
\end{equation}

Therefore

\begin{equation}
\frac{\partial c}{\partial\theta}=0.
\tag{TRISO-SPH-023}
\label{eq:triso-sph-023}
\end{equation}

and

\begin{equation}
\frac{\partial c}{\partial\varphi}=0.
\tag{TRISO-SPH-024}
\label{eq:triso-sph-024}
\end{equation}

Substitute these zero angular derivatives into the gradient:

\begin{equation}
\nabla c=\mathbf e_r\frac{\partial c}{\partial r}.
\tag{TRISO-SPH-025}
\label{eq:triso-sph-025}
\end{equation}

\textbf{[CONSTITUTIVE]} Fick's law therefore becomes radial:

\begin{equation}
\boxed{\mathbf J=-D(r,t)\frac{\partial c}{\partial r}\mathbf e_r.}
\tag{TRISO-SPH-026}
\label{eq:triso-sph-026}
\end{equation}

Thus

\begin{equation}
J_\theta=0.
\tag{TRISO-SPH-027}
\label{eq:triso-sph-027}
\end{equation}

and

\begin{equation}
J_\varphi=0.
\tag{TRISO-SPH-028}
\label{eq:triso-sph-028}
\end{equation}

\textbf{[EXACT]} Insert the zero angular fluxes into the divergence:

\begin{equation}
\nabla\cdot\mathbf J=\frac1{r^2}\frac{\partial}{\partial r}(r^2J_r).
\tag{TRISO-SPH-029}
\label{eq:triso-sph-029}
\end{equation}

The radial flux component is

\begin{equation}
J_r=-D(r,t)\frac{\partial c}{\partial r}.
\tag{TRISO-SPH-030}
\label{eq:triso-sph-030}
\end{equation}

Substitution gives

\begin{equation}
\nabla\cdot\mathbf J=\frac1{r^2}\frac{\partial}{\partial r}\left(-r^2D(r,t)\frac{\partial c}{\partial r}\right).
\tag{TRISO-SPH-031}
\label{eq:triso-sph-031}
\end{equation}

Insert this into conservation:

\begin{equation}
\frac{\partial c}{\partial t}+\frac1{r^2}\frac{\partial}{\partial r}\left(-r^2D(r,t)\frac{\partial c}{\partial r}\right)=S(r,t).
\tag{TRISO-SPH-032}
\label{eq:triso-sph-032}
\end{equation}

\textbf{[EXACT]} Move the negative term to the right:

\begin{equation}
\boxed{\frac{\partial c}{\partial t}=\frac1{r^2}\frac{\partial}{\partial r}\left(r^2D(r,t)\frac{\partial c}{\partial r}\right)+S(r,t).}
\tag{TRISO-SPH-033}
\label{eq:triso-sph-033}
\end{equation}


\subsection*{4.1 Specialisation to one homogeneous layer}

\textbf{[ASSUMPTION]} In one material layer $i$, the benchmark assumes the diffusivity is constant with respect to radius and time during the analysis:

\begin{equation}
D(r,t)=D_i.
\tag{TRISO-SPH-034}
\label{eq:triso-sph-034}
\end{equation}

\textbf{[EXACT]} Substitute $D_i$ into the conservative equation:

\begin{equation}
\frac{\partial c_i}{\partial t}=\frac1{r^2}\frac{\partial}{\partial r}\left(r^2D_i\frac{\partial c_i}{\partial r}\right)+S_i.
\tag{TRISO-SPH-035}
\label{eq:triso-sph-035}
\end{equation}

\textbf{[EXACT]} Because $D_i$ is constant with respect to $r$, take it outside the derivative:

\begin{equation}
\frac{\partial c_i}{\partial t}=\frac{D_i}{r^2}\frac{\partial}{\partial r}\left(r^2\frac{\partial c_i}{\partial r}\right)+S_i.
\tag{TRISO-SPH-036}
\label{eq:triso-sph-036}
\end{equation}

\textbf{[EXACT]} Apply the product rule:

\begin{equation}
\frac{\partial}{\partial r}\left(r^2\frac{\partial c_i}{\partial r}\right)=\frac{\partial r^2}{\partial r}\frac{\partial c_i}{\partial r}+r^2\frac{\partial^2c_i}{\partial r^2}.
\tag{TRISO-SPH-037}
\label{eq:triso-sph-037}
\end{equation}

Differentiate $r^2$:

\begin{equation}
\frac{\partial r^2}{\partial r}=2r.
\tag{TRISO-SPH-038}
\label{eq:triso-sph-038}
\end{equation}

Substitute:

\begin{equation}
\frac{\partial}{\partial r}\left(r^2\frac{\partial c_i}{\partial r}\right)=2r\frac{\partial c_i}{\partial r}+r^2\frac{\partial^2c_i}{\partial r^2}.
\tag{TRISO-SPH-039}
\label{eq:triso-sph-039}
\end{equation}

Divide by $r^2$:

\begin{equation}
\frac1{r^2}\frac{\partial}{\partial r}\left(r^2\frac{\partial c_i}{\partial r}\right)=\frac2r\frac{\partial c_i}{\partial r}+\frac{\partial^2c_i}{\partial r^2}.
\tag{TRISO-SPH-040}
\label{eq:triso-sph-040}
\end{equation}

Therefore:

\begin{equation}
\boxed{\frac{\partial c_i}{\partial t}=D_i\left(\frac{\partial^2c_i}{\partial r^2}+\frac2r\frac{\partial c_i}{\partial r}\right)+S_i.}
\tag{TRISO-SPH-041}
\label{eq:triso-sph-041}
\end{equation}

\textbf{[IMPORTANT]} Equation TRISO-SPH-041 is a within-layer constant-diffusivity equation. It must not be differentiated through a discontinuous material interface.

\subsection*{5. Source and reaction terms}

The conservation equation contains a net volumetric source term. We now separate the physical processes that may contribute to that term.

\textbf{[EXACT]} Define the net source in material layer i as

\begin{equation}
S_i=S_{i,\mathrm{gen}}+S_{i,\mathrm{other}}+S_{i,\mathrm{release}}-S_{i,\mathrm{decay}}-S_{i,\mathrm{trap}}.
\tag{TRISO-GOV-100}
\label{eq:triso-gov-100}
\end{equation}

Each term has units

\begin{equation}
[S_{i,\mathrm{gen}}]
=
[S_{i,\mathrm{other}}]
=
[S_{i,\mathrm{release}}]
=
[S_{i,\mathrm{decay}}]
=
[S_{i,\mathrm{trap}}]
=
\mathrm{mol\,m^{-3}\,s^{-1}}.
\tag{TRISO-GOV-101}
\label{eq:triso-gov-101}
\end{equation}

This equation is a bookkeeping definition. It does not yet choose a constitutive model for trapping, release, or generation.

\subsubsection*{5.1 Fission-product generation}

Let $S_{i,\mathrm{gen}}$ denote the local rate at which the tracked species is created.

\textbf{[ASSUMPTION]} For the simplest TRISO source benchmark, generation is confined to the fuel kernel:

\begin{equation}
S_{i,\mathrm{gen}}=
\begin{cases}
S_0,&0\le r<r_1,\\
0,&r_1<r<R.
\end{cases}
\tag{TRISO-GOV-102}
\label{eq:triso-gov-102}
\end{equation}

Here

\begin{equation}
[S_0]=\mathrm{mol\,m^{-3}\,s^{-1}}.
\tag{TRISO-GOV-103}
\label{eq:triso-gov-103}
\end{equation}

A microscopic fission-based expression such as a fission rate multiplied by a product yield would require a selected species, fission cross section, neutron flux, and yield correlation. Those inputs are not fixed by the current project evidence.

\textbf{[SOURCE NEEDED]} Species-specific generation correlation and parameter values.

\subsubsection*{5.2 Radioactive decay}

Suppose the tracked atoms disappear by radioactive decay independently with a constant decay probability per unit time.

Let $lambda_d$ be the decay constant:

\begin{equation}
[\lambda_d]=\mathrm{s^{-1}}.
\tag{TRISO-GOV-104}
\label{eq:triso-gov-104}
\end{equation}

Consider an amount $N$ of the tracked species.

During a short time interval (dt), the expected fraction that decays is proportional to $\lambda_d dt$:

\begin{equation}
dN_{\mathrm{decay}}=\lambda_d N\,dt.
\tag{TRISO-GOV-105}
\label{eq:triso-gov-105}
\end{equation}

Because decay removes atoms from the tracked species, the change in tracked inventory is negative:

\begin{equation}
dN=-dN_{\mathrm{decay}}.
\tag{TRISO-GOV-106}
\label{eq:triso-gov-106}
\end{equation}

Substitute the decay amount:

\begin{equation}
dN=-\lambda_dN\,dt.
\tag{TRISO-GOV-107}
\label{eq:triso-gov-107}
\end{equation}

Divide by (dt):

\begin{equation}
\frac{dN}{dt}=-\lambda_dN.
\tag{TRISO-GOV-108}
\label{eq:triso-gov-108}
\end{equation}

For a fixed volume element, $N=c\,dV$. Therefore

\begin{equation}
\frac{d(c\,dV)}{dt}=-\lambda_dc\,dV.
\tag{TRISO-GOV-109}
\label{eq:triso-gov-109}
\end{equation}

For a fixed volume element, (dV) is constant in time:

\begin{equation}
\frac{\partial c}{\partial t}=-\lambda_dc.
\tag{TRISO-GOV-110}
\label{eq:triso-gov-110}
\end{equation}

Thus the local decay sink has the form

\begin{equation}
\boxed{
S_{i,\mathrm{decay}}=\lambda_{d,i}c_i
}
\tag{TRISO-GOV-111}
\label{eq:triso-gov-111}
\end{equation}

and it enters the conservation equation with a minus sign.

Dimensional check:

\begin{equation}
[\lambda_dc]
=
\mathrm{s^{-1}}\times\mathrm{mol\,m^{-3}}
=
\mathrm{mol\,m^{-3}\,s^{-1}}.
\tag{TRISO-GOV-112}
\label{eq:triso-gov-112}
\end{equation}

\textbf{[CONSTITUTIVE]} The first-order decay assumption is the mathematical statement that each tracked atom has the same constant decay hazard $lambda_d$, independent of concentration.

\textbf{[ASSUMPTION]} The current base benchmark sets this decay contribution to zero:

\begin{equation}
S_{i,\mathrm{decay}}=0.
\tag{TRISO-GOV-113}
\label{eq:triso-gov-113}
\end{equation}

\subsubsection*{5.3 Trapping and release from traps}

Trapping is a transfer of the tracked species from the mobile population into a trapped population. Release is the reverse transfer.

Define

\begin{equation}
S_{i,\mathrm{trap}}
\tag{TRISO-GOV-114}
\label{eq:triso-gov-114}
\end{equation}

as the positive rate at which mobile species are removed into traps, and

\begin{equation}
S_{i,\mathrm{release}}
\tag{TRISO-GOV-115}
\label{eq:triso-gov-115}
\end{equation}

as the positive rate at which trapped species are returned to the mobile population.

Their units are

\begin{equation}
[S_{i,\mathrm{trap}}]
=
[S_{i,\mathrm{release}}]
=
\mathrm{mol\,m^{-3}\,s^{-1}}.
\tag{TRISO-GOV-116}
\label{eq:triso-gov-116}
\end{equation}

The mobile-species equation therefore contains

\begin{equation}
S_{i,\mathrm{release}}-S_{i,\mathrm{trap}}.
\tag{TRISO-GOV-117}
\label{eq:triso-gov-117}
\end{equation}

No first-order, saturation, occupancy, or irradiation-dependent trapping law is assumed here.

\textbf{[QUESTION FOR SUPERVISOR]} Select the physical trapping/release constitutive model, if trapping is required in the final species model.

\subsubsection*{5.4 General reaction-inclusive equation}

Substitute the separated source terms into the conservation equation:

\begin{equation}
\frac{\partial c_i}{\partial t}
=
\frac{1}{r^2}
\frac{\partial}{\partial r}
\left(
r^2D_i\frac{\partial c_i}{\partial r}
\right)
+
S_{i,\mathrm{gen}}
+
S_{i,\mathrm{other}}
+
S_{i,\mathrm{release}}
-
S_{i,\mathrm{decay}}
-
S_{i,\mathrm{trap}}.
\tag{TRISO-GOV-118}
\label{eq:triso-gov-118}
\end{equation}

This is the general layer-wise source/reaction form under spherical symmetry.

\textbf{[QUESTION FOR SUPERVISOR]} Confirm whether the final physical model should contain decay and/or trapping in addition to diffusion.

\subsubsection*{5.5 Two source problems must remain distinct}

There are two different physical initial-value problems in this project.

\textbf{Problem A: initially empty particle with continuing generation}

The initial inventory is zero:

\begin{equation}
c_i(r,0)=0.
\tag{TRISO-IC-100}
\label{eq:triso-ic-100}
\end{equation}

The kernel generation remains active for $t>0$:

\begin{equation}
S_{1,\mathrm{gen}}=S_0.
\tag{TRISO-IC-101}
\label{eq:triso-ic-101}
\end{equation}

This is the source-driven problem used by the homogeneous Part-I steady analytical benchmark.

\textbf{Problem B: initially loaded kernel with no continuing generation}

The kernel initially contains mobile species:

\begin{equation}
c_1(r,0)=c_0
\qquad
0\le r<r_1,
\tag{TRISO-IC-102}
\label{eq:triso-ic-102}
\end{equation}

while the coatings initially contain none:

\begin{equation}
c_i(r,0)=0
\qquad
r_1<r<R.
\tag{TRISO-IC-103}
\label{eq:triso-ic-103}
\end{equation}

After $t=0$, the benchmark source is zero:

\begin{equation}
S_{i,\mathrm{gen}}=0.
\tag{TRISO-IC-104}
\label{eq:triso-ic-104}
\end{equation}

This is the initial-inventory release problem used by the production WOS verification contract.

These problems have different source terms and different physical histories. They must not be merged into one equation by notation alone.

\subsection*{6. Initial, centre, interface and outer conditions}

\subsubsection*{6.1 Centre condition from spherical symmetry}

At the exact centre there is no preferred radial direction.

\textbf{[ASSUMPTION]} Extend the radial concentration profile evenly through the mathematical origin for the purpose of the local limit:

\begin{equation}
c(-r,t)=c(r,t).
\tag{TRISO-BC-100}
\label{eq:triso-bc-100}
\end{equation}

Differentiate this relation with respect to $r$:

\begin{equation}
-c_r(-r,t)=c_r(r,t).
\tag{TRISO-BC-101}
\label{eq:triso-bc-101}
\end{equation}

Set $r=0$:

\begin{equation}
-c_r(0,t)=c_r(0,t).
\tag{TRISO-BC-102}
\label{eq:triso-bc-102}
\end{equation}

Therefore

\begin{equation}
\boxed{c_r(0,t)=0.}
\tag{TRISO-BC-103}
\label{eq:triso-bc-103}
\end{equation}

This is the mathematical expression of spherical symmetry at the centre.

It also agrees with the physical interpretation: a nonzero radial derivative at the centre would select one direction as different from the opposite direction.

\subsubsection*{6.2 Apparent singularity at the origin}

The spherical operator contains

\begin{equation}
\frac{2}{r}c_r.
\tag{TRISO-BC-104}
\label{eq:triso-bc-104}
\end{equation}

At $r=0$, this expression is of the form (0/0) for a smooth symmetric field, so it must not be evaluated by direct substitution.

Because $c_r(0,t)=0$, apply L'Hôpital's rule:

\begin{equation}
\lim_{r\to0}\frac{c_r(r,t)}{r}
=
\lim_{r\to0}\frac{c_{rr}(r,t)}{1}.
\tag{TRISO-BC-105}
\label{eq:triso-bc-105}
\end{equation}

Therefore

\begin{equation}
\lim_{r\to0}\frac{c_r(r,t)}{r}
=
c_{rr}(0,t).
\tag{TRISO-BC-106}
\label{eq:triso-bc-106}
\end{equation}

Multiply by 2:

\begin{equation}
\lim_{r\to0}\frac{2}{r}c_r(r,t)
=
2c_{rr}(0,t).
\tag{TRISO-BC-107}
\label{eq:triso-bc-107}
\end{equation}

The full spherical diffusion operator therefore has the centre limit

\begin{equation}
\lim_{r\to0}
\left(
c_{rr}+\frac2r c_r
\right)
=
c_{rr}(0,t)+2c_{rr}(0,t).
\tag{TRISO-BC-108}
\label{eq:triso-bc-108}
\end{equation}

Collect the two identical curvature contributions:

\begin{equation}
\boxed{
\lim_{r\to0}
\left(
c_{rr}+\frac2r c_r
\right)
=
3c_{rr}(0,t).
}
\tag{TRISO-BC-109}
\label{eq:triso-bc-109}
\end{equation}

This is the continuum origin of the factor 3 that later becomes the factor 6 in the central FTCS stencil.

\textbf{[ASSUMPTION]} The limit requires sufficient smoothness of the radial field near the origin.

\subsubsection*{6.3 Material-interface conservation}

Consider an infinitesimally thin spherical control volume surrounding interface $r=r_k$.

Let its inner radius be $r_k-\varepsilon$ and its outer radius be $r_k+\varepsilon$.

The volume is

\begin{equation}
V_\varepsilon
=
\frac{4\pi}{3}
\left[
(r_k+\varepsilon)^3-(r_k-\varepsilon)^3
\right].
\tag{TRISO-INT-100}
\label{eq:triso-int-100}
\end{equation}

As $\varepsilon\to0$,

\begin{equation}
V_\varepsilon\to0.
\tag{TRISO-INT-101}
\label{eq:triso-int-101}
\end{equation}

The inward diffusive amount rate crossing the inner surface is

\begin{equation}
4\pi(r_k-\varepsilon)^2J_{r,k}^{-}.
\tag{TRISO-INT-102}
\label{eq:triso-int-102}
\end{equation}

The outward diffusive amount rate crossing the outer surface is

\begin{equation}
4\pi(r_k+\varepsilon)^2J_{r,k}^{+}.
\tag{TRISO-INT-103}
\label{eq:triso-int-103}
\end{equation}

Let $\Gamma_k$ be any explicitly modelled interfacial inventory per unit area, and let $g_k$ be any explicitly modelled interfacial production rate per unit area.

Their units are

\begin{equation}
[\Gamma_k]=\mathrm{mol\,m^{-2}},
\qquad
[g_k]=\mathrm{mol\,m^{-2}\,s^{-1}}.
\tag{TRISO-INT-104}
\label{eq:triso-int-104}
\end{equation}

The interface balance is

\begin{equation}
\frac{d}{dt}
\left(
4\pi r_k^2\Gamma_k
\right)
=
4\pi(r_k-\varepsilon)^2J_{r,k}^{-}
-
4\pi(r_k+\varepsilon)^2J_{r,k}^{+}
+
4\pi r_k^2g_k
+
o(1).
\tag{TRISO-INT-105}
\label{eq:triso-int-105}
\end{equation}

Divide by $4\pi r_k^2$:

\begin{equation}
\frac{d\Gamma_k}{dt}
=
\left(
\frac{r_k-\varepsilon}{r_k}
\right)^2
J_{r,k}^{-}
-
\left(
\frac{r_k+\varepsilon}{r_k}
\right)^2
J_{r,k}^{+}
+
g_k
+
o(1).
\tag{TRISO-INT-106}
\label{eq:triso-int-106}
\end{equation}

Take the zero-thickness limit:

\begin{equation}
\frac{d\Gamma_k}{dt}
=
J_{r,k}^{-}
-
J_{r,k}^{+}
+
g_k.
\tag{TRISO-INT-107}
\label{eq:triso-int-107}
\end{equation}

Therefore the general interface jump condition is

\begin{equation}
\boxed{
J_{r,k}^{-}-J_{r,k}^{+}
=
\frac{d\Gamma_k}{dt}-g_k.
}
\tag{TRISO-INT-108}
\label{eq:triso-int-108}
\end{equation}

For an ideal interface with no interfacial storage and no interfacial generation,

\begin{equation}
\Gamma_k=0,
\qquad
g_k=0.
\tag{TRISO-INT-109}
\label{eq:triso-int-109}
\end{equation}

Hence

\begin{equation}
J_{r,k}^{-}=J_{r,k}^{+}.
\tag{TRISO-INT-110}
\label{eq:triso-int-110}
\end{equation}

Thus

\begin{equation}
\boxed{
J_{r,k}^{-}=J_{r,k}^{+}.
}
\tag{TRISO-INT-111}
\label{eq:triso-int-111}
\end{equation}

Now substitute Fick's law on the inner side:

\begin{equation}
J_{r,k}^{-}
=
-D_k
\frac{\partial c_k}{\partial r}
\bigg|_{r_k^-}.
\tag{TRISO-INT-112}
\label{eq:triso-int-112}
\end{equation}

Substitute Fick's law on the outer side:

\begin{equation}
J_{r,k}^{+}
=
-D_{k+1}
\frac{\partial c_{k+1}}{\partial r}
\bigg|_{r_k^+}.
\tag{TRISO-INT-113}
\label{eq:triso-int-113}
\end{equation}

Equate them:

\begin{equation}
\boxed{
-D_k
\frac{\partial c_k}{\partial r}
\bigg|_{r_k^-}
=
-D_{k+1}
\frac{\partial c_{k+1}}{\partial r}
\bigg|_{r_k^+}.
}
\tag{TRISO-INT-114}
\label{eq:triso-int-114}
\end{equation}

This condition came from conservation. It did not require concentration continuity.

\subsubsection*{6.4 Concentration continuity is a separate interface assumption}

Flux continuity answers the question:

> Is species amount conserved across the zero-thickness interface?

It does not answer:

> What equilibrium relation connects the two concentrations at the interface?

For an ideal perfectly equilibrated interface, one may impose concentration continuity:

\begin{equation}
\boxed{
c_k(r_k,t)=c_{k+1}(r_k,t).
}
\tag{TRISO-INT-115}
\label{eq:triso-int-115}
\end{equation}

\textbf{[ASSUMPTION]} This is an ideal-interface constitutive/equilibrium assumption.

A species-specific partition coefficient $K_k$ instead gives a different relation:

\begin{equation}
\boxed{
c_{k+1}(r_k,t)=K_kc_k(r_k,t).
}
\tag{TRISO-INT-116}
\label{eq:triso-int-116}
\end{equation}

The value and definition of $K_k$ depend on the species and the two materials.

\textbf{[SOURCE NEEDED]} Species-specific partition/solubility data if $K_k\ne1$ is required physically.

\subsubsection*{6.5 Interfacial resistance is a third, distinct model}

Partitioning and interfacial resistance are not the same statement.

A finite interfacial mass-transfer coefficient $h_{\mathrm{int}}$ can instead be used in a constitutive resistance law such as

\begin{equation}
J_{r,k}
=
h_{\mathrm{int}}
\left(
c_k-\frac{c_{k+1}}{K_k}
\right).
\tag{TRISO-INT-117}
\label{eq:triso-int-117}
\end{equation}

The units are

\begin{equation}
[h_{\mathrm{int}}]=\mathrm{m\,s^{-1}}.
\tag{TRISO-INT-118}
\label{eq:triso-int-118}
\end{equation}

This relation introduces a finite concentration jump for finite resistance.

\textbf{[SOURCE NEEDED]} The physical interfacial-resistance law and coefficient, if such resistance is required.

\textbf{[ASSUMPTION]} The frozen numerical benchmark uses $K_k=1$ and no explicit interfacial resistance, so (TRISO-INT-115) and (TRISO-INT-117) are not simultaneously imposed.

\subsubsection*{6.6 Outer boundary conditions}

At the outer surface $r=R$, define the outward radial flux as

\begin{equation}
J_R=J_r(R,t).
\tag{TRISO-BC-110}
\label{eq:triso-bc-110}
\end{equation}

By Fick's law,

\begin{equation}
J_R
=
-D_5
\frac{\partial c_5}{\partial r}
\bigg|_{R}.
\tag{TRISO-BC-111}
\label{eq:triso-bc-111}
\end{equation}

\#\#\#\# Dirichlet: absorbing or prescribed surface concentration

The simplest absorbing boundary is

\begin{equation}
\boxed{
c_5(R,t)=0.
}
\tag{TRISO-BC-112}
\label{eq:triso-bc-112}
\end{equation}

More generally, a prescribed surface concentration $c_b(t)$ is

\begin{equation}
c_5(R,t)=c_b(t).
\tag{TRISO-BC-113}
\label{eq:triso-bc-113}
\end{equation}

The absorbing case is $c_b=0$.

\textbf{[ASSUMPTION]} The production WOS verification benchmark uses the absorbing case.

\#\#\#\# Neumann: prescribed outward flux

A prescribed outward flux is written

\begin{equation}
\boxed{
J_R=J_b(t).
}
\tag{TRISO-BC-114}
\label{eq:triso-bc-114}
\end{equation}

Substitute the diffusive flux:

\begin{equation}
-D_5
\frac{\partial c_5}{\partial r}
\bigg|_R
=
J_b(t).
\tag{TRISO-BC-115}
\label{eq:triso-bc-115}
\end{equation}

The units are

\begin{equation}
[J_b]=\mathrm{mol\,m^{-2}\,s^{-1}}.
\tag{TRISO-BC-116}
\label{eq:triso-bc-116}
\end{equation}

\#\#\#\# Robin: finite external mass transfer

Let the external coolant concentration be $c_\infty(t)$.

\textbf{[CONSTITUTIVE]} A linear external mass-transfer law is

\begin{equation}
J_R
=
h
\left[
c_5(R,t)-c_\infty(t)
\right].
\tag{TRISO-BC-117}
\label{eq:triso-bc-117}
\end{equation}

The units of $h$ are

\begin{equation}
[h]=\mathrm{m\,s^{-1}}.
\tag{TRISO-BC-118}
\label{eq:triso-bc-118}
\end{equation}

Substitute the diffusive surface flux:

\begin{equation}
-D_5
\frac{\partial c_5}{\partial r}
\bigg|_R
=
h
\left[
c_5(R,t)-c_\infty(t)
\right].
\tag{TRISO-BC-119}
\label{eq:triso-bc-119}
\end{equation}

Both sides have units

\begin{equation}
\mathrm{mol\,m^{-2}\,s^{-1}}.
\tag{TRISO-BC-120}
\label{eq:triso-bc-120}
\end{equation}

For zero bulk concentration,

\begin{equation}
c_\infty=0,
\tag{TRISO-BC-121}
\label{eq:triso-bc-121}
\end{equation}

so

\begin{equation}
-D_5c_5'(R,t)=hc_5(R,t).
\tag{TRISO-BC-122}
\label{eq:triso-bc-122}
\end{equation}

For very large $h$, a finite flux requires

\begin{equation}
c_5(R,t)-c_\infty(t)\to0.
\tag{TRISO-BC-123}
\label{eq:triso-bc-123}
\end{equation}

Thus the Robin condition approaches the Dirichlet condition

\begin{equation}
c_5(R,t)=c_\infty(t).
\tag{TRISO-BC-124}
\label{eq:triso-bc-124}
\end{equation}

For $h\to0$,

\begin{equation}
J_R\to0,
\tag{TRISO-BC-125}
\label{eq:triso-bc-125}
\end{equation}

which approaches the zero-flux Neumann condition.

The three outer-boundary models are therefore mathematically distinct:

\begin{equation}
\text{Dirichlet: concentration prescribed},
\end{equation}

\begin{equation}
\text{Neumann: flux prescribed},
\end{equation}

\begin{equation}
\text{Robin: flux responds to concentration difference}.
\end{equation}

\textbf{[ASSUMPTION]} The Part-I analytical benchmark uses Robin with $c_\infty=0$.

\textbf{[ASSUMPTION]} The production WOS verification benchmark uses absorbing Dirichlet (c\_5(R,t)=0).

\subsubsection*{6.7 Initial-condition summary}

The initial condition must be selected together with the source model.

Problem A:

\begin{equation}
c_i(r,0)=0,
\qquad
S_{1,\mathrm{gen}}=S_0.
\tag{TRISO-IC-105}
\label{eq:triso-ic-105}
\end{equation}

Problem B:

\begin{equation}
c_1(r,0)=c_0,
\qquad
c_{2..5}(r,0)=0,
\qquad
S_{i,\mathrm{gen}}=0.
\tag{TRISO-IC-106}
\label{eq:triso-ic-106}
\end{equation}

The first homogeneous analytical Part-I benchmark belongs to Problem A.

The production WOS release benchmark belongs to Problem B.

\subsection*{7. Part I homogeneous analytical benchmark}

The following solution belongs to \textbf{Problem A}: an initially empty homogeneous sphere with a continuing uniform source and a finite-transfer Robin boundary.

\textbf{[ASSUMPTION]} Replace the five-layer particle temporarily by one homogeneous sphere of radius $R$ and constant diffusivity $D$.

The governing equation is

\begin{equation}
\frac{\partial c}{\partial t}
=
D
\left(
\frac{\partial^2c}{\partial r^2}
+
\frac2r
\frac{\partial c}{\partial r}
\right)
+
S_0.
\tag{TRISO-ANA-100}
\label{eq:triso-ana-100}
\end{equation}

At steady state the concentration no longer changes with time:

\begin{equation}
\frac{\partial w}{\partial t}=0.
\tag{TRISO-ANA-101}
\label{eq:triso-ana-101}
\end{equation}

Therefore

\begin{equation}
0
=
D
\left(
\frac{d^2w}{dr^2}
+
\frac2r\frac{dw}{dr}
\right)
+
S_0.
\tag{TRISO-ANA-102}
\label{eq:triso-ana-102}
\end{equation}

Move the source term to the other side:

\begin{equation}
D
\left(
w''
+
\frac2r w'
\right)
=
-S_0.
\tag{TRISO-ANA-103}
\label{eq:triso-ana-103}
\end{equation}

Divide by $D$:

\begin{equation}
w''
+
\frac2r w'
=
-\frac{S_0}{D}.
\tag{TRISO-ANA-104}
\label{eq:triso-ana-104}
\end{equation}

Multiply by $r^2$:

\begin{equation}
r^2w''
+
2rw'
=
-\frac{S_0}{D}r^2.
\tag{TRISO-ANA-105}
\label{eq:triso-ana-105}
\end{equation}

Recognise the product derivative on the left. Verify it explicitly using the product rule:

\begin{equation}
\frac{d}{dr}(r^2w')
=
\frac{d r^2}{dr}w'
+
r^2w''.
\tag{TRISO-ANA-106}
\label{eq:triso-ana-106}
\end{equation}

Differentiate $r^2$:

\begin{equation}
\frac{d r^2}{dr}=2r.
\tag{TRISO-ANA-107}
\label{eq:triso-ana-107}
\end{equation}

Therefore

\begin{equation}
\frac{d}{dr}(r^2w')
=
2rw'+r^2w''.
\tag{TRISO-ANA-108}
\label{eq:triso-ana-108}
\end{equation}

Hence the steady equation becomes

\begin{equation}
\frac{d}{dr}(r^2w')
=
-\frac{S_0}{D}r^2.
\tag{TRISO-ANA-109}
\label{eq:triso-ana-109}
\end{equation}

Integrate both sides with respect to $r$:

\begin{equation}
\int
\frac{d}{dr}(r^2w')\,dr
=
-\frac{S_0}{D}
\int r^2\,dr.
\tag{TRISO-ANA-110}
\label{eq:triso-ana-110}
\end{equation}

The left-hand integral is

\begin{equation}
r^2w'.
\tag{TRISO-ANA-111}
\label{eq:triso-ana-111}
\end{equation}

The right-hand integral is

\begin{equation}
-\frac{S_0}{D}\frac{r^3}{3}.
\tag{TRISO-ANA-112}
\label{eq:triso-ana-112}
\end{equation}

Introduce the integration constant $A$:

\begin{equation}
r^2w'
=
-\frac{S_0r^3}{3D}
+
A.
\tag{TRISO-ANA-113}
\label{eq:triso-ana-113}
\end{equation}

At the centre, regularity requires $w'(0)$ to remain finite.

If $A\ne0$, then division by $r^2$ gives a term proportional to $1/r^2$, which diverges.

Therefore

\begin{equation}
A=0.
\tag{TRISO-ANA-114}
\label{eq:triso-ana-114}
\end{equation}

Substitute $A=0$:

\begin{equation}
r^2w'
=
-\frac{S_0r^3}{3D}.
\tag{TRISO-ANA-115}
\label{eq:triso-ana-115}
\end{equation}

For $r>0$, divide by $r^2$:

\begin{equation}
w'
=
-\frac{S_0r}{3D}.
\tag{TRISO-ANA-116}
\label{eq:triso-ana-116}
\end{equation}

Integrate again:

\begin{equation}
\int dw
=
-\frac{S_0}{3D}\int r\,dr.
\tag{TRISO-ANA-117}
\label{eq:triso-ana-117}
\end{equation}

The left-hand integral is

\begin{equation}
w.
\tag{TRISO-ANA-118}
\label{eq:triso-ana-118}
\end{equation}

The right-hand integral is

\begin{equation}
-\frac{S_0}{3D}\frac{r^2}{2}.
\tag{TRISO-ANA-119}
\label{eq:triso-ana-119}
\end{equation}

Introduce the second integration constant $B$:

\begin{equation}
w
=
B
-
\frac{S_0r^2}{6D}.
\tag{TRISO-ANA-120}
\label{eq:triso-ana-120}
\end{equation}

Now apply the Robin boundary condition at $r=R$:

\begin{equation}
-Dw'(R)=hw(R)
\tag{TRISO-ANA-121}
\label{eq:triso-ana-121}
\end{equation}

because $c_\infty=0$ for this benchmark.

Evaluate the derivative at $R$:

\begin{equation}
w'(R)
=
-\frac{S_0R}{3D}.
\tag{TRISO-ANA-122}
\label{eq:triso-ana-122}
\end{equation}

Multiply by (-D):

\begin{equation}
-Dw'(R)
=
\frac{S_0R}{3}.
\tag{TRISO-ANA-123}
\label{eq:triso-ana-123}
\end{equation}

Evaluate the concentration at $R$:

\begin{equation}
w(R)
=
B
-
\frac{S_0R^2}{6D}.
\tag{TRISO-ANA-124}
\label{eq:triso-ana-124}
\end{equation}

Substitute both expressions into the Robin condition:

\begin{equation}
\frac{S_0R}{3}
=
h
\left(
B
-
\frac{S_0R^2}{6D}
\right).
\tag{TRISO-ANA-125}
\label{eq:triso-ana-125}
\end{equation}

Divide by $h$:

\begin{equation}
\frac{S_0R}{3h}
=
B
-
\frac{S_0R^2}{6D}.
\tag{TRISO-ANA-126}
\label{eq:triso-ana-126}
\end{equation}

Add the quadratic term to both sides:

\begin{equation}
B
=
\frac{S_0R}{3h}
+
\frac{S_0R^2}{6D}.
\tag{TRISO-ANA-127}
\label{eq:triso-ana-127}
\end{equation}

Substitute $B$ into the profile:

\begin{equation}
w(r)
=
\frac{S_0R}{3h}
+
\frac{S_0R^2}{6D}
-
\frac{S_0r^2}{6D}.
\tag{TRISO-ANA-128}
\label{eq:triso-ana-128}
\end{equation}

Collect the two quadratic terms:

\begin{equation}
\boxed{
w(r)
=
\frac{S_0R}{3h}
+
\frac{S_0}{6D}
\left(
R^2-r^2
\right).
}
\tag{TRISO-ANA-129}
\label{eq:triso-ana-129}
\end{equation}

\subsubsection*{7.1 Dimensional check}

The first term has units

\begin{equation}
\left[
\frac{S_0R}{h}
\right]
=
\frac{
\mathrm{mol\,m^{-3}\,s^{-1}}
\times
\mathrm m
}{
\mathrm{m\,s^{-1}}
}
=
\mathrm{mol\,m^{-3}}.
\tag{TRISO-ANA-130}
\label{eq:triso-ana-130}
\end{equation}

The second term has units

\begin{equation}
\left[
\frac{S_0R^2}{D}
\right]
=
\frac{
\mathrm{mol\,m^{-3}\,s^{-1}}
\times
\mathrm{m^2}
}{
\mathrm{m^2\,s^{-1}}
}
=
\mathrm{mol\,m^{-3}}.
\tag{TRISO-ANA-131}
\label{eq:triso-ana-131}
\end{equation}

Both terms therefore have concentration units.

\subsubsection*{7.2 Independent global generation/release balance}

The total generation rate in the homogeneous sphere is

\begin{equation}
\dot N_{\mathrm{gen}}
=
\int_0^R
S_0\,4\pi r^2\,dr.
\tag{TRISO-ANA-132}
\label{eq:triso-ana-132}
\end{equation}

Pull out constants:

\begin{equation}
\dot N_{\mathrm{gen}}
=
4\pi S_0
\int_0^Rr^2\,dr.
\tag{TRISO-ANA-133}
\label{eq:triso-ana-133}
\end{equation}

Evaluate the integral:

\begin{equation}
\int_0^Rr^2\,dr
=
\frac{R^3}{3}.
\tag{TRISO-ANA-134}
\label{eq:triso-ana-134}
\end{equation}

Therefore

\begin{equation}
\boxed{
\dot N_{\mathrm{gen}}
=
\frac{4\pi R^3S_0}{3}.
}
\tag{TRISO-ANA-135}
\label{eq:triso-ana-135}
\end{equation}

At steady state, this must equal the outward surface release rate:

\begin{equation}
\dot N_{\mathrm{out}}
=
4\pi R^2
\left[
h w(R)
\right].
\tag{TRISO-ANA-136}
\label{eq:triso-ana-136}
\end{equation}

Set generation equal to release:

\begin{equation}
\frac{4\pi R^3S_0}{3}
=
4\pi R^2h w(R).
\tag{TRISO-ANA-137}
\label{eq:triso-ana-137}
\end{equation}

Cancel $4\pi R^2$:

\begin{equation}
\frac{S_0R}{3}
=
h w(R).
\tag{TRISO-ANA-138}
\label{eq:triso-ana-138}
\end{equation}

Divide by $h$:

\begin{equation}
w(R)=\frac{S_0R}{3h}.
\tag{TRISO-ANA-139}
\label{eq:triso-ana-139}
\end{equation}

Evaluate the analytical profile at $r=R$:

\begin{equation}
w(R)
=
\frac{S_0R}{3h}
+
\frac{S_0}{6D}(R^2-R^2).
\tag{TRISO-ANA-140}
\label{eq:triso-ana-140}
\end{equation}

The second term is zero:

\begin{equation}
w(R)
=
\frac{S_0R}{3h}.
\tag{TRISO-ANA-141}
\label{eq:triso-ana-141}
\end{equation}

The independently derived balance and the analytical solution therefore agree.

\textbf{[VERIFIED]} The homogeneous steady benchmark is internally consistent under its stated assumptions.

\subsection*{7.3 Scope of this benchmark}

This analytical solution does \textbf{not} represent the production WOS initial-kernel release problem.

It represents Problem A:

\noindent\textbullet\ initially empty particle;\\
\noindent\textbullet\ continuing uniform source;\\
\noindent\textbullet\ homogeneous diffusivity;\\
\noindent\textbullet\ Robin outer boundary.\\

The production WOS verification problem is Problem B and uses an initial kernel inventory with an absorbing outer boundary.

Keeping these benchmarks separate prevents source and boundary semantics from being mixed.

\subsection*{8. Homogeneous transient deviation problem}

The following transient benchmark belongs to Problem A: an initially empty homogeneous sphere with a continuing uniform source and a Robin outer boundary.

The steady solution $w(r)$ has already been obtained. We now remove the steady part so that the remaining transient problem has no source term.

\subsubsection*{8.1 Define the transient deviation}

Define

\begin{equation}
v(r,t)=c(r,t)-w(r).
\tag{TRISO-ANA-200}
\label{eq:triso-ana-200}
\end{equation}

Rearrange this definition:

\begin{equation}
c(r,t)=v(r,t)+w(r).
\tag{TRISO-ANA-201}
\label{eq:triso-ana-201}
\end{equation}

Because $w$ is a steady solution, it does not depend on time:

\begin{equation}
\frac{\partial w}{\partial t}=0.
\tag{TRISO-ANA-202}
\label{eq:triso-ana-202}
\end{equation}

Differentiate $c=v+w$ with respect to time:

\begin{equation}
\frac{\partial c}{\partial t}
=
\frac{\partial v}{\partial t}
+
\frac{\partial w}{\partial t}.
\tag{TRISO-ANA-203}
\label{eq:triso-ana-203}
\end{equation}

Substitute (TRISO-ANA-202):

\begin{equation}
\boxed{
\frac{\partial c}{\partial t}
=
\frac{\partial v}{\partial t}.
}
\tag{TRISO-ANA-204}
\label{eq:triso-ana-204}
\end{equation}

Differentiate $c=v+w$ with respect to radius:

\begin{equation}
\frac{\partial c}{\partial r}
=
\frac{\partial v}{\partial r}
+
\frac{dw}{dr}.
\tag{TRISO-ANA-205}
\label{eq:triso-ana-205}
\end{equation}

Differentiate once more:

\begin{equation}
\frac{\partial^2 c}{\partial r^2}
=
\frac{\partial^2v}{\partial r^2}
+
\frac{d^2w}{dr^2}.
\tag{TRISO-ANA-206}
\label{eq:triso-ana-206}
\end{equation}

Start from the full homogeneous source-driven PDE:

\begin{equation}
\frac{\partial c}{\partial t}
=
D
\left(
\frac{\partial^2c}{\partial r^2}
+
\frac2r\frac{\partial c}{\partial r}
\right)
+
S_0.
\tag{TRISO-ANA-207}
\label{eq:triso-ana-207}
\end{equation}

Substitute the time derivative from (TRISO-ANA-204):

\begin{equation}
\frac{\partial v}{\partial t}
=
D
\left(
\frac{\partial^2c}{\partial r^2}
+
\frac2r\frac{\partial c}{\partial r}
\right)
+
S_0.
\tag{TRISO-ANA-208}
\label{eq:triso-ana-208}
\end{equation}

Substitute the second spatial derivative from (TRISO-ANA-206):

\begin{equation}
\frac{\partial v}{\partial t}
=
D
\left(
\frac{\partial^2v}{\partial r^2}
+
\frac{d^2w}{dr^2}
+
\frac2r\frac{\partial c}{\partial r}
\right)
+
S_0.
\tag{TRISO-ANA-209}
\label{eq:triso-ana-209}
\end{equation}

Substitute the first spatial derivative from (TRISO-ANA-205):

\begin{equation}
\frac{\partial v}{\partial t}
=
D
\left(
\frac{\partial^2v}{\partial r^2}
+
\frac{d^2w}{dr^2}
+
\frac2r
\left[
\frac{\partial v}{\partial r}
+
\frac{dw}{dr}
\right]
\right)
+
S_0.
\tag{TRISO-ANA-210}
\label{eq:triso-ana-210}
\end{equation}

Distribute the factor $D$:

\begin{equation}
\frac{\partial v}{\partial t}
=
D\frac{\partial^2v}{\partial r^2}
+
D\frac{d^2w}{dr^2}
+
\frac{2D}{r}\frac{\partial v}{\partial r}
+
\frac{2D}{r}\frac{dw}{dr}
+
S_0.
\tag{TRISO-ANA-211}
\label{eq:triso-ana-211}
\end{equation}

Rearrange the terms into transient and steady groups:

\begin{equation}
\frac{\partial v}{\partial t}
=
D
\left(
\frac{\partial^2v}{\partial r^2}
+
\frac2r\frac{\partial v}{\partial r}
\right)
+
\left[
D
\left(
\frac{d^2w}{dr^2}
+
\frac2r\frac{dw}{dr}
\right)
+
S_0
\right].
\tag{TRISO-ANA-212}
\label{eq:triso-ana-212}
\end{equation}

The steady solution satisfies

\begin{equation}
D
\left(
\frac{d^2w}{dr^2}
+
\frac2r\frac{dw}{dr}
\right)
+
S_0
=
0.
\tag{TRISO-ANA-213}
\label{eq:triso-ana-213}
\end{equation}

Substitute (TRISO-ANA-213) into (TRISO-ANA-212):

\begin{equation}
\boxed{
\frac{\partial v}{\partial t}
=
D
\left(
\frac{\partial^2v}{\partial r^2}
+
\frac2r\frac{\partial v}{\partial r}
\right).
}
\tag{TRISO-ANA-214}
\label{eq:triso-ana-214}
\end{equation}

The source has disappeared because the steady part $w$ already accounts for the long-time source balance.

\subsubsection*{8.2 Transform the centre condition}

The original centre condition is

\begin{equation}
\frac{\partial c}{\partial r}(0,t)=0.
\tag{TRISO-ANA-215}
\label{eq:triso-ana-215}
\end{equation}

Substitute (TRISO-ANA-205):

\begin{equation}
\frac{\partial v}{\partial r}(0,t)+w'(0)=0.
\tag{TRISO-ANA-216}
\label{eq:triso-ana-216}
\end{equation}

The steady solution has

\begin{equation}
w'(r)=-\frac{S_0r}{3D}.
\tag{TRISO-ANA-217}
\label{eq:triso-ana-217}
\end{equation}

Evaluate it at the centre:

\begin{equation}
w'(0)=0.
\tag{TRISO-ANA-218}
\label{eq:triso-ana-218}
\end{equation}

Therefore

\begin{equation}
\boxed{
\frac{\partial v}{\partial r}(0,t)=0.
}
\tag{TRISO-ANA-219}
\label{eq:triso-ana-219}
\end{equation}

\subsubsection*{8.3 Transform the outer Robin condition}

The original Robin condition is

\begin{equation}
-Dc_r(R,t)=hw(R,t).
\tag{TRISO-ANA-220}
\label{eq:triso-ana-220}
\end{equation}

Substitute $c=v+w$:

\begin{equation}
-D
\left[
v_r(R,t)+w'(R)
\right]
=
h
\left[
v(R,t)+w(R)
\right].
\tag{TRISO-ANA-221}
\label{eq:triso-ana-221}
\end{equation}

Rearrange the transient and steady terms:

\begin{equation}
-Dv_r(R,t)-hw(R)
=
hv(R,t)+Dw'(R).
\tag{TRISO-ANA-222}
\label{eq:triso-ana-222}
\end{equation}

The steady solution satisfies

\begin{equation}
-Dw'(R)=hw(R).
\tag{TRISO-ANA-223}
\label{eq:triso-ana-223}
\end{equation}

Therefore

\begin{equation}
-Dv_r(R,t)=hv(R,t).
\tag{TRISO-ANA-224}
\label{eq:triso-ana-224}
\end{equation}

Hence the transient Robin condition is

\begin{equation}
\boxed{
-Dv_r(R,t)=hv(R,t).
}
\tag{TRISO-ANA-225}
\label{eq:triso-ana-225}
\end{equation}

\subsubsection*{8.4 Transform the initial condition}

The original initial condition for Problem A is

\begin{equation}
c(r,0)=0.
\tag{TRISO-ANA-226}
\label{eq:triso-ana-226}
\end{equation}

Apply the definition $v=c-w$:

\begin{equation}
v(r,0)=c(r,0)-w(r).
\tag{TRISO-ANA-227}
\label{eq:triso-ana-227}
\end{equation}

Substitute $c(r,0)=0$:

\begin{equation}
\boxed{
v(r,0)=-w(r).
}
\tag{TRISO-ANA-228}
\label{eq:triso-ana-228}
\end{equation}

Thus the complete transient deviation problem is

\begin{equation}
v_t
=
D
\left(
v_{rr}+\frac2r v_r
\right),
\end{equation}

with

\begin{equation}
v_r(0,t)=0,
\end{equation}

\begin{equation}
-Dv_r(R,t)=hv(R,t),
\end{equation}

and

\begin{equation}
v(r,0)=-w(r).
\end{equation}

\subsection*{9. Separation of variables}

\subsubsection*{9.1 Assume a separated solution}

Seek a non-zero transient mode in the form

\begin{equation}
v(r,t)=\phi(r)T(t).
\tag{TRISO-ANA-229}
\label{eq:triso-ana-229}
\end{equation}

Differentiate with respect to time:

\begin{equation}
v_t=\phi(r)T'(t).
\tag{TRISO-ANA-230}
\label{eq:triso-ana-230}
\end{equation}

Differentiate with respect to radius:

\begin{equation}
v_r=\phi'(r)T(t).
\tag{TRISO-ANA-231}
\label{eq:triso-ana-231}
\end{equation}

Differentiate once more:

\begin{equation}
v_{rr}=\phi''(r)T(t).
\tag{TRISO-ANA-232}
\label{eq:triso-ana-232}
\end{equation}

Substitute these three expressions into the transient PDE:

\begin{equation}
\phi T'
=
D
\left[
\phi''T
+
\frac2r\phi'T
\right].
\tag{TRISO-ANA-233}
\label{eq:triso-ana-233}
\end{equation}

Factor out $T$ on the right:

\begin{equation}
\phi T'
=
DT
\left(
\phi''+\frac2r\phi'
\right).
\tag{TRISO-ANA-234}
\label{eq:triso-ana-234}
\end{equation}

Divide by $D\phi T$, assuming the separated factors are non-zero at the point considered:

\begin{equation}
\frac{T'}{DT}
=
\frac{\phi''+2\phi'/r}{\phi}.
\tag{TRISO-ANA-235}
\label{eq:triso-ana-235}
\end{equation}

The left side depends only on $t$, while the right side depends only on $r$.

For one separated mode to satisfy the equation for every $r$ and $t$, both sides must equal the same constant.

Choose the separation constant as $-k^2$:

\begin{equation}
\frac{T'}{DT}=-k^2.
\tag{TRISO-ANA-236}
\label{eq:triso-ana-236}
\end{equation}

Then the time equation is

\begin{equation}
T'=-Dk^2T.
\tag{TRISO-ANA-237}
\label{eq:triso-ana-237}
\end{equation}

The radial equation is

\begin{equation}
\boxed{
\phi''
+
\frac2r\phi'
+
k^2\phi
=
0.
}
\tag{TRISO-ANA-238}
\label{eq:triso-ana-238}
\end{equation}

\subsubsection*{9.2 Dimensions of the separation constant}

The left side of (TRISO-ANA-236) has units

\begin{equation}
\left[
\frac{T'}{DT}
\right]
=
\frac{\mathrm{s^{-1}}}{\mathrm{m^2\,s^{-1}}}
=
\mathrm{m^{-2}}.
\tag{TRISO-ANA-239}
\label{eq:triso-ana-239}
\end{equation}

Therefore

\begin{equation}
[k^2]=\mathrm{m^{-2}}.
\tag{TRISO-ANA-240}
\label{eq:triso-ana-240}
\end{equation}

Hence

\begin{equation}
[k]=\mathrm{m^{-1}}.
\tag{TRISO-ANA-241}
\label{eq:triso-ana-241}
\end{equation}

Define the temporal decay rate

\begin{equation}
\boxed{
\Lambda=Dk^2.
}
\tag{TRISO-ANA-242}
\label{eq:triso-ana-242}
\end{equation}

Then

\begin{equation}
[\Lambda]
=
\mathrm{m^2\,s^{-1}}
\times
\mathrm{m^{-2}}
=
\mathrm{s^{-1}}.
\tag{TRISO-ANA-243}
\label{eq:triso-ana-243}
\end{equation}

Thus $k$ is a spatial wave number, while $\Lambda$ is a temporal decay rate.

\subsubsection*{9.3 Solve the temporal equation}

Starting from

\begin{equation}
T'=-\Lambda T,
\tag{TRISO-ANA-244}
\label{eq:triso-ana-244}
\end{equation}

divide by $T$:

\begin{equation}
\frac{T'}{T}=-\Lambda.
\tag{TRISO-ANA-245}
\label{eq:triso-ana-245}
\end{equation}

Write the derivative as a differential:

\begin{equation}
\frac{dT}{T}=-\Lambda\,dt.
\tag{TRISO-ANA-246}
\label{eq:triso-ana-246}
\end{equation}

Integrate:

\begin{equation}
\int\frac{dT}{T}
=
-\Lambda\int dt.
\tag{TRISO-ANA-247}
\label{eq:triso-ana-247}
\end{equation}

Therefore

\begin{equation}
\ln|T|
=
-\Lambda t+C.
\tag{TRISO-ANA-248}
\label{eq:triso-ana-248}
\end{equation}

Exponentiate:

\begin{equation}
|T|=e^{C}e^{-\Lambda t}.
\tag{TRISO-ANA-249}
\label{eq:triso-ana-249}
\end{equation}

Absorb the constant into an arbitrary amplitude $C_T$:

\begin{equation}
T(t)=C_Te^{-\Lambda t}.
\tag{TRISO-ANA-250}
\label{eq:triso-ana-250}
\end{equation}

The constant $C_T$ can be absorbed into the spatial amplitude, so take

\begin{equation}
\boxed{
T(t)=e^{-\Lambda t}.
}
\tag{TRISO-ANA-251}
\label{eq:triso-ana-251}
\end{equation}

Therefore each separated mode has the form

\begin{equation}
v(r,t)=\phi(r)e^{-\Lambda t}.
\tag{TRISO-ANA-252}
\label{eq:triso-ana-252}
\end{equation}

\subsection*{10. Radial eigenproblem, Robin condition, and modal expansion}

\subsubsection*{10.1 Transform the radial eigenproblem with $u=r\phi$}

Start from

\begin{equation}
\phi''
+
\frac2r\phi'
+
k^2\phi
=
0.
\tag{TRISO-ANA-253}
\label{eq:triso-ana-253}
\end{equation}

Introduce

\begin{equation}
u(r)=r\phi(r).
\tag{TRISO-ANA-254}
\label{eq:triso-ana-254}
\end{equation}

Solve the definition for $\phi$:

\begin{equation}
\phi(r)=\frac{u(r)}{r}.
\tag{TRISO-ANA-255}
\label{eq:triso-ana-255}
\end{equation}

Differentiate using the quotient rule:

\begin{equation}
\phi'
=
\frac{r u'-u}{r^2}.
\tag{TRISO-ANA-256}
\label{eq:triso-ana-256}
\end{equation}

Differentiate again. Write the numerator as $n=ru'-u$:

\begin{equation}
n'=u'+ru''-u'.
\tag{TRISO-ANA-257}
\label{eq:triso-ana-257}
\end{equation}

Therefore

\begin{equation}
n'=ru''.
\tag{TRISO-ANA-258}
\label{eq:triso-ana-258}
\end{equation}

Apply the quotient rule to $n/r^2$:

\begin{equation}
\phi''
=
\frac{n'r^2-n(2r)}{r^4}.
\tag{TRISO-ANA-259}
\label{eq:triso-ana-259}
\end{equation}

Substitute $n'=ru''$ and $n=ru'-u$:

\begin{equation}
\phi''
=
\frac{r^3u''-2r(ru'-u)}{r^4}.
\tag{TRISO-ANA-260}
\label{eq:triso-ana-260}
\end{equation}

Expand the numerator:

\begin{equation}
\phi''
=
\frac{r^3u''-2r^2u'+2ru}{r^4}.
\tag{TRISO-ANA-261}
\label{eq:triso-ana-261}
\end{equation}

Divide each term by $r^4$:

\begin{equation}
\phi''
=
\frac{u''}{r}
-
\frac{2u'}{r^2}
+
\frac{2u}{r^3}.
\tag{TRISO-ANA-262}
\label{eq:triso-ana-262}
\end{equation}

Now substitute (TRISO-ANA-256), (TRISO-ANA-255), and (TRISO-ANA-262) into (TRISO-ANA-253):

\begin{equation}
\left(
\frac{u''}{r}
-
\frac{2u'}{r^2}
+
\frac{2u}{r^3}
\right)
+
\frac2r
\left(
\frac{ru'-u}{r^2}
\right)
+
k^2\frac{u}{r}
=
0.
\tag{TRISO-ANA-263}
\label{eq:triso-ana-263}
\end{equation}

Expand the second term:

\begin{equation}
\frac{u''}{r}
-
\frac{2u'}{r^2}
+
\frac{2u}{r^3}
+
\frac{2u'}{r^2}
-
\frac{2u}{r^3}
+
k^2\frac{u}{r}
=
0.
\tag{TRISO-ANA-264}
\label{eq:triso-ana-264}
\end{equation}

Cancel the (u') terms explicitly:

\begin{equation}
\frac{u''}{r}
+
k^2\frac{u}{r}
=
0.
\tag{TRISO-ANA-265}
\label{eq:triso-ana-265}
\end{equation}

Multiply by $r$:

\begin{equation}
\boxed{
u''+k^2u=0.
}
\tag{TRISO-ANA-266}
\label{eq:triso-ana-266}
\end{equation}

\subsubsection*{10.2 Solve the transformed radial equation}

The characteristic equation is

\begin{equation}
m^2+k^2=0.
\tag{TRISO-ANA-267}
\label{eq:triso-ana-267}
\end{equation}

Its roots are

\begin{equation}
m=\pm ik.
\tag{TRISO-ANA-268}
\label{eq:triso-ana-268}
\end{equation}

Therefore the real-valued solution is

\begin{equation}
\boxed{
u(r)=A\sin(kr)+B\cos(kr).
}
\tag{TRISO-ANA-269}
\label{eq:triso-ana-269}
\end{equation}

Substitute into $\phi=u/r$:

\begin{equation}
\phi(r)
=
\frac{A\sin(kr)+B\cos(kr)}{r}.
\tag{TRISO-ANA-270}
\label{eq:triso-ana-270}
\end{equation}

\subsubsection*{10.3 Centre regularity}

Examine the cosine contribution as (r	o0).

Use the known limit

\begin{equation}
\lim_{r\to0}\cos(kr)=1.
\tag{TRISO-ANA-271}
\label{eq:triso-ana-271}
\end{equation}

Therefore

\begin{equation}
\lim_{r\to0}\frac{B\cos(kr)}{r}
=
\lim_{r\to0}\frac{B}{r}.
\tag{TRISO-ANA-272}
\label{eq:triso-ana-272}
\end{equation}

For $B\ne0$, this diverges.

A physical concentration perturbation must remain finite at the particle centre.

Therefore

\begin{equation}
B=0.
\tag{TRISO-ANA-273}
\label{eq:triso-ana-273}
\end{equation}

The eigenfunction becomes

\begin{equation}
\phi(r)=A\frac{\sin(kr)}{r}.
\tag{TRISO-ANA-274}
\label{eq:triso-ana-274}
\end{equation}

For the sine term, use

\begin{equation}
\sin(kr)=kr+O(r^3)
\qquad
(r\to0).
\tag{TRISO-ANA-275}
\label{eq:triso-ana-275}
\end{equation}

Divide by $r$:

\begin{equation}
\frac{\sin(kr)}{r}
=
k+O(r^2).
\tag{TRISO-ANA-276}
\label{eq:triso-ana-276}
\end{equation}

Therefore

\begin{equation}
\boxed{
\lim_{r\to0}\phi(r)=Ak.
}
\tag{TRISO-ANA-277}
\label{eq:triso-ana-277}
\end{equation}

The apparent (1/r) singularity is removable for the sine branch.

It is often convenient to absorb $k$ into the modal amplitude. Define

\begin{equation}
C=A k.
\tag{TRISO-ANA-278}
\label{eq:triso-ana-278}
\end{equation}

Then the same mode may be written as

\begin{equation}
\phi(r)=C\frac{\sin(kr)}{kr}.
\tag{TRISO-ANA-279}
\label{eq:triso-ana-279}
\end{equation}

The normalization is arbitrary; only the relative spatial shape matters for the eigenvalue problem.

\subsubsection*{10.4 Derive the Robin eigencondition}

Start from the normalized form

\begin{equation}
\phi(r)=C\frac{\sin(kr)}{kr}.
\tag{TRISO-ANA-280}
\label{eq:triso-ana-280}
\end{equation}

The derivative is easier to obtain by treating (C/k) as a constant:

\begin{equation}
\phi(r)=\frac{C}{k}\frac{\sin(kr)}{r}.
\tag{TRISO-ANA-281}
\label{eq:triso-ana-281}
\end{equation}

Differentiate $\sin(kr)/r$ using the quotient rule:

\begin{equation}
\frac{d}{dr}
\left(
\frac{\sin(kr)}{r}
\right)
=
\frac{
kr\cos(kr)-\sin(kr)
}{
r^2
}.
\tag{TRISO-ANA-282}
\label{eq:triso-ana-282}
\end{equation}

Therefore

\begin{equation}
\boxed{
\phi'(r)
=
\frac{C}{k}
\frac{
kr\cos(kr)-\sin(kr)
}{
r^2
}.
}
\tag{TRISO-ANA-283}
\label{eq:triso-ana-283}
\end{equation}

At the outer boundary,

\begin{equation}
-D\phi'(R)=h\phi(R).
\tag{TRISO-ANA-284}
\label{eq:triso-ana-284}
\end{equation}

Substitute $phi'(R)$:

\begin{equation}
-D
\frac{C}{k}
\frac{
kR\cos(kR)-\sin(kR)
}{
R^2
}
=
h\phi(R).
\tag{TRISO-ANA-285}
\label{eq:triso-ana-285}
\end{equation}

Substitute $phi(R)=C\sin(kR$/(kR)):

\begin{equation}
-D
\frac{C}{k}
\frac{
kR\cos(kR)-\sin(kR)
}{
R^2
}
=
h
\frac{C}{kR}
\sin(kR).
\tag{TRISO-ANA-286}
\label{eq:triso-ana-286}
\end{equation}

Multiply both sides by $kR^2/C$, assuming $C\ne0$:

\begin{equation}
-D
\left[
kR\cos(kR)-\sin(kR)
\right]
=
hR\sin(kR).
\tag{TRISO-ANA-287}
\label{eq:triso-ana-287}
\end{equation}

Multiply by (-1):

\begin{equation}
D
\left[
\sin(kR)-kR\cos(kR)
\right]
=
hR\sin(kR).
\tag{TRISO-ANA-288}
\label{eq:triso-ana-288}
\end{equation}

Define the dimensionless eigenvariable

\begin{equation}
\mu=kR.
\tag{TRISO-ANA-289}
\label{eq:triso-ana-289}
\end{equation}

Its dimensions are

\begin{equation}
[\mu]=[k][R]=\mathrm{m^{-1}}\times\mathrm m=1.
\tag{TRISO-ANA-290}
\label{eq:triso-ana-290}
\end{equation}

Define the Biot number

\begin{equation}
\mathrm{Bi}=\frac{hR}{D}.
\tag{TRISO-ANA-291}
\label{eq:triso-ana-291}
\end{equation}

Its dimensions are

\begin{equation}
[\mathrm{Bi}]
=
\frac{
\mathrm{m\,s^{-1}}\times\mathrm m
}{
\mathrm{m^2\,s^{-1}}
}
=
1.
\tag{TRISO-ANA-292}
\label{eq:triso-ana-292}
\end{equation}

Substitute $\mu=kR$ and $hR/D=\mathrm{Bi}$ into (TRISO-ANA-288):

\begin{equation}
\boxed{
\sin\mu-\mu\cos\mu
=
\mathrm{Bi}\sin\mu.
}
\tag{TRISO-ANA-293}
\label{eq:triso-ana-293}
\end{equation}

Rearrange:

\begin{equation}
(1-\mathrm{Bi})\sin\mu
=
\mu\cos\mu.
\tag{TRISO-ANA-294}
\label{eq:triso-ana-294}
\end{equation}

For $\sin\mu\ne0$, divide by $\sin\mu$:

\begin{equation}
1-\mathrm{Bi}
=
\mu\frac{\cos\mu}{\sin\mu}.
\tag{TRISO-ANA-295}
\label{eq:triso-ana-295}
\end{equation}

Use $\cot\mu=\cos\mu/\sin\mu$:

\begin{equation}
\boxed{
\mu\cot\mu=1-\mathrm{Bi}.
}
\tag{TRISO-ANA-296}
\label{eq:triso-ana-296}
\end{equation}

\subsubsection*{10.5 Check whether division by (sinmu) loses roots}

The undivided equation (TRISO-ANA-293) must be used for this check.

Suppose

\begin{equation}
\sin\mu=0.
\tag{TRISO-ANA-297}
\label{eq:triso-ana-297}
\end{equation}

Then $mu=n\pi$ for integer $n$.

Substitute into (TRISO-ANA-293):

\begin{equation}
0-n\pi\cos(n\pi)=0.
\tag{TRISO-ANA-298}
\label{eq:triso-ana-298}
\end{equation}

Because

\begin{equation}
\cos(n\pi)=(-1)^n,
\tag{TRISO-ANA-299}
\label{eq:triso-ana-299}
\end{equation}

this becomes

\begin{equation}
-n\pi(-1)^n=0.
\tag{TRISO-ANA-300}
\label{eq:triso-ana-300}
\end{equation}

For positive $n$, this is not zero.

Therefore no positive eigenvalue is lost when dividing by (sinmu).

The only simultaneous zero is $mu=0$, which does not satisfy the positive transient-mode condition for the Robin problem with $h>0$.

\subsubsection*{10.6 Eigenvalue definitions}

Let $mu_n$ denote the positive roots of (TRISO-ANA-293).

Then

\begin{equation}
k_nR=\mu_n.
\tag{TRISO-ANA-301}
\label{eq:triso-ana-301}
\end{equation}

Therefore

\begin{equation}
\boxed{
k_n=\frac{\mu_n}{R}.
}
\tag{TRISO-ANA-302}
\label{eq:triso-ana-302}
\end{equation}

Since

\begin{equation}
\Lambda_n=Dk_n^2,
\tag{TRISO-ANA-303}
\label{eq:triso-ana-303}
\end{equation}

substitute (TRISO-ANA-302):

\begin{equation}
\Lambda_n
=
D
\left(
\frac{\mu_n}{R}
\right)^2.
\tag{TRISO-ANA-304}
\label{eq:triso-ana-304}
\end{equation}

Thus

\begin{equation}
\boxed{
\Lambda_n
=
D\frac{\mu_n^2}{R^2}.
}
\tag{TRISO-ANA-305}
\label{eq:triso-ana-305}
\end{equation}

The dimensions are

\begin{equation}
[\Lambda_n]
=
\mathrm{m^2\,s^{-1}}
\times
\mathrm{m^{-2}}
=
\mathrm{s^{-1}}.
\tag{TRISO-ANA-306}
\label{eq:triso-ana-306}
\end{equation}

The separated transient mode is therefore

\begin{equation}
v_n(r,t)=\phi_n(r)e^{-\Lambda_nt}.
\tag{TRISO-ANA-307}
\label{eq:triso-ana-307}
\end{equation}

\subsubsection*{10.7 Sturm–Liouville form}

Start from the radial eigenproblem:

\begin{equation}
\phi_n''
+
\frac2r\phi_n'
+
k_n^2\phi_n
=
0.
\tag{TRISO-SL-200}
\label{eq:triso-sl-200}
\end{equation}

Multiply by $r^2$:

\begin{equation}
r^2\phi_n''
+
2r\phi_n'
+
k_n^2r^2\phi_n
=
0.
\tag{TRISO-SL-201}
\label{eq:triso-sl-201}
\end{equation}

The first two terms are a product derivative because

\begin{equation}
\frac{d}{dr}(r^2\phi_n')
=
2r\phi_n'
+
r^2\phi_n''.
\tag{TRISO-SL-202}
\label{eq:triso-sl-202}
\end{equation}

Therefore

\begin{equation}
\frac{d}{dr}(r^2\phi_n')
+
k_n^2r^2\phi_n
=
0.
\tag{TRISO-SL-203}
\label{eq:triso-sl-203}
\end{equation}

Move the eigenvalue term to the other side:

\begin{equation}
-\frac{d}{dr}(r^2\phi_n')
=
k_n^2r^2\phi_n.
\tag{TRISO-SL-204}
\label{eq:triso-sl-204}
\end{equation}

This is the self-adjoint Sturm–Liouville form

\begin{equation}
-\frac{d}{dr}
\left(
p(r)\frac{d\phi_n}{dr}
\right)
+
q(r)\phi_n
=
\lambda_n w(r)\phi_n
\end{equation}

with the identifications

\begin{equation}
p(r)=r^2,
\tag{TRISO-SL-205}
\label{eq:triso-sl-205}
\end{equation}

\begin{equation}
q(r)=0,
\tag{TRISO-SL-206}
\label{eq:triso-sl-206}
\end{equation}

\begin{equation}
w(r)=r^2,
\tag{TRISO-SL-207}
\label{eq:triso-sl-207}
\end{equation}

and

\begin{equation}
\lambda_n=k_n^2.
\tag{TRISO-SL-208}
\label{eq:triso-sl-208}
\end{equation}

The eigenvalue in this Sturm–Liouville problem is therefore $k_n^2$, with units $\mathrm{m^{-2}}$, not the temporal decay rate $\Lambda_n$.

The interval is $0<r<R$.

The centre condition is regularity of $\phi_n$, equivalent for these modes to a finite $phi_n(0)$ and zero radial derivative at the centre.

The outer boundary is the homogeneous Robin condition

\begin{equation}
-D\phi_n'(R)=h\phi_n(R).
\tag{TRISO-SL-209}
\label{eq:triso-sl-209}
\end{equation}

[THEOREM / STANDARD FORM] Sturm–Liouville theory provides the framework for eigenvalues and eigenfunctions of self-adjoint second-order problems. See the NIST Digital Library of Mathematical Functions, §1.13(viii), which identifies Sturm–Liouville eigenvalues/eigenfunctions and the Liouville form. \cite{NIST_DLMF}.

Because the centre endpoint has $p(0)=0$, it is more precise to call this a radial \textbf{singular} Sturm–Liouville endpoint rather than an ordinary regular endpoint. The orthogonality used below can nevertheless be derived directly for the present eigenfunctions, so no stronger theorem is needed.

\subsubsection*{10.8 Derive orthogonality directly}

Take two distinct eigenfunctions $\phi_m$ and $\phi_n$ with eigenvalues $k_m^2$ and $k_n^2$:

\begin{equation}
-\frac{d}{dr}
(r^2\phi_m')
=
k_m^2r^2\phi_m,
\tag{TRISO-SL-210}
\label{eq:triso-sl-210}
\end{equation}

and

\begin{equation}
-\frac{d}{dr}
(r^2\phi_n')
=
k_n^2r^2\phi_n.
\tag{TRISO-SL-211}
\label{eq:triso-sl-211}
\end{equation}

Multiply the first equation by $\phi_n$:

\begin{equation}
-\phi_n\frac{d}{dr}(r^2\phi_m')
=
k_m^2r^2\phi_m\phi_n.
\tag{TRISO-SL-212}
\label{eq:triso-sl-212}
\end{equation}

Multiply the second equation by $\phi_m$:

\begin{equation}
-\phi_m\frac{d}{dr}(r^2\phi_n')
=
k_n^2r^2\phi_m\phi_n.
\tag{TRISO-SL-213}
\label{eq:triso-sl-213}
\end{equation}

Subtract the second equation from the first:

\begin{equation}
-\phi_n\frac{d}{dr}(r^2\phi_m')
+
\phi_m\frac{d}{dr}(r^2\phi_n')
=
(k_m^2-k_n^2)r^2\phi_m\phi_n.
\tag{TRISO-SL-214}
\label{eq:triso-sl-214}
\end{equation}

Recognise the left side as a derivative:

\begin{equation}
\frac{d}{dr}
\left[
r^2
(
\phi_m\phi_n'
-
\phi_n\phi_m'
)
\right]
=
(k_n^2-k_m^2)r^2\phi_m\phi_n.
\tag{TRISO-SL-215}
\label{eq:triso-sl-215}
\end{equation}

Integrate from (0) to $R$:

\begin{equation}
\int_0^R
\frac{d}{dr}
\left[
r^2
(
\phi_m\phi_n'
-
\phi_n\phi_m'
)
\right]dr
=
(k_n^2-k_m^2)
\int_0^R
r^2\phi_m\phi_n\,dr.
\tag{TRISO-SL-216}
\label{eq:triso-sl-216}
\end{equation}

Evaluate the left-hand integral:

\begin{equation}
\left[
r^2
(
\phi_m\phi_n'
-
\phi_n\phi_m'
)
\right]_0^R
=
(k_n^2-k_m^2)
\int_0^R
r^2\phi_m\phi_n\,dr.
\tag{TRISO-SL-217}
\label{eq:triso-sl-217}
\end{equation}

At $r=R$, both eigenfunctions satisfy the same Robin condition:

\begin{equation}
\phi_m'(R)=-\frac{h}{D}\phi_m(R),
\tag{TRISO-SL-218}
\label{eq:triso-sl-218}
\end{equation}

and

\begin{equation}
\phi_n'(R)=-\frac{h}{D}\phi_n(R).
\tag{TRISO-SL-219}
\label{eq:triso-sl-219}
\end{equation}

Therefore the outer boundary term is zero:

\begin{equation}
R^2
[
\phi_m(R)\phi_n'(R)
-
\phi_n(R)\phi_m'(R)
]
=0.
\tag{TRISO-SL-220}
\label{eq:triso-sl-220}
\end{equation}

At the centre, the regular eigenfunctions are finite and their derivatives remain bounded, while $r^2\to0$.

Hence

\begin{equation}
\lim_{r\to0}
r^2
[
\phi_m\phi_n'
-
\phi_n\phi_m'
]
=0.
\tag{TRISO-SL-221}
\label{eq:triso-sl-221}
\end{equation}

Therefore the complete boundary term is zero:

\begin{equation}
(k_n^2-k_m^2)
\int_0^R
r^2\phi_m\phi_n\,dr
=
0.
\tag{TRISO-SL-222}
\label{eq:triso-sl-222}
\end{equation}

For distinct eigenvalues,

\begin{equation}
k_n^2\ne k_m^2,
\tag{TRISO-SL-223}
\label{eq:triso-sl-223}
\end{equation}

so

\begin{equation}
\boxed{
\int_0^R
r^2\phi_m(r)\phi_n(r)\,dr
=
0,
\qquad m\ne n.
}
\tag{TRISO-SL-224}
\label{eq:triso-sl-224}
\end{equation}

The weight is therefore

\begin{equation}
\boxed{
w(r)=r^2.
}
\tag{TRISO-SL-225}
\label{eq:triso-sl-225}
\end{equation}

The same $r^2$ weight also follows directly from spherical volume $dV=4\pi r^2dr$.

\subsubsection*{10.9 Modal coefficient projection}

At $t=0$, (TRISO-ANA-228) gives

\begin{equation}
v(r,0)=-w(r).
\tag{TRISO-SL-226}
\label{eq:triso-sl-226}
\end{equation}

Represent the initial transient as an eigenfunction series:

\begin{equation}
-w(r)
=
\sum_{n=1}^{\infty}
A_n\phi_n(r).
\tag{TRISO-SL-227}
\label{eq:triso-sl-227}
\end{equation}

Multiply both sides by $r^2\phi_m(r)$:

\begin{equation}
-r^2w(r)\phi_m(r)
=
\sum_{n=1}^{\infty}
A_n r^2\phi_n(r)\phi_m(r).
\tag{TRISO-SL-228}
\label{eq:triso-sl-228}
\end{equation}

Integrate from (0) to $R$:

\begin{equation}
-\int_0^R
r^2w(r)\phi_m(r)\,dr
=
\sum_{n=1}^{\infty}
A_n
\int_0^R
r^2\phi_n(r)\phi_m(r)\,dr.
\tag{TRISO-SL-229}
\label{eq:triso-sl-229}
\end{equation}

For $n\ne m$, orthogonality makes the corresponding integrals zero:

\begin{equation}
\int_0^R
r^2\phi_n\phi_m\,dr
=
0.
\tag{TRISO-SL-230}
\label{eq:triso-sl-230}
\end{equation}

The remaining $n=m$ term is

\begin{equation}
-\int_0^R
r^2w(r)\phi_m(r)\,dr
=
A_m
\int_0^R
r^2\phi_m(r)^2\,dr.
\tag{TRISO-SL-231}
\label{eq:triso-sl-231}
\end{equation}

Divide by the non-zero mode norm:

\begin{equation}
\boxed{
A_m
=
-
\frac{
\int_0^R
r^2w(r)\phi_m(r)\,dr
}{
\int_0^R
r^2\phi_m(r)^2\,dr
}.
}
\tag{TRISO-SL-232}
\label{eq:triso-sl-232}
\end{equation}

Rename $m$ to $n$:

\begin{equation}
\boxed{
A_n
=
-
\frac{
\int_0^R
r^2w(r)\phi_n(r)\,dr
}{
\int_0^R
r^2\phi_n(r)^2\,dr
}.
}
\tag{TRISO-SL-233}
\label{eq:triso-sl-233}
\end{equation}

Each mode evolves with $e^{-\Lambda_nt}$, so

\begin{equation}
v(r,t)
=
\sum_{n=1}^{\infty}
A_n\phi_n(r)e^{-\Lambda_nt}.
\tag{TRISO-SL-234}
\label{eq:triso-sl-234}
\end{equation}

Since $c=v+w$,

\begin{equation}
\boxed{
c(r,t)
=
w(r)
+
\sum_{n=1}^{\infty}
A_n\phi_n(r)e^{-\Lambda_nt}.
}
\tag{TRISO-SL-235}
\label{eq:triso-sl-235}
\end{equation}

\subsubsection*{10.10 Checks on the transient series}

Each mode has decay factor

\begin{equation}
e^{-\Lambda_nt}.
\tag{TRISO-SL-236}
\label{eq:triso-sl-236}
\end{equation}

At $t=0$,

\begin{equation}
e^{-\Lambda_n\cdot0}=1.
\tag{TRISO-SL-237}
\label{eq:triso-sl-237}
\end{equation}

Therefore

\begin{equation}
c(r,0)
=
w(r)
+
\sum_{n=1}^{\infty}A_n\phi_n(r).
\tag{TRISO-SL-238}
\label{eq:triso-sl-238}
\end{equation}

Using the defining expansion (TRISO-SL-227),

\begin{equation}
\sum_{n=1}^{\infty}A_n\phi_n(r)
=
-w(r).
\tag{TRISO-SL-239}
\label{eq:triso-sl-239}
\end{equation}

Hence, formally,

\begin{equation}
\boxed{
c(r,0)=0.
}
\tag{TRISO-SL-240}
\label{eq:triso-sl-240}
\end{equation}

As $t\to\infty$, every mode with $\Lambda_n>0$ satisfies

\begin{equation}
e^{-\Lambda_nt}\to0.
\tag{TRISO-SL-241}
\label{eq:triso-sl-241}
\end{equation}

Therefore, provided the modal expansion has the required convergence,

\begin{equation}
\boxed{
c(r,t)\to w(r).
}
\tag{TRISO-SL-242}
\label{eq:triso-sl-242}
\end{equation}

The convergence of the infinite series itself has not been numerically established here. The statements above are the formal consequences of the eigen-expansion framework.

\subsection*{11. Five-layer steady analytical formulation}

\subsubsection*{11.1 Geometry, domains, and layer quantities}

Define the concentric radii

\begin{equation}
r_0=0<r_1<r_2<r_3<r_4<r_5=R.
\tag{TRISO-ML-300}
\label{eq:triso-ml-300}
\end{equation}

The five material regions are

\begin{equation}
\Omega_1=(r_0,r_1)
\quad\text{fuel kernel},
\tag{TRISO-ML-301}
\label{eq:triso-ml-301}
\end{equation}

\begin{equation}
\Omega_2=(r_1,r_2)
\quad\text{buffer},
\tag{TRISO-ML-302}
\label{eq:triso-ml-302}
\end{equation}

\begin{equation}
\Omega_3=(r_2,r_3)
\quad\text{IPyC},
\tag{TRISO-ML-303}
\label{eq:triso-ml-303}
\end{equation}

\begin{equation}
\Omega_4=(r_3,r_4)
\quad\text{SiC},
\tag{TRISO-ML-304}
\label{eq:triso-ml-304}
\end{equation}

and

\begin{equation}
\Omega_5=(r_4,r_5)
\quad\text{OPyC}.
\tag{TRISO-ML-305}
\label{eq:triso-ml-305}
\end{equation}

In material layer $i$, define

\begin{equation}
c_i(r,t),
\qquad
r_{i-1}<r<r_i.
\tag{TRISO-ML-306}
\label{eq:triso-ml-306}
\end{equation}

The concentration units are

\begin{equation}
[c_i]=\mathrm{mol\,m^{-3}}.
\tag{TRISO-ML-307}
\label{eq:triso-ml-307}
\end{equation}

Let $D_i$ be the diffusivity in layer $i$:

\begin{equation}
[D_i]=\mathrm{m^2\,s^{-1}}.
\tag{TRISO-ML-308}
\label{eq:triso-ml-308}
\end{equation}

Let $S_i$ be the net volumetric source in layer $i$:

\begin{equation}
[S_i]=\mathrm{mol\,m^{-3}\,s^{-1}}.
\tag{TRISO-ML-309}
\label{eq:triso-ml-309}
\end{equation}

The coordinates $r,t$, outer radius $R$, transfer coefficient $h$, and external concentration $c_\infty$ are global quantities. The fields $c_i$, diffusivities $D_i$, and sources $S_i$ are material-layer quantities.

For the source-driven five-layer benchmark,

\begin{equation}
S_1=S_0,
\tag{TRISO-ML-310}
\label{eq:triso-ml-310}
\end{equation}

while

\begin{equation}
S_i=0,
\qquad
i=2,3,4,5.
\tag{TRISO-ML-311}
\label{eq:triso-ml-311}
\end{equation}

\textbf{[ASSUMPTION]} Each $D_i>0$ is constant within its material layer for this analytical benchmark.

\textbf{[ASSUMPTION]} Interfaces have zero storage, zero interfacial source, $K_i=1$, and no explicit interfacial resistance.

\subsubsection*{11.2 Kernel steady solution}

The steady conservative equation in the kernel is

\begin{equation}
0=
\frac1{r^2}
\frac{d}{dr}
\left(
r^2D_1\frac{dc_1}{dr}
\right)
+S_0.
\tag{TRISO-ML-312}
\label{eq:triso-ml-312}
\end{equation}

Move the source term to the right:

\begin{equation}
\frac1{r^2}
\frac{d}{dr}
\left(
r^2D_1\frac{dc_1}{dr}
\right)
=-S_0.
\tag{TRISO-ML-313}
\label{eq:triso-ml-313}
\end{equation}

Multiply by $r^2$:

\begin{equation}
\frac{d}{dr}
\left(
r^2D_1\frac{dc_1}{dr}
\right)
=-S_0r^2.
\tag{TRISO-ML-314}
\label{eq:triso-ml-314}
\end{equation}

Because $D_1$ is constant in the kernel,

\begin{equation}
D_1
\frac{d}{dr}
\left(
r^2\frac{dc_1}{dr}
\right)
=-S_0r^2.
\tag{TRISO-ML-315}
\label{eq:triso-ml-315}
\end{equation}

Divide by $D_1$:

\begin{equation}
\frac{d}{dr}
\left(
r^2\frac{dc_1}{dr}
\right)
=-\frac{S_0}{D_1}r^2.
\tag{TRISO-ML-316}
\label{eq:triso-ml-316}
\end{equation}

Integrate:

\begin{equation}
r^2\frac{dc_1}{dr}
=
-\frac{S_0r^3}{3D_1}
+C_1.
\tag{TRISO-ML-317}
\label{eq:triso-ml-317}
\end{equation}

Divide by $r^2$, for $r>0$:

\begin{equation}
\frac{dc_1}{dr}
=
-\frac{S_0r}{3D_1}
+\frac{C_1}{r^2}.
\tag{TRISO-ML-318}
\label{eq:triso-ml-318}
\end{equation}

Centre regularity requires $dc_1/dr$ to remain finite as $r\to0$. The term $C_1/r^2$ diverges unless

\begin{equation}
C_1=0.
\tag{TRISO-ML-319}
\label{eq:triso-ml-319}
\end{equation}

Therefore

\begin{equation}
\frac{dc_1}{dr}
=
-\frac{S_0r}{3D_1}.
\tag{TRISO-ML-320}
\label{eq:triso-ml-320}
\end{equation}

Integrate again:

\begin{equation}
c_1(r)
=
-\frac{S_0r^2}{6D_1}
+A_1.
\tag{TRISO-ML-321}
\label{eq:triso-ml-321}
\end{equation}

Hence the regular kernel profile is

\begin{equation}
\boxed{
c_1(r)
=
A_1-\frac{S_0r^2}{6D_1}.
}
\tag{TRISO-ML-322}
\label{eq:triso-ml-322}
\end{equation}

\subsubsection*{11.3 Source-free coating solutions}

For $i=2,3,4,5$,

\begin{equation}
S_i=0.
\tag{TRISO-ML-323}
\label{eq:triso-ml-323}
\end{equation}

The steady equation is

\begin{equation}
0=
\frac1{r^2}
\frac{d}{dr}
\left(
r^2D_i\frac{dc_i}{dr}
\right).
\tag{TRISO-ML-324}
\label{eq:triso-ml-324}
\end{equation}

Multiply by $r^2$:

\begin{equation}
\frac{d}{dr}
\left(
r^2D_i\frac{dc_i}{dr}
\right)=0.
\tag{TRISO-ML-325}
\label{eq:triso-ml-325}
\end{equation}

Use constant $D_i$:

\begin{equation}
D_i
\frac{d}{dr}
\left(
r^2\frac{dc_i}{dr}
\right)=0.
\tag{TRISO-ML-326}
\label{eq:triso-ml-326}
\end{equation}

Divide by $D_i>0$:

\begin{equation}
\frac{d}{dr}
\left(
r^2\frac{dc_i}{dr}
\right)=0.
\tag{TRISO-ML-327}
\label{eq:triso-ml-327}
\end{equation}

Integrate:

\begin{equation}
r^2\frac{dc_i}{dr}=C_i.
\tag{TRISO-ML-328}
\label{eq:triso-ml-328}
\end{equation}

Divide by $r^2$:

\begin{equation}
\frac{dc_i}{dr}=\frac{C_i}{r^2}.
\tag{TRISO-ML-329}
\label{eq:triso-ml-329}
\end{equation}

Integrate:

\begin{equation}
c_i(r)=C_i\int r^{-2}dr+A_i.
\tag{TRISO-ML-330}
\label{eq:triso-ml-330}
\end{equation}

Since

\begin{equation}
\int r^{-2}dr=-\frac1r,
\tag{TRISO-ML-331}
\label{eq:triso-ml-331}
\end{equation}

we obtain

\begin{equation}
c_i(r)=A_i-\frac{C_i}{r}.
\tag{TRISO-ML-332}
\label{eq:triso-ml-332}
\end{equation}

Define $B_i=-C_i$. Then

\begin{equation}
\boxed{
c_i(r)=A_i+\frac{B_i}{r},
\qquad i=2,3,4,5.
}
\tag{TRISO-ML-333}
\label{eq:triso-ml-333}
\end{equation}

The coating layers do not include $r=0$, so their $1/r$ terms are finite within their own domains and are not removed by centre regularity.

\subsubsection*{11.4 Total kernel generation and common steady flux}

The total kernel generation rate is

\begin{equation}
\dot N_{\mathrm{gen}}
=
4\pi
\int_0^{r_1}
S_0r^2\,dr.
\tag{TRISO-ML-334}
\label{eq:triso-ml-334}
\end{equation}

Pull the constants outside:

\begin{equation}
\dot N_{\mathrm{gen}}
=
4\pi S_0
\int_0^{r_1}r^2\,dr.
\tag{TRISO-ML-335}
\label{eq:triso-ml-335}
\end{equation}

Evaluate the integral:

\begin{equation}
\int_0^{r_1}r^2\,dr
=
\left[\frac{r^3}{3}\right]_0^{r_1}.
\tag{TRISO-ML-336}
\label{eq:triso-ml-336}
\end{equation}

Therefore

\begin{equation}
\int_0^{r_1}r^2\,dr
=
\frac{r_1^3}{3}.
\tag{TRISO-ML-337}
\label{eq:triso-ml-337}
\end{equation}

Hence

\begin{equation}
\boxed{
\dot N_{\mathrm{gen}}
=
\frac{4\pi S_0r_1^3}{3}.
}
\tag{TRISO-ML-338}
\label{eq:triso-ml-338}
\end{equation}

At steady state, with no coating source, reaction, or storage, the same total amount rate crosses every sphere outside the kernel:

\begin{equation}
4\pi r^2J_r(r)
=
\dot N_{\mathrm{gen}},
\qquad r>r_1.
\tag{TRISO-ML-339}
\label{eq:triso-ml-339}
\end{equation}

Substitute (TRISO-ML-338):

\begin{equation}
4\pi r^2J_r(r)
=
\frac{4\pi S_0r_1^3}{3}.
\tag{TRISO-ML-340}
\label{eq:triso-ml-340}
\end{equation}

Cancel $4\pi$:

\begin{equation}
r^2J_r(r)=\frac{S_0r_1^3}{3}.
\tag{TRISO-ML-341}
\label{eq:triso-ml-341}
\end{equation}

Divide by $r^2$:

\begin{equation}
\boxed{
J_r(r)=\frac{S_0r_1^3}{3r^2}.
}
\tag{TRISO-ML-342}
\label{eq:triso-ml-342}
\end{equation}

The positive sign is outward.

\subsubsection*{11.5 Recover shell gradients from Fick's law}

In shell $i$,

\begin{equation}
J_r=-D_i\frac{dc_i}{dr}.
\tag{TRISO-ML-343}
\label{eq:triso-ml-343}
\end{equation}

Substitute (TRISO-ML-342):

\begin{equation}
-D_i\frac{dc_i}{dr}
=
\frac{S_0r_1^3}{3r^2}.
\tag{TRISO-ML-344}
\label{eq:triso-ml-344}
\end{equation}

Divide by $-D_i$:

\begin{equation}
\boxed{
\frac{dc_i}{dr}
=
-\frac{S_0r_1^3}{3D_ir^2}.
}
\tag{TRISO-ML-345}
\label{eq:triso-ml-345}
\end{equation}

Integrate from $r$ to the outer radius $r_i$ of that shell:

\begin{equation}
\int_{c_i(r)}^{c_i(r_i)}dc_i
=
-\frac{S_0r_1^3}{3D_i}
\int_r^{r_i}\rho^{-2}d\rho.
\tag{TRISO-ML-346}
\label{eq:triso-ml-346}
\end{equation}

The radial integral is

\begin{equation}
\int_r^{r_i}\rho^{-2}d\rho
=
\frac1r-\frac1{r_i}.
\tag{TRISO-ML-347}
\label{eq:triso-ml-347}
\end{equation}

Therefore

\begin{equation}
c_i(r_i)-c_i(r)
=
-\frac{S_0r_1^3}{3D_i}
\left(
\frac1r-\frac1{r_i}
\right).
\tag{TRISO-ML-348}
\label{eq:triso-ml-348}
\end{equation}

Multiply by $-1$:

\begin{equation}
\boxed{
c_i(r)-c_i(r_i)
=
\frac{S_0r_1^3}{3D_i}
\left(
\frac1r-\frac1{r_i}
\right).
}
\tag{TRISO-ML-349}
\label{eq:triso-ml-349}
\end{equation}

Differentiating $A_i+B_i/r$ gives

\begin{equation}
\frac{dc_i}{dr}=-\frac{B_i}{r^2}.
\tag{TRISO-ML-350}
\label{eq:triso-ml-350}
\end{equation}

Compare with (TRISO-ML-345):

\begin{equation}
-\frac{B_i}{r^2}
=
-\frac{S_0r_1^3}{3D_ir^2}.
\tag{TRISO-ML-351}
\label{eq:triso-ml-351}
\end{equation}

Hence

\begin{equation}
\boxed{
B_i=\frac{S_0r_1^3}{3D_i}.
}
\tag{TRISO-ML-352}
\label{eq:triso-ml-352}
\end{equation}

This independently agrees with the direct shell ODE solution.

\subsubsection*{11.6 Interface matching}

At $r=r_1$, flux continuity is

\begin{equation}
-D_1c_1'(r_1)
=
-D_2c_2'(r_1),
\tag{TRISO-ML-353}
\label{eq:triso-ml-353}
\end{equation}

and ideal concentration continuity is

\begin{equation}
c_1(r_1)=c_2(r_1).
\tag{TRISO-ML-354}
\label{eq:triso-ml-354}
\end{equation}

At $r=r_2$,

\begin{equation}
-D_2c_2'(r_2)
=
-D_3c_3'(r_2),
\tag{TRISO-ML-355}
\label{eq:triso-ml-355}
\end{equation}

and

\begin{equation}
c_2(r_2)=c_3(r_2).
\tag{TRISO-ML-356}
\label{eq:triso-ml-356}
\end{equation}

At $r=r_3$,

\begin{equation}
-D_3c_3'(r_3)
=
-D_4c_4'(r_3),
\tag{TRISO-ML-357}
\label{eq:triso-ml-357}
\end{equation}

and

\begin{equation}
c_3(r_3)=c_4(r_3).
\tag{TRISO-ML-358}
\label{eq:triso-ml-358}
\end{equation}

At $r=r_4$,

\begin{equation}
-D_4c_4'(r_4)
=
-D_5c_5'(r_4),
\tag{TRISO-ML-359}
\label{eq:triso-ml-359}
\end{equation}

and

\begin{equation}
c_4(r_4)=c_5(r_4).
\tag{TRISO-ML-360}
\label{eq:triso-ml-360}
\end{equation}

The flux equations are already satisfied by the common steady amount rate. The concentration equations relate the additive constants.

Across shell $i$,

\begin{equation}
c_i(r_{i-1})-c_i(r_i)
=
\frac{S_0r_1^3}{3D_i}
\left(
\frac1{r_{i-1}}-\frac1{r_i}
\right),
\qquad i=2,3,4,5.
\tag{TRISO-ML-361}
\label{eq:triso-ml-361}
\end{equation}

Thus each inner interface concentration is obtained from the next outer interface concentration by adding that shell's concentration drop.

\subsubsection*{11.7 Outer Robin condition and inward propagation}

At $R=r_5$,

\begin{equation}
-D_5c_5'(R)
=
h[c_5(R)-c_\infty].
\tag{TRISO-ML-362}
\label{eq:triso-ml-362}
\end{equation}

For the benchmark,

\begin{equation}
c_\infty=0.
\tag{TRISO-ML-363}
\label{eq:triso-ml-363}
\end{equation}

Therefore

\begin{equation}
-D_5c_5'(R)=hc_5(R).
\tag{TRISO-ML-364}
\label{eq:triso-ml-364}
\end{equation}

The common flux gives

\begin{equation}
-D_5c_5'(R)
=
\frac{S_0r_1^3}{3R^2}.
\tag{TRISO-ML-365}
\label{eq:triso-ml-365}
\end{equation}

Equate the two expressions:

\begin{equation}
hc_5(R)
=
\frac{S_0r_1^3}{3R^2}.
\tag{TRISO-ML-366}
\label{eq:triso-ml-366}
\end{equation}

Divide by $h$:

\begin{equation}
\boxed{
c_5(R)
=
\frac{S_0r_1^3}{3hR^2}.
}
\tag{TRISO-ML-367}
\label{eq:triso-ml-367}
\end{equation}

Move inward through OPyC:

\begin{equation}
c_5(r_4)
=
c_5(R)
+
\frac{S_0r_1^3}{3D_5}
\left(
\frac1{r_4}-\frac1R
\right).
\tag{TRISO-ML-368}
\label{eq:triso-ml-368}
\end{equation}

Apply continuity:

\begin{equation}
c_4(r_4)=c_5(r_4).
\tag{TRISO-ML-369}
\label{eq:triso-ml-369}
\end{equation}

Move inward through SiC:

\begin{equation}
c_4(r_3)
=
c_4(r_4)
+
\frac{S_0r_1^3}{3D_4}
\left(
\frac1{r_3}-\frac1{r_4}
\right).
\tag{TRISO-ML-370}
\label{eq:triso-ml-370}
\end{equation}

Apply continuity:

\begin{equation}
c_3(r_3)=c_4(r_3).
\tag{TRISO-ML-371}
\label{eq:triso-ml-371}
\end{equation}

Move inward through IPyC:

\begin{equation}
c_3(r_2)
=
c_3(r_3)
+
\frac{S_0r_1^3}{3D_3}
\left(
\frac1{r_2}-\frac1{r_3}
\right).
\tag{TRISO-ML-372}
\label{eq:triso-ml-372}
\end{equation}

Apply continuity:

\begin{equation}
c_2(r_2)=c_3(r_2).
\tag{TRISO-ML-373}
\label{eq:triso-ml-373}
\end{equation}

Move inward through the buffer:

\begin{equation}
c_2(r_1)
=
c_2(r_2)
+
\frac{S_0r_1^3}{3D_2}
\left(
\frac1{r_1}-\frac1{r_2}
\right).
\tag{TRISO-ML-374}
\label{eq:triso-ml-374}
\end{equation}

Apply continuity:

\begin{equation}
c_1(r_1)=c_2(r_1).
\tag{TRISO-ML-375}
\label{eq:triso-ml-375}
\end{equation}

From the kernel profile,

\begin{equation}
c_1(r)-c_1(r_1)
=
\frac{S_0}{6D_1}(r_1^2-r^2).
\tag{TRISO-ML-376}
\label{eq:triso-ml-376}
\end{equation}

Hence

\begin{equation}
\boxed{
c_1(r)
=
c_2(r_1)
+
\frac{S_0}{6D_1}(r_1^2-r^2).
}
\tag{TRISO-ML-377}
\label{eq:triso-ml-377}
\end{equation}

Equations (TRISO-ML-367) through (TRISO-ML-377) give the complete steady five-layer solution recursively.

\subsubsection*{11.8 Spherical resistance formulation}

Define the common steady amount rate outside the kernel:

\begin{equation}
\dot N=4\pi r^2J_r.
\tag{TRISO-ML-378}
\label{eq:triso-ml-378}
\end{equation}

In shell $i$,

\begin{equation}
\dot N
=
-4\pi r^2D_i\frac{dc_i}{dr}.
\tag{TRISO-ML-379}
\label{eq:triso-ml-379}
\end{equation}

Rearrange:

\begin{equation}
dc_i
=
-\frac{\dot N}{4\pi D_i}\frac{dr}{r^2}.
\tag{TRISO-ML-380}
\label{eq:triso-ml-380}
\end{equation}

Integrate from $r_{i-1}$ to $r_i$:

\begin{equation}
c_i(r_i)-c_i(r_{i-1})
=
-\frac{\dot N}{4\pi D_i}
\int_{r_{i-1}}^{r_i}r^{-2}dr.
\tag{TRISO-ML-381}
\label{eq:triso-ml-381}
\end{equation}

Evaluate the radial integral:

\begin{equation}
\int_{r_{i-1}}^{r_i}r^{-2}dr
=
\frac1{r_{i-1}}-\frac1{r_i}.
\tag{TRISO-ML-382}
\label{eq:triso-ml-382}
\end{equation}

Therefore

\begin{equation}
c_i(r_{i-1})-c_i(r_i)
=
\dot N
\frac1{4\pi D_i}
\left(
\frac1{r_{i-1}}-\frac1{r_i}
\right).
\tag{TRISO-ML-383}
\label{eq:triso-ml-383}
\end{equation}

Define

\begin{equation}
\boxed{
\mathcal R_i
=
\frac1{4\pi D_i}
\left(
\frac1{r_{i-1}}-\frac1{r_i}
\right),
\qquad i=2,3,4,5.
}
\tag{TRISO-ML-384}
\label{eq:triso-ml-384}
\end{equation}

Then

\begin{equation}
c_i(r_{i-1})-c_i(r_i)=\dot N\mathcal R_i.
\tag{TRISO-ML-385}
\label{eq:triso-ml-385}
\end{equation}

Its units are

\begin{equation}
[\mathcal R_i]
=
\mathrm{s\,m^{-3}}.
\tag{TRISO-ML-386}
\label{eq:triso-ml-386}
\end{equation}

For external transfer,

\begin{equation}
\dot N
=
4\pi R^2h[c_5(R)-c_\infty].
\tag{TRISO-ML-387}
\label{eq:triso-ml-387}
\end{equation}

Rearrange:

\begin{equation}
c_5(R)-c_\infty
=
\dot N
\frac1{4\pi R^2h}.
\tag{TRISO-ML-388}
\label{eq:triso-ml-388}
\end{equation}

Define

\begin{equation}
\boxed{
\mathcal R_h
=
\frac1{4\pi R^2h}.
}
\tag{TRISO-ML-389}
\label{eq:triso-ml-389}
\end{equation}

Its units are also

\begin{equation}
[\mathcal R_h]=\mathrm{s\,m^{-3}}.
\tag{TRISO-ML-390}
\label{eq:triso-ml-390}
\end{equation}

Because the same $\dot N$ passes through every coating and the external film, the concentration drops add:

\begin{equation}
c_2(r_1)-c_\infty
=
\dot N
\left(
\mathcal R_2+\mathcal R_3+\mathcal R_4+\mathcal R_5+\mathcal R_h
\right).
\tag{TRISO-ML-391}
\label{eq:triso-ml-391}
\end{equation}

Thus

\begin{equation}
\boxed{
c_2(r_1)
=
c_\infty
+
\dot N
\left(
\mathcal R_2+\mathcal R_3+\mathcal R_4+\mathcal R_5+\mathcal R_h
\right).
}
\tag{TRISO-ML-392}
\label{eq:triso-ml-392}
\end{equation}

For the benchmark $c_\infty=0$ and

\begin{equation}
\dot N=\frac{4\pi S_0r_1^3}{3}.
\tag{TRISO-ML-393}
\label{eq:triso-ml-393}
\end{equation}

Substituting (TRISO-ML-384), (TRISO-ML-389), and (TRISO-ML-393) into (TRISO-ML-392) reproduces exactly the inward-recursion concentration at $r_1$.

The kernel itself is source-containing, so it is not represented by the same source-free shell resistance. Its centre-to-interface concentration rise is instead

\begin{equation}
c_1(0)-c_1(r_1)
=
\frac{S_0r_1^2}{6D_1}.
\tag{TRISO-ML-394}
\label{eq:triso-ml-394}
\end{equation}

This distinction prevents a source-containing kernel from being incorrectly treated as an ordinary source-free series resistance.

\subsection*{12. Five-layer transient analytical formulation}

\subsubsection*{12.1 Define the steady reference and transient deviation}

Let

\begin{equation}
c_{i,\mathrm{ss}}(r)
\tag{TRISO-ML-400}
\label{eq:triso-ml-400}
\end{equation}

denote the steady five-layer solution derived in Section 11.

Define the transient deviation in each layer:

\begin{equation}
\boxed{
v_i(r,t)
=
c_i(r,t)-c_{i,\mathrm{ss}}(r).
}
\tag{TRISO-ML-401}
\label{eq:triso-ml-401}
\end{equation}

Rearrange:

\begin{equation}
c_i(r,t)
=
v_i(r,t)+c_{i,\mathrm{ss}}(r).
\tag{TRISO-ML-402}
\label{eq:triso-ml-402}
\end{equation}

Because the steady reference is time independent,

\begin{equation}
\frac{\partial c_{i,\mathrm{ss}}}{\partial t}=0.
\tag{TRISO-ML-403}
\label{eq:triso-ml-403}
\end{equation}

Therefore

\begin{equation}
\frac{\partial c_i}{\partial t}
=
\frac{\partial v_i}{\partial t}.
\tag{TRISO-ML-404}
\label{eq:triso-ml-404}
\end{equation}

Similarly,

\begin{equation}
\frac{\partial c_i}{\partial r}
=
\frac{\partial v_i}{\partial r}
+
\frac{dc_{i,\mathrm{ss}}}{dr},
\tag{TRISO-ML-405}
\label{eq:triso-ml-405}
\end{equation}

and

\begin{equation}
\frac{\partial^2c_i}{\partial r^2}
=
\frac{\partial^2v_i}{\partial r^2}
+
\frac{d^2c_{i,\mathrm{ss}}}{dr^2}.
\tag{TRISO-ML-406}
\label{eq:triso-ml-406}
\end{equation}

The full layer equation is

\begin{equation}
\frac{\partial c_i}{\partial t}
=
D_i
\left(
\frac{\partial^2c_i}{\partial r^2}
+
\frac2r\frac{\partial c_i}{\partial r}
\right)
+
S_i.
\tag{TRISO-ML-407}
\label{eq:triso-ml-407}
\end{equation}

Substitute (TRISO-ML-404) through (TRISO-ML-406):

\begin{equation}
\frac{\partial v_i}{\partial t}
=
D_i
\left(
v_{i,rr}
+
c_{i,\mathrm{ss}}''
+
\frac2r v_{i,r}
+
\frac2r c_{i,\mathrm{ss}}'
\right)
+
S_i.
\tag{TRISO-ML-408}
\label{eq:triso-ml-408}
\end{equation}

Group transient and steady terms:

\begin{equation}
\frac{\partial v_i}{\partial t}
=
D_i
\left(
v_{i,rr}+\frac2r v_{i,r}
\right)
+
\left[
D_i
\left(
c_{i,\mathrm{ss}}''
+\frac2r c_{i,\mathrm{ss}}'
\right)
+
S_i
\right].
\tag{TRISO-ML-409}
\label{eq:triso-ml-409}
\end{equation}

The steady solution satisfies

\begin{equation}
D_i
\left(
c_{i,\mathrm{ss}}''
+\frac2r c_{i,\mathrm{ss}}'
\right)
+
S_i
=
0.
\tag{TRISO-ML-410}
\label{eq:triso-ml-410}
\end{equation}

Therefore

\begin{equation}
\boxed{
\frac{\partial v_i}{\partial t}
=
D_i
\left(
\frac{\partial^2v_i}{\partial r^2}
+
\frac2r\frac{\partial v_i}{\partial r}
\right).
}
\tag{TRISO-ML-411}
\label{eq:triso-ml-411}
\end{equation}

\subsubsection*{12.2 Homogeneous transient conditions}

At the centre, both $c_1$ and $c_{1,\mathrm{ss}}$ satisfy zero radial derivative. Therefore

\begin{equation}
\boxed{
v_{1,r}(0,t)=0.
}
\tag{TRISO-ML-412}
\label{eq:triso-ml-412}
\end{equation}

At interface $r=r_i$, both the full and steady solutions satisfy concentration continuity. Subtracting the steady relation from the full relation gives

\begin{equation}
\boxed{
v_i(r_i,t)=v_{i+1}(r_i,t).
}
\tag{TRISO-ML-413}
\label{eq:triso-ml-413}
\end{equation}

Both the full and steady solutions also satisfy flux continuity. Subtraction gives

\begin{equation}
\boxed{
-D_iv_i'(r_i,t)
=
-D_{i+1}v_{i+1}'(r_i,t).
}
\tag{TRISO-ML-414}
\label{eq:triso-ml-414}
\end{equation}

At $R$, the full Robin condition is

\begin{equation}
-D_5c_5'(R,t)=h[c_5(R,t)-c_\infty].
\tag{TRISO-ML-415}
\label{eq:triso-ml-415}
\end{equation}

The steady solution satisfies

\begin{equation}
-D_5c_{5,\mathrm{ss}}'(R)
=
h[c_{5,\mathrm{ss}}(R)-c_\infty].
\tag{TRISO-ML-416}
\label{eq:triso-ml-416}
\end{equation}

Subtract (TRISO-ML-416) from (TRISO-ML-415):

\begin{equation}
\boxed{
-D_5v_5'(R,t)=hv_5(R,t).
}
\tag{TRISO-ML-417}
\label{eq:triso-ml-417}
\end{equation}

For the initially empty source-driven problem,

\begin{equation}
c_i(r,0)=0.
\tag{TRISO-ML-418}
\label{eq:triso-ml-418}
\end{equation}

Hence

\begin{equation}
\boxed{
v_i(r,0)
=
-c_{i,\mathrm{ss}}(r).
}
\tag{TRISO-ML-419}
\label{eq:triso-ml-419}
\end{equation}

\subsubsection*{12.3 Separation in each layer}

For one global transient mode, assume

\begin{equation}
v_i(r,t)
=
\phi_i(r)e^{-\Lambda t}.
\tag{TRISO-ML-420}
\label{eq:triso-ml-420}
\end{equation}

The same temporal factor must apply in every layer because the interface conditions couple the layer amplitudes at the same physical time. A single global eigenmode cannot use independent exponential time factors on the two sides of one interface and still satisfy the interface equations for all $t$, except in a degenerate zero-amplitude case.

Differentiate with respect to time:

\begin{equation}
\frac{\partial v_i}{\partial t}
=
-\Lambda\phi_i e^{-\Lambda t}.
\tag{TRISO-ML-421}
\label{eq:triso-ml-421}
\end{equation}

Differentiate with respect to radius:

\begin{equation}
\frac{\partial v_i}{\partial r}
=
\phi_i'e^{-\Lambda t}.
\tag{TRISO-ML-422}
\label{eq:triso-ml-422}
\end{equation}

Differentiate again:

\begin{equation}
\frac{\partial^2v_i}{\partial r^2}
=
\phi_i''e^{-\Lambda t}.
\tag{TRISO-ML-423}
\label{eq:triso-ml-423}
\end{equation}

Substitute into (TRISO-ML-411):

\begin{equation}
-\Lambda\phi_i e^{-\Lambda t}
=
D_i
\left(
\phi_i''
+\frac2r\phi_i'
\right)
e^{-\Lambda t}.
\tag{TRISO-ML-424}
\label{eq:triso-ml-424}
\end{equation}

Cancel the non-zero exponential factor:

\begin{equation}
-\Lambda\phi_i
=
D_i
\left(
\phi_i''
+\frac2r\phi_i'
\right).
\tag{TRISO-ML-425}
\label{eq:triso-ml-425}
\end{equation}

Divide by $D_i$:

\begin{equation}
\phi_i''
+\frac2r\phi_i'
+
\frac{\Lambda}{D_i}\phi_i
=
0.
\tag{TRISO-ML-426}
\label{eq:triso-ml-426}
\end{equation}

Define

\begin{equation}
\boxed{
k_i^2=\frac{\Lambda}{D_i}.
}
\tag{TRISO-ML-427}
\label{eq:triso-ml-427}
\end{equation}

Then

\begin{equation}
[k_i^2]
=
\frac{\mathrm{s^{-1}}}{\mathrm{m^2\,s^{-1}}}
=
\mathrm{m^{-2}},
\tag{TRISO-ML-428}
\label{eq:triso-ml-428}
\end{equation}

so

\begin{equation}
[k_i]=\mathrm{m^{-1}}.
\tag{TRISO-ML-429}
\label{eq:triso-ml-429}
\end{equation}

Each layer generally has a different $k_i$, because each layer has a different $D_i$, even though all layers in one global mode share the same $\Lambda$.

The radial equation is

\begin{equation}
\boxed{
\phi_i''
+\frac2r\phi_i'
+k_i^2\phi_i
=
0.
}
\tag{TRISO-ML-430}
\label{eq:triso-ml-430}
\end{equation}

\subsubsection*{12.4 Apply the proven transformation $u_i=r\phi_i$}

Section 10 proved that the transformation

\begin{equation}
u_i=r\phi_i
\tag{TRISO-ML-431}
\label{eq:triso-ml-431}
\end{equation}

maps (TRISO-ML-430) to

\begin{equation}
\boxed{
u_i''+k_i^2u_i=0.
}
\tag{TRISO-ML-432}
\label{eq:triso-ml-432}
\end{equation}

Therefore

\begin{equation}
\boxed{
u_i(r)
=
A_i\sin(k_ir)+B_i\cos(k_ir).
}
\tag{TRISO-ML-433}
\label{eq:triso-ml-433}
\end{equation}

In the kernel,

\begin{equation}
\phi_1(r)=\frac{u_1(r)}{r}.
\tag{TRISO-ML-434}
\label{eq:triso-ml-434}
\end{equation}

The cosine contribution $B_1\cos(k_1r)/r$ diverges as $r\to0$, exactly as proved in Section 10.

Therefore

\begin{equation}
\boxed{
B_1=0.
}
\tag{TRISO-ML-435}
\label{eq:triso-ml-435}
\end{equation}

Thus

\begin{equation}
u_1(r)=A_1\sin(k_1r).
\tag{TRISO-ML-436}
\label{eq:triso-ml-436}
\end{equation}

No coating layer contains the origin, so $B_i$ is not forced to zero for $i=2,3,4,5$.

\subsubsection*{12.5 Transform concentration continuity}

At interface $r=r_i$,

\begin{equation}
\phi_i(r_i)=\phi_{i+1}(r_i).
\tag{TRISO-ML-437}
\label{eq:triso-ml-437}
\end{equation}

Use $\phi_i=u_i/r$:

\begin{equation}
\frac{u_i(r_i)}{r_i}
=
\frac{u_{i+1}(r_i)}{r_i}.
\tag{TRISO-ML-438}
\label{eq:triso-ml-438}
\end{equation}

Multiply by the common non-zero radius $r_i$:

\begin{equation}
\boxed{
u_i(r_i)=u_{i+1}(r_i).
}
\tag{TRISO-ML-439}
\label{eq:triso-ml-439}
\end{equation}

\subsubsection*{12.6 Transform flux continuity}

Start from

\begin{equation}
\phi_i(r)=\frac{u_i(r)}{r}.
\tag{TRISO-ML-440}
\label{eq:triso-ml-440}
\end{equation}

Differentiate using the quotient rule:

\begin{equation}
\phi_i'(r)
=
\frac{ru_i'(r)-u_i(r)}{r^2}.
\tag{TRISO-ML-441}
\label{eq:triso-ml-441}
\end{equation}

Separate the two terms:

\begin{equation}
\boxed{
\phi_i'(r)
=
\frac{u_i'(r)}{r}
-
\frac{u_i(r)}{r^2}.
}
\tag{TRISO-ML-442}
\label{eq:triso-ml-442}
\end{equation}

Flux continuity at $r=r_i$ is

\begin{equation}
-D_i\phi_i'(r_i)
=
-D_{i+1}\phi_{i+1}'(r_i).
\tag{TRISO-ML-443}
\label{eq:triso-ml-443}
\end{equation}

Cancel the common minus sign:

\begin{equation}
D_i\phi_i'(r_i)
=
D_{i+1}\phi_{i+1}'(r_i).
\tag{TRISO-ML-444}
\label{eq:triso-ml-444}
\end{equation}

Substitute (TRISO-ML-442) on both sides:

\begin{equation}
D_i
\left[
\frac{u_i'(r_i)}{r_i}
-
\frac{u_i(r_i)}{r_i^2}
\right]
=
D_{i+1}
\left[
\frac{u_{i+1}'(r_i)}{r_i}
-
\frac{u_{i+1}(r_i)}{r_i^2}
\right].
\tag{TRISO-ML-445}
\label{eq:triso-ml-445}
\end{equation}

Multiply by the common factor $r_i$:

\begin{equation}
\boxed{
D_i
\left[
u_i'(r_i)-\frac{u_i(r_i)}{r_i}
\right]
=
D_{i+1}
\left[
u_{i+1}'(r_i)-\frac{u_{i+1}(r_i)}{r_i}
\right].
}
\tag{TRISO-ML-446}
\label{eq:triso-ml-446}
\end{equation}

This is the transformed ideal flux-continuity condition.

\subsubsection*{12.7 Transform the outer Robin condition}

The transient Robin condition is

\begin{equation}
-D_5\phi_5'(R)=h\phi_5(R).
\tag{TRISO-ML-447}
\label{eq:triso-ml-447}
\end{equation}

Use

\begin{equation}
\phi_5'(R)
=
\frac{u_5'(R)}{R}
-
\frac{u_5(R)}{R^2},
\tag{TRISO-ML-448}
\label{eq:triso-ml-448}
\end{equation}

and

\begin{equation}
\phi_5(R)=\frac{u_5(R)}{R}.
\tag{TRISO-ML-449}
\label{eq:triso-ml-449}
\end{equation}

Substitute both:

\begin{equation}
-D_5
\left[
\frac{u_5'(R)}{R}
-
\frac{u_5(R)}{R^2}
\right]
=
h\frac{u_5(R)}{R}.
\tag{TRISO-ML-450}
\label{eq:triso-ml-450}
\end{equation}

Multiply by $R$:

\begin{equation}
\boxed{
-D_5
\left[
u_5'(R)-\frac{u_5(R)}{R}
\right]
=
hu_5(R).
}
\tag{TRISO-ML-451}
\label{eq:triso-ml-451}
\end{equation}

Equivalently,

\begin{equation}
D_5u_5'(R)
+
\left(
h-\frac{D_5}{R}
\right)u_5(R)
=
0.
\tag{TRISO-ML-452}
\label{eq:triso-ml-452}
\end{equation}

\subsubsection*{12.8 Count the unknown coefficients}

Before centre regularity, five layers would provide ten coefficients:

\begin{equation}
(A_1,B_1,A_2,B_2,A_3,B_3,A_4,B_4,A_5,B_5).
\tag{TRISO-ML-453}
\label{eq:triso-ml-453}
\end{equation}

Centre regularity fixes

\begin{equation}
B_1=0.
\tag{TRISO-ML-454}
\label{eq:triso-ml-454}
\end{equation}

Therefore nine independent coefficients remain.

Define

\begin{equation}
\boxed{
\mathbf a
=
(A_1,A_2,B_2,A_3,B_3,A_4,B_4,A_5,B_5)^T.
}
\tag{TRISO-ML-455}
\label{eq:triso-ml-455}
\end{equation}

Thus

\begin{equation}
\mathbf a\in\mathbb R^9
\tag{TRISO-ML-456}
\label{eq:triso-ml-456}
\end{equation}

for real $\Lambda>0$.

\subsubsection*{12.9 Define interface shorthand}

For compact matrix notation, define at interface $r=r_j$

\begin{equation}
s_{ij}=\sin(k_ir_j),
\qquad
c_{ij}=\cos(k_ir_j).
\tag{TRISO-ML-457}
\label{eq:triso-ml-457}
\end{equation}

For a sine basis term,

\begin{equation}
u_i=A_i\sin(k_ir),
\tag{TRISO-ML-458}
\label{eq:triso-ml-458}
\end{equation}

and

\begin{equation}
u_i'=A_ik_i\cos(k_ir).
\tag{TRISO-ML-459}
\label{eq:triso-ml-459}
\end{equation}

Therefore the transformed flux factor for the sine basis at $r_j$ is

\begin{equation}
F^{(s)}_{ij}
=
D_i
\left(
k_ic_{ij}-\frac{s_{ij}}{r_j}
\right).
\tag{TRISO-ML-460}
\label{eq:triso-ml-460}
\end{equation}

For a cosine basis term,

\begin{equation}
u_i=B_i\cos(k_ir),
\tag{TRISO-ML-461}
\label{eq:triso-ml-461}
\end{equation}

and

\begin{equation}
u_i'=-B_ik_i\sin(k_ir).
\tag{TRISO-ML-462}
\label{eq:triso-ml-462}
\end{equation}

Therefore its transformed flux factor is

\begin{equation}
F^{(c)}_{ij}
=
D_i
\left(
-k_is_{ij}-\frac{c_{ij}}{r_j}
\right).
\tag{TRISO-ML-463}
\label{eq:triso-ml-463}
\end{equation}

These quantities depend on $\Lambda$ through $k_i=\sqrt{\Lambda/D_i}$.

\subsubsection*{12.10 Assemble the global homogeneous coefficient system}

There are two equations at each of the four internal interfaces:

\noindent\textbullet\ one concentration-continuity equation;\\
\noindent\textbullet\ one flux-continuity equation.\\

This gives

\begin{equation}
4\times2=8
\tag{TRISO-ML-464}
\label{eq:triso-ml-464}
\end{equation}

interface equations.

The outer Robin boundary supplies one more equation:

\begin{equation}
8+1=9.
\tag{TRISO-ML-465}
\label{eq:triso-ml-465}
\end{equation}

These nine equations determine the nine coefficients up to an arbitrary overall modal normalization when $\Lambda$ is an eigenvalue.

Write

\begin{equation}
\boxed{
\mathbf M(\Lambda)\mathbf a=\mathbf0.
}
\tag{TRISO-ML-466}
\label{eq:triso-ml-466}
\end{equation}

The matrix dimensions are

\begin{equation}
\boxed{
\mathbf M(\Lambda)\in\mathbb R^{9\times9}.
}
\tag{TRISO-ML-467}
\label{eq:triso-ml-467}
\end{equation}

Using the coefficient order in (TRISO-ML-455), rows 1–2 correspond to $r_1$, rows 3–4 to $r_2$, rows 5–6 to $r_3$, rows 7–8 to $r_4$, and row 9 to the outer Robin condition.

The explicit matrix is

\begin{equation}
\mathbf M(\Lambda)=
\begin{pmatrix}
s_{11} & -s_{21} & -c_{21} & 0 & 0 & 0 & 0 & 0 & 0\\
F^{(s)}_{11} & -F^{(s)}_{21} & -F^{(c)}_{21} & 0 & 0 & 0 & 0 & 0 & 0\\
0 & s_{22} & c_{22} & -s_{32} & -c_{32} & 0 & 0 & 0 & 0\\
0 & F^{(s)}_{22} & F^{(c)}_{22} & -F^{(s)}_{32} & -F^{(c)}_{32} & 0 & 0 & 0 & 0\\
0 & 0 & 0 & s_{33} & c_{33} & -s_{43} & -c_{43} & 0 & 0\\
0 & 0 & 0 & F^{(s)}_{33} & F^{(c)}_{33} & -F^{(s)}_{43} & -F^{(c)}_{43} & 0 & 0\\
0 & 0 & 0 & 0 & 0 & s_{44} & c_{44} & -s_{54} & -c_{54}\\
0 & 0 & 0 & 0 & 0 & F^{(s)}_{44} & F^{(c)}_{44} & -F^{(s)}_{54} & -F^{(c)}_{54}\\
0 & 0 & 0 & 0 & 0 & 0 & 0 & G_s & G_c
\end{pmatrix}.
\tag{TRISO-ML-468}
\label{eq:triso-ml-468}
\end{equation}

The outer-row coefficients follow from (TRISO-ML-452).

For the sine basis,

\begin{equation}
G_s
=
D_5k_5\cos(k_5R)
+
\left(
h-\frac{D_5}{R}
\right)
\sin(k_5R).
\tag{TRISO-ML-469}
\label{eq:triso-ml-469}
\end{equation}

For the cosine basis,

\begin{equation}
G_c
=
-D_5k_5\sin(k_5R)
+
\left(
h-\frac{D_5}{R}
\right)
\cos(k_5R).
\tag{TRISO-ML-470}
\label{eq:triso-ml-470}
\end{equation}

Centre regularity does not appear as a matrix row because it has already been used to eliminate $B_1$ from the unknown vector.

\subsubsection*{12.11 Global eigenvalue condition}

For a generic value of $\Lambda$, the homogeneous system

\begin{equation}
\mathbf M(\Lambda)\mathbf a=\mathbf0
\tag{TRISO-ML-471}
\label{eq:triso-ml-471}
\end{equation}

has only the trivial solution

\begin{equation}
\mathbf a=\mathbf0
\tag{TRISO-ML-472}
\label{eq:triso-ml-472}
\end{equation}

when $\mathbf M$ is nonsingular.

A non-zero global eigenmode requires a non-trivial coefficient vector:

\begin{equation}
\mathbf a\ne\mathbf0.
\tag{TRISO-ML-473}
\label{eq:triso-ml-473}
\end{equation}

A square homogeneous linear system has a non-trivial solution only if its matrix is singular.

Therefore

\begin{equation}
\boxed{
\det\mathbf M(\Lambda)=0.
}
\tag{TRISO-ML-474}
\label{eq:triso-ml-474}
\end{equation}

Define

\begin{equation}
\boxed{
F(\Lambda)=\det\mathbf M(\Lambda).
}
\tag{TRISO-ML-475}
\label{eq:triso-ml-475}
\end{equation}

The global modal decay rates are the positive roots

\begin{equation}
F(\Lambda_n)=0.
\tag{TRISO-ML-476}
\label{eq:triso-ml-476}
\end{equation}

For each root $\Lambda_n$, the layer wave numbers are

\begin{equation}
\boxed{
k_{i,n}
=
\sqrt{\frac{\Lambda_n}{D_i}}.
}
\tag{TRISO-ML-477}
\label{eq:triso-ml-477}
\end{equation}

Thus one global decay rate $\Lambda_n$ generates five material-dependent spatial wave numbers.

\subsubsection*{12.12 Global conservative self-adjoint structure}

Return to the separated transient equation before the substitution $u_i=r\phi_i$.

Equation (TRISO-ML-425) is

\begin{equation}
-\Lambda\phi_i
=
D_i
\left(
\phi_i''
+
\frac2r\phi_i'
\right).
\tag{TRISO-ML-478}
\label{eq:triso-ml-478}
\end{equation}

Multiply by $r^2$:

\begin{equation}
-\Lambda r^2\phi_i
=
D_i
\left(
r^2\phi_i''
+
2r\phi_i'
\right).
\tag{TRISO-ML-479}
\label{eq:triso-ml-479}
\end{equation}

Because $D_i$ is constant inside layer $i$,

\begin{equation}
D_i
\left(
r^2\phi_i''
+
2r\phi_i'
\right)
=
\frac{d}{dr}
\left(
r^2D_i\phi_i'
\right).
\tag{TRISO-ML-480}
\label{eq:triso-ml-480}
\end{equation}

Therefore

\begin{equation}
-\Lambda r^2\phi_i
=
\frac{d}{dr}
\left(
r^2D_i\phi_i'
\right).
\tag{TRISO-ML-481}
\label{eq:triso-ml-481}
\end{equation}

Multiply by $-1$:

\begin{equation}
\boxed{
-\frac{d}{dr}
\left(
r^2D_i\phi_i'
\right)
=
\Lambda r^2\phi_i.
}
\tag{TRISO-ML-482}
\label{eq:triso-ml-482}
\end{equation}

This is the conservative eigen-equation in layer $i$.

Define the piecewise diffusivity

\begin{equation}
D(r)=D_i,
\qquad
r_{i-1}<r<r_i.
\tag{TRISO-ML-483}
\label{eq:triso-ml-483}
\end{equation}

Then the piecewise Sturm–Liouville coefficient is

\begin{equation}
\boxed{
p(r)=r^2D(r).
}
\tag{TRISO-ML-484}
\label{eq:triso-ml-484}
\end{equation}

There is no zeroth-order potential term:

\begin{equation}
\boxed{
q(r)=0.
}
\tag{TRISO-ML-485}
\label{eq:triso-ml-485}
\end{equation}

Comparing

\begin{equation}
-\frac{d}{dr}
\left(
p(r)\phi'
\right)
+
q(r)\phi
=
\Lambda w(r)\phi
\tag{TRISO-ML-486}
\label{eq:triso-ml-486}
\end{equation}

with (TRISO-ML-482) shows that the weight is

\begin{equation}
\boxed{
w(r)=r^2.
}
\tag{TRISO-ML-487}
\label{eq:triso-ml-487}
\end{equation}

Thus the weight $r^2$ follows from the physical conservative eigen-equation; it is not imported by analogy with the homogeneous sphere.

\subsubsection*{12.13 Layerwise Lagrange identity for two global modes}

Let global mode $m$ have eigenvalue $\Lambda_m$ and layer functions $\phi_i^{(m)}$.

In layer $i$,

\begin{equation}
-\frac{d}{dr}
\left(
r^2D_i\frac{d\phi_i^{(m)}}{dr}
\right)
=
\Lambda_m r^2\phi_i^{(m)}.
\tag{TRISO-ML-488}
\label{eq:triso-ml-488}
\end{equation}

Let global mode $n$ have eigenvalue $\Lambda_n$:

\begin{equation}
-\frac{d}{dr}
\left(
r^2D_i\frac{d\phi_i^{(n)}}{dr}
\right)
=
\Lambda_n r^2\phi_i^{(n)}.
\tag{TRISO-ML-489}
\label{eq:triso-ml-489}
\end{equation}

Multiply the $m$-equation by $\phi_i^{(n)}$:

\begin{equation}
-\phi_i^{(n)}
\frac{d}{dr}
\left(
r^2D_i\phi_i^{(m)\prime}
\right)
=
\Lambda_m r^2
\phi_i^{(m)}
\phi_i^{(n)}.
\tag{TRISO-ML-490}
\label{eq:triso-ml-490}
\end{equation}

Multiply the $n$-equation by $\phi_i^{(m)}$:

\begin{equation}
-\phi_i^{(m)}
\frac{d}{dr}
\left(
r^2D_i\phi_i^{(n)\prime}
\right)
=
\Lambda_n r^2
\phi_i^{(m)}
\phi_i^{(n)}.
\tag{TRISO-ML-491}
\label{eq:triso-ml-491}
\end{equation}

Subtract (TRISO-ML-491) from (TRISO-ML-490):

\begin{equation}
-\phi_i^{(n)}
\frac{d}{dr}
\left(
r^2D_i\phi_i^{(m)\prime}
\right)
+
\phi_i^{(m)}
\frac{d}{dr}
\left(
r^2D_i\phi_i^{(n)\prime}
\right)
=
(\Lambda_m-\Lambda_n)
r^2\phi_i^{(m)}\phi_i^{(n)}.
\tag{TRISO-ML-492}
\label{eq:triso-ml-492}
\end{equation}

Define

\begin{equation}
P_i(r)=r^2D_i.
\tag{TRISO-ML-493}
\label{eq:triso-ml-493}
\end{equation}

Use the product rule:

\begin{equation}
\frac{d}{dr}
\left[
P_i
\left(
\phi_i^{(m)}\phi_i^{(n)\prime}
-
\phi_i^{(n)}\phi_i^{(m)\prime}
\right)
\right]
\end{equation}

\begin{equation}
=
\phi_i^{(m)}
\frac{d}{dr}
\left(
P_i\phi_i^{(n)\prime}
\right)
-
\phi_i^{(n)}
\frac{d}{dr}
\left(
P_i\phi_i^{(m)\prime}
\right).
\tag{TRISO-ML-494}
\label{eq:triso-ml-494}
\end{equation}

Therefore (TRISO-ML-492) becomes

\begin{equation}
\frac{d}{dr}
\left[
r^2D_i
\left(
\phi_i^{(m)}\phi_i^{(n)\prime}
-
\phi_i^{(n)}\phi_i^{(m)\prime}
\right)
\right]
=
(\Lambda_m-\Lambda_n)
r^2\phi_i^{(m)}\phi_i^{(n)}.
\tag{TRISO-ML-495}
\label{eq:triso-ml-495}
\end{equation}

Integrate over layer $i$:

\begin{equation}
\int_{r_{i-1}}^{r_i}
\frac{d}{dr}
\left[
r^2D_i
\left(
\phi_i^{(m)}\phi_i^{(n)\prime}
-
\phi_i^{(n)}\phi_i^{(m)\prime}
\right)
\right]dr
\end{equation}

\begin{equation}
=
(\Lambda_m-\Lambda_n)
\int_{r_{i-1}}^{r_i}
r^2\phi_i^{(m)}\phi_i^{(n)}\,dr.
\tag{TRISO-ML-496}
\label{eq:triso-ml-496}
\end{equation}

Evaluate the derivative integral:

\begin{equation}
\left[
r^2D_i
\left(
\phi_i^{(m)}\phi_i^{(n)\prime}
-
\phi_i^{(n)}\phi_i^{(m)\prime}
\right)
\right]_{r_{i-1}}^{r_i}
\end{equation}

\begin{equation}
=
(\Lambda_m-\Lambda_n)
\int_{r_{i-1}}^{r_i}
r^2\phi_i^{(m)}\phi_i^{(n)}\,dr.
\tag{TRISO-ML-497}
\label{eq:triso-ml-497}
\end{equation}

\subsubsection*{12.14 Sum over all five layers}

Define the layer boundary expression

\begin{equation}
\mathcal B_i(r)
=
r^2D_i
\left(
\phi_i^{(m)}\phi_i^{(n)\prime}
-
\phi_i^{(n)}\phi_i^{(m)\prime}
\right).
\tag{TRISO-ML-498}
\label{eq:triso-ml-498}
\end{equation}

Equation (TRISO-ML-497) is

\begin{equation}
\mathcal B_i(r_i)
-
\mathcal B_i(r_{i-1})
=
(\Lambda_m-\Lambda_n)
\int_{r_{i-1}}^{r_i}
r^2\phi_i^{(m)}\phi_i^{(n)}\,dr.
\tag{TRISO-ML-499}
\label{eq:triso-ml-499}
\end{equation}

Sum from $i=1$ to $5$:

\begin{equation}
\sum_{i=1}^{5}
\left[
\mathcal B_i(r_i)
-
\mathcal B_i(r_{i-1})
\right]
\end{equation}

\begin{equation}
=
(\Lambda_m-\Lambda_n)
\sum_{i=1}^{5}
\int_{r_{i-1}}^{r_i}
r^2\phi_i^{(m)}\phi_i^{(n)}\,dr.
\tag{TRISO-ML-500}
\label{eq:triso-ml-500}
\end{equation}

Write the left side explicitly:

\begin{equation}
\mathcal B_1(r_1)-\mathcal B_1(0)
+
\mathcal B_2(r_2)-\mathcal B_2(r_1)
\end{equation}

\begin{equation}
+
\mathcal B_3(r_3)-\mathcal B_3(r_2)
+
\mathcal B_4(r_4)-\mathcal B_4(r_3)
+
\mathcal B_5(R)-\mathcal B_5(r_4).
\tag{TRISO-ML-501}
\label{eq:triso-ml-501}
\end{equation}

Group the four internal-interface contributions:

\begin{equation}
-\mathcal B_1(0)
+
\left[
\mathcal B_1(r_1)-\mathcal B_2(r_1)
\right]
+
\left[
\mathcal B_2(r_2)-\mathcal B_3(r_2)
\right]
\end{equation}

\begin{equation}
+
\left[
\mathcal B_3(r_3)-\mathcal B_4(r_3)
\right]
+
\left[
\mathcal B_4(r_4)-\mathcal B_5(r_4)
\right]
+
\mathcal B_5(R).
\tag{TRISO-ML-502}
\label{eq:triso-ml-502}
\end{equation}

\subsubsection*{12.15 Explicit cancellation at one internal interface}

Consider interface $r=r_j$ between layers $j$ and $j+1$.

The contribution from the left layer is

\begin{equation}
\mathcal B_j(r_j)
=
r_j^2D_j
\left[
\phi_j^{(m)}\phi_j^{(n)\prime}
-
\phi_j^{(n)}\phi_j^{(m)\prime}
\right]_{r_j}.
\tag{TRISO-ML-503}
\label{eq:triso-ml-503}
\end{equation}

The contribution from the right layer enters with a minus sign:

\begin{equation}
-\mathcal B_{j+1}(r_j)
=
-r_j^2D_{j+1}
\left[
\phi_{j+1}^{(m)}\phi_{j+1}^{(n)\prime}
-
\phi_{j+1}^{(n)}\phi_{j+1}^{(m)\prime}
\right]_{r_j}.
\tag{TRISO-ML-504}
\label{eq:triso-ml-504}
\end{equation}

For ideal concentration continuity,

\begin{equation}
\phi_j^{(m)}(r_j)
=
\phi_{j+1}^{(m)}(r_j),
\tag{TRISO-ML-505}
\label{eq:triso-ml-505}
\end{equation}

and

\begin{equation}
\phi_j^{(n)}(r_j)
=
\phi_{j+1}^{(n)}(r_j).
\tag{TRISO-ML-506}
\label{eq:triso-ml-506}
\end{equation}

For ideal flux continuity,

\begin{equation}
D_j\phi_j^{(m)\prime}(r_j)
=
D_{j+1}\phi_{j+1}^{(m)\prime}(r_j),
\tag{TRISO-ML-507}
\label{eq:triso-ml-507}
\end{equation}

and

\begin{equation}
D_j\phi_j^{(n)\prime}(r_j)
=
D_{j+1}\phi_{j+1}^{(n)\prime}(r_j).
\tag{TRISO-ML-508}
\label{eq:triso-ml-508}
\end{equation}

Use (TRISO-ML-505) and (TRISO-ML-508) in the first product of (TRISO-ML-503):

\begin{equation}
D_j
\phi_j^{(m)}
\phi_j^{(n)\prime}
=
\phi_{j+1}^{(m)}
D_{j+1}\phi_{j+1}^{(n)\prime}.
\tag{TRISO-ML-509}
\label{eq:triso-ml-509}
\end{equation}

Use (TRISO-ML-506) and (TRISO-ML-507) in the second product:

\begin{equation}
D_j
\phi_j^{(n)}
\phi_j^{(m)\prime}
=
\phi_{j+1}^{(n)}
D_{j+1}\phi_{j+1}^{(m)\prime}.
\tag{TRISO-ML-510}
\label{eq:triso-ml-510}
\end{equation}

Therefore

\begin{equation}
\mathcal B_j(r_j)
=
r_j^2D_{j+1}
\left[
\phi_{j+1}^{(m)}
\phi_{j+1}^{(n)\prime}
-
\phi_{j+1}^{(n)}
\phi_{j+1}^{(m)\prime}
\right].
\tag{TRISO-ML-511}
\label{eq:triso-ml-511}
\end{equation}

The right side of (TRISO-ML-511) is exactly

\begin{equation}
\mathcal B_{j+1}(r_j).
\tag{TRISO-ML-512}
\label{eq:triso-ml-512}
\end{equation}

Hence

\begin{equation}
\boxed{
\mathcal B_j(r_j)-\mathcal B_{j+1}(r_j)=0.
}
\tag{TRISO-ML-513}
\label{eq:triso-ml-513}
\end{equation}

The same argument applies independently at $r_1,r_2,r_3,r_4$, including when adjacent diffusivities are unequal.

Thus all four internal-interface terms in (TRISO-ML-502) cancel pairwise.

\subsubsection*{12.16 Centre boundary contribution}

At the centre,

\begin{equation}
\mathcal B_1(0)
=
\lim_{r\to0}
r^2D_1
\left(
\phi_1^{(m)}\phi_1^{(n)\prime}
-
\phi_1^{(n)}\phi_1^{(m)\prime}
\right).
\tag{TRISO-ML-514}
\label{eq:triso-ml-514}
\end{equation}

Regularity gives finite centre values for both eigenfunctions:

\begin{equation}
|\phi_1^{(m)}(0)|<\infty,
\qquad
|\phi_1^{(n)}(0)|<\infty.
\tag{TRISO-ML-515}
\label{eq:triso-ml-515}
\end{equation}

Spherical symmetry gives

\begin{equation}
\phi_1^{(m)\prime}(0)=0,
\qquad
\phi_1^{(n)\prime}(0)=0.
\tag{TRISO-ML-516}
\label{eq:triso-ml-516}
\end{equation}

For regular eigenfunctions, the bracketed quantity remains bounded as $r\to0$.

Since

\begin{equation}
r^2D_1\to0
\qquad
(r\to0),
\tag{TRISO-ML-517}
\label{eq:triso-ml-517}
\end{equation}

we obtain

\begin{equation}
\boxed{
\mathcal B_1(0)=0.
}
\tag{TRISO-ML-518}
\label{eq:triso-ml-518}
\end{equation}

\subsubsection*{12.17 Outer Robin boundary contribution}

At $r=R$,

\begin{equation}
\mathcal B_5(R)
=
R^2D_5
\left[
\phi_5^{(m)}(R)\phi_5^{(n)\prime}(R)
-
\phi_5^{(n)}(R)\phi_5^{(m)\prime}(R)
\right].
\tag{TRISO-ML-519}
\label{eq:triso-ml-519}
\end{equation}

Both modes satisfy the same homogeneous Robin condition:

\begin{equation}
-D_5\phi_5^{(m)\prime}(R)
=
h\phi_5^{(m)}(R),
\tag{TRISO-ML-520}
\label{eq:triso-ml-520}
\end{equation}

and

\begin{equation}
-D_5\phi_5^{(n)\prime}(R)
=
h\phi_5^{(n)}(R).
\tag{TRISO-ML-521}
\label{eq:triso-ml-521}
\end{equation}

Solve the first condition for the derivative:

\begin{equation}
D_5\phi_5^{(m)\prime}(R)
=
-h\phi_5^{(m)}(R).
\tag{TRISO-ML-522}
\label{eq:triso-ml-522}
\end{equation}

Similarly,

\begin{equation}
D_5\phi_5^{(n)\prime}(R)
=
-h\phi_5^{(n)}(R).
\tag{TRISO-ML-523}
\label{eq:triso-ml-523}
\end{equation}

Substitute into (TRISO-ML-519):

\begin{equation}
\mathcal B_5(R)
=
R^2
\left[
-h\phi_5^{(m)}(R)\phi_5^{(n)}(R)
+
h\phi_5^{(n)}(R)\phi_5^{(m)}(R)
\right].
\tag{TRISO-ML-524}
\label{eq:triso-ml-524}
\end{equation}

The two products are identical and have opposite signs:

\begin{equation}
\boxed{
\mathcal B_5(R)=0.
}
\tag{TRISO-ML-525}
\label{eq:triso-ml-525}
\end{equation}

\subsubsection*{12.18 Global multilayer orthogonality}

All terms on the left side of (TRISO-ML-500) now vanish:

\noindent\textbullet\ centre term by (TRISO-ML-518);\\
\noindent\textbullet\ four interface pairs by (TRISO-ML-513);\\
\noindent\textbullet\ outer Robin term by (TRISO-ML-525).\\

Therefore

\begin{equation}
0
=
(\Lambda_m-\Lambda_n)
\sum_{i=1}^{5}
\int_{r_{i-1}}^{r_i}
r^2
\phi_i^{(m)}(r)
\phi_i^{(n)}(r)
\,dr.
\tag{TRISO-ML-526}
\label{eq:triso-ml-526}
\end{equation}

For distinct eigenvalues,

\begin{equation}
\Lambda_m\ne\Lambda_n.
\tag{TRISO-ML-527}
\label{eq:triso-ml-527}
\end{equation}

Divide by $\Lambda_m-\Lambda_n$:

\begin{equation}
\boxed{
\sum_{i=1}^{5}
\int_{r_{i-1}}^{r_i}
r^2
\phi_i^{(m)}(r)
\phi_i^{(n)}(r)
\,dr
=
0,
\qquad
m\ne n.
}
\tag{TRISO-ML-528}
\label{eq:triso-ml-528}
\end{equation}

Thus the correct global weight is

\begin{equation}
\boxed{
w(r)=r^2.
}
\tag{TRISO-ML-529}
\label{eq:triso-ml-529}
\end{equation}

Define the global piecewise eigenfunction

\begin{equation}
\Phi_n(r)
=
\phi_i^{(n)}(r),
\qquad
r_{i-1}<r<r_i.
\tag{TRISO-ML-530}
\label{eq:triso-ml-530}
\end{equation}

Then define the weighted inner product

\begin{equation}
\boxed{
\langle f,g\rangle_w
=
\sum_{i=1}^{5}
\int_{r_{i-1}}^{r_i}
r^2f_i(r)g_i(r)\,dr.
}
\tag{TRISO-ML-531}
\label{eq:triso-ml-531}
\end{equation}

Global orthogonality is

\begin{equation}
\boxed{
\langle\Phi_m,\Phi_n\rangle_w=0,
\qquad
m\ne n.
}
\tag{TRISO-ML-532}
\label{eq:triso-ml-532}
\end{equation}

\subsubsection*{12.19 Modal norm and normalization}

Define the norm of mode $n$:

\begin{equation}
\boxed{
N_n
=
\langle\Phi_n,\Phi_n\rangle_w
=
\sum_{i=1}^{5}
\int_{r_{i-1}}^{r_i}
r^2
\left[
\phi_i^{(n)}(r)
\right]^2
dr.
}
\tag{TRISO-ML-533}
\label{eq:triso-ml-533}
\end{equation}

If $\phi_i$ carries concentration units, then

\begin{equation}
[N_n]
=
\mathrm{m^3}
[\phi]^2.
\tag{TRISO-ML-534}
\label{eq:triso-ml-534}
\end{equation}

If instead the eigenfunctions are chosen dimensionless, then

\begin{equation}
[N_n]=\mathrm{m^3}.
\tag{TRISO-ML-535}
\label{eq:triso-ml-535}
\end{equation}

The eigenvalue problem determines each coefficient vector only up to an arbitrary non-zero scale.

If

\begin{equation}
\Phi_n\to\alpha_n\Phi_n,
\tag{TRISO-ML-536}
\label{eq:triso-ml-536}
\end{equation}

then

\begin{equation}
N_n\to\alpha_n^2N_n.
\tag{TRISO-ML-537}
\label{eq:triso-ml-537}
\end{equation}

The corresponding modal amplitude transforms inversely:

\begin{equation}
A_n\to\frac{A_n}{\alpha_n}.
\tag{TRISO-ML-538}
\label{eq:triso-ml-538}
\end{equation}

Therefore the physical product

\begin{equation}
A_n\Phi_n
\tag{TRISO-ML-539}
\label{eq:triso-ml-539}
\end{equation}

is unchanged.

One convenient convention is unit weighted norm:

\begin{equation}
N_n=1.
\tag{TRISO-ML-540}
\label{eq:triso-ml-540}
\end{equation}

No such normalization is required for the projection formula below.

\subsubsection*{12.20 Project the initial transient state}

For the source-driven benchmark,

\begin{equation}
v_i(r,0)
=
-c_{i,\mathrm{ss}}(r).
\tag{TRISO-ML-541}
\label{eq:triso-ml-541}
\end{equation}

Assume the global modal representation

\begin{equation}
v_i(r,t)
=
\sum_{n=1}^{\infty}
A_n
\phi_i^{(n)}(r)
e^{-\Lambda_nt}.
\tag{TRISO-ML-542}
\label{eq:triso-ml-542}
\end{equation}

At $t=0$,

\begin{equation}
e^{-\Lambda_n0}=1.
\tag{TRISO-ML-543}
\label{eq:triso-ml-543}
\end{equation}

Therefore

\begin{equation}
-c_{i,\mathrm{ss}}(r)
=
\sum_{n=1}^{\infty}
A_n\phi_i^{(n)}(r),
\qquad
r_{i-1}<r<r_i.
\tag{TRISO-ML-544}
\label{eq:triso-ml-544}
\end{equation}

Multiply the equation in layer $i$ by

\begin{equation}
r^2\phi_i^{(m)}(r).
\tag{TRISO-ML-545}
\label{eq:triso-ml-545}
\end{equation}

This gives

\begin{equation}
-r^2
c_{i,\mathrm{ss}}(r)
\phi_i^{(m)}(r)
=
\sum_{n=1}^{\infty}
A_n
r^2
\phi_i^{(n)}(r)
\phi_i^{(m)}(r).
\tag{TRISO-ML-546}
\label{eq:triso-ml-546}
\end{equation}

Integrate over layer $i$:

\begin{equation}
-\int_{r_{i-1}}^{r_i}
r^2
c_{i,\mathrm{ss}}
\phi_i^{(m)}
\,dr
=
\sum_{n=1}^{\infty}
A_n
\int_{r_{i-1}}^{r_i}
r^2
\phi_i^{(n)}
\phi_i^{(m)}
\,dr.
\tag{TRISO-ML-547}
\label{eq:triso-ml-547}
\end{equation}

Sum all five layers:

\begin{equation}
-\sum_{i=1}^{5}
\int_{r_{i-1}}^{r_i}
r^2
c_{i,\mathrm{ss}}
\phi_i^{(m)}
\,dr
\end{equation}

\begin{equation}
=
\sum_{n=1}^{\infty}
A_n
\sum_{i=1}^{5}
\int_{r_{i-1}}^{r_i}
r^2
\phi_i^{(n)}
\phi_i^{(m)}
\,dr.
\tag{TRISO-ML-548}
\label{eq:triso-ml-548}
\end{equation}

For every $n\ne m$, global orthogonality gives

\begin{equation}
\sum_{i=1}^{5}
\int_{r_{i-1}}^{r_i}
r^2
\phi_i^{(n)}
\phi_i^{(m)}
\,dr
=
0.
\tag{TRISO-ML-549}
\label{eq:triso-ml-549}
\end{equation}

Therefore only the $n=m$ term remains:

\begin{equation}
-\sum_{i=1}^{5}
\int_{r_{i-1}}^{r_i}
r^2
c_{i,\mathrm{ss}}
\phi_i^{(m)}
\,dr
=
A_m
\sum_{i=1}^{5}
\int_{r_{i-1}}^{r_i}
r^2
\left[
\phi_i^{(m)}
\right]^2
dr.
\tag{TRISO-ML-550}
\label{eq:triso-ml-550}
\end{equation}

The denominator is $N_m$:

\begin{equation}
-\sum_{i=1}^{5}
\int_{r_{i-1}}^{r_i}
r^2
c_{i,\mathrm{ss}}
\phi_i^{(m)}
\,dr
=
A_mN_m.
\tag{TRISO-ML-551}
\label{eq:triso-ml-551}
\end{equation}

Divide by $N_m>0$:

\begin{equation}
\boxed{
A_m
=
-
\frac{
\displaystyle
\sum_{i=1}^{5}
\int_{r_{i-1}}^{r_i}
r^2
c_{i,\mathrm{ss}}(r)
\phi_i^{(m)}(r)
\,dr
}{
\displaystyle
\sum_{i=1}^{5}
\int_{r_{i-1}}^{r_i}
r^2
\left[
\phi_i^{(m)}(r)
\right]^2
dr
}.
}
\tag{TRISO-ML-552}
\label{eq:triso-ml-552}
\end{equation}

This is the multilayer modal coefficient formula, conditional on the assumed completeness of the global eigenfunction family for representing the initial transient state.

\subsubsection*{12.21 Complete formal five-layer transient solution}

Because

\begin{equation}
c_i=v_i+c_{i,\mathrm{ss}},
\tag{TRISO-ML-553}
\label{eq:triso-ml-553}
\end{equation}

substitute the modal expansion (TRISO-ML-542):

\begin{equation}
\boxed{
c_i(r,t)
=
c_{i,\mathrm{ss}}(r)
+
\sum_{n=1}^{\infty}
A_n
\phi_i^{(n)}(r)
e^{-\Lambda_nt},
\qquad
r_{i-1}<r<r_i.
}
\tag{TRISO-ML-554}
\label{eq:triso-ml-554}
\end{equation}

\#\#\#\# Initial-time check

At $t=0$,

\begin{equation}
c_i(r,0)
=
c_{i,\mathrm{ss}}(r)
+
\sum_{n=1}^{\infty}
A_n\phi_i^{(n)}(r).
\tag{TRISO-ML-555}
\label{eq:triso-ml-555}
\end{equation}

If the eigenfunction expansion represents the initial transient,

\begin{equation}
\sum_{n=1}^{\infty}
A_n\phi_i^{(n)}(r)
=
-c_{i,\mathrm{ss}}(r).
\tag{TRISO-ML-556}
\label{eq:triso-ml-556}
\end{equation}

Therefore

\begin{equation}
c_i(r,0)=0.
\tag{TRISO-ML-557}
\label{eq:triso-ml-557}
\end{equation}

This recovery is formal and depends on the completeness/convergence statement discussed below.

\#\#\#\# Long-time check

For every positive decay rate,

\begin{equation}
\Lambda_n>0,
\tag{TRISO-ML-558}
\label{eq:triso-ml-558}
\end{equation}

so

\begin{equation}
e^{-\Lambda_nt}\to0
\qquad
(t\to\infty).
\tag{TRISO-ML-559}
\label{eq:triso-ml-559}
\end{equation}

Formally,

\begin{equation}
\boxed{
c_i(r,t)\to c_{i,\mathrm{ss}}(r)
\qquad
(t\to\infty).
}
\tag{TRISO-ML-560}
\label{eq:triso-ml-560}
\end{equation}

\#\#\#\# Interface check

Every global eigenmode separately satisfies

\begin{equation}
\phi_i^{(n)}(r_i)
=
\phi_{i+1}^{(n)}(r_i),
\tag{TRISO-ML-561}
\label{eq:triso-ml-561}
\end{equation}

and

\begin{equation}
D_i\phi_i^{(n)\prime}(r_i)
=
D_{i+1}\phi_{i+1}^{(n)\prime}(r_i).
\tag{TRISO-ML-562}
\label{eq:triso-ml-562}
\end{equation}

A linear combination of such modes therefore preserves both homogeneous interface conditions.

Adding the steady solution restores the corresponding full interface conditions.

\#\#\#\# Centre check

Every kernel eigenfunction has $B_1=0$, so it is regular at $r=0$.

Therefore each transient mode is finite at the centre and satisfies the centre symmetry condition.

\#\#\#\# Outer-boundary check

Every mode satisfies

\begin{equation}
-D_5\phi_5^{(n)\prime}(R)
=
h\phi_5^{(n)}(R).
\tag{TRISO-ML-563}
\label{eq:triso-ml-563}
\end{equation}

Multiplication by the scalar factor $A_ne^{-\Lambda_nt}$ preserves this relation.

Therefore the complete transient sum satisfies the homogeneous Robin condition whenever termwise boundary evaluation is justified.

Adding the steady solution restores the full Robin boundary condition.

\subsubsection*{12.22 Completeness and convergence status}

The orthogonality result (TRISO-ML-528) was derived directly from the conservative differential equations, interface conditions, centre regularity, and outer Robin condition.

It does not require a completeness theorem.

\textbf{[VERIFIED]} Distinct global eigenmodes are orthogonal in the weighted inner product with $w(r)=r^2$.

The projection algebra leading to (TRISO-ML-552) is also valid once an expansion in the eigenfunctions is admitted.

[DERIVED CONDITIONALLY] The modal projection formula follows from orthogonality under the assumption that the initial transient lies in the closure of the global eigenfunction span.

Completeness is a separate spectral statement.

Standard regular Sturm–Liouville completeness theorems provide completeness in an appropriate weighted $L^2$ space for regular self-adjoint problems on finite intervals. The present TRISO problem is more delicate because:

1. $p(r)=r^2D(r)$ vanishes at $r=0$, so the centre is a singular endpoint;
2. $D(r)$ is piecewise constant and discontinuous at four internal interfaces;
3. the operator domain includes transmission conditions enforcing continuity of concentration and flux.

The current derivation has explicitly demonstrated the self-adjoint boundary/interface cancellation needed for symmetry of the operator.

However, this project has not yet supplied a theorem specifically covering the singular endpoint together with these piecewise transmission conditions and proving completeness of the resulting eigenfunctions.

[SOURCE / THEOREM NEEDED] A rigorous completeness theorem for this self-adjoint singular Sturm–Liouville/transmission problem, or an equivalent self-adjoint compact-resolvent operator formulation.

Accordingly, the current statuses are:

\noindent\textbullet\ orthogonality: \textbf{VERIFIED by direct derivation};\\
\noindent\textbullet\ modal projection formula: \textbf{DERIVED CONDITIONALLY};\\
\noindent\textbullet\ completeness: \textbf{THEOREM-DEPENDENT / NOT YET PROVED IN THIS PROJECT};\\
\noindent\textbullet\ numerical convergence of truncated modal sums: \textbf{UNVERIFIED}.\\

The standard Sturm–Liouville framework and regular-problem completeness results provide mathematical context, but they are not being silently promoted to a proof for this singular piecewise problem.

\subsubsection*{12.23 Homogeneous-diffusivity reduction check}

Set

\begin{equation}
D_1=D_2=D_3=D_4=D_5=D.
\tag{TRISO-ML-564}
\label{eq:triso-ml-564}
\end{equation}

Then

\begin{equation}
p(r)=r^2D
\tag{TRISO-ML-565}
\label{eq:triso-ml-565}
\end{equation}

throughout the sphere.

Flux continuity becomes

\begin{equation}
D\phi_i'(r_i)=D\phi_{i+1}'(r_i).
\tag{TRISO-ML-566}
\label{eq:triso-ml-566}
\end{equation}

Cancel $D>0$:

\begin{equation}
\phi_i'(r_i)=\phi_{i+1}'(r_i).
\tag{TRISO-ML-567}
\label{eq:triso-ml-567}
\end{equation}

Together with concentration continuity,

\begin{equation}
\phi_i(r_i)=\phi_{i+1}(r_i),
\tag{TRISO-ML-568}
\label{eq:triso-ml-568}
\end{equation}

the artificial internal material boundaries become transparent to the homogeneous eigenfunction.

The global inner product becomes

\begin{equation}
\sum_{i=1}^{5}
\int_{r_{i-1}}^{r_i}
r^2\phi_m\phi_n\,dr.
\tag{TRISO-ML-569}
\label{eq:triso-ml-569}
\end{equation}

Because the five intervals partition $(0,R)$,

\begin{equation}
\sum_{i=1}^{5}
\int_{r_{i-1}}^{r_i}
r^2\phi_m\phi_n\,dr
=
\int_0^R
r^2\phi_m\phi_n\,dr.
\tag{TRISO-ML-570}
\label{eq:triso-ml-570}
\end{equation}

Therefore the multilayer orthogonality relation reduces to

\begin{equation}
\boxed{
\int_0^R
r^2\phi_m(r)\phi_n(r)\,dr
=
0,
\qquad
m\ne n,
}
\tag{TRISO-ML-571}
\label{eq:triso-ml-571}
\end{equation}

which is exactly the homogeneous spherical weight derived in Section 10.

\subsubsection*{12.24 Unequal-diffusivity interface check}

The interface cancellation did not require

\begin{equation}
D_j=D_{j+1}.
\tag{TRISO-ML-572}
\label{eq:triso-ml-572}
\end{equation}

Instead it used the physical transmission condition

\begin{equation}
D_j\phi_j'
=
D_{j+1}\phi_{j+1}'.
\tag{TRISO-ML-573}
\label{eq:triso-ml-573}
\end{equation}

Therefore unequal adjacent diffusivities remain fully compatible with the global orthogonality proof.

The diffusivity discontinuity is carried by $p(r)=r^2D(r)$, while the weight remains $r^2$.

\subsubsection*{12.25 Updated analytical status}

\textbf{[VERIFIED]} Five-layer steady analytical solution under the frozen ideal benchmark assumptions.

\textbf{[VERIFIED]} Explicit global eigenvalue system and determinant condition.

\textbf{[VERIFIED]} Piecewise conservative self-adjoint/Lagrange structure.

\textbf{[VERIFIED]} Pairwise cancellation of all four ideal-interface boundary terms.

\textbf{[VERIFIED]} Centre and Robin boundary cancellation.

\textbf{[VERIFIED]} Global orthogonality with weight $r^2$.

[DERIVED CONDITIONALLY] Global modal coefficient projection.

[FORMAL / CONDITIONALLY VERIFIED] Infinite five-layer transient modal representation.

[THEOREM-DEPENDENT] Completeness of the global eigenfunction family.

[UNVERIFIED] Numerical convergence rate and truncation error of the modal series.

\subsection*{13. Original FTCS deterministic benchmark discretisation}

This section derives the original Ray notebook's FTCS method as a transparent deterministic benchmark.

It is \textbf{not} the supervisor repository's production Walk-on-Spheres method.

The present subsection is restricted to an interior node lying entirely inside one homogeneous material region, so that the local diffusivity is constant and no material interface lies inside the stencil.

\subsubsection*{13.1 Starting continuous equation}

Inside one homogeneous layer, begin from

\begin{equation}
\frac{\partial c}{\partial t}
=
D
\left(
\frac{\partial^2c}{\partial r^2}
+
\frac2r\frac{\partial c}{\partial r}
\right)
+
S.
\tag{TRISO-DIS-100}
\label{eq:triso-dis-100}
\end{equation}

\textbf{[ASSUMPTION]} $D$ is constant over the local stencil.

\textbf{[ASSUMPTION]} The mesh is uniform for this Part-I benchmark.

\subsubsection*{13.2 Define the spatial mesh}

Divide the interval

\begin{equation}
0\le r\le R
\tag{TRISO-DIS-101}
\label{eq:triso-dis-101}
\end{equation}

into $N$ equal intervals.

Define

\begin{equation}
\Delta r=\frac{R}{N}.
\tag{TRISO-DIS-102}
\label{eq:triso-dis-102}
\end{equation}

The node locations are

\begin{equation}
r_i=i\Delta r,
\qquad
i=0,1,\ldots,N.
\tag{TRISO-DIS-103}
\label{eq:triso-dis-103}
\end{equation}

Therefore

\begin{equation}
r_0=0,
\tag{TRISO-DIS-104}
\label{eq:triso-dis-104}
\end{equation}

and

\begin{equation}
r_N=R.
\tag{TRISO-DIS-105}
\label{eq:triso-dis-105}
\end{equation}

\subsubsection*{13.3 Define the temporal mesh}

Let the time step be

\begin{equation}
\Delta t>0.
\tag{TRISO-DIS-106}
\label{eq:triso-dis-106}
\end{equation}

Define

\begin{equation}
t_j=j\Delta t,
\qquad
j=0,1,2,\ldots.
\tag{TRISO-DIS-107}
\label{eq:triso-dis-107}
\end{equation}

Let

\begin{equation}
C_i^j
\tag{TRISO-DIS-108}
\label{eq:triso-dis-108}
\end{equation}

denote the numerical approximation to

\begin{equation}
c(r_i,t_j).
\tag{TRISO-DIS-109}
\label{eq:triso-dis-109}
\end{equation}

Thus

\begin{equation}
C_i^j\approx c(r_i,t_j).
\tag{TRISO-DIS-110}
\label{eq:triso-dis-110}
\end{equation}

The units remain

\begin{equation}
[C_i^j]=\mathrm{mol\,m^{-3}}.
\tag{TRISO-DIS-111}
\label{eq:triso-dis-111}
\end{equation}

\subsubsection*{13.4 Forward approximation of the time derivative}

At fixed radius $r_i$, Taylor-expand the exact solution from $t_j$ to $t_j+\Delta t$:

\begin{equation}
c(r_i,t_j+\Delta t)
=
c(r_i,t_j)
+
\Delta t
\frac{\partial c}{\partial t}(r_i,t_j)
+
O(\Delta t^2).
\tag{TRISO-DIS-112}
\label{eq:triso-dis-112}
\end{equation}

Subtract $c(r_i,t_j)$ from both sides:

\begin{equation}
c(r_i,t_j+\Delta t)-c(r_i,t_j)
=
\Delta t
\frac{\partial c}{\partial t}(r_i,t_j)
+
O(\Delta t^2).
\tag{TRISO-DIS-113}
\label{eq:triso-dis-113}
\end{equation}

Divide by $\Delta t$:

\begin{equation}
\frac{
c(r_i,t_j+\Delta t)-c(r_i,t_j)
}{
\Delta t
}
=
\frac{\partial c}{\partial t}(r_i,t_j)
+
O(\Delta t).
\tag{TRISO-DIS-114}
\label{eq:triso-dis-114}
\end{equation}

Rearrange:

\begin{equation}
\frac{\partial c}{\partial t}(r_i,t_j)
=
\frac{
c(r_i,t_j+\Delta t)-c(r_i,t_j)
}{
\Delta t
}
+
O(\Delta t).
\tag{TRISO-DIS-115}
\label{eq:triso-dis-115}
\end{equation}

Replace the exact nodal values by the numerical unknowns:

\begin{equation}
\boxed{
\frac{\partial c}{\partial t}(r_i,t_j)
\approx
\frac{C_i^{j+1}-C_i^j}{\Delta t}.
}
\tag{TRISO-DIS-116}
\label{eq:triso-dis-116}
\end{equation}

The forward-time approximation is first-order accurate in time.

\subsubsection*{13.5 Centred approximation of the first radial derivative}

At fixed $t_j$, Taylor-expand about $r_i$ toward $r_i+\Delta r$:

\begin{equation}
c(r_i+\Delta r,t_j)
=
c_i
+
\Delta r\,c_{r,i}
+
\frac{\Delta r^2}{2}c_{rr,i}
+
\frac{\Delta r^3}{6}c_{rrr,i}
+
O(\Delta r^4).
\tag{TRISO-DIS-117}
\label{eq:triso-dis-117}
\end{equation}

Taylor-expand toward $r_i-\Delta r$:

\begin{equation}
c(r_i-\Delta r,t_j)
=
c_i
-
\Delta r\,c_{r,i}
+
\frac{\Delta r^2}{2}c_{rr,i}
-
\frac{\Delta r^3}{6}c_{rrr,i}
+
O(\Delta r^4).
\tag{TRISO-DIS-118}
\label{eq:triso-dis-118}
\end{equation}

Subtract the backward expansion from the forward expansion:

\begin{equation}
c(r_i+\Delta r,t_j)
-
c(r_i-\Delta r,t_j)
=
2\Delta r\,c_{r,i}
+
\frac{\Delta r^3}{3}c_{rrr,i}
+
O(\Delta r^5).
\tag{TRISO-DIS-119}
\label{eq:triso-dis-119}
\end{equation}

Divide by $2\Delta r$:

\begin{equation}
\frac{
c(r_i+\Delta r,t_j)-c(r_i-\Delta r,t_j)
}{
2\Delta r
}
=
c_{r,i}
+
O(\Delta r^2).
\tag{TRISO-DIS-120}
\label{eq:triso-dis-120}
\end{equation}

Therefore

\begin{equation}
\boxed{
\frac{\partial c}{\partial r}(r_i,t_j)
\approx
\frac{
C_{i+1}^j-C_{i-1}^j
}{
2\Delta r
}.
}
\tag{TRISO-DIS-121}
\label{eq:triso-dis-121}
\end{equation}

\subsubsection*{13.6 Centred approximation of the second radial derivative}

Add the two Taylor expansions (TRISO-DIS-117) and (TRISO-DIS-118):

\begin{equation}
c(r_i+\Delta r,t_j)
+
c(r_i-\Delta r,t_j)
=
2c_i
+
\Delta r^2c_{rr,i}
+
O(\Delta r^4).
\tag{TRISO-DIS-122}
\label{eq:triso-dis-122}
\end{equation}

Subtract $2c_i$:

\begin{equation}
c(r_i+\Delta r,t_j)
-
2c_i
+
c(r_i-\Delta r,t_j)
=
\Delta r^2c_{rr,i}
+
O(\Delta r^4).
\tag{TRISO-DIS-123}
\label{eq:triso-dis-123}
\end{equation}

Divide by $\Delta r^2$:

\begin{equation}
\frac{
c(r_i+\Delta r,t_j)-2c_i+c(r_i-\Delta r,t_j)
}{
\Delta r^2
}
=
c_{rr,i}
+
O(\Delta r^2).
\tag{TRISO-DIS-124}
\label{eq:triso-dis-124}
\end{equation}

Therefore

\begin{equation}
\boxed{
\frac{\partial^2c}{\partial r^2}(r_i,t_j)
\approx
\frac{
C_{i-1}^j-2C_i^j+C_{i+1}^j
}{
\Delta r^2
}.
}
\tag{TRISO-DIS-125}
\label{eq:triso-dis-125}
\end{equation}

\subsubsection*{13.7 Substitute the discrete derivatives into the PDE}

Evaluate the continuous equation at $(r_i,t_j)$:

\begin{equation}
c_{t,i}^j
=
D
\left(
c_{rr,i}^j
+
\frac2{r_i}c_{r,i}^j
\right)
+
S_i^j.
\tag{TRISO-DIS-126}
\label{eq:triso-dis-126}
\end{equation}

Substitute the forward-time approximation:

\begin{equation}
\frac{
C_i^{j+1}-C_i^j
}{
\Delta t
}
=
D
\left(
c_{rr,i}^j
+
\frac2{r_i}c_{r,i}^j
\right)
+
S_i^j.
\tag{TRISO-DIS-127}
\label{eq:triso-dis-127}
\end{equation}

Substitute the centred second derivative:

\begin{equation}
\frac{
C_i^{j+1}-C_i^j
}{
\Delta t
}
=
D
\left[
\frac{
C_{i-1}^j-2C_i^j+C_{i+1}^j
}{
\Delta r^2
}
+
\frac2{r_i}c_{r,i}^j
\right]
+
S_i^j.
\tag{TRISO-DIS-128}
\label{eq:triso-dis-128}
\end{equation}

Substitute the centred first derivative:

\begin{equation}
\frac{
C_i^{j+1}-C_i^j
}{
\Delta t
}
=
D
\left[
\frac{
C_{i-1}^j-2C_i^j+C_{i+1}^j
}{
\Delta r^2
}
+
\frac2{r_i}
\frac{
C_{i+1}^j-C_{i-1}^j
}{
2\Delta r
}
\right]
+
S_i^j.
\tag{TRISO-DIS-129}
\label{eq:triso-dis-129}
\end{equation}

Cancel the factor $2$ in the radial first-derivative term:

\begin{equation}
\frac{
C_i^{j+1}-C_i^j
}{
\Delta t
}
=
D
\left[
\frac{
C_{i-1}^j-2C_i^j+C_{i+1}^j
}{
\Delta r^2
}
+
\frac{
C_{i+1}^j-C_{i-1}^j
}{
r_i\Delta r
}
\right]
+
S_i^j.
\tag{TRISO-DIS-130}
\label{eq:triso-dis-130}
\end{equation}

Use the uniform-mesh identity

\begin{equation}
r_i=i\Delta r.
\tag{TRISO-DIS-131}
\label{eq:triso-dis-131}
\end{equation}

Then

\begin{equation}
r_i\Delta r
=
i\Delta r^2.
\tag{TRISO-DIS-132}
\label{eq:triso-dis-132}
\end{equation}

Therefore

\begin{equation}
\frac{
C_i^{j+1}-C_i^j
}{
\Delta t
}
=
\frac{D}{\Delta r^2}
\left[
C_{i-1}^j
-
2C_i^j
+
C_{i+1}^j
+
\frac1i
\left(
C_{i+1}^j-C_{i-1}^j
\right)
\right]
+
S_i^j.
\tag{TRISO-DIS-133}
\label{eq:triso-dis-133}
\end{equation}

\subsubsection*{13.8 Collect neighbour coefficients}

Expand the $1/i$ term:

\begin{equation}
\frac{
C_i^{j+1}-C_i^j
}{
\Delta t
}
=
\frac{D}{\Delta r^2}
\left[
C_{i-1}^j
-
2C_i^j
+
C_{i+1}^j
+
\frac1iC_{i+1}^j
-
\frac1iC_{i-1}^j
\right]
+
S_i^j.
\tag{TRISO-DIS-134}
\label{eq:triso-dis-134}
\end{equation}

Collect the $C_{i-1}^j$ terms:

\begin{equation}
C_{i-1}^j
-
\frac1iC_{i-1}^j
=
\left(
1-\frac1i
\right)
C_{i-1}^j.
\tag{TRISO-DIS-135}
\label{eq:triso-dis-135}
\end{equation}

Collect the $C_{i+1}^j$ terms:

\begin{equation}
C_{i+1}^j
+
\frac1iC_{i+1}^j
=
\left(
1+\frac1i
\right)
C_{i+1}^j.
\tag{TRISO-DIS-136}
\label{eq:triso-dis-136}
\end{equation}

Thus

\begin{equation}
\frac{
C_i^{j+1}-C_i^j
}{
\Delta t
}
=
\frac{D}{\Delta r^2}
\left[
\left(
1-\frac1i
\right)C_{i-1}^j
-
2C_i^j
+
\left(
1+\frac1i
\right)C_{i+1}^j
\right]
+
S_i^j.
\tag{TRISO-DIS-137}
\label{eq:triso-dis-137}
\end{equation}

Multiply by $\Delta t$:

\begin{equation}
C_i^{j+1}-C_i^j
=
\frac{D\Delta t}{\Delta r^2}
\left[
\left(
1-\frac1i
\right)C_{i-1}^j
-
2C_i^j
+
\left(
1+\frac1i
\right)C_{i+1}^j
\right]
+
S_i^j\Delta t.
\tag{TRISO-DIS-138}
\label{eq:triso-dis-138}
\end{equation}

Add $C_i^j$ to both sides:

\begin{equation}
C_i^{j+1}
=
C_i^j
+
\frac{D\Delta t}{\Delta r^2}
\left[
\left(
1-\frac1i
\right)C_{i-1}^j
-
2C_i^j
+
\left(
1+\frac1i
\right)C_{i+1}^j
\right]
+
S_i^j\Delta t.
\tag{TRISO-DIS-139}
\label{eq:triso-dis-139}
\end{equation}

\subsubsection*{13.9 Define the Fourier number}

Define

\begin{equation}
\boxed{
\mathrm{Fo}
=
\frac{D\Delta t}{\Delta r^2}.
}
\tag{TRISO-DIS-140}
\label{eq:triso-dis-140}
\end{equation}

Its units are

\begin{equation}
[\mathrm{Fo}]
=
\frac{
\mathrm{m^2\,s^{-1}}\mathrm{s}
}{
\mathrm{m^2}
}
=
1.
\tag{TRISO-DIS-141}
\label{eq:triso-dis-141}
\end{equation}

Therefore $\mathrm{Fo}$ is dimensionless.

Substitute the definition into (TRISO-DIS-139):

\begin{equation}
C_i^{j+1}
=
C_i^j
+
\mathrm{Fo}
\left[
\left(
1-\frac1i
\right)C_{i-1}^j
-
2C_i^j
+
\left(
1+\frac1i
\right)C_{i+1}^j
\right]
+
S_i^j\Delta t.
\tag{TRISO-DIS-142}
\label{eq:triso-dis-142}
\end{equation}

Distribute $\mathrm{Fo}$:

\begin{equation}
C_i^{j+1}
=
C_i^j
+
\mathrm{Fo}
\left(
1-\frac1i
\right)C_{i-1}^j
-
2\mathrm{Fo}C_i^j
+
\mathrm{Fo}
\left(
1+\frac1i
\right)C_{i+1}^j
+
S_i^j\Delta t.
\tag{TRISO-DIS-143}
\label{eq:triso-dis-143}
\end{equation}

Collect the central coefficient:

\begin{equation}
C_i^j-2\mathrm{Fo}C_i^j
=
(1-2\mathrm{Fo})C_i^j.
\tag{TRISO-DIS-144}
\label{eq:triso-dis-144}
\end{equation}

Hence the original Ray interior FTCS update is

\begin{equation}
\boxed{
C_i^{j+1}
=
\mathrm{Fo}
\left(
1-\frac1i
\right)C_{i-1}^j
+
(1-2\mathrm{Fo})C_i^j
+
\mathrm{Fo}
\left(
1+\frac1i
\right)C_{i+1}^j
+
S_i^j\Delta t.
}
\tag{TRISO-DIS-145}
\label{eq:triso-dis-145}
\end{equation}

This is algebraically equivalent to the ordering used in the original notebook.

\subsubsection*{13.10 Dimensional check of the source increment}

The source contribution is

\begin{equation}
S_i^j\Delta t.
\tag{TRISO-DIS-146}
\label{eq:triso-dis-146}
\end{equation}

Its units are

\begin{equation}
[S_i^j\Delta t]
=
\mathrm{mol\,m^{-3}\,s^{-1}}
\times
\mathrm{s}.
\tag{TRISO-DIS-147}
\label{eq:triso-dis-147}
\end{equation}

Therefore

\begin{equation}
[S_i^j\Delta t]
=
\mathrm{mol\,m^{-3}},
\tag{TRISO-DIS-148}
\label{eq:triso-dis-148}
\end{equation}

which matches the concentration units of every other term in (TRISO-DIS-145).

\subsubsection*{13.11 Local consistency order}

The forward-time derivative has truncation error

\begin{equation}
O(\Delta t).
\tag{TRISO-DIS-149}
\label{eq:triso-dis-149}
\end{equation}

The centred first derivative has truncation error

\begin{equation}
O(\Delta r^2).
\tag{TRISO-DIS-150}
\label{eq:triso-dis-150}
\end{equation}

The centred second derivative has truncation error

\begin{equation}
O(\Delta r^2).
\tag{TRISO-DIS-151}
\label{eq:triso-dis-151}
\end{equation}

Therefore, away from $r=0$, material interfaces, and the outer boundary, the local differential approximation is formally

\begin{equation}
\boxed{
O(\Delta t)+O(\Delta r^2).
}
\tag{TRISO-DIS-152}
\label{eq:triso-dis-152}
\end{equation}

This is a consistency statement only. It is not by itself a proof of stability or convergence.

\subsubsection*{13.12 Scope and limitations of the interior stencil}

Equation (TRISO-DIS-145) assumes:

1. $i\ge1$, so the $1/i$ factor is defined;
2. the stencil lies inside one homogeneous constant-$D$ region;
3. the spatial mesh is uniform;
4. the time step is uniform;
5. the source value $S_i^j$ is known explicitly at time level $j$.

It must \textbf{not} be applied unchanged:

\noindent\textbullet\ at $r=0$;\\
\noindent\textbullet\ across a discontinuous material interface;\\
\noindent\textbullet\ at the outer boundary.\\

Those cases require separate derivations.

The original notebook's FTCS method is retained as a transparent deterministic benchmark. It does not redefine the production WOS method and it does not yet establish the preferred discretisation for the full five-layer discontinuous-$D$ problem.

\subsection*{14. Centre discretisation}

The interior stencil in Section 13 cannot be evaluated at $i=0$ because it contains the factor $1/i$.

The centre must therefore be derived from the regular spherical limit.

\subsubsection*{14.1 Continuous centre operator}

Inside the homogeneous kernel, the radial diffusion operator is

\begin{equation}
\mathcal L[c]
=
\frac{\partial^2c}{\partial r^2}
+
\frac2r\frac{\partial c}{\partial r}.
\tag{TRISO-DIS-200}
\label{eq:triso-dis-200}
\end{equation}

For a sufficiently smooth spherically symmetric field,

\begin{equation}
\frac{\partial c}{\partial r}(0,t)=0.
\tag{TRISO-DIS-201}
\label{eq:triso-dis-201}
\end{equation}

Consider the apparently singular term

\begin{equation}
\lim_{r\to0}
\frac2r
\frac{\partial c}{\partial r}.
\tag{TRISO-DIS-202}
\label{eq:triso-dis-202}
\end{equation}

Because the numerator tends to zero,

\begin{equation}
\lim_{r\to0}
\frac{\partial c/\partial r}{r}
\end{equation}

has the indeterminate form $0/0$.

Apply L'Hôpital's rule with respect to $r$:

\begin{equation}
\lim_{r\to0}
\frac{\partial c/\partial r}{r}
=
\lim_{r\to0}
\frac{\partial^2c/\partial r^2}{1}.
\tag{TRISO-DIS-203}
\label{eq:triso-dis-203}
\end{equation}

Therefore

\begin{equation}
\lim_{r\to0}
\frac1r
\frac{\partial c}{\partial r}
=
\frac{\partial^2c}{\partial r^2}(0,t).
\tag{TRISO-DIS-204}
\label{eq:triso-dis-204}
\end{equation}

Multiply by $2$:

\begin{equation}
\lim_{r\to0}
\frac2r
\frac{\partial c}{\partial r}
=
2
\frac{\partial^2c}{\partial r^2}(0,t).
\tag{TRISO-DIS-205}
\label{eq:triso-dis-205}
\end{equation}

Hence

\begin{equation}
\mathcal L[c](0,t)
=
\frac{\partial^2c}{\partial r^2}(0,t)
+
2\frac{\partial^2c}{\partial r^2}(0,t).
\tag{TRISO-DIS-206}
\label{eq:triso-dis-206}
\end{equation}

Collect the terms:

\begin{equation}
\boxed{
\mathcal L[c](0,t)
=
3
\frac{\partial^2c}{\partial r^2}(0,t).
}
\tag{TRISO-DIS-207}
\label{eq:triso-dis-207}
\end{equation}

This is the exact smooth-origin limit of the spherical radial operator.

\subsubsection*{14.2 Introduce a symmetric ghost point}

The uniform mesh has

\begin{equation}
r_0=0,
\qquad
r_1=\Delta r.
\tag{TRISO-DIS-208}
\label{eq:triso-dis-208}
\end{equation}

For the purpose of constructing a centred derivative at the origin, introduce a mathematical ghost point

\begin{equation}
r_{-1}=-\Delta r.
\tag{TRISO-DIS-209}
\label{eq:triso-dis-209}
\end{equation}

Spherical symmetry corresponds to an even extension of the radial concentration:

\begin{equation}
c(-r,t)=c(r,t).
\tag{TRISO-DIS-210}
\label{eq:triso-dis-210}
\end{equation}

Evaluate this at $r=\Delta r$:

\begin{equation}
c(-\Delta r,t)=c(\Delta r,t).
\tag{TRISO-DIS-211}
\label{eq:triso-dis-211}
\end{equation}

In nodal notation,

\begin{equation}
\boxed{
C_{-1}^j=C_1^j.
}
\tag{TRISO-DIS-212}
\label{eq:triso-dis-212}
\end{equation}

The ghost point is a mathematical device. It does not represent a physical negative-radius material region.

\subsubsection*{14.3 Centre second derivative}

Use the centred second-derivative formula at $i=0$:

\begin{equation}
\frac{\partial^2c}{\partial r^2}(0,t_j)
\approx
\frac{
C_{-1}^j-2C_0^j+C_1^j
}{
\Delta r^2
}.
\tag{TRISO-DIS-213}
\label{eq:triso-dis-213}
\end{equation}

Substitute the symmetry relation $C_{-1}^j=C_1^j$:

\begin{equation}
\frac{\partial^2c}{\partial r^2}(0,t_j)
\approx
\frac{
C_1^j-2C_0^j+C_1^j
}{
\Delta r^2
}.
\tag{TRISO-DIS-214}
\label{eq:triso-dis-214}
\end{equation}

Add the two $C_1^j$ terms:

\begin{equation}
\boxed{
\frac{\partial^2c}{\partial r^2}(0,t_j)
\approx
\frac{
2(C_1^j-C_0^j)
}{
\Delta r^2
}.
}
\tag{TRISO-DIS-215}
\label{eq:triso-dis-215}
\end{equation}

\subsubsection*{14.4 Discrete spherical operator at the centre}

From the exact centre limit,

\begin{equation}
\mathcal L[c](0,t)
=
3c_{rr}(0,t).
\tag{TRISO-DIS-216}
\label{eq:triso-dis-216}
\end{equation}

Substitute the discrete second derivative (TRISO-DIS-215):

\begin{equation}
\mathcal L[c](0,t_j)
\approx
3
\frac{
2(C_1^j-C_0^j)
}{
\Delta r^2
}.
\tag{TRISO-DIS-217}
\label{eq:triso-dis-217}
\end{equation}

Multiply the factors $3$ and $2$:

\begin{equation}
\boxed{
\mathcal L[c](0,t_j)
\approx
\frac{
6(C_1^j-C_0^j)
}{
\Delta r^2
}.
}
\tag{TRISO-DIS-218}
\label{eq:triso-dis-218}
\end{equation}

This is the origin of the factor $6$ in the Ray notebook's centre update.

\subsubsection*{14.5 Apply the centre PDE}

In the homogeneous kernel, the PDE at the centre is interpreted through the regular limit:

\begin{equation}
\frac{\partial c}{\partial t}(0,t)
=
D_1\mathcal L[c](0,t)
+
S_0.
\tag{TRISO-DIS-219}
\label{eq:triso-dis-219}
\end{equation}

Use the forward-time approximation:

\begin{equation}
\frac{
C_0^{j+1}-C_0^j
}{
\Delta t
}
=
D_1\mathcal L[c](0,t_j)
+
S_0.
\tag{TRISO-DIS-220}
\label{eq:triso-dis-220}
\end{equation}

Substitute (TRISO-DIS-218):

\begin{equation}
\frac{
C_0^{j+1}-C_0^j
}{
\Delta t
}
=
D_1
\frac{
6(C_1^j-C_0^j)
}{
\Delta r^2
}
+
S_0.
\tag{TRISO-DIS-221}
\label{eq:triso-dis-221}
\end{equation}

Multiply by $\Delta t$:

\begin{equation}
C_0^{j+1}-C_0^j
=
\frac{
6D_1\Delta t
}{
\Delta r^2
}
(C_1^j-C_0^j)
+
S_0\Delta t.
\tag{TRISO-DIS-222}
\label{eq:triso-dis-222}
\end{equation}

Define the kernel Fourier number

\begin{equation}
\boxed{
\mathrm{Fo}_1
=
\frac{D_1\Delta t}{\Delta r^2}.
}
\tag{TRISO-DIS-223}
\label{eq:triso-dis-223}
\end{equation}

Substitute it:

\begin{equation}
C_0^{j+1}-C_0^j
=
6\mathrm{Fo}_1
(C_1^j-C_0^j)
+
S_0\Delta t.
\tag{TRISO-DIS-224}
\label{eq:triso-dis-224}
\end{equation}

Add $C_0^j$ to both sides:

\begin{equation}
\boxed{
C_0^{j+1}
=
C_0^j
+
6\mathrm{Fo}_1
(C_1^j-C_0^j)
+
S_0\Delta t.
}
\tag{TRISO-DIS-225}
\label{eq:triso-dis-225}
\end{equation}

Expand the difference:

\begin{equation}
C_0^{j+1}
=
C_0^j
+
6\mathrm{Fo}_1C_1^j
-
6\mathrm{Fo}_1C_0^j
+
S_0\Delta t.
\tag{TRISO-DIS-226}
\label{eq:triso-dis-226}
\end{equation}

Collect the centre coefficient:

\begin{equation}
\boxed{
C_0^{j+1}
=
(1-6\mathrm{Fo}_1)C_0^j
+
6\mathrm{Fo}_1C_1^j
+
S_0\Delta t.
}
\tag{TRISO-DIS-227}
\label{eq:triso-dis-227}
\end{equation}

\subsubsection*{14.6 Consistency of the centre approximation}

For a smooth even radial field, expand about $r=0$:

\begin{equation}
c(\Delta r,t)
=
c(0,t)
+
\frac{\Delta r^2}{2}c_{rr}(0,t)
+
\frac{\Delta r^4}{24}c_{rrrr}(0,t)
+
O(\Delta r^6).
\tag{TRISO-DIS-228}
\label{eq:triso-dis-228}
\end{equation}

Subtract $c(0,t)$:

\begin{equation}
c(\Delta r,t)-c(0,t)
=
\frac{\Delta r^2}{2}c_{rr}(0,t)
+
\frac{\Delta r^4}{24}c_{rrrr}(0,t)
+
O(\Delta r^6).
\tag{TRISO-DIS-229}
\label{eq:triso-dis-229}
\end{equation}

Multiply by $2/\Delta r^2$:

\begin{equation}
\frac{
2[c(\Delta r,t)-c(0,t)]
}{
\Delta r^2
}
=
c_{rr}(0,t)
+
\frac{\Delta r^2}{12}c_{rrrr}(0,t)
+
O(\Delta r^4).
\tag{TRISO-DIS-230}
\label{eq:triso-dis-230}
\end{equation}

Therefore the centre second derivative in (TRISO-DIS-215) is second-order accurate in space:

\begin{equation}
c_{rr}(0,t)
=
\frac{
2(C_1-C_0)
}{
\Delta r^2
}
+
O(\Delta r^2).
\tag{TRISO-DIS-231}
\label{eq:triso-dis-231}
\end{equation}

Multiplying by the exact factor $3$ does not change the spatial order:

\begin{equation}
\mathcal L[c](0,t)
=
\frac{
6(C_1-C_0)
}{
\Delta r^2
}
+
O(\Delta r^2).
\tag{TRISO-DIS-232}
\label{eq:triso-dis-232}
\end{equation}

Combined with forward Euler time stepping, the centre equation is locally

\begin{equation}
\boxed{
O(\Delta t)+O(\Delta r^2).
}
\tag{TRISO-DIS-233}
\label{eq:triso-dis-233}
\end{equation}

\subsubsection*{14.7 Scope of the centre formula}

The centre formula requires:

\noindent\textbullet\ spherical symmetry;\\
\noindent\textbullet\ sufficient smoothness at $r=0$;\\
\noindent\textbullet\ the centre to lie inside one homogeneous kernel material;\\
\noindent\textbullet\ constant $D_1$ over the centre stencil.\\

It does not determine the treatment of material interfaces or the outer surface.

The coefficient $1-6\mathrm{Fo}_1$ will later enter the FTCS monotonicity/stability discussion, but no stability conclusion is drawn here.

\subsection*{15. Material-interface discretisation for discontinuous diffusivity}

The homogeneous interior FTCS stencil from Section 13 must not be centred across a material interface because the diffusivity is discontinuous there.

This section derives an interface-aligned discrete constraint directly from the two physical interface conditions already established in the continuous model:

1. ideal concentration continuity;
2. diffusive flux continuity.

The derivation is for the frozen ideal benchmark with $K=1$ and no explicit interfacial resistance.

\subsubsection*{15.1 Place a mesh node at the material interface}

Let a material interface occur at

\begin{equation}
r=r_I.
\tag{TRISO-DIS-300}
\label{eq:triso-dis-300}
\end{equation}

Choose the mesh so that one numerical node lies exactly at the interface:

\begin{equation}
r_I=r_{i}.
\tag{TRISO-DIS-301}
\label{eq:triso-dis-301}
\end{equation}

Let the material immediately inside the interface have diffusivity

\begin{equation}
D^-,
\tag{TRISO-DIS-302}
\label{eq:triso-dis-302}
\end{equation}

and the material immediately outside have diffusivity

\begin{equation}
D^+.
\tag{TRISO-DIS-303}
\label{eq:triso-dis-303}
\end{equation}

For a uniform mesh,

\begin{equation}
r_{i-1}=r_I-\Delta r,
\tag{TRISO-DIS-304}
\label{eq:triso-dis-304}
\end{equation}

and

\begin{equation}
r_{i+1}=r_I+\Delta r.
\tag{TRISO-DIS-305}
\label{eq:triso-dis-305}
\end{equation}

The interface concentration is represented by one nodal unknown:

\begin{equation}
C_I^j.
\tag{TRISO-DIS-306}
\label{eq:triso-dis-306}
\end{equation}

\subsubsection*{15.2 Discrete concentration continuity}

For the ideal $K=1$ interface, the continuous condition is

\begin{equation}
c^-(r_I,t)=c^+(r_I,t).
\tag{TRISO-DIS-307}
\label{eq:triso-dis-307}
\end{equation}

Represent both limiting concentrations by the same interface unknown:

\begin{equation}
c^-(r_I,t_j)\approx C_I^j,
\tag{TRISO-DIS-308}
\label{eq:triso-dis-308}
\end{equation}

and

\begin{equation}
c^+(r_I,t_j)\approx C_I^j.
\tag{TRISO-DIS-309}
\label{eq:triso-dis-309}
\end{equation}

Thus concentration continuity is built directly into the interface-node representation.

No averaging of the two material concentrations is required because there is only one ideal-interface concentration degree of freedom.

\subsubsection*{15.3 Continuous flux continuity}

The exact ideal-interface condition is

\begin{equation}
-D^-
\left.
\frac{\partial c^-}{\partial r}
\right|_{r_I^-}
=
-D^+
\left.
\frac{\partial c^+}{\partial r}
\right|_{r_I^+}.
\tag{TRISO-DIS-310}
\label{eq:triso-dis-310}
\end{equation}

The minus signs occur because the outward radial diffusive flux is

\begin{equation}
J_r=-D\frac{\partial c}{\partial r}.
\tag{TRISO-DIS-311}
\label{eq:triso-dis-311}
\end{equation}

\subsubsection*{15.4 Approximate the inner-side gradient}

On the inner material side, use the interface node and its inner neighbour.

The radial distance is

\begin{equation}
r_I-r_{i-1}=\Delta r.
\tag{TRISO-DIS-312}
\label{eq:triso-dis-312}
\end{equation}

A first-order one-sided approximation is

\begin{equation}
\left.
\frac{\partial c^-}{\partial r}
\right|_{r_I^-}
\approx
\frac{
C_I^j-C_{i-1}^j
}{
\Delta r
}.
\tag{TRISO-DIS-313}
\label{eq:triso-dis-313}
\end{equation}

Therefore the inner-side outward flux is approximated by

\begin{equation}
J_I^-
\approx
-D^-
\frac{
C_I^j-C_{i-1}^j
}{
\Delta r
}.
\tag{TRISO-DIS-314}
\label{eq:triso-dis-314}
\end{equation}

\subsubsection*{15.5 Approximate the outer-side gradient}

On the outer material side,

\begin{equation}
r_{i+1}-r_I=\Delta r.
\tag{TRISO-DIS-315}
\label{eq:triso-dis-315}
\end{equation}

The one-sided gradient is

\begin{equation}
\left.
\frac{\partial c^+}{\partial r}
\right|_{r_I^+}
\approx
\frac{
C_{i+1}^j-C_I^j
}{
\Delta r
}.
\tag{TRISO-DIS-316}
\label{eq:triso-dis-316}
\end{equation}

Therefore

\begin{equation}
J_I^+
\approx
-D^+
\frac{
C_{i+1}^j-C_I^j
}{
\Delta r
}.
\tag{TRISO-DIS-317}
\label{eq:triso-dis-317}
\end{equation}

\subsubsection*{15.6 Enforce discrete flux continuity}

Set the two approximated fluxes equal:

\begin{equation}
-D^-
\frac{
C_I^j-C_{i-1}^j
}{
\Delta r
}
=
-D^+
\frac{
C_{i+1}^j-C_I^j
}{
\Delta r
}.
\tag{TRISO-DIS-318}
\label{eq:triso-dis-318}
\end{equation}

Cancel the common factor $-1/\Delta r$:

\begin{equation}
D^-
\left(
C_I^j-C_{i-1}^j
\right)
=
D^+
\left(
C_{i+1}^j-C_I^j
\right).
\tag{TRISO-DIS-319}
\label{eq:triso-dis-319}
\end{equation}

Expand both sides:

\begin{equation}
D^-C_I^j-D^-C_{i-1}^j
=
D^+C_{i+1}^j-D^+C_I^j.
\tag{TRISO-DIS-320}
\label{eq:triso-dis-320}
\end{equation}

Add $D^+C_I^j$ to both sides:

\begin{equation}
(D^-+D^+)C_I^j-D^-C_{i-1}^j
=
D^+C_{i+1}^j.
\tag{TRISO-DIS-321}
\label{eq:triso-dis-321}
\end{equation}

Add $D^-C_{i-1}^j$ to both sides:

\begin{equation}
(D^-+D^+)C_I^j
=
D^-C_{i-1}^j
+
D^+C_{i+1}^j.
\tag{TRISO-DIS-322}
\label{eq:triso-dis-322}
\end{equation}

Divide by $D^-+D^+>0$:

\begin{equation}
\boxed{
C_I^j
=
\frac{
D^-C_{i-1}^j
+
D^+C_{i+1}^j
}{
D^-+D^+
}.
}
\tag{TRISO-DIS-323}
\label{eq:triso-dis-323}
\end{equation}

This is an algebraic interface constraint, not a homogeneous-material FTCS update.

\subsubsection*{15.7 Unequal grid spacing}

The same derivation can be retained if the interface is not equally spaced from its neighbouring nodes.

Define

\begin{equation}
\Delta r^-=r_I-r_{i-1},
\tag{TRISO-DIS-324}
\label{eq:triso-dis-324}
\end{equation}

and

\begin{equation}
\Delta r^+=r_{i+1}-r_I.
\tag{TRISO-DIS-325}
\label{eq:triso-dis-325}
\end{equation}

Then the two flux approximations are

\begin{equation}
J_I^-
\approx
-D^-
\frac{
C_I-C_{i-1}
}{
\Delta r^-
},
\tag{TRISO-DIS-326}
\label{eq:triso-dis-326}
\end{equation}

and

\begin{equation}
J_I^+
\approx
-D^+
\frac{
C_{i+1}-C_I
}{
\Delta r^+
}.
\tag{TRISO-DIS-327}
\label{eq:triso-dis-327}
\end{equation}

Flux continuity gives

\begin{equation}
\frac{D^-}{\Delta r^-}
(C_I-C_{i-1})
=
\frac{D^+}{\Delta r^+}
(C_{i+1}-C_I).
\tag{TRISO-DIS-328}
\label{eq:triso-dis-328}
\end{equation}

Expand:

\begin{equation}
\frac{D^-}{\Delta r^-}C_I
-
\frac{D^-}{\Delta r^-}C_{i-1}
=
\frac{D^+}{\Delta r^+}C_{i+1}
-
\frac{D^+}{\Delta r^+}C_I.
\tag{TRISO-DIS-329}
\label{eq:triso-dis-329}
\end{equation}

Collect the interface unknown:

\begin{equation}
\left(
\frac{D^-}{\Delta r^-}
+
\frac{D^+}{\Delta r^+}
\right)
C_I
=
\frac{D^-}{\Delta r^-}C_{i-1}
+
\frac{D^+}{\Delta r^+}C_{i+1}.
\tag{TRISO-DIS-330}
\label{eq:triso-dis-330}
\end{equation}

Therefore

\begin{equation}
\boxed{
C_I
=
\frac{
\dfrac{D^-}{\Delta r^-}C_{i-1}
+
\dfrac{D^+}{\Delta r^+}C_{i+1}
}{
\dfrac{D^-}{\Delta r^-}
+
\dfrac{D^+}{\Delta r^+}
}.
}
\tag{TRISO-DIS-331}
\label{eq:triso-dis-331}
\end{equation}

Equation (TRISO-DIS-323) is recovered when

\begin{equation}
\Delta r^-=\Delta r^+=\Delta r.
\tag{TRISO-DIS-332}
\label{eq:triso-dis-332}
\end{equation}

\subsubsection*{15.8 Equivalent two-node conductance and harmonic diffusivity}

Sometimes the interface concentration is eliminated so that the flux is written directly between the two neighbouring material nodes.

Start from the inner-side flux relation:

\begin{equation}
J_I
=
-D^-
\frac{
C_I-C_{i-1}
}{
\Delta r^-
}.
\tag{TRISO-DIS-333}
\label{eq:triso-dis-333}
\end{equation}

Rearrange for the inner concentration drop:

\begin{equation}
C_{i-1}-C_I
=
J_I
\frac{\Delta r^-}{D^-}.
\tag{TRISO-DIS-334}
\label{eq:triso-dis-334}
\end{equation}

From the outer-side relation,

\begin{equation}
J_I
=
-D^+
\frac{
C_{i+1}-C_I
}{
\Delta r^+
}.
\tag{TRISO-DIS-335}
\label{eq:triso-dis-335}
\end{equation}

Rearrange:

\begin{equation}
C_I-C_{i+1}
=
J_I
\frac{\Delta r^+}{D^+}.
\tag{TRISO-DIS-336}
\label{eq:triso-dis-336}
\end{equation}

Add the two concentration drops:

\begin{equation}
C_{i-1}-C_{i+1}
=
J_I
\left(
\frac{\Delta r^-}{D^-}
+
\frac{\Delta r^+}{D^+}
\right).
\tag{TRISO-DIS-337}
\label{eq:triso-dis-337}
\end{equation}

Solve for the flux:

\begin{equation}
\boxed{
J_I
=
\frac{
C_{i-1}-C_{i+1}
}{
\dfrac{\Delta r^-}{D^-}
+
\dfrac{\Delta r^+}{D^+}
}.
}
\tag{TRISO-DIS-338}
\label{eq:triso-dis-338}
\end{equation}

The denominator is the sum of the two local diffusion resistances per unit area.

Define the total node-to-node distance

\begin{equation}
\Delta r_{\mathrm{tot}}
=
\Delta r^-+\Delta r^+.
\tag{TRISO-DIS-339}
\label{eq:triso-dis-339}
\end{equation}

Define an effective diffusivity $D_{\mathrm{eff}}$ by

\begin{equation}
J_I
=
D_{\mathrm{eff}}
\frac{
C_{i-1}-C_{i+1}
}{
\Delta r_{\mathrm{tot}}
}.
\tag{TRISO-DIS-340}
\label{eq:triso-dis-340}
\end{equation}

Equate (TRISO-DIS-338) and (TRISO-DIS-340):

\begin{equation}
\frac{D_{\mathrm{eff}}}{\Delta r_{\mathrm{tot}}}
=
\frac1{
\dfrac{\Delta r^-}{D^-}
+
\dfrac{\Delta r^+}{D^+}
}.
\tag{TRISO-DIS-341}
\label{eq:triso-dis-341}
\end{equation}

Multiply by $\Delta r_{\mathrm{tot}}$:

\begin{equation}
\boxed{
D_{\mathrm{eff}}
=
\frac{
\Delta r^-+\Delta r^+
}{
\dfrac{\Delta r^-}{D^-}
+
\dfrac{\Delta r^+}{D^+}
}.
}
\tag{TRISO-DIS-342}
\label{eq:triso-dis-342}
\end{equation}

For equal half-distances,

\begin{equation}
\Delta r^-=\Delta r^+,
\tag{TRISO-DIS-343}
\label{eq:triso-dis-343}
\end{equation}

the effective diffusivity becomes

\begin{equation}
D_{\mathrm{eff}}
=
\frac{2}{
\dfrac1{D^-}
+
\dfrac1{D^+}
}.
\tag{TRISO-DIS-344}
\label{eq:triso-dis-344}
\end{equation}

Thus

\begin{equation}
\boxed{
D_{\mathrm{eff}}
=
\frac{
2D^-D^+
}{
D^-+D^+
}.
}
\tag{TRISO-DIS-345}
\label{eq:triso-dis-345}
\end{equation}

This is the harmonic mean of the adjacent diffusivities.

It appears because diffusion resistances add in series; it is not an arbitrary averaging rule.

\subsubsection*{15.9 Limiting checks}

If

\begin{equation}
D^-=D^+=D,
\tag{TRISO-DIS-346}
\label{eq:triso-dis-346}
\end{equation}

then (TRISO-DIS-323) becomes

\begin{equation}
C_I
=
\frac{
DC_{i-1}+DC_{i+1}
}{
2D
}.
\tag{TRISO-DIS-347}
\label{eq:triso-dis-347}
\end{equation}

Cancel $D$:

\begin{equation}
\boxed{
C_I
=
\frac{
C_{i-1}+C_{i+1}
}{2}.
}
\tag{TRISO-DIS-348}
\label{eq:triso-dis-348}
\end{equation}

This is the expected linear interpolation for equal diffusivity and equal spacing.

If

\begin{equation}
D^+\ll D^-,
\tag{TRISO-DIS-349}
\label{eq:triso-dis-349}
\end{equation}

then the low-diffusivity outer material contributes the dominant diffusion resistance in (TRISO-DIS-338).

The interface concentration correspondingly approaches the concentration on the high-diffusivity side only according to the resistance-weighted relation; it is not obtained from an arithmetic diffusivity average.

\subsubsection*{15.10 Accuracy and scope}

The one-sided interface gradients in (TRISO-DIS-313) and (TRISO-DIS-316) are first-order approximations to the limiting interface derivatives when used with only one neighbouring node on each side.

Therefore the algebraic interface relation derived here is conservative with respect to the approximated flux, but this particular gradient construction is not automatically second-order accurate at the interface.

A higher-order interface treatment would require additional same-material nodes or a finite-volume formulation with carefully defined face fluxes.

\textbf{[IMPORTANT]} The present result establishes the correct \textbf{discrete transmission logic}:

\noindent\textbullet\ one ideal-interface concentration;\\
\noindent\textbullet\ no differentiation of $D$ through its jump;\\
\noindent\textbullet\ equal discrete flux on both sides;\\
\noindent\textbullet\ resistance-weighted, rather than arithmetic, diffusivity coupling.\\

It does not yet establish the globally preferred five-layer spatial discretisation.

\subsection*{16. Outer Robin boundary: ghost-point elimination and surface update}

The outer surface is

\begin{equation}
r_N=R.
\tag{TRISO-DIS-400}
\label{eq:triso-dis-400}
\end{equation}

For the homogeneous Part-I FTCS benchmark, the outer material has constant diffusivity $D$ and the external bulk concentration is

\begin{equation}
c_\infty=0.
\tag{TRISO-DIS-401}
\label{eq:triso-dis-401}
\end{equation}

The continuous Robin condition is therefore

\begin{equation}
-D
\frac{\partial c}{\partial r}(R,t)
=
h c(R,t).
\tag{TRISO-DIS-402}
\label{eq:triso-dis-402}
\end{equation}

The left side is the outward diffusive flux from the particle. The right side is the outward external mass-transfer flux.

\subsubsection*{16.1 Introduce the outer ghost point}

The last physical node is

\begin{equation}
r_N=R.
\tag{TRISO-DIS-403}
\label{eq:triso-dis-403}
\end{equation}

The adjacent physical interior node is

\begin{equation}
r_{N-1}=R-\Delta r.
\tag{TRISO-DIS-404}
\label{eq:triso-dis-404}
\end{equation}

Introduce a mathematical ghost point outside the particle:

\begin{equation}
r_{N+1}=R+\Delta r.
\tag{TRISO-DIS-405}
\label{eq:triso-dis-405}
\end{equation}

Its numerical value is denoted

\begin{equation}
C_{N+1}^j.
\tag{TRISO-DIS-406}
\label{eq:triso-dis-406}
\end{equation}

The ghost value is not an external physical concentration. It is an algebraic device used to retain a centred derivative at $r=R$.

\subsubsection*{16.2 Centred approximation of the surface gradient}

At time $t_j$, approximate the radial derivative by

\begin{equation}
\frac{\partial c}{\partial r}(R,t_j)
\approx
\frac{
C_{N+1}^j-C_{N-1}^j
}{
2\Delta r
}.
\tag{TRISO-DIS-407}
\label{eq:triso-dis-407}
\end{equation}

Substitute this approximation into the Robin condition:

\begin{equation}
-D
\frac{
C_{N+1}^j-C_{N-1}^j
}{
2\Delta r
}
=
hC_N^j.
\tag{TRISO-DIS-408}
\label{eq:triso-dis-408}
\end{equation}

Multiply both sides by $2\Delta r$:

\begin{equation}
-D
\left(
C_{N+1}^j-C_{N-1}^j
\right)
=
2h\Delta r\,C_N^j.
\tag{TRISO-DIS-409}
\label{eq:triso-dis-409}
\end{equation}

Divide by $-D$:

\begin{equation}
C_{N+1}^j-C_{N-1}^j
=
-\frac{2h\Delta r}{D}C_N^j.
\tag{TRISO-DIS-410}
\label{eq:triso-dis-410}
\end{equation}

Add $C_{N-1}^j$ to both sides:

\begin{equation}
C_{N+1}^j
=
C_{N-1}^j
-
\frac{2h\Delta r}{D}C_N^j.
\tag{TRISO-DIS-411}
\label{eq:triso-dis-411}
\end{equation}

Define the dimensionless mesh transfer parameter

\begin{equation}
\boxed{
\kappa
=
\frac{h\Delta r}{D}.
}
\tag{TRISO-DIS-412}
\label{eq:triso-dis-412}
\end{equation}

Its units are

\begin{equation}
[\kappa]
=
\frac{
\mathrm{m\,s^{-1}}\mathrm m
}{
\mathrm{m^2\,s^{-1}}
}
=
1.
\tag{TRISO-DIS-413}
\label{eq:triso-dis-413}
\end{equation}

Therefore the ghost relation is

\begin{equation}
\boxed{
C_{N+1}^j
=
C_{N-1}^j
-
2\kappa C_N^j.
}
\tag{TRISO-DIS-414}
\label{eq:triso-dis-414}
\end{equation}

\subsubsection*{16.3 Surface approximation of the second radial derivative}

Use the centred second derivative at node $N$:

\begin{equation}
\frac{\partial^2c}{\partial r^2}(R,t_j)
\approx
\frac{
C_{N-1}^j
-
2C_N^j
+
C_{N+1}^j
}{
\Delta r^2
}.
\tag{TRISO-DIS-415}
\label{eq:triso-dis-415}
\end{equation}

Substitute the ghost relation (TRISO-DIS-414):

\begin{equation}
\frac{\partial^2c}{\partial r^2}(R,t_j)
\approx
\frac{
C_{N-1}^j
-
2C_N^j
+
C_{N-1}^j
-
2\kappa C_N^j
}{
\Delta r^2
}.
\tag{TRISO-DIS-416}
\label{eq:triso-dis-416}
\end{equation}

Collect the two interior-neighbour terms:

\begin{equation}
C_{N-1}^j+C_{N-1}^j
=
2C_{N-1}^j.
\tag{TRISO-DIS-417}
\label{eq:triso-dis-417}
\end{equation}

Collect the two surface terms:

\begin{equation}
-2C_N^j-2\kappa C_N^j
=
-2(1+\kappa)C_N^j.
\tag{TRISO-DIS-418}
\label{eq:triso-dis-418}
\end{equation}

Therefore

\begin{equation}
\boxed{
\frac{\partial^2c}{\partial r^2}(R,t_j)
\approx
\frac{
2C_{N-1}^j
-
2(1+\kappa)C_N^j
}{
\Delta r^2
}.
}
\tag{TRISO-DIS-419}
\label{eq:triso-dis-419}
\end{equation}

\subsubsection*{16.4 Surface approximation of the first radial derivative}

The centred derivative is

\begin{equation}
\frac{\partial c}{\partial r}(R,t_j)
\approx
\frac{
C_{N+1}^j-C_{N-1}^j
}{
2\Delta r
}.
\tag{TRISO-DIS-420}
\label{eq:triso-dis-420}
\end{equation}

Substitute (TRISO-DIS-414):

\begin{equation}
\frac{\partial c}{\partial r}(R,t_j)
\approx
\frac{
C_{N-1}^j
-
2\kappa C_N^j
-
C_{N-1}^j
}{
2\Delta r
}.
\tag{TRISO-DIS-421}
\label{eq:triso-dis-421}
\end{equation}

Cancel the two $C_{N-1}^j$ terms:

\begin{equation}
\frac{\partial c}{\partial r}(R,t_j)
\approx
-\frac{
2\kappa C_N^j
}{
2\Delta r
}.
\tag{TRISO-DIS-422}
\label{eq:triso-dis-422}
\end{equation}

Cancel the factor $2$:

\begin{equation}
\boxed{
\frac{\partial c}{\partial r}(R,t_j)
\approx
-\frac{\kappa}{\Delta r}C_N^j.
}
\tag{TRISO-DIS-423}
\label{eq:triso-dis-423}
\end{equation}

Using $\kappa=h\Delta r/D$,

\begin{equation}
-\frac{\kappa}{\Delta r}C_N^j
=
-\frac hD C_N^j.
\tag{TRISO-DIS-424}
\label{eq:triso-dis-424}
\end{equation}

Thus the eliminated ghost relation reproduces the discrete Robin gradient exactly within the chosen centred boundary approximation.

\subsubsection*{16.5 Discrete spherical operator at the outer surface}

The spherical operator is

\begin{equation}
\mathcal L[c](R,t)
=
c_{rr}(R,t)
+
\frac2R c_r(R,t).
\tag{TRISO-DIS-425}
\label{eq:triso-dis-425}
\end{equation}

Substitute the discrete second derivative (TRISO-DIS-419):

\begin{equation}
\mathcal L[c](R,t_j)
\approx
\frac{
2C_{N-1}^j
-
2(1+\kappa)C_N^j
}{
\Delta r^2
}
+
\frac2R c_r(R,t_j).
\tag{TRISO-DIS-426}
\label{eq:triso-dis-426}
\end{equation}

Substitute the discrete first derivative (TRISO-DIS-423):

\begin{equation}
\mathcal L[c](R,t_j)
\approx
\frac{
2C_{N-1}^j
-
2(1+\kappa)C_N^j
}{
\Delta r^2
}
-
\frac{2\kappa}{R\Delta r}C_N^j.
\tag{TRISO-DIS-427}
\label{eq:triso-dis-427}
\end{equation}

Use

\begin{equation}
R=N\Delta r.
\tag{TRISO-DIS-428}
\label{eq:triso-dis-428}
\end{equation}

Therefore

\begin{equation}
R\Delta r
=
N\Delta r^2.
\tag{TRISO-DIS-429}
\label{eq:triso-dis-429}
\end{equation}

Hence

\begin{equation}
\frac{2\kappa}{R\Delta r}
=
\frac{2\kappa}{N\Delta r^2}.
\tag{TRISO-DIS-430}
\label{eq:triso-dis-430}
\end{equation}

Substitute:

\begin{equation}
\mathcal L[c](R,t_j)
\approx
\frac{
2C_{N-1}^j
-
2(1+\kappa)C_N^j
}{
\Delta r^2
}
-
\frac{
2\kappa C_N^j
}{
N\Delta r^2
}.
\tag{TRISO-DIS-431}
\label{eq:triso-dis-431}
\end{equation}

Put the terms over the common denominator:

\begin{equation}
\mathcal L[c](R,t_j)
\approx
\frac2{\Delta r^2}
\left[
C_{N-1}^j
-
(1+\kappa)C_N^j
-
\frac{\kappa}{N}C_N^j
\right].
\tag{TRISO-DIS-432}
\label{eq:triso-dis-432}
\end{equation}

Collect the surface coefficient:

\begin{equation}
(1+\kappa)
+
\frac{\kappa}{N}
=
1+\kappa
\left(
1+\frac1N
\right).
\tag{TRISO-DIS-433}
\label{eq:triso-dis-433}
\end{equation}

Therefore

\begin{equation}
\boxed{
\mathcal L[c](R,t_j)
\approx
\frac2{\Delta r^2}
\left[
C_{N-1}^j
-
\left(
1+\kappa\left(1+\frac1N\right)
\right)
C_N^j
\right].
}
\tag{TRISO-DIS-434}
\label{eq:triso-dis-434}
\end{equation}

\subsubsection*{16.6 Apply the surface PDE}

For the homogeneous benchmark,

\begin{equation}
\frac{\partial c}{\partial t}(R,t)
=
D\mathcal L[c](R,t)
+
S_R.
\tag{TRISO-DIS-435}
\label{eq:triso-dis-435}
\end{equation}

For the original homogeneous source benchmark,

\begin{equation}
S_R=S_0.
\tag{TRISO-DIS-436}
\label{eq:triso-dis-436}
\end{equation}

For the physical five-layer kernel-confined source problem, the OPyC source would instead be zero. This distinction must be preserved when the boundary formula is reused outside the Part-I benchmark.

Use forward Euler:

\begin{equation}
\frac{
C_N^{j+1}-C_N^j
}{
\Delta t
}
=
D\mathcal L[c](R,t_j)
+
S_R^j.
\tag{TRISO-DIS-437}
\label{eq:triso-dis-437}
\end{equation}

Substitute (TRISO-DIS-434):

\begin{equation}
\frac{
C_N^{j+1}-C_N^j
}{
\Delta t
}
=
\frac{2D}{\Delta r^2}
\left[
C_{N-1}^j
-
\left(
1+\kappa\left(1+\frac1N\right)
\right)
C_N^j
\right]
+
S_R^j.
\tag{TRISO-DIS-438}
\label{eq:triso-dis-438}
\end{equation}

Multiply by $\Delta t$:

\begin{equation}
C_N^{j+1}-C_N^j
=
\frac{2D\Delta t}{\Delta r^2}
\left[
C_{N-1}^j
-
\left(
1+\kappa\left(1+\frac1N\right)
\right)
C_N^j
\right]
+
S_R^j\Delta t.
\tag{TRISO-DIS-439}
\label{eq:triso-dis-439}
\end{equation}

Use

\begin{equation}
\mathrm{Fo}
=
\frac{D\Delta t}{\Delta r^2}.
\tag{TRISO-DIS-440}
\label{eq:triso-dis-440}
\end{equation}

Then

\begin{equation}
C_N^{j+1}-C_N^j
=
2\mathrm{Fo}
\left[
C_{N-1}^j
-
\left(
1+\kappa\left(1+\frac1N\right)
\right)
C_N^j
\right]
+
S_R^j\Delta t.
\tag{TRISO-DIS-441}
\label{eq:triso-dis-441}
\end{equation}

Add $C_N^j$ to both sides:

\begin{equation}
C_N^{j+1}
=
C_N^j
+
2\mathrm{Fo}C_{N-1}^j
-
2\mathrm{Fo}
\left(
1+\kappa\left(1+\frac1N\right)
\right)
C_N^j
+
S_R^j\Delta t.
\tag{TRISO-DIS-442}
\label{eq:triso-dis-442}
\end{equation}

Collect the surface coefficient:

\begin{equation}
\boxed{
C_N^{j+1}
=
2\mathrm{Fo}C_{N-1}^j
+
\left[
1
-
2\mathrm{Fo}
\left(
1+\kappa\left(1+\frac1N\right)
\right)
\right]
C_N^j
+
S_R^j\Delta t.
}
\tag{TRISO-DIS-443}
\label{eq:triso-dis-443}
\end{equation}

This recovers the surface update written in the original Ray derivation.

\subsubsection*{16.7 Dimensional checks}

The Fourier number is dimensionless:

\begin{equation}
[\mathrm{Fo}]=1.
\tag{TRISO-DIS-444}
\label{eq:triso-dis-444}
\end{equation}

The mesh transfer parameter is dimensionless:

\begin{equation}
[\kappa]=1.
\tag{TRISO-DIS-445}
\label{eq:triso-dis-445}
\end{equation}

Therefore every coefficient multiplying a concentration in (TRISO-DIS-443) is dimensionless.

The source increment has units

\begin{equation}
[S_R\Delta t]
=
\mathrm{mol\,m^{-3}}.
\tag{TRISO-DIS-446}
\label{eq:triso-dis-446}
\end{equation}

Thus every term in the update has concentration units.

\subsubsection*{16.8 Limiting checks}

If

\begin{equation}
h=0,
\tag{TRISO-DIS-447}
\label{eq:triso-dis-447}
\end{equation}

then

\begin{equation}
\kappa=0.
\tag{TRISO-DIS-448}
\label{eq:triso-dis-448}
\end{equation}

The ghost relation becomes

\begin{equation}
C_{N+1}^j=C_{N-1}^j.
\tag{TRISO-DIS-449}
\label{eq:triso-dis-449}
\end{equation}

This is the expected symmetric zero-gradient ghost condition for a zero-flux Neumann boundary.

The surface update becomes

\begin{equation}
C_N^{j+1}
=
2\mathrm{Fo}C_{N-1}^j
+
(1-2\mathrm{Fo})C_N^j
+
S_R^j\Delta t.
\tag{TRISO-DIS-450}
\label{eq:triso-dis-450}
\end{equation}

For finite $h>0$, increasing $h$ increases $\kappa$, which strengthens the outward-transfer contribution in the surface coefficient.

The formal limit $h\to\infty$ is more delicate for this explicit ghost formulation because

\begin{equation}
\kappa=\frac{h\Delta r}{D}\to\infty
\tag{TRISO-DIS-451}
\label{eq:triso-dis-451}
\end{equation}

at fixed $\Delta r$.

The continuum Robin condition approaches the absorbing Dirichlet condition $c(R,t)=0$, but the explicit ghost update becomes increasingly stiff rather than automatically turning into a numerically well-conditioned Dirichlet update.

Therefore an absorbing Dirichlet boundary should be imposed directly when that is the intended numerical model, rather than obtained by taking $\kappa\to\infty$ in (TRISO-DIS-443).

\subsubsection*{16.9 Complete truncation error of the Robin ghost surface closure}

The centred Robin derivative by itself is second-order accurate, but the complete surface PDE closure also inserts the ghost value into a second-derivative stencil.

The complete boundary operator must therefore be expanded directly.

Define

\begin{equation}
\beta=\frac{h}{D}.
\tag{TRISO-DIS-452}
\label{eq:triso-dis-452}
\end{equation}

The exact Robin condition gives

\begin{equation}
c_r(R)=-\beta c(R).
\tag{TRISO-DIS-453}
\label{eq:triso-dis-453}
\end{equation}

The ghost construction is

\begin{equation}
c_g(R+\Delta r)
=
c(R-\Delta r)
-
2\beta\Delta r\,c(R).
\tag{TRISO-DIS-454}
\label{eq:triso-dis-454}
\end{equation}

Taylor-expand the exact interior value:

\begin{equation}
c(R-\Delta r)
=
c(R)
-
\Delta r\,c_r(R)
+
\frac{\Delta r^2}{2}c_{rr}(R)
-
\frac{\Delta r^3}{6}c_{rrr}(R)
+
O(\Delta r^4).
\tag{TRISO-DIS-455}
\label{eq:triso-dis-455}
\end{equation}

Substitute the Robin derivative $c_r(R)=-\beta c(R)$:

\begin{equation}
c(R-\Delta r)
=
c(R)
+
\beta\Delta r\,c(R)
+
\frac{\Delta r^2}{2}c_{rr}(R)
-
\frac{\Delta r^3}{6}c_{rrr}(R)
+
O(\Delta r^4).
\tag{TRISO-DIS-456}
\label{eq:triso-dis-456}
\end{equation}

Substitute this expansion into the ghost construction:

\begin{equation}
c_g(R+\Delta r)
=
c(R)
-
\beta\Delta r\,c(R)
+
\frac{\Delta r^2}{2}c_{rr}(R)
-
\frac{\Delta r^3}{6}c_{rrr}(R)
+
O(\Delta r^4).
\tag{TRISO-DIS-457}
\label{eq:triso-dis-457}
\end{equation}

The exact smooth continuation would be

\begin{equation}
c(R+\Delta r)
=
c(R)
+
\Delta r\,c_r(R)
+
\frac{\Delta r^2}{2}c_{rr}(R)
+
\frac{\Delta r^3}{6}c_{rrr}(R)
+
O(\Delta r^4).
\tag{TRISO-DIS-458}
\label{eq:triso-dis-458}
\end{equation}

Apply $c_r(R)=-\beta c(R)$:

\begin{equation}
c(R+\Delta r)
=
c(R)
-
\beta\Delta r\,c(R)
+
\frac{\Delta r^2}{2}c_{rr}(R)
+
\frac{\Delta r^3}{6}c_{rrr}(R)
+
O(\Delta r^4).
\tag{TRISO-DIS-459}
\label{eq:triso-dis-459}
\end{equation}

Subtract the exact continuation from the ghost continuation:

\begin{equation}
c_g(R+\Delta r)-c(R+\Delta r)
=
-\frac{\Delta r^3}{3}c_{rrr}(R)
+
O(\Delta r^4).
\tag{TRISO-DIS-460}
\label{eq:triso-dis-460}
\end{equation}

Thus the ghost value error is

\begin{equation}
\boxed{
c_g-c_{\mathrm{exact}}
=
O(\Delta r^3).
}
\tag{TRISO-DIS-461}
\label{eq:triso-dis-461}
\end{equation}

The ghost value enters the centred second derivative divided by $\Delta r^2$:

\begin{equation}
c_{rr}^{\,g}(R)
=
\frac{
c(R-\Delta r)-2c(R)+c_g(R+\Delta r)
}{
\Delta r^2
}.
\tag{TRISO-DIS-462}
\label{eq:triso-dis-462}
\end{equation}

Write the ghost value as

\begin{equation}
c_g(R+\Delta r)
=
c(R+\Delta r)
+
\varepsilon_g,
\tag{TRISO-DIS-463}
\label{eq:triso-dis-463}
\end{equation}

where

\begin{equation}
\varepsilon_g
=
-\frac{\Delta r^3}{3}c_{rrr}(R)
+
O(\Delta r^4).
\tag{TRISO-DIS-464}
\label{eq:triso-dis-464}
\end{equation}

Substitute:

\begin{equation}
c_{rr}^{\,g}(R)
=
\frac{
c(R-\Delta r)-2c(R)+c(R+\Delta r)
}{
\Delta r^2
}
+
\frac{\varepsilon_g}{\Delta r^2}.
\tag{TRISO-DIS-465}
\label{eq:triso-dis-465}
\end{equation}

The ordinary centred second derivative contributes

\begin{equation}
c_{rr}(R)+O(\Delta r^2).
\tag{TRISO-DIS-466}
\label{eq:triso-dis-466}
\end{equation}

The ghost-error contribution is

\begin{equation}
\frac{\varepsilon_g}{\Delta r^2}
=
-\frac{\Delta r}{3}c_{rrr}(R)
+
O(\Delta r^2).
\tag{TRISO-DIS-467}
\label{eq:triso-dis-467}
\end{equation}

Therefore

\begin{equation}
\boxed{
c_{rr}^{\,g}(R)
=
c_{rr}(R)
-
\frac{\Delta r}{3}c_{rrr}(R)
+
O(\Delta r^2).
}
\tag{TRISO-DIS-468}
\label{eq:triso-dis-468}
\end{equation}

The first-derivative Robin closure remains second-order, but the second-derivative part is generically first-order at the boundary.

Hence the complete spherical surface operator has generic local spatial truncation

\begin{equation}
\boxed{
\mathcal L_h[c](R)
=
\mathcal L[c](R)
+
O(\Delta r).
}
\tag{TRISO-DIS-469}
\label{eq:triso-dis-469}
\end{equation}

With forward Euler time stepping, the boundary local consistency is therefore generically

\begin{equation}
\boxed{
O(\Delta t)+O(\Delta r),
}
\tag{TRISO-DIS-470}
\label{eq:triso-dis-470}
\end{equation}

unless additional cancellation or superconvergence is demonstrated.

[CORRECTION] The previous claim of $O(\Delta t)$+$O(\Delta r^2)$ for the complete Robin surface update was overstated. Only the centred first-derivative approximation was second-order.

[OPEN REVIEW FINDING R2-D01] The analytical overclaim is corrected, but the independent reviewer requires an actual grid-refinement study before this finding is closed. No global FTCS convergence order is claimed here.

\subsubsection*{16.10 Scope}

Equation (TRISO-DIS-443) belongs to the homogeneous Part-I FTCS benchmark.

For the physical five-layer problem, the same derivational structure may be applied to the OPyC layer by replacing $D$ with $D_5$ and using the physically selected outer source and boundary parameters.

The formula must not be confused with the production WOS absorbing-boundary treatment.

No stability conclusion is drawn in this section.

\subsection*{17. FTCS monotonicity and stability analysis}

The original Ray notebook inferred a stability restriction from non-negative update coefficients.

That argument is useful, but the mathematical statement must be made precisely.

Non-negative coefficients with an appropriate row sum establish a monotonicity/maximum-principle style property for the homogeneous update. They are not automatically a necessary-and-sufficient characterization of every possible matrix stability notion.

This section first derives the coefficient restrictions and then states exactly what they prove.

\subsubsection*{17.1 Remove forcing for stability analysis}

The complete explicit update contains source terms.

To study propagation of perturbations, compare two numerical solutions subject to the same prescribed source.

Let their difference be

\begin{equation}
E_i^j
=
C_i^j-\widetilde C_i^j.
\tag{TRISO-DIS-500}
\label{eq:triso-dis-500}
\end{equation}

Because both solutions have the same additive source, subtraction cancels that source.

Thus the error/perturbation equation is homogeneous.

Stability of the linear time-marching operator can therefore be analysed from the source-free amplification step

\begin{equation}
\mathbf E^{j+1}
=
\mathbf A\mathbf E^j.
\tag{TRISO-DIS-501}
\label{eq:triso-dis-501}
\end{equation}

\subsubsection*{17.2 Interior-row coefficients}

For an interior homogeneous node, Section 13 gives

\begin{equation}
E_i^{j+1}
=
\mathrm{Fo}
\left(
1-\frac1i
\right)
E_{i-1}^j
+
(1-2\mathrm{Fo})E_i^j
+
\mathrm{Fo}
\left(
1+\frac1i
\right)
E_{i+1}^j.
\tag{TRISO-DIS-502}
\label{eq:triso-dis-502}
\end{equation}

Define

\begin{equation}
a_i
=
\mathrm{Fo}
\left(
1-\frac1i
\right),
\tag{TRISO-DIS-503}
\label{eq:triso-dis-503}
\end{equation}

\begin{equation}
b_i
=
1-2\mathrm{Fo},
\tag{TRISO-DIS-504}
\label{eq:triso-dis-504}
\end{equation}

and

\begin{equation}
d_i
=
\mathrm{Fo}
\left(
1+\frac1i
\right).
\tag{TRISO-DIS-505}
\label{eq:triso-dis-505}
\end{equation}

For $i\ge1$,

\begin{equation}
1-\frac1i\ge0.
\tag{TRISO-DIS-506}
\label{eq:triso-dis-506}
\end{equation}

Since $\mathrm{Fo}\ge0$,

\begin{equation}
a_i\ge0.
\tag{TRISO-DIS-507}
\label{eq:triso-dis-507}
\end{equation}

Similarly,

\begin{equation}
1+\frac1i>0,
\tag{TRISO-DIS-508}
\label{eq:triso-dis-508}
\end{equation}

so

\begin{equation}
d_i\ge0.
\tag{TRISO-DIS-509}
\label{eq:triso-dis-509}
\end{equation}

The central coefficient is non-negative when

\begin{equation}
1-2\mathrm{Fo}\ge0.
\tag{TRISO-DIS-510}
\label{eq:triso-dis-510}
\end{equation}

Rearrange:

\begin{equation}
2\mathrm{Fo}\le1.
\tag{TRISO-DIS-511}
\label{eq:triso-dis-511}
\end{equation}

Therefore

\begin{equation}
\boxed{
\mathrm{Fo}\le\frac12.
}
\tag{TRISO-DIS-512}
\label{eq:triso-dis-512}
\end{equation}

Now add the three interior coefficients:

\begin{equation}
a_i+b_i+d_i
=
\mathrm{Fo}\left(1-\frac1i\right)
+
1-2\mathrm{Fo}
+
\mathrm{Fo}\left(1+\frac1i\right).
\tag{TRISO-DIS-513}
\label{eq:triso-dis-513}
\end{equation}

Expand:

\begin{equation}
a_i+b_i+d_i
=
\mathrm{Fo}
-\frac{\mathrm{Fo}}i
+
1
-
2\mathrm{Fo}
+
\mathrm{Fo}
+
\frac{\mathrm{Fo}}i.
\tag{TRISO-DIS-514}
\label{eq:triso-dis-514}
\end{equation}

Cancel the radial terms:

\begin{equation}
-\frac{\mathrm{Fo}}i
+
\frac{\mathrm{Fo}}i
=
0.
\tag{TRISO-DIS-515}
\label{eq:triso-dis-515}
\end{equation}

Cancel the diffusion contributions:

\begin{equation}
\mathrm{Fo}-2\mathrm{Fo}+\mathrm{Fo}=0.
\tag{TRISO-DIS-516}
\label{eq:triso-dis-516}
\end{equation}

Hence

\begin{equation}
\boxed{
a_i+b_i+d_i=1.
}
\tag{TRISO-DIS-517}
\label{eq:triso-dis-517}
\end{equation}

When (TRISO-DIS-512) holds, an interior update is therefore a convex combination of the previous-time neighbouring values.

\subsubsection*{17.3 Centre-row coefficients}

Section 14 gives the homogeneous centre error update

\begin{equation}
E_0^{j+1}
=
(1-6\mathrm{Fo})E_0^j
+
6\mathrm{Fo}E_1^j.
\tag{TRISO-DIS-518}
\label{eq:triso-dis-518}
\end{equation}

The neighbour coefficient satisfies

\begin{equation}
6\mathrm{Fo}\ge0.
\tag{TRISO-DIS-519}
\label{eq:triso-dis-519}
\end{equation}

The centre coefficient is non-negative when

\begin{equation}
1-6\mathrm{Fo}\ge0.
\tag{TRISO-DIS-520}
\label{eq:triso-dis-520}
\end{equation}

Therefore

\begin{equation}
6\mathrm{Fo}\le1.
\tag{TRISO-DIS-521}
\label{eq:triso-dis-521}
\end{equation}

Hence

\begin{equation}
\boxed{
\mathrm{Fo}\le\frac16.
}
\tag{TRISO-DIS-522}
\label{eq:triso-dis-522}
\end{equation}

The centre-row sum is

\begin{equation}
(1-6\mathrm{Fo})+6\mathrm{Fo}.
\tag{TRISO-DIS-523}
\label{eq:triso-dis-523}
\end{equation}

Therefore

\begin{equation}
\boxed{
(1-6\mathrm{Fo})+6\mathrm{Fo}=1.
}
\tag{TRISO-DIS-524}
\label{eq:triso-dis-524}
\end{equation}

Under (TRISO-DIS-522), the centre row is also a convex combination.

\subsubsection*{17.4 Robin surface-row coefficients}

Section 16 gives the homogeneous surface error update

\begin{equation}
E_N^{j+1}
=
2\mathrm{Fo}E_{N-1}^j
+
\left[
1
-
2\mathrm{Fo}
\left(
1+\kappa\left(1+\frac1N\right)
\right)
\right]
E_N^j.
\tag{TRISO-DIS-525}
\label{eq:triso-dis-525}
\end{equation}

The interior-neighbour coefficient is

\begin{equation}
2\mathrm{Fo}\ge0.
\tag{TRISO-DIS-526}
\label{eq:triso-dis-526}
\end{equation}

The surface coefficient is non-negative when

\begin{equation}
1
-
2\mathrm{Fo}
\left(
1+\kappa\left(1+\frac1N\right)
\right)
\ge0.
\tag{TRISO-DIS-527}
\label{eq:triso-dis-527}
\end{equation}

Move the second term to the other side:

\begin{equation}
1
\ge
2\mathrm{Fo}
\left(
1+\kappa\left(1+\frac1N\right)
\right).
\tag{TRISO-DIS-528}
\label{eq:triso-dis-528}
\end{equation}

For $h\ge0$, $D>0$, and $\Delta r>0$,

\begin{equation}
\kappa\ge0.
\tag{TRISO-DIS-529}
\label{eq:triso-dis-529}
\end{equation}

Therefore the denominator below is positive.

Divide:

\begin{equation}
\boxed{
\mathrm{Fo}
\le
\frac{
1
}{
2\left[
1+\kappa\left(1+\frac1N\right)
\right]
}.
}
\tag{TRISO-DIS-530}
\label{eq:triso-dis-530}
\end{equation}

Now calculate the surface-row sum:

\begin{equation}
2\mathrm{Fo}
+
1
-
2\mathrm{Fo}
\left(
1+\kappa\left(1+\frac1N\right)
\right).
\tag{TRISO-DIS-531}
\label{eq:triso-dis-531}
\end{equation}

Expand the last term:

\begin{equation}
2\mathrm{Fo}
+
1
-
2\mathrm{Fo}
-
2\mathrm{Fo}\kappa
\left(
1+\frac1N
\right).
\tag{TRISO-DIS-532}
\label{eq:triso-dis-532}
\end{equation}

Cancel $2\mathrm{Fo}-2\mathrm{Fo}$:

\begin{equation}
\boxed{
\text{surface row sum}
=
1
-
2\mathrm{Fo}\kappa
\left(
1+\frac1N
\right).
}
\tag{TRISO-DIS-533}
\label{eq:triso-dis-533}
\end{equation}

For $\kappa\ge0$,

\begin{equation}
\text{surface row sum}\le1.
\tag{TRISO-DIS-534}
\label{eq:triso-dis-534}
\end{equation}

Under the non-negativity restriction (TRISO-DIS-530),

\begin{equation}
\text{surface row sum}\ge0.
\tag{TRISO-DIS-535}
\label{eq:triso-dis-535}
\end{equation}

Thus the Robin row is sub-convex: part of the previous concentration can leave through the external boundary.

\subsubsection*{17.5 Combined coefficient-non-negativity condition}

The three restrictions are

\begin{equation}
\mathrm{Fo}\le\frac12
\tag{TRISO-DIS-536}
\label{eq:triso-dis-536}
\end{equation}

for ordinary interior nodes,

\begin{equation}
\mathrm{Fo}\le\frac16
\tag{TRISO-DIS-537}
\label{eq:triso-dis-537}
\end{equation}

for the centre,

and

\begin{equation}
\mathrm{Fo}
\le
\frac{
1
}{
2\left[
1+\kappa\left(1+\frac1N\right)
\right]
}
\tag{TRISO-DIS-538}
\label{eq:triso-dis-538}
\end{equation}

for the Robin surface.

Because

\begin{equation}
\frac16<\frac12,
\tag{TRISO-DIS-539}
\label{eq:triso-dis-539}
\end{equation}

the interior bound is never the controlling restriction once the centre node is included.

Therefore a sufficient coefficient-non-negativity condition for this homogeneous benchmark is

\begin{equation}
\boxed{
\mathrm{Fo}
\le
\min
\left\{
\frac16,
\,
\frac{
1
}{
2\left[
1+\kappa\left(1+\frac1N\right)
\right]
}
\right\}.
}
\tag{TRISO-DIS-540}
\label{eq:triso-dis-540}
\end{equation}

This is the corrected form of the original notebook restriction.

\subsubsection*{17.6 Amplification matrix}

Collect the nodal errors into

\begin{equation}
\mathbf E^j
=
(E_0^j,E_1^j,\ldots,E_N^j)^T.
\tag{TRISO-DIS-541}
\label{eq:triso-dis-541}
\end{equation}

The homogeneous update is

\begin{equation}
\boxed{
\mathbf E^{j+1}
=
\mathbf A\mathbf E^j.
}
\tag{TRISO-DIS-542}
\label{eq:triso-dis-542}
\end{equation}

The matrix has size

\begin{equation}
\mathbf A\in\mathbb R^{(N+1)\times(N+1)}.
\tag{TRISO-DIS-543}
\label{eq:triso-dis-543}
\end{equation}

The centre row contains

\begin{equation}
A_{0,0}=1-6\mathrm{Fo},
\tag{TRISO-DIS-544}
\label{eq:triso-dis-544}
\end{equation}

and

\begin{equation}
A_{0,1}=6\mathrm{Fo}.
\tag{TRISO-DIS-545}
\label{eq:triso-dis-545}
\end{equation}

For an ordinary interior row $i$,

\begin{equation}
A_{i,i-1}
=
\mathrm{Fo}
\left(
1-\frac1i
\right),
\tag{TRISO-DIS-546}
\label{eq:triso-dis-546}
\end{equation}

\begin{equation}
A_{i,i}
=
1-2\mathrm{Fo},
\tag{TRISO-DIS-547}
\label{eq:triso-dis-547}
\end{equation}

and

\begin{equation}
A_{i,i+1}
=
\mathrm{Fo}
\left(
1+\frac1i
\right).
\tag{TRISO-DIS-548}
\label{eq:triso-dis-548}
\end{equation}

The Robin surface row contains

\begin{equation}
A_{N,N-1}=2\mathrm{Fo},
\tag{TRISO-DIS-549}
\label{eq:triso-dis-549}
\end{equation}

and

\begin{equation}
A_{N,N}
=
1
-
2\mathrm{Fo}
\left(
1+\kappa\left(1+\frac1N\right)
\right).
\tag{TRISO-DIS-550}
\label{eq:triso-dis-550}
\end{equation}

All other entries are zero for this homogeneous tridiagonal benchmark.

\subsubsection*{17.7 $\ell_\infty$ stability under the monotonicity restriction}

The induced infinity norm of a matrix is

\begin{equation}
\|\mathbf A\|_\infty
=
\max_i
\sum_j
|A_{ij}|.
\tag{TRISO-DIS-551}
\label{eq:triso-dis-551}
\end{equation}

Under (TRISO-DIS-540), every non-zero entry of $\mathbf A$ is non-negative.

Therefore

\begin{equation}
|A_{ij}|=A_{ij}.
\tag{TRISO-DIS-552}
\label{eq:triso-dis-552}
\end{equation}

For the centre row, the row sum is exactly

\begin{equation}
1.
\tag{TRISO-DIS-553}
\label{eq:triso-dis-553}
\end{equation}

For each ordinary interior row, the row sum is exactly

\begin{equation}
1.
\tag{TRISO-DIS-554}
\label{eq:triso-dis-554}
\end{equation}

For the Robin row, the row sum is at most

\begin{equation}
1.
\tag{TRISO-DIS-555}
\label{eq:triso-dis-555}
\end{equation}

Therefore

\begin{equation}
\boxed{
\|\mathbf A\|_\infty\le1.
}
\tag{TRISO-DIS-556}
\label{eq:triso-dis-556}
\end{equation}

Apply the matrix norm inequality:

\begin{equation}
\|\mathbf E^{j+1}\|_\infty
=
\|\mathbf A\mathbf E^j\|_\infty
\le
\|\mathbf A\|_\infty
\|\mathbf E^j\|_\infty.
\tag{TRISO-DIS-557}
\label{eq:triso-dis-557}
\end{equation}

Use (TRISO-DIS-556):

\begin{equation}
\|\mathbf E^{j+1}\|_\infty
\le
\|\mathbf E^j\|_\infty.
\tag{TRISO-DIS-558}
\label{eq:triso-dis-558}
\end{equation}

Repeat the inequality over $j$ steps:

\begin{equation}
\boxed{
\|\mathbf E^j\|_\infty
\le
\|\mathbf E^0\|_\infty.
}
\tag{TRISO-DIS-559}
\label{eq:triso-dis-559}
\end{equation}

Thus the coefficient restriction (TRISO-DIS-540) is not merely a heuristic: for this assembled homogeneous benchmark it is a sufficient condition for non-amplification in the discrete $\ell_\infty$ norm.

\subsubsection*{17.8 Positivity and discrete maximum-principle interpretation}

Suppose

\begin{equation}
E_i^j\ge0
\tag{TRISO-DIS-560}
\label{eq:triso-dis-560}
\end{equation}

for every node.

Under (TRISO-DIS-540), every amplification coefficient is non-negative.

Therefore every component of

\begin{equation}
\mathbf E^{j+1}
=
\mathbf A\mathbf E^j
\tag{TRISO-DIS-561}
\label{eq:triso-dis-561}
\end{equation}

is also non-negative.

Hence the homogeneous update preserves non-negativity.

For an interior or centre row whose coefficients sum to one, the new value lies between the minimum and maximum of the contributing old values.

At the Robin boundary, the row sum is less than or equal to one because concentration can leave the domain.

This is the precise monotonicity/maximum-principle content of the notebook's coefficient argument.

\subsubsection*{17.9 Spectral-radius consequence}

For every square matrix,

\begin{equation}
\rho(\mathbf A)
\le
\|\mathbf A\|
\tag{TRISO-DIS-562}
\label{eq:triso-dis-562}
\end{equation}

for any induced matrix norm.

Using the infinity norm,

\begin{equation}
\rho(\mathbf A)
\le
\|\mathbf A\|_\infty.
\tag{TRISO-DIS-563}
\label{eq:triso-dis-563}
\end{equation}

Under (TRISO-DIS-540),

\begin{equation}
\|\mathbf A\|_\infty\le1.
\tag{TRISO-DIS-564}
\label{eq:triso-dis-564}
\end{equation}

Therefore

\begin{equation}
\boxed{
\rho(\mathbf A)\le1.
}
\tag{TRISO-DIS-565}
\label{eq:triso-dis-565}
\end{equation}

So the monotonicity restriction also provides a sufficient spectral-radius bound for this homogeneous assembled amplification matrix.

This does \textbf{not} prove that (TRISO-DIS-540) is necessary for spectral stability.

There may be parameter values with some negative coefficients for which

\begin{equation}
\rho(\mathbf A)\le1.
\tag{TRISO-DIS-566}
\label{eq:triso-dis-566}
\end{equation}

Determining the exact necessary-and-sufficient spectral stability region would require analysis of the eigenvalues of the specific amplification matrix.

\subsubsection*{17.10 Relation to the standard explicit-Euler eigenvalue condition}

Write a semi-discrete diffusion system abstractly as

\begin{equation}
\frac{d\mathbf C}{dt}
=
\mathbf L\mathbf C.
\tag{TRISO-DIS-567}
\label{eq:triso-dis-567}
\end{equation}

Forward Euler gives

\begin{equation}
\mathbf C^{j+1}
=
\left(
\mathbf I+\Delta t\,\mathbf L
\right)
\mathbf C^j.
\tag{TRISO-DIS-568}
\label{eq:triso-dis-568}
\end{equation}

Therefore

\begin{equation}
\mathbf A
=
\mathbf I+\Delta t\,\mathbf L.
\tag{TRISO-DIS-569}
\label{eq:triso-dis-569}
\end{equation}

If $\lambda_\ell$ is an eigenvalue of $\mathbf L$, then the corresponding amplification eigenvalue is

\begin{equation}
g_\ell
=
1+\Delta t\,\lambda_\ell.
\tag{TRISO-DIS-570}
\label{eq:triso-dis-570}
\end{equation}

For a real non-positive diffusion eigenvalue,

\begin{equation}
\lambda_\ell\le0,
\tag{TRISO-DIS-571}
\label{eq:triso-dis-571}
\end{equation}

the scalar forward-Euler stability requirement is

\begin{equation}
|1+\Delta t\,\lambda_\ell|\le1.
\tag{TRISO-DIS-572}
\label{eq:triso-dis-572}
\end{equation}

For real $\lambda_\ell\le0$, this is equivalent to

\begin{equation}
-1
\le
1+\Delta t\,\lambda_\ell
\le
1.
\tag{TRISO-DIS-573}
\label{eq:triso-dis-573}
\end{equation}

Subtract $1$:

\begin{equation}
-2
\le
\Delta t\,\lambda_\ell
\le
0.
\tag{TRISO-DIS-574}
\label{eq:triso-dis-574}
\end{equation}

Because $\lambda_\ell<0$ for a decaying mode, the lower inequality gives

\begin{equation}
\Delta t
\le
\frac{2}{|\lambda_\ell|}.
\tag{TRISO-DIS-575}
\label{eq:triso-dis-575}
\end{equation}

For all modes,

\begin{equation}
\boxed{
\Delta t
\le
\frac{2}{
\max_\ell|\lambda_\ell|
}
}
\tag{TRISO-DIS-576}
\label{eq:triso-dis-576}
\end{equation}

would be the exact scalar forward-Euler restriction if the relevant semi-discrete operator has a real non-positive spectrum and is diagonalizable in the norm under consideration.

The present section does not compute $\max|\lambda_\ell|$ for the spherical matrix.

Therefore (TRISO-DIS-576) is a framework for a sharper spectral analysis, not a completed numerical bound for this benchmark.

\subsubsection*{17.11 Important limitation: material interfaces}

The amplification matrix in Sections 17.6–17.10 corresponds to the homogeneous benchmark using:

\noindent\textbullet\ the centre row from Section 14;\\
\noindent\textbullet\ homogeneous interior rows from Section 13;\\
\noindent\textbullet\ the Robin surface row from Section 16.\\

The full five-layer problem contains discontinuous diffusivities and the interface transmission treatment from Section 15.

Its assembled operator is therefore different.

The bound

\begin{equation}
\mathrm{Fo}
\le
\min
\left\{
\frac16,
\frac1{
2[1+\kappa(1+1/N)]
}
\right\}
\tag{TRISO-DIS-577}
\label{eq:triso-dis-577}
\end{equation}

must \textbf{not} be presented as a proved stability bound for an arbitrary five-layer discretisation.

A five-layer stability condition depends on the final chosen conservative spatial discretisation, the layer-specific diffusivities, mesh spacings, and interface treatment.

\subsubsection*{17.12 Correct status of the original notebook claim}

The original notebook's coefficient argument is therefore classified as follows.

\textbf{[VERIFIED]} Interior coefficient non-negativity requires

\begin{equation}
\mathrm{Fo}\le\frac12.
\tag{TRISO-DIS-578}
\label{eq:triso-dis-578}
\end{equation}

\textbf{[VERIFIED]} Centre coefficient non-negativity requires

\begin{equation}
\mathrm{Fo}\le\frac16.
\tag{TRISO-DIS-579}
\label{eq:triso-dis-579}
\end{equation}

\textbf{[VERIFIED]} Robin-surface coefficient non-negativity requires

\begin{equation}
\mathrm{Fo}
\le
\frac1{
2[1+\kappa(1+1/N)]
}.
\tag{TRISO-DIS-580}
\label{eq:triso-dis-580}
\end{equation}

\textbf{[VERIFIED]} Their minimum is a sufficient monotonicity condition for the homogeneous benchmark.

\textbf{[VERIFIED]} Under that same condition, the assembled homogeneous amplification matrix satisfies

\begin{equation}
\|\mathbf A\|_\infty\le1
\tag{TRISO-DIS-581}
\label{eq:triso-dis-581}
\end{equation}

and consequently

\begin{equation}
\rho(\mathbf A)\le1.
\tag{TRISO-DIS-582}
\label{eq:triso-dis-582}
\end{equation}

[NOT PROVED] The condition is necessary for spectral stability.

[NOT APPLICABLE WITHOUT RE-DERIVATION] The same bound is the exact stability condition for the final discontinuous-$D$, five-layer discretisation.

\subsection*{18. Conservative finite-volume discretisation of the five-layer PDE}

The original FTCS derivation is a useful benchmark, but the actual five-layer PDE contains discontinuous material diffusivities.

A conservative finite-volume formulation is therefore derived directly from

\begin{equation}
\frac{\partial c}{\partial t}
=
\frac1{r^2}
\frac{\partial}{\partial r}
\left(
r^2D\frac{\partial c}{\partial r}
\right)
+
S.
\tag{TRISO-FV-100}
\label{eq:triso-fv-100}
\end{equation}

This choice is a mathematical discretisation of the verified continuum model. It does not replace or redefine the production WOS algorithm.

\subsubsection*{18.1 Spherical control-volume geometry}

Let cell $P$ occupy

\begin{equation}
r_{P-\frac12}
<
r
<
r_{P+\frac12}.
\tag{TRISO-FV-101}
\label{eq:triso-fv-101}
\end{equation}

Its west face is

\begin{equation}
r_w=r_{P-\frac12},
\tag{TRISO-FV-102}
\label{eq:triso-fv-102}
\end{equation}

and its east face is

\begin{equation}
r_e=r_{P+\frac12}.
\tag{TRISO-FV-103}
\label{eq:triso-fv-103}
\end{equation}

The corresponding spherical face areas are

\begin{equation}
\boxed{
A_w=4\pi r_w^2
}
\tag{TRISO-FV-104}
\label{eq:triso-fv-104}
\end{equation}

and

\begin{equation}
\boxed{
A_e=4\pi r_e^2.
}
\tag{TRISO-FV-105}
\label{eq:triso-fv-105}
\end{equation}

The exact cell volume is

\begin{equation}
V_P
=
\int_{r_w}^{r_e}4\pi r^2\,dr.
\tag{TRISO-FV-106}
\label{eq:triso-fv-106}
\end{equation}

Evaluate the integral:

\begin{equation}
V_P
=
4\pi
\left[
\frac{r^3}{3}
\right]_{r_w}^{r_e}.
\tag{TRISO-FV-107}
\label{eq:triso-fv-107}
\end{equation}

Therefore

\begin{equation}
\boxed{
V_P
=
\frac{4\pi}{3}
\left(
r_e^3-r_w^3
\right).
}
\tag{TRISO-FV-108}
\label{eq:triso-fv-108}
\end{equation}

\subsubsection*{18.1.1 Frozen finite-volume unknown and representative coordinate}

The canonical finite-volume unknown remains the \textbf{exact spherical cell average} defined later in (TRISO-FV-118). It is not redefined as a point value.

For geometry and two-point flux reconstruction, assign each cell a representative radial coordinate equal to its spherical volume centroid:

\begin{equation}
\boxed{
r_P
=
\frac{
\displaystyle
\int_{r_w}^{r_e}
r\,4\pi r^2\,dr
}{
\displaystyle
\int_{r_w}^{r_e}
4\pi r^2\,dr
}.
}
\tag{TRISO-FV-109A}
\label{eq:triso-fv-109a}
\end{equation}

Evaluate the numerator:

\begin{equation}
\int_{r_w}^{r_e}
4\pi r^3\,dr
=
\pi
\left(
r_e^4-r_w^4
\right).
\tag{TRISO-FV-109B}
\label{eq:triso-fv-109b}
\end{equation}

Use the exact volume (TRISO-FV-108):

\begin{equation}
V_P
=
\frac{4\pi}{3}
\left(
r_e^3-r_w^3
\right).
\tag{TRISO-FV-109C}
\label{eq:triso-fv-109c}
\end{equation}

Therefore

\begin{equation}
\boxed{
r_P
=
\frac34
\frac{
r_e^4-r_w^4
}{
r_e^3-r_w^3
}.
}
\tag{TRISO-FV-109D}
\label{eq:triso-fv-109d}
\end{equation}

The canonical representation is therefore:

\noindent\textbullet\ $C_P$: exact spherical volume average over cell $P$;\\
\noindent\textbullet\ $r_P$: spherical volume-centroid coordinate used as the representative location of that average in two-point reconstruction;\\
\noindent\textbullet\ $r_{P+1/2}$: physical face coordinate;\\
\noindent\textbullet\ $r_{P+1}-r_P$: distance between representative cell coordinates;\\
\noindent\textbullet\ $\delta r_P=r_{P+1/2}-r_P$: distance from cell $P$'s representative coordinate to its east face;\\
\noindent\textbullet\ $\delta r_{P+1}=r_{P+1}-r_{P+1/2}$: distance from the shared face to the neighbouring representative coordinate.\\

At a material interface, the physical interface is aligned with a face $r_{P+1/2}$.

At the outer boundary,

\begin{equation}
\delta r_R=R-r_{M-1}.
\tag{TRISO-FV-109E}
\label{eq:triso-fv-109e}
\end{equation}

\textbf{[IMPORTANT]} The face-gradient formulas treat the exact cell averages as reconstructed values located at their volume centroids. That reconstruction is an approximation whose spatial order must be established in the subsequent accuracy study; it is not part of the exact control-volume balance.

\subsubsection*{18.2 Integrate conservation over one spherical cell}

Multiply (TRISO-FV-100) by the spherical volume element

\begin{equation}
4\pi r^2\,dr.
\tag{TRISO-FV-109}
\label{eq:triso-fv-109}
\end{equation}

This gives

\begin{equation}
4\pi r^2
\frac{\partial c}{\partial t}\,dr
=
4\pi
\frac{\partial}{\partial r}
\left(
r^2D\frac{\partial c}{\partial r}
\right)dr
+
4\pi r^2S\,dr.
\tag{TRISO-FV-110}
\label{eq:triso-fv-110}
\end{equation}

Integrate from $r_w$ to $r_e$:

\begin{equation}
\int_{r_w}^{r_e}
4\pi r^2
\frac{\partial c}{\partial t}\,dr
=
\int_{r_w}^{r_e}
4\pi
\frac{\partial}{\partial r}
\left(
r^2D\frac{\partial c}{\partial r}
\right)dr
+
\int_{r_w}^{r_e}
4\pi r^2S\,dr.
\tag{TRISO-FV-111}
\label{eq:triso-fv-111}
\end{equation}

Evaluate the derivative integral:

\begin{equation}
\int_{r_w}^{r_e}
4\pi
\frac{\partial}{\partial r}
\left(
r^2D\frac{\partial c}{\partial r}
\right)dr
=
4\pi
\left[
r^2D\frac{\partial c}{\partial r}
\right]_{r_w}^{r_e}.
\tag{TRISO-FV-112}
\label{eq:triso-fv-112}
\end{equation}

Expand the boundary evaluation:

\begin{equation}
4\pi
\left[
r^2D\frac{\partial c}{\partial r}
\right]_{r_w}^{r_e}
=
A_eD_e
\left.\frac{\partial c}{\partial r}\right|_e
-
A_wD_w
\left.\frac{\partial c}{\partial r}\right|_w.
\tag{TRISO-FV-113}
\label{eq:triso-fv-113}
\end{equation}

Define the outward radial Fickian flux

\begin{equation}
J_r=-D\frac{\partial c}{\partial r}.
\tag{TRISO-FV-114}
\label{eq:triso-fv-114}
\end{equation}

Then

\begin{equation}
D\frac{\partial c}{\partial r}=-J_r.
\tag{TRISO-FV-115}
\label{eq:triso-fv-115}
\end{equation}

Substitute into (TRISO-FV-113):

\begin{equation}
A_eD_e c_r|_e-A_wD_wc_r|_w
=
-A_eJ_e+A_wJ_w.
\tag{TRISO-FV-116}
\label{eq:triso-fv-116}
\end{equation}

Thus the integrated conservation equation is

\begin{equation}
\int_{r_w}^{r_e}
4\pi r^2
\frac{\partial c}{\partial t}\,dr
=
A_wJ_w-A_eJ_e
+
\int_{r_w}^{r_e}4\pi r^2S\,dr.
\tag{TRISO-FV-117}
\label{eq:triso-fv-117}
\end{equation}

This is the discrete starting point: accumulation equals inward face flow minus outward face flow plus generation.

\subsubsection*{18.3 Cell-average unknown}

Define the volume-averaged concentration

\begin{equation}
\boxed{
C_P(t)
=
\frac1{V_P}
\int_{r_w}^{r_e}
c(r,t)\,4\pi r^2\,dr.
}
\tag{TRISO-FV-118}
\label{eq:triso-fv-118}
\end{equation}

Multiply by $V_P$:

\begin{equation}
V_PC_P
=
\int_{r_w}^{r_e}
c\,4\pi r^2\,dr.
\tag{TRISO-FV-119}
\label{eq:triso-fv-119}
\end{equation}

For a fixed mesh, $V_P$ is constant in time.

Differentiate:

\begin{equation}
V_P\frac{dC_P}{dt}
=
\int_{r_w}^{r_e}
4\pi r^2
\frac{\partial c}{\partial t}\,dr.
\tag{TRISO-FV-120}
\label{eq:triso-fv-120}
\end{equation}

Define the volume-averaged source

\begin{equation}
\boxed{
S_P
=
\frac1{V_P}
\int_{r_w}^{r_e}
S(r,t)\,4\pi r^2\,dr.
}
\tag{TRISO-FV-121}
\label{eq:triso-fv-121}
\end{equation}

Therefore

\begin{equation}
S_PV_P
=
\int_{r_w}^{r_e}
S(r,t)\,4\pi r^2\,dr.
\tag{TRISO-FV-122}
\label{eq:triso-fv-122}
\end{equation}

Substitute (TRISO-FV-120) and (TRISO-FV-122) into (TRISO-FV-117):

\begin{equation}
\boxed{
V_P\frac{dC_P}{dt}
=
A_wJ_w-A_eJ_e+S_PV_P.
}
\tag{TRISO-FV-123}
\label{eq:triso-fv-123}
\end{equation}

Equation (TRISO-FV-123) is an exact control-volume balance before face-flux approximation.

\subsubsection*{18.4 Interior face flux inside one material}

Let cells $P$ and $E$ share east face $e$.

Assume the face lies inside one material with constant diffusivity $D_e$.

Let the cell-centre distance be

\begin{equation}
\delta r_{PE}=r_E-r_P.
\tag{TRISO-FV-124}
\label{eq:triso-fv-124}
\end{equation}

Approximate the face gradient by

\begin{equation}
\left.
\frac{\partial c}{\partial r}
\right|_e
\approx
\frac{C_E-C_P}{\delta r_{PE}}.
\tag{TRISO-FV-125}
\label{eq:triso-fv-125}
\end{equation}

Fick's law gives

\begin{equation}
J_e
\approx
-D_e
\frac{C_E-C_P}{\delta r_{PE}}.
\tag{TRISO-FV-126}
\label{eq:triso-fv-126}
\end{equation}

Reverse the numerator:

\begin{equation}
\boxed{
J_e
\approx
D_e
\frac{C_P-C_E}{\delta r_{PE}}.
}
\tag{TRISO-FV-127}
\label{eq:triso-fv-127}
\end{equation}

Similarly, for west neighbour $W$,

\begin{equation}
\boxed{
J_w
\approx
D_w
\frac{C_W-C_P}{\delta r_{WP}}.
}
\tag{TRISO-FV-128}
\label{eq:triso-fv-128}
\end{equation}

The sign convention is consistent with $J_r>0$ meaning outward radial transport.

\subsubsection*{18.5 Face crossing a material interface}

Now let face $e$ coincide with an interface between two materials.

Let the distance from centre $P$ to the interface be

\begin{equation}
\delta r_P.
\tag{TRISO-FV-129}
\label{eq:triso-fv-129}
\end{equation}

Let the distance from the interface to centre $E$ be

\begin{equation}
\delta r_E.
\tag{TRISO-FV-130}
\label{eq:triso-fv-130}
\end{equation}

Let the diffusivities be

\begin{equation}
D_P
\tag{TRISO-FV-131}
\label{eq:triso-fv-131}
\end{equation}

and

\begin{equation}
D_E.
\tag{TRISO-FV-132}
\label{eq:triso-fv-132}
\end{equation}

Let the ideal interface concentration be $C_I$.

The flux from $P$ to the interface is

\begin{equation}
J_e
=
D_P
\frac{C_P-C_I}{\delta r_P}.
\tag{TRISO-FV-133}
\label{eq:triso-fv-133}
\end{equation}

The flux from the interface to $E$ is

\begin{equation}
J_e
=
D_E
\frac{C_I-C_E}{\delta r_E}.
\tag{TRISO-FV-134}
\label{eq:triso-fv-134}
\end{equation}

Solve the first equation for the concentration drop:

\begin{equation}
C_P-C_I
=
J_e\frac{\delta r_P}{D_P}.
\tag{TRISO-FV-135}
\label{eq:triso-fv-135}
\end{equation}

Solve the second equation for its drop:

\begin{equation}
C_I-C_E
=
J_e\frac{\delta r_E}{D_E}.
\tag{TRISO-FV-136}
\label{eq:triso-fv-136}
\end{equation}

Add the two equations:

\begin{equation}
C_P-C_E
=
J_e
\left(
\frac{\delta r_P}{D_P}
+
\frac{\delta r_E}{D_E}
\right).
\tag{TRISO-FV-137}
\label{eq:triso-fv-137}
\end{equation}

Solve for $J_e$:

\begin{equation}
\boxed{
J_e
=
\frac{
C_P-C_E
}{
\dfrac{\delta r_P}{D_P}
+
\dfrac{\delta r_E}{D_E}
}.
}
\tag{TRISO-FV-138}
\label{eq:triso-fv-138}
\end{equation}

Define the face conductance per unit area

\begin{equation}
\boxed{
g_e
=
\left(
\frac{\delta r_P}{D_P}
+
\frac{\delta r_E}{D_E}
\right)^{-1}.
}
\tag{TRISO-FV-139}
\label{eq:triso-fv-139}
\end{equation}

Then

\begin{equation}
\boxed{
J_e=g_e(C_P-C_E).
}
\tag{TRISO-FV-140}
\label{eq:triso-fv-140}
\end{equation}

The interface law is therefore generated by adding diffusion resistances, exactly as in the continuous steady resistance derivation.

\subsubsection*{18.6 Effective face diffusivity}

Define the centre-to-centre distance

\begin{equation}
\delta r_{PE}
=
\delta r_P+\delta r_E.
\tag{TRISO-FV-141}
\label{eq:triso-fv-141}
\end{equation}

Define $D_e^{\mathrm{eff}}$ by

\begin{equation}
J_e
=
D_e^{\mathrm{eff}}
\frac{C_P-C_E}{\delta r_{PE}}.
\tag{TRISO-FV-142}
\label{eq:triso-fv-142}
\end{equation}

Compare with (TRISO-FV-138):

\begin{equation}
\frac{D_e^{\mathrm{eff}}}{\delta r_{PE}}
=
\left(
\frac{\delta r_P}{D_P}
+
\frac{\delta r_E}{D_E}
\right)^{-1}.
\tag{TRISO-FV-143}
\label{eq:triso-fv-143}
\end{equation}

Multiply by $\delta r_{PE}$:

\begin{equation}
\boxed{
D_e^{\mathrm{eff}}
=
\frac{
\delta r_P+\delta r_E
}{
\dfrac{\delta r_P}{D_P}
+
\dfrac{\delta r_E}{D_E}
}.
}
\tag{TRISO-FV-144}
\label{eq:triso-fv-144}
\end{equation}

For equal half-cell distances,

\begin{equation}
\delta r_P=\delta r_E,
\tag{TRISO-FV-145}
\label{eq:triso-fv-145}
\end{equation}

this reduces to

\begin{equation}
\boxed{
D_e^{\mathrm{eff}}
=
\frac{2D_PD_E}{D_P+D_E}.
}
\tag{TRISO-FV-146}
\label{eq:triso-fv-146}
\end{equation}

Thus the harmonic mean arises naturally from the conservative face-flux derivation.

\subsubsection*{18.7 Semi-discrete conservative equation}

Define the west conductance

\begin{equation}
G_w
=
\frac{A_wD_w^{\mathrm{eff}}}{\delta r_{WP}},
\tag{TRISO-FV-147}
\label{eq:triso-fv-147}
\end{equation}

and the east conductance

\begin{equation}
G_e
=
\frac{A_eD_e^{\mathrm{eff}}}{\delta r_{PE}}.
\tag{TRISO-FV-148}
\label{eq:triso-fv-148}
\end{equation}

Their units are

\begin{equation}
[G_w]=[G_e]=\mathrm{m^3\,s^{-1}}.
\tag{TRISO-FV-149}
\label{eq:triso-fv-149}
\end{equation}

Using (TRISO-FV-127) and (TRISO-FV-128),

\begin{equation}
A_wJ_w
=
G_w(C_W-C_P),
\tag{TRISO-FV-150}
\label{eq:triso-fv-150}
\end{equation}

and

\begin{equation}
A_eJ_e
=
G_e(C_P-C_E).
\tag{TRISO-FV-151}
\label{eq:triso-fv-151}
\end{equation}

Substitute these into (TRISO-FV-123):

\begin{equation}
V_P\frac{dC_P}{dt}
=
G_w(C_W-C_P)
-
G_e(C_P-C_E)
+
S_PV_P.
\tag{TRISO-FV-152}
\label{eq:triso-fv-152}
\end{equation}

Expand:

\begin{equation}
V_P\frac{dC_P}{dt}
=
G_wC_W
-
G_wC_P
-
G_eC_P
+
G_eC_E
+
S_PV_P.
\tag{TRISO-FV-153}
\label{eq:triso-fv-153}
\end{equation}

Collect the central concentration:

\begin{equation}
\boxed{
V_P\frac{dC_P}{dt}
=
G_wC_W
-
(G_w+G_e)C_P
+
G_eC_E
+
S_PV_P.
}
\tag{TRISO-FV-154}
\label{eq:triso-fv-154}
\end{equation}

Divide by $V_P$:

\begin{equation}
\boxed{
\frac{dC_P}{dt}
=
\frac{G_w}{V_P}C_W
-
\frac{G_w+G_e}{V_P}C_P
+
\frac{G_e}{V_P}C_E
+
S_P.
}
\tag{TRISO-FV-155}
\label{eq:triso-fv-155}
\end{equation}

This is the conservative semi-discrete equation for an ordinary spherical control volume.

\subsubsection*{18.8 Explicit Euler time discretisation}

Apply forward Euler:

\begin{equation}
\frac{
C_P^{j+1}-C_P^j
}{
\Delta t
}
=
\frac{G_w}{V_P}C_W^j
-
\frac{G_w+G_e}{V_P}C_P^j
+
\frac{G_e}{V_P}C_E^j
+
S_P^j.
\tag{TRISO-FV-156}
\label{eq:triso-fv-156}
\end{equation}

Multiply by $\Delta t$:

\begin{equation}
C_P^{j+1}-C_P^j
=
\frac{\Delta t\,G_w}{V_P}C_W^j
-
\frac{\Delta t(G_w+G_e)}{V_P}C_P^j
+
\frac{\Delta t\,G_e}{V_P}C_E^j
+
S_P^j\Delta t.
\tag{TRISO-FV-157}
\label{eq:triso-fv-157}
\end{equation}

Add $C_P^j$:

\begin{equation}
\boxed{
C_P^{j+1}
=
\frac{\Delta t\,G_w}{V_P}C_W^j
+
\left[
1-
\frac{\Delta t(G_w+G_e)}{V_P}
\right]C_P^j
+
\frac{\Delta t\,G_e}{V_P}C_E^j
+
S_P^j\Delta t.
}
\tag{TRISO-FV-158}
\label{eq:triso-fv-158}
\end{equation}

This is an explicit conservative finite-volume update.

\subsubsection*{18.9 Exact discrete conservation over multiple cells}

Sum (TRISO-FV-123) over all control volumes $P=1,\ldots,M$:

\begin{equation}
\sum_{P=1}^{M}
V_P\frac{dC_P}{dt}
=
\sum_{P=1}^{M}
(A_wJ_w-A_eJ_e)
+
\sum_{P=1}^{M}S_PV_P.
\tag{TRISO-FV-159}
\label{eq:triso-fv-159}
\end{equation}

At a shared internal face, the east flux of one cell is the west flux of the next cell.

For example,

\begin{equation}
-A_{e,P}J_{e,P}
+
A_{w,P+1}J_{w,P+1}
=
0
\tag{TRISO-FV-160}
\label{eq:triso-fv-160}
\end{equation}

because

\begin{equation}
A_{e,P}=A_{w,P+1}
\tag{TRISO-FV-161}
\label{eq:triso-fv-161}
\end{equation}

and the same single face flux is used:

\begin{equation}
J_{e,P}=J_{w,P+1}.
\tag{TRISO-FV-162}
\label{eq:triso-fv-162}
\end{equation}

Therefore every internal-face contribution cancels pairwise.

Only the physical domain boundaries remain:

\begin{equation}
\boxed{
\frac{d}{dt}
\left(
\sum_PV_PC_P
\right)
=
\text{boundary inflow}
-
\text{boundary outflow}
+
\sum_PS_PV_P.
}
\tag{TRISO-FV-163}
\label{eq:triso-fv-163}
\end{equation}

This is the principal conservation advantage of the finite-volume construction.

\subsubsection*{18.10 Centre control volume}

For the central cell,

\begin{equation}
r_w=0.
\tag{TRISO-FV-164}
\label{eq:triso-fv-164}
\end{equation}

Therefore its west-face area is

\begin{equation}
A_w=4\pi(0)^2.
\tag{TRISO-FV-165}
\label{eq:triso-fv-165}
\end{equation}

Hence

\begin{equation}
\boxed{
A_w=0.
}
\tag{TRISO-FV-166}
\label{eq:triso-fv-166}
\end{equation}

The centre requires no artificial inward flux condition in the control-volume balance because the spherical face at $r=0$ has zero area.

The central balance becomes

\begin{equation}
\boxed{
V_0\frac{dC_0}{dt}
=
-A_eJ_e
+
S_0V_0.
}
\tag{TRISO-FV-167}
\label{eq:triso-fv-167}
\end{equation}

Using the outward face-flux convention,

\begin{equation}
J_e
=
G_e^{(A=1)}(C_0-C_1),
\tag{TRISO-FV-168}
\label{eq:triso-fv-168}
\end{equation}

where $G_e^{(A=1)}$ denotes the conductance per unit area.

Equivalently, using total conductance $G_e$,

\begin{equation}
\boxed{
V_0\frac{dC_0}{dt}
=
G_e(C_1-C_0)
+
S_0V_0.
}
\tag{TRISO-FV-169}
\label{eq:triso-fv-169}
\end{equation}

Thus centre regularity is built geometrically into the zero-area inner face.

\subsubsection*{18.11 Kernel-confined source}

For a cell entirely inside the kernel,

\begin{equation}
S_P=S_0
\tag{TRISO-FV-170}
\label{eq:triso-fv-170}
\end{equation}

for the constant-source benchmark.

For a cell entirely outside the kernel,

\begin{equation}
S_P=0.
\tag{TRISO-FV-171}
\label{eq:triso-fv-171}
\end{equation}

If a control-volume face is aligned with the kernel boundary $r_1$, no cell straddles the source discontinuity.

Then the discrete total generation is

\begin{equation}
\dot N_{\mathrm{gen}}^{\,h}
=
\sum_{P\in\mathrm{kernel}}
S_0V_P.
\tag{TRISO-FV-172}
\label{eq:triso-fv-172}
\end{equation}

Because the kernel control volumes exactly partition $0<r<r_1$,

\begin{equation}
\sum_{P\in\mathrm{kernel}}V_P
=
\frac{4\pi r_1^3}{3}.
\tag{TRISO-FV-173}
\label{eq:triso-fv-173}
\end{equation}

Therefore

\begin{equation}
\boxed{
\dot N_{\mathrm{gen}}^{\,h}
=
\frac{4\pi S_0r_1^3}{3}.
}
\tag{TRISO-FV-174}
\label{eq:triso-fv-174}
\end{equation}

The aligned finite-volume source inventory exactly reproduces the continuous total generation for constant $S_0$.

\subsubsection*{18.12 Cell-centred Robin boundary closure}

At the physical outer surface,

\begin{equation}
r=R.
\tag{TRISO-FV-175}
\label{eq:triso-fv-175}
\end{equation}

Let the centre of the outermost OPyC control volume be at

\begin{equation}
r_P<R.
\tag{TRISO-FV-176}
\label{eq:triso-fv-176}
\end{equation}

Define the centre-to-surface distance

\begin{equation}
\boxed{
\delta r_R=R-r_P.
}
\tag{TRISO-FV-177}
\label{eq:triso-fv-177}
\end{equation}

Let the outer OPyC diffusivity be

\begin{equation}
D_5.
\tag{TRISO-FV-178}
\label{eq:triso-fv-178}
\end{equation}

Let the physical surface concentration be

\begin{equation}
C_R.
\tag{TRISO-FV-179}
\label{eq:triso-fv-179}
\end{equation}

The outer-cell unknown is the cell-centred or cell-average concentration

\begin{equation}
C_P.
\tag{TRISO-FV-180}
\label{eq:triso-fv-180}
\end{equation}

These two concentrations are not silently identified.

\#\#\#\# 18.12.1 Half-cell diffusion relation

Approximate the OPyC concentration gradient between the outer cell centre and the physical surface by

\begin{equation}
\left.
\frac{\partial c}{\partial r}
\right|_{P\rightarrow R}
\approx
\frac{C_R-C_P}{\delta r_R}.
\tag{TRISO-FV-181}
\label{eq:triso-fv-181}
\end{equation}

Fick's law gives the outward radial flux

\begin{equation}
J_R
=
-D_5
\frac{C_R-C_P}{\delta r_R}.
\tag{TRISO-FV-182}
\label{eq:triso-fv-182}
\end{equation}

Reverse the concentration difference:

\begin{equation}
\boxed{
J_R
=
D_5
\frac{C_P-C_R}{\delta r_R}.
}
\tag{TRISO-FV-183}
\label{eq:triso-fv-183}
\end{equation}

Solve for the half-cell concentration drop:

\begin{equation}
C_P-C_R
=
J_R\frac{\delta r_R}{D_5}.
\tag{TRISO-FV-184}
\label{eq:triso-fv-184}
\end{equation}

The quantity

\begin{equation}
\frac{\delta r_R}{D_5}
\tag{TRISO-FV-185}
\label{eq:triso-fv-185}
\end{equation}

is the diffusion resistance per unit area of the outer half-cell.

\#\#\#\# 18.12.2 External film relation

The physical Robin law is

\begin{equation}
\boxed{
J_R
=
h(C_R-c_\infty).
}
\tag{TRISO-FV-186}
\label{eq:triso-fv-186}
\end{equation}

Solve for the film concentration drop:

\begin{equation}
C_R-c_\infty
=
\frac{J_R}{h}.
\tag{TRISO-FV-187}
\label{eq:triso-fv-187}
\end{equation}

The external-film resistance per unit area is therefore

\begin{equation}
\frac1h.
\tag{TRISO-FV-188}
\label{eq:triso-fv-188}
\end{equation}

\#\#\#\# 18.12.3 Add the two series concentration drops

Write the total cell-centre-to-bulk concentration difference as

\begin{equation}
C_P-c_\infty
=
(C_P-C_R)
+
(C_R-c_\infty).
\tag{TRISO-FV-189}
\label{eq:triso-fv-189}
\end{equation}

Substitute the half-cell drop (TRISO-FV-184):

\begin{equation}
C_P-c_\infty
=
J_R\frac{\delta r_R}{D_5}
+
(C_R-c_\infty).
\tag{TRISO-FV-190}
\label{eq:triso-fv-190}
\end{equation}

Substitute the film drop (TRISO-FV-187):

\begin{equation}
C_P-c_\infty
=
J_R\frac{\delta r_R}{D_5}
+
\frac{J_R}{h}.
\tag{TRISO-FV-191}
\label{eq:triso-fv-191}
\end{equation}

Factor out $J_R$:

\begin{equation}
C_P-c_\infty
=
J_R
\left(
\frac{\delta r_R}{D_5}
+
\frac1h
\right).
\tag{TRISO-FV-192}
\label{eq:triso-fv-192}
\end{equation}

Solve for the outward boundary flux:

\begin{equation}
\boxed{
J_R
=
\frac{
C_P-c_\infty
}{
\dfrac{\delta r_R}{D_5}
+
\dfrac1h
}.
}
\tag{TRISO-FV-193}
\label{eq:triso-fv-193}
\end{equation}

Thus the half-cell diffusion resistance and external-film resistance add in series.

\#\#\#\# 18.12.4 Effective boundary transfer coefficient

Define the effective cell-centre-to-bulk transfer coefficient

\begin{equation}
\boxed{
h_{\mathrm{eff}}
=
\left(
\frac{\delta r_R}{D_5}
+
\frac1h
\right)^{-1}.
}
\tag{TRISO-FV-194}
\label{eq:triso-fv-194}
\end{equation}

Then

\begin{equation}
\boxed{
J_R
=
h_{\mathrm{eff}}
(C_P-c_\infty).
}
\tag{TRISO-FV-195}
\label{eq:triso-fv-195}
\end{equation}

Multiply numerator and denominator of (TRISO-FV-194) by $hD_5$:

\begin{equation}
h_{\mathrm{eff}}
=
\frac{hD_5}{
h\delta r_R+D_5
}.
\tag{TRISO-FV-196}
\label{eq:triso-fv-196}
\end{equation}

Therefore

\begin{equation}
\boxed{
h_{\mathrm{eff}}
=
\frac{hD_5}{
D_5+h\delta r_R
}.
}
\tag{TRISO-FV-197}
\label{eq:triso-fv-197}
\end{equation}

The units are

\begin{equation}
[h_{\mathrm{eff}}]
=
\mathrm{m\,s^{-1}}.
\tag{TRISO-FV-198}
\label{eq:triso-fv-198}
\end{equation}

\#\#\#\# 18.12.5 Recover the physical surface concentration

The surface concentration can also be obtained explicitly.

From (TRISO-FV-187),

\begin{equation}
C_R
=
c_\infty+\frac{J_R}{h}.
\tag{TRISO-FV-199}
\label{eq:triso-fv-199}
\end{equation}

Substitute (TRISO-FV-193):

\begin{equation}
C_R
=
c_\infty
+
\frac1h
\frac{
C_P-c_\infty
}{
\dfrac{\delta r_R}{D_5}
+
\dfrac1h
}.
\tag{TRISO-FV-200}
\label{eq:triso-fv-200}
\end{equation}

Multiply the second term's denominator by $h$:

\begin{equation}
C_R
=
c_\infty
+
\frac{
C_P-c_\infty
}{
1+\dfrac{h\delta r_R}{D_5}
}.
\tag{TRISO-FV-201}
\label{eq:triso-fv-201}
\end{equation}

Define the half-cell boundary Biot number

\begin{equation}
\boxed{
\mathrm{Bi}_R
=
\frac{h\delta r_R}{D_5}.
}
\tag{TRISO-FV-202}
\label{eq:triso-fv-202}
\end{equation}

Then

\begin{equation}
\boxed{
C_R
=
c_\infty
+
\frac{
C_P-c_\infty
}{
1+\mathrm{Bi}_R
}.
}
\tag{TRISO-FV-203}
\label{eq:triso-fv-203}
\end{equation}

This expression keeps the cell-centre and physical surface concentrations distinct.

\#\#\#\# 18.12.6 Limiting checks

If

\begin{equation}
h\to0,
\tag{TRISO-FV-204}
\label{eq:triso-fv-204}
\end{equation}

then

\begin{equation}
\frac1h\to\infty.
\tag{TRISO-FV-205}
\label{eq:triso-fv-205}
\end{equation}

Therefore

\begin{equation}
h_{\mathrm{eff}}\to0.
\tag{TRISO-FV-206}
\label{eq:triso-fv-206}
\end{equation}

Hence

\begin{equation}
J_R\to0.
\tag{TRISO-FV-207}
\label{eq:triso-fv-207}
\end{equation}

This recovers the insulating Neumann limit.

If

\begin{equation}
\delta r_R\to0,
\tag{TRISO-FV-208}
\label{eq:triso-fv-208}
\end{equation}

then the half-cell resistance vanishes:

\begin{equation}
\frac{\delta r_R}{D_5}\to0.
\tag{TRISO-FV-209}
\label{eq:triso-fv-209}
\end{equation}

Therefore

\begin{equation}
h_{\mathrm{eff}}\to h.
\tag{TRISO-FV-210}
\label{eq:triso-fv-210}
\end{equation}

Thus the cell-centred closure approaches the physical Robin law as the outer cell centre approaches the surface.

If

\begin{equation}
h\to\infty,
\tag{TRISO-FV-211}
\label{eq:triso-fv-211}
\end{equation}

then the film resistance vanishes:

\begin{equation}
\frac1h\to0.
\tag{TRISO-FV-212}
\label{eq:triso-fv-212}
\end{equation}

Therefore

\begin{equation}
\boxed{
h_{\mathrm{eff}}
\to
\frac{D_5}{\delta r_R}.
}
\tag{TRISO-FV-213}
\label{eq:triso-fv-213}
\end{equation}

The resulting flux is

\begin{equation}
J_R
\to
\frac{D_5}{\delta r_R}
(C_P-c_\infty).
\tag{TRISO-FV-214}
\label{eq:triso-fv-214}
\end{equation}

This is the expected half-cell diffusion flux to a prescribed Dirichlet surface concentration $C_R=c_\infty$.

Unlike the ghost-point formula, this cell-centred resistance closure remains finite in the $h\to\infty$ limit.

\#\#\#\# 18.12.7 Outer-face amount conductance

The physical outer area is

\begin{equation}
\boxed{
A_R=4\pi R^2.
}
\tag{TRISO-FV-215}
\label{eq:triso-fv-215}
\end{equation}

Define the total outer-boundary conductance

\begin{equation}
\boxed{
G_R=A_Rh_{\mathrm{eff}}.
}
\tag{TRISO-FV-216}
\label{eq:triso-fv-216}
\end{equation}

Its units are

\begin{equation}
[G_R]
=
\mathrm{m^2}
\times
\mathrm{m\,s^{-1}}
=
\mathrm{m^3\,s^{-1}}.
\tag{TRISO-FV-217}
\label{eq:triso-fv-217}
\end{equation}

The outward amount rate is

\begin{equation}
A_RJ_R
=
G_R(C_P-c_\infty).
\tag{TRISO-FV-218}
\label{eq:triso-fv-218}
\end{equation}

\#\#\#\# 18.12.8 Final outer-cell semi-discrete balance

For the outermost control volume, the exact balance is

\begin{equation}
V_P\frac{dC_P}{dt}
=
A_wJ_w
-
A_RJ_R
+
S_PV_P.
\tag{TRISO-FV-219}
\label{eq:triso-fv-219}
\end{equation}

The west-face amount rate is

\begin{equation}
A_wJ_w
=
G_w(C_W-C_P).
\tag{TRISO-FV-220}
\label{eq:triso-fv-220}
\end{equation}

The outer amount rate is

\begin{equation}
A_RJ_R
=
G_R(C_P-c_\infty).
\tag{TRISO-FV-221}
\label{eq:triso-fv-221}
\end{equation}

Substitute both:

\begin{equation}
V_P\frac{dC_P}{dt}
=
G_w(C_W-C_P)
-
G_R(C_P-c_\infty)
+
S_PV_P.
\tag{TRISO-FV-222}
\label{eq:triso-fv-222}
\end{equation}

Expand the west term:

\begin{equation}
V_P\frac{dC_P}{dt}
=
G_wC_W
-
G_wC_P
-
G_R(C_P-c_\infty)
+
S_PV_P.
\tag{TRISO-FV-223}
\label{eq:triso-fv-223}
\end{equation}

Expand the boundary term:

\begin{equation}
V_P\frac{dC_P}{dt}
=
G_wC_W
-
G_wC_P
-
G_RC_P
+
G_Rc_\infty
+
S_PV_P.
\tag{TRISO-FV-224}
\label{eq:triso-fv-224}
\end{equation}

Collect the cell-centre concentration:

\begin{equation}
\boxed{
V_P\frac{dC_P}{dt}
=
G_wC_W
-
(G_w+G_R)C_P
+
G_Rc_\infty
+
S_PV_P.
}
\tag{TRISO-FV-225}
\label{eq:triso-fv-225}
\end{equation}

Divide by $V_P$:

\begin{equation}
\boxed{
\frac{dC_P}{dt}
=
\frac{G_w}{V_P}C_W
-
\frac{G_w+G_R}{V_P}C_P
+
\frac{G_R}{V_P}c_\infty
+
S_P.
}
\tag{TRISO-FV-226}
\label{eq:triso-fv-226}
\end{equation}

For the benchmark

\begin{equation}
c_\infty=0,
\tag{TRISO-FV-227}
\label{eq:triso-fv-227}
\end{equation}

this reduces to

\begin{equation}
\boxed{
\frac{dC_P}{dt}
=
\frac{G_w}{V_P}C_W
-
\frac{G_w+G_R}{V_P}C_P
+
S_P.
}
\tag{TRISO-FV-228}
\label{eq:triso-fv-228}
\end{equation}

For the physical five-layer kernel-confined source model, the outer OPyC cell has

\begin{equation}
S_P=0.
\tag{TRISO-FV-229}
\label{eq:triso-fv-229}
\end{equation}

\#\#\#\# 18.12.9 Explicit Euler outer-cell update

Apply forward Euler to (TRISO-FV-226):

\begin{equation}
\frac{
C_P^{j+1}-C_P^j
}{
\Delta t
}
=
\frac{G_w}{V_P}C_W^j
-
\frac{G_w+G_R}{V_P}C_P^j
+
\frac{G_R}{V_P}c_\infty^j
+
S_P^j.
\tag{TRISO-FV-230}
\label{eq:triso-fv-230}
\end{equation}

Multiply by $\Delta t$:

\begin{equation}
C_P^{j+1}-C_P^j
=
\frac{\Delta tG_w}{V_P}C_W^j
-
\frac{\Delta t(G_w+G_R)}{V_P}C_P^j
+
\frac{\Delta tG_R}{V_P}c_\infty^j
+
S_P^j\Delta t.
\tag{TRISO-FV-231}
\label{eq:triso-fv-231}
\end{equation}

Add $C_P^j$:

\begin{equation}
\boxed{
C_P^{j+1}
=
\frac{\Delta tG_w}{V_P}C_W^j
+
\left[
1-\frac{\Delta t(G_w+G_R)}{V_P}
\right]C_P^j
+
\frac{\Delta tG_R}{V_P}c_\infty^j
+
S_P^j\Delta t.
}
\tag{TRISO-FV-232}
\label{eq:triso-fv-232}
\end{equation}

This closes the cell-centred Robin boundary without identifying the cell-centre concentration with the physical surface concentration.

\subsubsection*{18.13 Explicit-Euler positivity condition for an ordinary cell}

From (TRISO-FV-158), the neighbour coefficients are

\begin{equation}
\frac{\Delta t\,G_w}{V_P}\ge0
\tag{TRISO-FV-600}
\label{eq:triso-fv-600}
\end{equation}

and

\begin{equation}
\frac{\Delta t\,G_e}{V_P}\ge0.
\tag{TRISO-FV-601}
\label{eq:triso-fv-601}
\end{equation}

The central coefficient is non-negative when

\begin{equation}
1-
\frac{\Delta t(G_w+G_e)}{V_P}
\ge0.
\tag{TRISO-FV-602}
\label{eq:triso-fv-602}
\end{equation}

Rearrange:

\begin{equation}
\Delta t(G_w+G_e)
\le
V_P.
\tag{TRISO-FV-603}
\label{eq:triso-fv-603}
\end{equation}

Therefore

\begin{equation}
\boxed{
\Delta t
\le
\frac{V_P}{G_w+G_e}.
}
\tag{TRISO-FV-604}
\label{eq:triso-fv-604}
\end{equation}

For every ordinary cell, a sufficient global coefficient-positivity restriction is

\begin{equation}
\boxed{
\Delta t
\le
\min_P
\frac{V_P}{G_w+G_e}.
}
\tag{TRISO-FV-605}
\label{eq:triso-fv-605}
\end{equation}

This is the finite-volume analogue of the earlier homogeneous FTCS monotonicity restriction.

The outer Robin closure is now available in Section 18.12. The final global positivity bound must include its boundary conductance.

\subsubsection*{18.14 Global five-layer semi-discrete matrix}

Let the finite-volume mesh contain $M$ spherical cells.

Index the cell-average concentrations by

\begin{equation}
\mathbf C(t)
=
(C_0,C_1,\ldots,C_{M-1})^T.
\tag{TRISO-FV-233}
\label{eq:triso-fv-233}
\end{equation}

Define the diagonal volume matrix

\begin{equation}
\boxed{
\mathbf V
=
\operatorname{diag}
(V_0,V_1,\ldots,V_{M-1}).
}
\tag{TRISO-FV-234}
\label{eq:triso-fv-234}
\end{equation}

Every cell volume is positive:

\begin{equation}
V_P>0.
\tag{TRISO-FV-235}
\label{eq:triso-fv-235}
\end{equation}

Let $G_{P+\frac12}$ denote the total conductance of the face shared by cells $P$ and $P+1$.

At an ordinary same-material face,

\begin{equation}
G_{P+\frac12}
=
\frac{
A_{P+\frac12}D
}{
r_{P+1}-r_P
}.
\tag{TRISO-FV-236}
\label{eq:triso-fv-236}
\end{equation}

At a material-interface face,

\begin{equation}
\boxed{
G_{P+\frac12}
=
A_{P+\frac12}
\left(
\frac{\delta r_P}{D_P}
+
\frac{\delta r_{P+1}}{D_{P+1}}
\right)^{-1}.
}
\tag{TRISO-FV-237}
\label{eq:triso-fv-237}
\end{equation}

Thus the matrix assembly does not require a separate interface unknown.

\subsubsection*{18.15 Central-cell row}

The central-cell balance is

\begin{equation}
V_0\frac{dC_0}{dt}
=
G_{\frac12}(C_1-C_0)
+
S_0V_0.
\tag{TRISO-FV-238}
\label{eq:triso-fv-238}
\end{equation}

Expand:

\begin{equation}
V_0\frac{dC_0}{dt}
=
-G_{\frac12}C_0
+
G_{\frac12}C_1
+
S_0V_0.
\tag{TRISO-FV-239}
\label{eq:triso-fv-239}
\end{equation}

Therefore the first row of the transport matrix contains

\begin{equation}
K_{0,0}
=
-G_{\frac12},
\tag{TRISO-FV-240}
\label{eq:triso-fv-240}
\end{equation}

and

\begin{equation}
K_{0,1}
=
G_{\frac12}.
\tag{TRISO-FV-241}
\label{eq:triso-fv-241}
\end{equation}

\subsubsection*{18.16 Ordinary-cell row}

For cell $P$, with

\begin{equation}
1\le P\le M-2,
\tag{TRISO-FV-242}
\label{eq:triso-fv-242}
\end{equation}

the conservative balance is

\begin{equation}
V_P\frac{dC_P}{dt}
=
G_{P-\frac12}(C_{P-1}-C_P)
+
G_{P+\frac12}(C_{P+1}-C_P)
+
S_PV_P.
\tag{TRISO-FV-243}
\label{eq:triso-fv-243}
\end{equation}

Expand the west-face term:

\begin{equation}
G_{P-\frac12}(C_{P-1}-C_P)
=
G_{P-\frac12}C_{P-1}
-
G_{P-\frac12}C_P.
\tag{TRISO-FV-244}
\label{eq:triso-fv-244}
\end{equation}

Expand the east-face term:

\begin{equation}
G_{P+\frac12}(C_{P+1}-C_P)
=
G_{P+\frac12}C_{P+1}
-
G_{P+\frac12}C_P.
\tag{TRISO-FV-245}
\label{eq:triso-fv-245}
\end{equation}

Substitute:

\begin{equation}
V_P\frac{dC_P}{dt}
=
G_{P-\frac12}C_{P-1}
-
G_{P-\frac12}C_P
+
G_{P+\frac12}C_{P+1}
-
G_{P+\frac12}C_P
+
S_PV_P.
\tag{TRISO-FV-246}
\label{eq:triso-fv-246}
\end{equation}

Collect the central coefficient:

\begin{equation}
\boxed{
V_P\frac{dC_P}{dt}
=
G_{P-\frac12}C_{P-1}
-
\left(
G_{P-\frac12}+G_{P+\frac12}
\right)C_P
+
G_{P+\frac12}C_{P+1}
+
S_PV_P.
}
\tag{TRISO-FV-247}
\label{eq:triso-fv-247}
\end{equation}

Therefore

\begin{equation}
K_{P,P-1}
=
G_{P-\frac12},
\tag{TRISO-FV-248}
\label{eq:triso-fv-248}
\end{equation}

\begin{equation}
K_{P,P}
=
-\left(
G_{P-\frac12}+G_{P+\frac12}
\right),
\tag{TRISO-FV-249}
\label{eq:triso-fv-249}
\end{equation}

and

\begin{equation}
K_{P,P+1}
=
G_{P+\frac12}.
\tag{TRISO-FV-250}
\label{eq:triso-fv-250}
\end{equation}

These formulas remain valid when either face is a material interface because the corresponding $G$ already contains the resistance-weighted discontinuous-$D$ coupling.

\subsubsection*{18.17 Outermost-cell row}

Let the outermost cell index be

\begin{equation}
P=M-1.
\tag{TRISO-FV-251}
\label{eq:triso-fv-251}
\end{equation}

Section 18.12 gives

\begin{equation}
V_{M-1}\frac{dC_{M-1}}{dt}
=
G_{M-\frac32}C_{M-2}
-
\left(
G_{M-\frac32}+G_R
\right)C_{M-1}
+
G_Rc_\infty
+
S_{M-1}V_{M-1}.
\tag{TRISO-FV-252}
\label{eq:triso-fv-252}
\end{equation}

Therefore

\begin{equation}
K_{M-1,M-2}
=
G_{M-\frac32},
\tag{TRISO-FV-253}
\label{eq:triso-fv-253}
\end{equation}

and

\begin{equation}
K_{M-1,M-1}
=
-\left(
G_{M-\frac32}+G_R
\right).
\tag{TRISO-FV-254}
\label{eq:triso-fv-254}
\end{equation}

The external concentration enters as a forcing term rather than as another particle unknown.

\subsubsection*{18.18 Assemble the transport matrix}

Define the source vector

\begin{equation}
\mathbf S
=
(S_0,S_1,\ldots,S_{M-1})^T.
\tag{TRISO-FV-255}
\label{eq:triso-fv-255}
\end{equation}

Define the external-boundary forcing vector

\begin{equation}
\mathbf b_\infty
=
(0,0,\ldots,0,G_Rc_\infty)^T.
\tag{TRISO-FV-256}
\label{eq:triso-fv-256}
\end{equation}

The complete semi-discrete system is

\begin{equation}
\boxed{
\mathbf V
\frac{d\mathbf C}{dt}
=
\mathbf K\mathbf C
+
\mathbf V\mathbf S
+
\mathbf b_\infty.
}
\tag{TRISO-FV-257}
\label{eq:triso-fv-257}
\end{equation}

The matrix $\mathbf K$ is tridiagonal:

\begin{equation}
\mathbf K
=
\begin{pmatrix}
-G_{\frac12}
&
G_{\frac12}
&
0
&
\cdots
&
0
\\
G_{\frac12}
&
-(G_{\frac12}+G_{\frac32})
&
G_{\frac32}
&
\ddots
&
\vdots
\\
0
&
G_{\frac32}
&
-(G_{\frac32}+G_{\frac52})
&
\ddots
&
0
\\
\vdots
&
\ddots
&
\ddots
&
\ddots
&
G_{M-\frac32}
\\
0
&
\cdots
&
0
&
G_{M-\frac32}
&
-(G_{M-\frac32}+G_R)
\end{pmatrix}.
\tag{TRISO-FV-258}
\label{eq:triso-fv-258}
\end{equation}

Because each shared-face conductance appears identically in the two neighbouring rows,

\begin{equation}
\boxed{
\mathbf K=\mathbf K^T.
}
\tag{TRISO-FV-259}
\label{eq:triso-fv-259}
\end{equation}

The time-evolution operator in concentration coordinates is

\begin{equation}
\boxed{
\mathbf L
=
\mathbf V^{-1}\mathbf K.
}
\tag{TRISO-FV-260}
\label{eq:triso-fv-260}
\end{equation}

The matrix $\mathbf L$ is generally not symmetric because cell volumes differ with radius.

\subsubsection*{18.19 Weighted self-adjoint structure of the semi-discrete operator}

Define the discrete volume-weighted inner product

\begin{equation}
\boxed{
\langle\mathbf x,\mathbf y\rangle_V
=
\mathbf x^T\mathbf V\mathbf y.
}
\tag{TRISO-FV-261}
\label{eq:triso-fv-261}
\end{equation}

Evaluate

\begin{equation}
\langle\mathbf x,\mathbf L\mathbf y\rangle_V
=
\mathbf x^T\mathbf V\mathbf L\mathbf y.
\tag{TRISO-FV-262}
\label{eq:triso-fv-262}
\end{equation}

Use $\mathbf L=\mathbf V^{-1}\mathbf K$:

\begin{equation}
\mathbf V\mathbf L
=
\mathbf K.
\tag{TRISO-FV-263}
\label{eq:triso-fv-263}
\end{equation}

Therefore

\begin{equation}
\langle\mathbf x,\mathbf L\mathbf y\rangle_V
=
\mathbf x^T\mathbf K\mathbf y.
\tag{TRISO-FV-264}
\label{eq:triso-fv-264}
\end{equation}

Because $\mathbf K=\mathbf K^T$,

\begin{equation}
\mathbf x^T\mathbf K\mathbf y
=
\mathbf y^T\mathbf K\mathbf x.
\tag{TRISO-FV-265}
\label{eq:triso-fv-265}
\end{equation}

Reverse the preceding steps:

\begin{equation}
\mathbf y^T\mathbf K\mathbf x
=
\langle\mathbf L\mathbf x,\mathbf y\rangle_V.
\tag{TRISO-FV-266}
\label{eq:triso-fv-266}
\end{equation}

Hence

\begin{equation}
\boxed{
\langle\mathbf x,\mathbf L\mathbf y\rangle_V
=
\langle\mathbf L\mathbf x,\mathbf y\rangle_V.
}
\tag{TRISO-FV-267}
\label{eq:triso-fv-267}
\end{equation}

Thus the conservative semi-discrete diffusion operator is self-adjoint in the volume-weighted discrete inner product.

This is the discrete analogue of the continuum $r^2$-weighted self-adjoint structure.

\subsubsection*{18.20 Negative-semidefinite diffusion form}

For any vector $\mathbf x$,

\begin{equation}
\mathbf x^T\mathbf K\mathbf x
\tag{TRISO-FV-268}
\label{eq:triso-fv-268}
\end{equation}

can be grouped face-by-face.

An internal face between $P$ and $P+1$ contributes

\begin{equation}
-G_{P+\frac12}x_P^2
+
2G_{P+\frac12}x_Px_{P+1}
-
G_{P+\frac12}x_{P+1}^2.
\tag{TRISO-FV-269}
\label{eq:triso-fv-269}
\end{equation}

Factor $-G_{P+\frac12}$:

\begin{equation}
-G_{P+\frac12}
\left(
x_P^2
-
2x_Px_{P+1}
+
x_{P+1}^2
\right).
\tag{TRISO-FV-270}
\label{eq:triso-fv-270}
\end{equation}

Recognise the square:

\begin{equation}
x_P^2
-
2x_Px_{P+1}
+
x_{P+1}^2
=
(x_{P+1}-x_P)^2.
\tag{TRISO-FV-271}
\label{eq:triso-fv-271}
\end{equation}

Therefore each internal face contributes

\begin{equation}
-G_{P+\frac12}
(x_{P+1}-x_P)^2.
\tag{TRISO-FV-272}
\label{eq:triso-fv-272}
\end{equation}

The Robin boundary contributes

\begin{equation}
-G_Rx_{M-1}^2.
\tag{TRISO-FV-273}
\label{eq:triso-fv-273}
\end{equation}

Thus

\begin{equation}
\boxed{
\mathbf x^T\mathbf K\mathbf x
=
-
\sum_{P=0}^{M-2}
G_{P+\frac12}
(x_{P+1}-x_P)^2
-
G_Rx_{M-1}^2.
}
\tag{TRISO-FV-274}
\label{eq:triso-fv-274}
\end{equation}

Since all conductances are non-negative,

\begin{equation}
\boxed{
\mathbf x^T\mathbf K\mathbf x\le0.
}
\tag{TRISO-FV-275}
\label{eq:triso-fv-275}
\end{equation}

For $G_R>0$, equality requires both

\begin{equation}
x_{P+1}=x_P
\tag{TRISO-FV-276}
\label{eq:triso-fv-276}
\end{equation}

for every internal face and

\begin{equation}
x_{M-1}=0.
\tag{TRISO-FV-277}
\label{eq:triso-fv-277}
\end{equation}

Therefore

\begin{equation}
\mathbf x=\mathbf0.
\tag{TRISO-FV-278}
\label{eq:triso-fv-278}
\end{equation}

Hence, for a finite-transfer or absorbing outer boundary with $G_R>0$,

\begin{equation}
\boxed{
\mathbf K
\text{ is negative definite.}
}
\tag{TRISO-FV-279}
\label{eq:triso-fv-279}
\end{equation}

If $G_R=0$, the constant vector is the expected zero mode of a closed no-flux particle.

\subsubsection*{18.21 Global discrete inventory balance}

Define the discrete total particle inventory

\begin{equation}
\boxed{
N_h(t)
=
\mathbf 1^T\mathbf V\mathbf C
=
\sum_{P=0}^{M-1}
V_PC_P.
}
\tag{TRISO-FV-280}
\label{eq:triso-fv-280}
\end{equation}

Differentiate:

\begin{equation}
\frac{dN_h}{dt}
=
\mathbf1^T
\mathbf V
\frac{d\mathbf C}{dt}.
\tag{TRISO-FV-281}
\label{eq:triso-fv-281}
\end{equation}

Use the semi-discrete system:

\begin{equation}
\frac{dN_h}{dt}
=
\mathbf1^T\mathbf K\mathbf C
+
\mathbf1^T\mathbf V\mathbf S
+
\mathbf1^T\mathbf b_\infty.
\tag{TRISO-FV-282}
\label{eq:triso-fv-282}
\end{equation}

All internal conductance contributions cancel in the row sum.

The only non-zero transport contribution is the Robin boundary:

\begin{equation}
\mathbf1^T\mathbf K\mathbf C
=
-G_RC_{M-1}.
\tag{TRISO-FV-283}
\label{eq:triso-fv-283}
\end{equation}

The boundary forcing is

\begin{equation}
\mathbf1^T\mathbf b_\infty
=
G_Rc_\infty.
\tag{TRISO-FV-284}
\label{eq:triso-fv-284}
\end{equation}

Therefore

\begin{equation}
\boxed{
\frac{dN_h}{dt}
=
\sum_{P=0}^{M-1}S_PV_P
-
G_R(C_{M-1}-c_\infty).
}
\tag{TRISO-FV-285}
\label{eq:triso-fv-285}
\end{equation}

This is exactly the discrete statement

\begin{equation}
\text{accumulation}
=
\text{generation}
-
\text{outward release}.
\tag{TRISO-FV-286}
\label{eq:triso-fv-286}
\end{equation}

For the kernel-confined constant source,

\begin{equation}
\sum_PS_PV_P
=
\frac{4\pi S_0r_1^3}{3}
\tag{TRISO-FV-287}
\label{eq:triso-fv-287}
\end{equation}

when the source/material interface is face-aligned.

Thus

\begin{equation}
\boxed{
\frac{dN_h}{dt}
=
\frac{4\pi S_0r_1^3}{3}
-
G_R(C_{M-1}-c_\infty).
}
\tag{TRISO-FV-288}
\label{eq:triso-fv-288}
\end{equation}

\subsubsection*{18.22 Explicit-Euler matrix update}

Apply forward Euler to (TRISO-FV-257):

\begin{equation}
\mathbf V
\frac{
\mathbf C^{j+1}-\mathbf C^j
}{
\Delta t
}
=
\mathbf K\mathbf C^j
+
\mathbf V\mathbf S^j
+
\mathbf b_\infty^j.
\tag{TRISO-FV-289}
\label{eq:triso-fv-289}
\end{equation}

Multiply by $\mathbf V^{-1}$:

\begin{equation}
\frac{
\mathbf C^{j+1}-\mathbf C^j
}{
\Delta t
}
=
\mathbf V^{-1}\mathbf K\mathbf C^j
+
\mathbf S^j
+
\mathbf V^{-1}\mathbf b_\infty^j.
\tag{TRISO-FV-290}
\label{eq:triso-fv-290}
\end{equation}

Multiply by $\Delta t$:

\begin{equation}
\mathbf C^{j+1}-\mathbf C^j
=
\Delta t\,\mathbf V^{-1}\mathbf K\mathbf C^j
+
\Delta t\,\mathbf S^j
+
\Delta t\,\mathbf V^{-1}\mathbf b_\infty^j.
\tag{TRISO-FV-291}
\label{eq:triso-fv-291}
\end{equation}

Add $\mathbf C^j$:

\begin{equation}
\boxed{
\mathbf C^{j+1}
=
\mathbf A_{\mathrm{FV}}\mathbf C^j
+
\Delta t\,\mathbf S^j
+
\Delta t\,\mathbf V^{-1}\mathbf b_\infty^j,
}
\tag{TRISO-FV-292}
\label{eq:triso-fv-292}
\end{equation}

where

\begin{equation}
\boxed{
\mathbf A_{\mathrm{FV}}
=
\mathbf I
+
\Delta t\,\mathbf V^{-1}\mathbf K.
}
\tag{TRISO-FV-293}
\label{eq:triso-fv-293}
\end{equation}

\subsubsection*{18.23 Completed coefficient-positivity bound}

For the central cell, the explicit update has central coefficient

\begin{equation}
1-
\frac{\Delta tG_{\frac12}}{V_0}.
\tag{TRISO-FV-294}
\label{eq:triso-fv-294}
\end{equation}

It is non-negative when

\begin{equation}
\boxed{
\Delta t
\le
\frac{V_0}{G_{\frac12}}.
}
\tag{TRISO-FV-295}
\label{eq:triso-fv-295}
\end{equation}

For an ordinary cell $P$,

\begin{equation}
1-
\frac{
\Delta t
\left(
G_{P-\frac12}+G_{P+\frac12}
\right)
}{
V_P
}
\ge0.
\tag{TRISO-FV-296}
\label{eq:triso-fv-296}
\end{equation}

Therefore

\begin{equation}
\boxed{
\Delta t
\le
\frac{
V_P
}{
G_{P-\frac12}+G_{P+\frac12}
}.
}
\tag{TRISO-FV-297}
\label{eq:triso-fv-297}
\end{equation}

For the outer cell,

\begin{equation}
1-
\frac{
\Delta t
\left(
G_{M-\frac32}+G_R
\right)
}{
V_{M-1}
}
\ge0.
\tag{TRISO-FV-298}
\label{eq:triso-fv-298}
\end{equation}

Therefore

\begin{equation}
\boxed{
\Delta t
\le
\frac{
V_{M-1}
}{
G_{M-\frac32}+G_R
}.
}
\tag{TRISO-FV-299}
\label{eq:triso-fv-299}
\end{equation}

Combine all cells:

\begin{equation}
\boxed{
\Delta t
\le
\min
\left\{
\frac{V_0}{G_{\frac12}},
\;
\min_{1\le P\le M-2}
\frac{
V_P
}{
G_{P-\frac12}+G_{P+\frac12}
},
\;
\frac{
V_{M-1}
}{
G_{M-\frac32}+G_R
}
\right\}.
}
\tag{TRISO-FV-300}
\label{eq:triso-fv-300}
\end{equation}

Under this condition, every off-diagonal amplification coefficient is non-negative and every diagonal amplification coefficient is non-negative.

\subsubsection*{18.24 Row sums and monotonicity}

For the centre row, the two homogeneous coefficients sum to

\begin{equation}
1.
\tag{TRISO-FV-301}
\label{eq:triso-fv-301}
\end{equation}

For every ordinary interior cell, the three homogeneous coefficients sum to

\begin{equation}
1.
\tag{TRISO-FV-302}
\label{eq:triso-fv-302}
\end{equation}

For the outer row, the particle-state coefficients sum to

\begin{equation}
1-
\frac{\Delta tG_R}{V_{M-1}}.
\tag{TRISO-FV-303}
\label{eq:triso-fv-303}
\end{equation}

For $G_R\ge0$,

\begin{equation}
1-
\frac{\Delta tG_R}{V_{M-1}}
\le1.
\tag{TRISO-FV-304}
\label{eq:triso-fv-304}
\end{equation}

Under (TRISO-FV-300), the row sum is also non-negative.

Therefore the homogeneous particle amplification matrix is substochastic.

Hence

\begin{equation}
\boxed{
\|\mathbf A_{\mathrm{FV}}\|_\infty\le1.
}
\tag{TRISO-FV-305}
\label{eq:triso-fv-305}
\end{equation}

Consequently,

\begin{equation}
\boxed{
\rho(\mathbf A_{\mathrm{FV}})\le1.
}
\tag{TRISO-FV-306}
\label{eq:triso-fv-306}
\end{equation}

Thus (TRISO-FV-300) is a sufficient explicit-Euler monotonicity and $\ell_\infty$-stability condition for the completed finite-volume system.

It is not asserted to be a necessary spectral-stability condition.

\subsubsection*{18.25 Semi-discrete spectral sign}

The generalized eigenproblem is

\begin{equation}
\mathbf K\mathbf x
=
\lambda
\mathbf V\mathbf x.
\tag{TRISO-FV-307}
\label{eq:triso-fv-307}
\end{equation}

Premultiply by $\mathbf x^T$:

\begin{equation}
\mathbf x^T\mathbf K\mathbf x
=
\lambda
\mathbf x^T\mathbf V\mathbf x.
\tag{TRISO-FV-308}
\label{eq:triso-fv-308}
\end{equation}

For non-zero $\mathbf x$,

\begin{equation}
\mathbf x^T\mathbf V\mathbf x>0.
\tag{TRISO-FV-309}
\label{eq:triso-fv-309}
\end{equation}

From (TRISO-FV-275),

\begin{equation}
\mathbf x^T\mathbf K\mathbf x\le0.
\tag{TRISO-FV-310}
\label{eq:triso-fv-310}
\end{equation}

Therefore

\begin{equation}
\boxed{
\lambda\le0.
}
\tag{TRISO-FV-311}
\label{eq:triso-fv-311}
\end{equation}

For $G_R>0$, $\mathbf K$ is negative definite, so

\begin{equation}
\boxed{
\lambda<0
}
\tag{TRISO-FV-312}
\label{eq:triso-fv-312}
\end{equation}

for every non-zero mode.

This establishes the expected diffusive sign of the semi-discrete spectrum without numerically enumerating eigenvalues.

\subsubsection*{18.26 Remaining accuracy and convergence questions}

The complete five-layer finite-volume algebra is now assembled.

The following statements are established:

\noindent\textbullet\ exact control-volume conservation;\\
\noindent\textbullet\ resistance-weighted interface fluxes;\\
\noindent\textbullet\ a closed cell-centred Robin boundary;\\
\noindent\textbullet\ a symmetric conductance matrix $\mathbf K$;\\
\noindent\textbullet\ volume-weighted self-adjointness;\\
\noindent\textbullet\ non-positive semi-discrete spectrum;\\
\noindent\textbullet\ a sufficient explicit-Euler monotonicity/$\ell_\infty$-stability bound.\\

The remaining mathematical questions are narrower:

[UNVERIFIED] Global spatial order of accuracy when $D(r)$ is discontinuous.

[UNVERIFIED] Global temporal/spatial convergence rate of the fully discrete scheme.

[UNVERIFIED] Numerical convergence study against the analytical benchmarks.

These require a dedicated consistency/convergence pass rather than further coefficient assembly.

\subsubsection*{18.27 Status of the five-layer deterministic discretisation}

\textbf{[VERIFIED]} The spherical control-volume geometry and exact integrated conservation balance.

\textbf{[VERIFIED]} Conservative internal face coupling.

\textbf{[VERIFIED]} Harmonic/resistance-weighted treatment of discontinuous diffusivity.

\textbf{[VERIFIED]} Pairwise cancellation of internal face fluxes.

\textbf{[VERIFIED]} Geometric centre treatment through the zero-area inner face.

\textbf{[VERIFIED]} Exact kernel generation inventory when material/source interfaces align with control-volume faces.

\textbf{[VERIFIED]} Cell-centred Robin surface closure through half-cell diffusion plus external-film resistance.

\textbf{[VERIFIED]} Final assembled five-layer coefficient matrix including the outer boundary.

[UNVERIFIED] Accuracy order of the complete multilayer finite-volume scheme.

[CONDITIONALLY VERIFIED] Sufficient explicit-Euler monotonicity/stability bound; convergence rate remains unverified.

The finite-volume derivation is therefore the current canonical deterministic route for the discontinuous-$D$ five-layer model, while the original FTCS scheme remains a transparent homogeneous benchmark.

\subsection*{18A. Finite-volume consistency: smooth same-material cells}

This section begins the accuracy/convergence stage authorised by the independent discrete-mathematics audit.

It treats only smooth cells whose two faces lie inside one material with constant diffusivity $D$. Material-interface and outer-boundary consistency are deferred to subsequent subsections.

\subsubsection*{18A.1 Exact cell average versus representative point value}

The canonical unknown is the exact spherical volume average

\begin{equation}
C_P
=
\frac1{V_P}
\int_{r_w}^{r_e}
c(r)\,4\pi r^2\,dr.
\tag{TRISO-ACC-100}
\label{eq:triso-acc-100}
\end{equation}

The representative coordinate $r_P$ is the spherical volume centroid:

\begin{equation}
r_P
=
\frac1{V_P}
\int_{r_w}^{r_e}
r\,4\pi r^2\,dr.
\tag{TRISO-ACC-101}
\label{eq:triso-acc-101}
\end{equation}

Subtract $r_P$ inside the weighted first moment:

\begin{equation}
\int_{r_w}^{r_e}
(r-r_P)\,4\pi r^2\,dr
=
\int_{r_w}^{r_e}
r\,4\pi r^2\,dr
-
r_P
\int_{r_w}^{r_e}
4\pi r^2\,dr.
\tag{TRISO-ACC-102}
\label{eq:triso-acc-102}
\end{equation}

Use the centroid definition in the first term:

\begin{equation}
\int_{r_w}^{r_e}
r\,4\pi r^2\,dr
=
r_PV_P.
\tag{TRISO-ACC-103}
\label{eq:triso-acc-103}
\end{equation}

Use the volume definition in the second term:

\begin{equation}
\int_{r_w}^{r_e}
4\pi r^2\,dr
=
V_P.
\tag{TRISO-ACC-104}
\label{eq:triso-acc-104}
\end{equation}

Therefore

\begin{equation}
\boxed{
\int_{r_w}^{r_e}
(r-r_P)\,4\pi r^2\,dr
=
0.
}
\tag{TRISO-ACC-105}
\label{eq:triso-acc-105}
\end{equation}

Taylor-expand a smooth concentration about $r_P$:

\begin{equation}
c(r)
=
c(r_P)
+
c_r(r_P)(r-r_P)
+
\frac12c_{rr}(r_P)(r-r_P)^2
+
O(h_P^3),
\tag{TRISO-ACC-106}
\label{eq:triso-acc-106}
\end{equation}

where

\begin{equation}
h_P=r_e-r_w.
\tag{TRISO-ACC-107}
\label{eq:triso-acc-107}
\end{equation}

Insert (TRISO-ACC-106) into the exact average:

\begin{equation}
C_P
=
\frac1{V_P}
\int_{r_w}^{r_e}
\left[
c(r_P)
+
c_r(r_P)(r-r_P)
+
\frac12c_{rr}(r_P)(r-r_P)^2
+
O(h_P^3)
\right]
4\pi r^2\,dr.
\tag{TRISO-ACC-108}
\label{eq:triso-acc-108}
\end{equation}

Separate the constant term:

\begin{equation}
\frac{c(r_P)}{V_P}
\int_{r_w}^{r_e}4\pi r^2\,dr
=
c(r_P).
\tag{TRISO-ACC-109}
\label{eq:triso-acc-109}
\end{equation}

The first-order term vanishes by (TRISO-ACC-105):

\begin{equation}
\frac{c_r(r_P)}{V_P}
\int_{r_w}^{r_e}
(r-r_P)4\pi r^2\,dr
=
0.
\tag{TRISO-ACC-110}
\label{eq:triso-acc-110}
\end{equation}

Define the weighted second central moment

\begin{equation}
\mu_{2,P}
=
\frac1{V_P}
\int_{r_w}^{r_e}
(r-r_P)^2\,4\pi r^2\,dr.
\tag{TRISO-ACC-111}
\label{eq:triso-acc-111}
\end{equation}

For a shape-regular refining radial mesh,

\begin{equation}
\mu_{2,P}=O(h_P^2).
\tag{TRISO-ACC-112}
\label{eq:triso-acc-112}
\end{equation}

Therefore

\begin{equation}
\boxed{
C_P
=
c(r_P)
+
\frac12c_{rr}(r_P)\mu_{2,P}
+
O(h_P^3).
}
\tag{TRISO-ACC-113}
\label{eq:triso-acc-113}
\end{equation}

In particular,

\begin{equation}
\boxed{
C_P-c(r_P)=O(h_P^2).
}
\tag{TRISO-ACC-114}
\label{eq:triso-acc-114}
\end{equation}

Thus locating the exact cell average at the spherical volume centroid is a second-order point-representation approximation for a smooth field. It is not exact equality.

\subsubsection*{18A.2 Exact integrated balance remains independent of reconstruction}

The control-volume identity

\begin{equation}
V_P\frac{dC_P}{dt}
=
A_wJ_w-A_eJ_e+S_PV_P
\tag{TRISO-ACC-115}
\label{eq:triso-acc-115}
\end{equation}

was obtained by exact integration.

No point-value approximation was used to obtain (TRISO-ACC-115).

Therefore the spatial consistency question is isolated to the approximation of the face fluxes.

\subsubsection*{18A.3 Two-point gradient on a smooth same-material face}

Consider a face $f=r_{P+\frac12}$ between cells $P$ and $E=P+1$.

Define

\begin{equation}
d_P=r_f-r_P,
\tag{TRISO-ACC-116}
\label{eq:triso-acc-116}
\end{equation}

and

\begin{equation}
d_E=r_E-r_f.
\tag{TRISO-ACC-117}
\label{eq:triso-acc-117}
\end{equation}

Hence

\begin{equation}
r_E-r_P=d_P+d_E.
\tag{TRISO-ACC-118}
\label{eq:triso-acc-118}
\end{equation}

Taylor-expand the exact point value at $r_P$ about the face:

\begin{equation}
c(r_P)
=
c_f
-
d_Pc_f'
+
\frac{d_P^2}{2}c_f''
-
\frac{d_P^3}{6}c_f'''
+
O(h^4).
\tag{TRISO-ACC-119}
\label{eq:triso-acc-119}
\end{equation}

Taylor-expand the exact point value at $r_E$:

\begin{equation}
c(r_E)
=
c_f
+
d_Ec_f'
+
\frac{d_E^2}{2}c_f''
+
\frac{d_E^3}{6}c_f'''
+
O(h^4).
\tag{TRISO-ACC-120}
\label{eq:triso-acc-120}
\end{equation}

Subtract (TRISO-ACC-119) from (TRISO-ACC-120):

\begin{equation}
c(r_E)-c(r_P)
=
(d_P+d_E)c_f'
+
\frac{d_E^2-d_P^2}{2}c_f''
+
\frac{d_E^3+d_P^3}{6}c_f'''
+
O(h^4).
\tag{TRISO-ACC-121}
\label{eq:triso-acc-121}
\end{equation}

Divide by $d_P+d_E$:

\begin{equation}
\frac{c(r_E)-c(r_P)}{r_E-r_P}
=
c_f'
+
\frac{d_E-d_P}{2}c_f''
+
\frac{d_E^3+d_P^3}{6(d_P+d_E)}c_f'''
+
O(h^3).
\tag{TRISO-ACC-122}
\label{eq:triso-acc-122}
\end{equation}

For a locally symmetric representative geometry,

\begin{equation}
d_E-d_P=O(h^2),
\tag{TRISO-ACC-123}
\label{eq:triso-acc-123}
\end{equation}

while

\begin{equation}
d_P=O(h),
\qquad
d_E=O(h).
\tag{TRISO-ACC-124}
\label{eq:triso-acc-124}
\end{equation}

Therefore

\begin{equation}
\frac{d_E^3+d_P^3}{d_P+d_E}
=
O(h^2).
\tag{TRISO-ACC-125}
\label{eq:triso-acc-125}
\end{equation}

Hence the point-value two-point gradient is

\begin{equation}
\boxed{
\frac{c(r_E)-c(r_P)}{r_E-r_P}
=
c_r(r_f)
+
O(h^2)
}
\tag{TRISO-ACC-126}
\label{eq:triso-acc-126}
\end{equation}

under the local-symmetry condition (TRISO-ACC-123).

\subsubsection*{18A.4 Effect of using exact cell averages}

Write the exact cell averages as

\begin{equation}
C_P=c(r_P)+\eta_P,
\tag{TRISO-ACC-127}
\label{eq:triso-acc-127}
\end{equation}

and

\begin{equation}
C_E=c(r_E)+\eta_E,
\tag{TRISO-ACC-128}
\label{eq:triso-acc-128}
\end{equation}

where

\begin{equation}
\eta_P=O(h^2),
\qquad
\eta_E=O(h^2).
\tag{TRISO-ACC-129}
\label{eq:triso-acc-129}
\end{equation}

The numerical gradient is

\begin{equation}
\frac{C_E-C_P}{r_E-r_P}.
\tag{TRISO-ACC-130}
\label{eq:triso-acc-130}
\end{equation}

Substitute (TRISO-ACC-127) and (TRISO-ACC-128):

\begin{equation}
\frac{C_E-C_P}{r_E-r_P}
=
\frac{c(r_E)-c(r_P)}{r_E-r_P}
+
\frac{\eta_E-\eta_P}{r_E-r_P}.
\tag{TRISO-ACC-131}
\label{eq:triso-acc-131}
\end{equation}

For a smoothly varying family of shape-regular cells, the cell-average representation error varies smoothly between adjacent cells, so

\begin{equation}
\eta_E-\eta_P=O(h^3).
\tag{TRISO-ACC-132}
\label{eq:triso-acc-132}
\end{equation}

Since

\begin{equation}
r_E-r_P=O(h),
\tag{TRISO-ACC-133}
\label{eq:triso-acc-133}
\end{equation}

we obtain

\begin{equation}
\frac{\eta_E-\eta_P}{r_E-r_P}
=
O(h^2).
\tag{TRISO-ACC-134}
\label{eq:triso-acc-134}
\end{equation}

Combine (TRISO-ACC-126) and (TRISO-ACC-134):

\begin{equation}
\boxed{
\frac{C_E-C_P}{r_E-r_P}
=
c_r(r_f)
+
O(h^2).
}
\tag{TRISO-ACC-135}
\label{eq:triso-acc-135}
\end{equation}

This result is conditional on smooth solution data, shape-regular refinement, and the local geometric relation (TRISO-ACC-123).

\subsubsection*{18A.5 Same-material face-flux consistency}

Inside one material,

\begin{equation}
J_f=-D\,c_r(r_f).
\tag{TRISO-ACC-136}
\label{eq:triso-acc-136}
\end{equation}

The two-point numerical flux is

\begin{equation}
J_f^h
=
-D
\frac{C_E-C_P}{r_E-r_P}.
\tag{TRISO-ACC-137}
\label{eq:triso-acc-137}
\end{equation}

Use (TRISO-ACC-135):

\begin{equation}
J_f^h
=
-D
\left[
c_r(r_f)+O(h^2)
\right].
\tag{TRISO-ACC-138}
\label{eq:triso-acc-138}
\end{equation}

Because $D$ is constant and finite,

\begin{equation}
\boxed{
J_f^h
=
J_f
+
O(h^2).
}
\tag{TRISO-ACC-139}
\label{eq:triso-acc-139}
\end{equation}

Thus an ordinary smooth same-material face flux is second-order consistent under the stated mesh assumptions.

\subsubsection*{18A.6 Complete smooth-cell divergence consistency}

The exact diffusion contribution to the cell-average evolution is

\begin{equation}
\mathcal D_P
=
\frac{
A_wJ_w-A_eJ_e
}{
V_P
}.
\tag{TRISO-ACC-140}
\label{eq:triso-acc-140}
\end{equation}

The numerical diffusion contribution is

\begin{equation}
\mathcal D_P^h
=
\frac{
A_wJ_w^h-A_eJ_e^h
}{
V_P
}.
\tag{TRISO-ACC-141}
\label{eq:triso-acc-141}
\end{equation}

Define the face-flux errors

\begin{equation}
\varepsilon_w
=
J_w^h-J_w,
\tag{TRISO-ACC-142}
\label{eq:triso-acc-142}
\end{equation}

and

\begin{equation}
\varepsilon_e
=
J_e^h-J_e.
\tag{TRISO-ACC-143}
\label{eq:triso-acc-143}
\end{equation}

Subtract the exact cell diffusion term from the numerical one:

\begin{equation}
\mathcal D_P^h-\mathcal D_P
=
\frac{
A_w(J_w^h-J_w)
-
A_e(J_e^h-J_e)
}{
V_P
}.
\tag{TRISO-ACC-144}
\label{eq:triso-acc-144}
\end{equation}

Use the error definitions:

\begin{equation}
\boxed{
\mathcal D_P^h-\mathcal D_P
=
\frac{
A_w\varepsilon_w-A_e\varepsilon_e
}{
V_P}.
}
\tag{TRISO-ACC-145}
\label{eq:triso-acc-145}
\end{equation}

The face analysis alone gives

\begin{equation}
\varepsilon_w=O(h^2),
\qquad
\varepsilon_e=O(h^2).
\tag{TRISO-ACC-146}
\label{eq:triso-acc-146}
\end{equation}

Since

\begin{equation}
V_P=O(h)
\tag{TRISO-ACC-147}
\label{eq:triso-acc-147}
\end{equation}

for a refining shell away from pathological mesh degeneration, the estimate (TRISO-ACC-146) by itself would permit only

\begin{equation}
\mathcal D_P^h-\mathcal D_P=O(h).
\tag{TRISO-ACC-148}
\label{eq:triso-acc-148}
\end{equation}

Therefore second-order face consistency alone is insufficient to establish second-order cell-divergence consistency.

A cancellation property of the leading face errors is required.

\#\#\#\# 18A.6.1 Smooth leading face-error field

For a smooth solution and a smoothly varying, locally symmetric mesh family, assume the face-flux truncation error admits the local expansion

\begin{equation}
\boxed{
\varepsilon_f
=
h^2E(r_f)
+
O(h^3),
}
\tag{TRISO-ACC-149}
\label{eq:triso-acc-149}
\end{equation}

where $E(r)$ is smooth within the material.

This is stronger than the statement $\varepsilon_f=\(O(h^2)$.

At the west face,

\begin{equation}
\varepsilon_w
=
h^2E(r_w)
+
O(h^3).
\tag{TRISO-ACC-150}
\label{eq:triso-acc-150}
\end{equation}

At the east face,

\begin{equation}
\varepsilon_e
=
h^2E(r_e)
+
O(h^3).
\tag{TRISO-ACC-151}
\label{eq:triso-acc-151}
\end{equation}

Substitute into the numerator of (TRISO-ACC-145):

\begin{equation}
A_w\varepsilon_w-A_e\varepsilon_e
=
h^2
\left[
A_wE(r_w)-A_eE(r_e)
\right]
+
O(h^4),
\tag{TRISO-ACC-152}
\label{eq:triso-acc-152}
\end{equation}

where the remainder scaling assumes a shape-regular cell with bounded spherical face areas.

Define

\begin{equation}
F(r)=A(r)E(r),
\tag{TRISO-ACC-153}
\label{eq:triso-acc-153}
\end{equation}

with

\begin{equation}
A(r)=4\pi r^2.
\tag{TRISO-ACC-154}
\label{eq:triso-acc-154}
\end{equation}

Then

\begin{equation}
A_wE(r_w)-A_eE(r_e)
=
F(r_w)-F(r_e).
\tag{TRISO-ACC-155}
\label{eq:triso-acc-155}
\end{equation}

Taylor-expand $F(r_e)$ about $r_w$:

\begin{equation}
F(r_e)
=
F(r_w)
+
(r_e-r_w)F'(r_w)
+
O(h^2).
\tag{TRISO-ACC-156}
\label{eq:triso-acc-156}
\end{equation}

Since

\begin{equation}
r_e-r_w=h_P=O(h),
\tag{TRISO-ACC-157}
\label{eq:triso-acc-157}
\end{equation}

subtracting gives

\begin{equation}
F(r_w)-F(r_e)
=
-h_PF'(r_w)
+
O(h^2).
\tag{TRISO-ACC-158}
\label{eq:triso-acc-158}
\end{equation}

Therefore

\begin{equation}
F(r_w)-F(r_e)=O(h).
\tag{TRISO-ACC-159}
\label{eq:triso-acc-159}
\end{equation}

Return to (TRISO-ACC-152):

\begin{equation}
A_w\varepsilon_w-A_e\varepsilon_e
=
h^2O(h)
+
O(h^4).
\tag{TRISO-ACC-160}
\label{eq:triso-acc-160}
\end{equation}

Hence

\begin{equation}
\boxed{
A_w\varepsilon_w-A_e\varepsilon_e
=
O(h^3).
}
\tag{TRISO-ACC-161}
\label{eq:triso-acc-161}
\end{equation}

Divide by

\begin{equation}
V_P=O(h).
\tag{TRISO-ACC-162}
\label{eq:triso-acc-162}
\end{equation}

Then

\begin{equation}
\boxed{
\mathcal D_P^h-\mathcal D_P
=
O(h^2).
}
\tag{TRISO-ACC-163}
\label{eq:triso-acc-163}
\end{equation}

Thus the complete smooth same-material finite-volume diffusion operator is second-order consistent \textbf{provided} the leading face-flux error varies smoothly from face to face as in (TRISO-ACC-149).

\#\#\#\# 18A.6.2 Relation to the exact spherical differential operator

The exact integrated diffusion term is

\begin{equation}
\mathcal D_P
=
\frac1{V_P}
\int_{r_w}^{r_e}
4\pi
\frac{\partial}{\partial r}
\left(
r^2D c_r
\right)dr.
\tag{TRISO-ACC-164}
\label{eq:triso-acc-164}
\end{equation}

For a smooth integrand, the exact cell average of the differential operator differs from its value at the volume centroid by the same centroid-moment mechanism used in Section 18A.1.

Therefore

\begin{equation}
\mathcal D_P
=
\left[
\frac1{r^2}
\frac{\partial}{\partial r}
\left(
r^2D c_r
\right)
\right]_{r=r_P}
+
O(h^2).
\tag{TRISO-ACC-165}
\label{eq:triso-acc-165}
\end{equation}

Combine (TRISO-ACC-163) and (TRISO-ACC-165):

\begin{equation}
\boxed{
\mathcal D_P^h
=
\left[
\frac1{r^2}
\frac{\partial}{\partial r}
\left(
r^2D c_r
\right)
\right]_{r=r_P}
+
O(h^2).
}
\tag{TRISO-ACC-166}
\label{eq:triso-acc-166}
\end{equation}

This result applies only to smooth same-material cells under the stated mesh/error-regularity assumptions.

\#\#\#\# 18A.6.3 Why the cancellation assumption matters

If the leading face-error coefficient is not smooth across the cell, then

\begin{equation}
A_w\varepsilon_w-A_e\varepsilon_e
\tag{TRISO-ACC-167}
\label{eq:triso-acc-167}
\end{equation}

need not be $O(h^3)$.

It may remain only

\begin{equation}
O(h^2).
\tag{TRISO-ACC-168}
\label{eq:triso-acc-168}
\end{equation}

Division by $V_P=\(O(h)$ would then give only

\begin{equation}
O(h)
\tag{TRISO-ACC-169}
\label{eq:triso-acc-169}
\end{equation}

cell-divergence consistency.

This is precisely why the smooth-cell result cannot be transferred automatically across a discontinuous material interface.

\subsubsection*{18A.7 Discontinuous-diffusivity interface-face consistency}

Consider a physical material interface at

\begin{equation}
r=r_I.
\tag{TRISO-ACC-170}
\label{eq:triso-acc-170}
\end{equation}

Let cell $P$ lie immediately to the left of the interface and cell $E$ immediately to the right.

Define

\begin{equation}
d_P=r_I-r_P,
\tag{TRISO-ACC-171}
\label{eq:triso-acc-171}
\end{equation}

and

\begin{equation}
d_E=r_E-r_I.
\tag{TRISO-ACC-172}
\label{eq:triso-acc-172}
\end{equation}

Let the diffusivities be

\begin{equation}
D^- \quad\text{for }r<r_I,
\tag{TRISO-ACC-173}
\label{eq:triso-acc-173}
\end{equation}

and

\begin{equation}
D^+ \quad\text{for }r>r_I.
\tag{TRISO-ACC-174}
\label{eq:triso-acc-174}
\end{equation}

For the frozen ideal interface,

\begin{equation}
c^-(r_I)=c^+(r_I)=c_I.
\tag{TRISO-ACC-175}
\label{eq:triso-acc-175}
\end{equation}

The exact outward flux is continuous:

\begin{equation}
J_I
=
-D^-c_r^-(r_I)
=
-D^+c_r^+(r_I).
\tag{TRISO-ACC-176}
\label{eq:triso-acc-176}
\end{equation}

\subsubsection*{18A.7.1 Expand the left-side concentration}

Taylor-expand the exact point value at $r_P=r_I-d_P$ about the interface from the left:

\begin{equation}
c(r_P)
=
c_I
-
d_Pc_r^-(r_I)
+
\frac{d_P^2}{2}c_{rr}^-(r_I)
-
\frac{d_P^3}{6}c_{rrr}^-(r_I)
+
O(h^4).
\tag{TRISO-ACC-177}
\label{eq:triso-acc-177}
\end{equation}

Use

\begin{equation}
c_r^-(r_I)
=
-\frac{J_I}{D^-}.
\tag{TRISO-ACC-178}
\label{eq:triso-acc-178}
\end{equation}

Substitute:

\begin{equation}
c(r_P)
=
c_I
+
\frac{d_P}{D^-}J_I
+
\frac{d_P^2}{2}c_{rr}^-(r_I)
-
\frac{d_P^3}{6}c_{rrr}^-(r_I)
+
O(h^4).
\tag{TRISO-ACC-179}
\label{eq:triso-acc-179}
\end{equation}

Rearrange the interface-to-left concentration difference:

\begin{equation}
c(r_P)-c_I
=
\frac{d_P}{D^-}J_I
+
\frac{d_P^2}{2}c_{rr}^-(r_I)
-
\frac{d_P^3}{6}c_{rrr}^-(r_I)
+
O(h^4).
\tag{TRISO-ACC-180}
\label{eq:triso-acc-180}
\end{equation}

\subsubsection*{18A.7.2 Expand the right-side concentration}

Taylor-expand from the right:

\begin{equation}
c(r_E)
=
c_I
+
d_Ec_r^+(r_I)
+
\frac{d_E^2}{2}c_{rr}^+(r_I)
+
\frac{d_E^3}{6}c_{rrr}^+(r_I)
+
O(h^4).
\tag{TRISO-ACC-181}
\label{eq:triso-acc-181}
\end{equation}

Use

\begin{equation}
c_r^+(r_I)
=
-\frac{J_I}{D^+}.
\tag{TRISO-ACC-182}
\label{eq:triso-acc-182}
\end{equation}

Therefore

\begin{equation}
c(r_E)
=
c_I
-
\frac{d_E}{D^+}J_I
+
\frac{d_E^2}{2}c_{rr}^+(r_I)
+
\frac{d_E^3}{6}c_{rrr}^+(r_I)
+
O(h^4).
\tag{TRISO-ACC-183}
\label{eq:triso-acc-183}
\end{equation}

Rearrange:

\begin{equation}
c_I-c(r_E)
=
\frac{d_E}{D^+}J_I
-
\frac{d_E^2}{2}c_{rr}^+(r_I)
-
\frac{d_E^3}{6}c_{rrr}^+(r_I)
+
O(h^4).
\tag{TRISO-ACC-184}
\label{eq:triso-acc-184}
\end{equation}

\subsubsection*{18A.7.3 Add the two exact concentration drops}

Add (TRISO-ACC-180) and (TRISO-ACC-184):

\begin{equation}
c(r_P)-c(r_E)
=
J_I
\left(
\frac{d_P}{D^-}
+
\frac{d_E}{D^+}
\right)
\end{equation}

\begin{equation}
+
\frac{d_P^2}{2}c_{rr}^-(r_I)
-
\frac{d_E^2}{2}c_{rr}^+(r_I)
-
\frac{d_P^3}{6}c_{rrr}^-(r_I)
-
\frac{d_E^3}{6}c_{rrr}^+(r_I)
+
O(h^4).
\tag{TRISO-ACC-185}
\label{eq:triso-acc-185}
\end{equation}

Define the two-half-cell resistance

\begin{equation}
\boxed{
R_I
=
\frac{d_P}{D^-}
+
\frac{d_E}{D^+}.
}
\tag{TRISO-ACC-186}
\label{eq:triso-acc-186}
\end{equation}

For a shape-regular mesh,

\begin{equation}
R_I=O(h).
\tag{TRISO-ACC-187}
\label{eq:triso-acc-187}
\end{equation}

Define the second-order remainder

\begin{equation}
Q_I
=
\frac{d_P^2}{2}c_{rr}^-(r_I)
-
\frac{d_E^2}{2}c_{rr}^+(r_I).
\tag{TRISO-ACC-188}
\label{eq:triso-acc-188}
\end{equation}

Then

\begin{equation}
Q_I=O(h^2).
\tag{TRISO-ACC-189}
\label{eq:triso-acc-189}
\end{equation}

Equation (TRISO-ACC-185) becomes

\begin{equation}
c(r_P)-c(r_E)
=
J_IR_I
+
Q_I
+
O(h^3).
\tag{TRISO-ACC-190}
\label{eq:triso-acc-190}
\end{equation}

Solve for the exact flux:

\begin{equation}
J_I
=
\frac{
c(r_P)-c(r_E)-Q_I+O(h^3)
}{
R_I
}.
\tag{TRISO-ACC-191}
\label{eq:triso-acc-191}
\end{equation}

Separate the resistance formula:

\begin{equation}
J_I
=
\frac{
c(r_P)-c(r_E)
}{
R_I
}
-
\frac{Q_I}{R_I}
+
O(h^2).
\tag{TRISO-ACC-192}
\label{eq:triso-acc-192}
\end{equation}

Because

\begin{equation}
Q_I=O(h^2)
\tag{TRISO-ACC-193}
\label{eq:triso-acc-193}
\end{equation}

and

\begin{equation}
R_I=O(h),
\tag{TRISO-ACC-194}
\label{eq:triso-acc-194}
\end{equation}

we have

\begin{equation}
\frac{Q_I}{R_I}=O(h).
\tag{TRISO-ACC-195}
\label{eq:triso-acc-195}
\end{equation}

Therefore the point-value resistance flux

\begin{equation}
J_I^{h,\mathrm{pt}}
=
\frac{
c(r_P)-c(r_E)
}{
R_I
}
\tag{TRISO-ACC-196}
\label{eq:triso-acc-196}
\end{equation}

satisfies, generically,

\begin{equation}
\boxed{
J_I^{h,\mathrm{pt}}
=
J_I
+
O(h).
}
\tag{TRISO-ACC-197}
\label{eq:triso-acc-197}
\end{equation}

Thus the basic two-point harmonic/resistance interface flux is generically first-order accurate at a discontinuity.

\subsubsection*{18A.7.4 Special cancellation condition}

The leading $O(h)$ flux error vanishes if

\begin{equation}
Q_I=O(h^3).
\tag{TRISO-ACC-198}
\label{eq:triso-acc-198}
\end{equation}

At leading order this requires

\begin{equation}
d_P^2c_{rr}^-(r_I)
-
d_E^2c_{rr}^+(r_I)
=
O(h^3).
\tag{TRISO-ACC-199}
\label{eq:triso-acc-199}
\end{equation}

For equal half-distances,

\begin{equation}
d_P=d_E,
\tag{TRISO-ACC-200}
\label{eq:triso-acc-200}
\end{equation}

a sufficient leading-order cancellation condition is

\begin{equation}
c_{rr}^-(r_I)
=
c_{rr}^+(r_I).
\tag{TRISO-ACC-201}
\label{eq:triso-acc-201}
\end{equation}

Such equality is not generally implied by concentration continuity and flux continuity when $D^-\ne D^+$.

Therefore second-order interface flux accuracy must not be assumed merely because the harmonic resistance is physically conservative.

\subsubsection*{18A.7.5 Effect of exact cell averages}

The numerical scheme uses exact cell averages rather than exact point values.

Write

\begin{equation}
C_P=c(r_P)+\eta_P,
\tag{TRISO-ACC-202}
\label{eq:triso-acc-202}
\end{equation}

and

\begin{equation}
C_E=c(r_E)+\eta_E.
\tag{TRISO-ACC-203}
\label{eq:triso-acc-203}
\end{equation}

For smooth one-sided fields within each material,

\begin{equation}
\eta_P=O(h^2),
\qquad
\eta_E=O(h^2).
\tag{TRISO-ACC-204}
\label{eq:triso-acc-204}
\end{equation}

The numerical interface flux is

\begin{equation}
J_I^h
=
\frac{
C_P-C_E
}{
R_I
}.
\tag{TRISO-ACC-205}
\label{eq:triso-acc-205}
\end{equation}

Substitute the average representations:

\begin{equation}
J_I^h
=
\frac{
c(r_P)-c(r_E)
}{
R_I
}
+
\frac{
\eta_P-\eta_E
}{
R_I
}.
\tag{TRISO-ACC-206}
\label{eq:triso-acc-206}
\end{equation}

Across a material discontinuity, the leading $O(h^2)$ average-representation coefficients on the two sides need not match smoothly.

Therefore, generically,

\begin{equation}
\eta_P-\eta_E=O(h^2).
\tag{TRISO-ACC-207}
\label{eq:triso-acc-207}
\end{equation}

Since

\begin{equation}
R_I=O(h),
\tag{TRISO-ACC-208}
\label{eq:triso-acc-208}
\end{equation}

the cell-average correction contributes

\begin{equation}
\frac{\eta_P-\eta_E}{R_I}
=
O(h).
\tag{TRISO-ACC-209}
\label{eq:triso-acc-209}
\end{equation}

Combine this with (TRISO-ACC-197):

\begin{equation}
\boxed{
J_I^h
=
J_I
+
O(h)
}
\tag{TRISO-ACC-210}
\label{eq:triso-acc-210}
\end{equation}

generically for the canonical cell-average, two-point resistance interface flux.

\subsubsection*{18A.7.6 Conservation remains exact despite first-order local accuracy}

The same numerical interface flux $J_I^h$ is used by both adjacent control volumes.

The left cell contains the outward interface amount rate

\begin{equation}
-A_IJ_I^h.
\tag{TRISO-ACC-211}
\label{eq:triso-acc-211}
\end{equation}

The right cell contains the corresponding inward amount rate

\begin{equation}
+A_IJ_I^h.
\tag{TRISO-ACC-212}
\label{eq:triso-acc-212}
\end{equation}

Add them:

\begin{equation}
-A_IJ_I^h+A_IJ_I^h=0.
\tag{TRISO-ACC-213}
\label{eq:triso-acc-213}
\end{equation}

Therefore

\begin{equation}
\boxed{
\text{interface conservation is exact at the discrete level}
}
\tag{TRISO-ACC-214}
\label{eq:triso-acc-214}
\end{equation}

even though the local interface flux is generically only first-order accurate.

Conservation and formal order are separate properties.

\subsubsection*{18A.7.7 Consequence for cells adjacent to the interface}

Let the interface face-flux error be

\begin{equation}
\varepsilon_I=O(h).
\tag{TRISO-ACC-215}
\label{eq:triso-acc-215}
\end{equation}

The interface amount-rate error is

\begin{equation}
A_I\varepsilon_I.
\tag{TRISO-ACC-216}
\label{eq:triso-acc-216}
\end{equation}

For an interface away from the origin,

\begin{equation}
A_I=O(1)
\tag{TRISO-ACC-217}
\label{eq:triso-acc-217}
\end{equation}

under radial refinement of a fixed physical geometry.

The adjacent cell volume satisfies

\begin{equation}
V_P=O(h).
\tag{TRISO-ACC-218}
\label{eq:triso-acc-218}
\end{equation}

Therefore the contribution of the interface flux error to the adjacent cell-average time derivative can scale as

\begin{equation}
\frac{
A_I\varepsilon_I
}{
V_P
}
=
O(1).
\tag{TRISO-ACC-219}
\label{eq:triso-acc-219}
\end{equation}

This does \textbf{not} by itself prove that the global solution fails to converge.

It shows that a pointwise local truncation-error argument at the interface-adjacent cell is insufficient for establishing a global order.

A global stability-plus-consistency argument in an appropriate integrated norm, or direct grid-refinement evidence, is required.

\subsubsection*{18A.7.8 Interface consistency status}

\textbf{[VERIFIED]} The harmonic/resistance interface flux exactly enforces one common discrete flux and therefore exact discrete conservation.

\textbf{[VERIFIED]} For piecewise smooth solutions satisfying ideal concentration and flux continuity, the basic two-point point-value interface flux is generically $O(h)$ accurate.

\textbf{[VERIFIED]} Using exact cell averages at volume centroids does not generically improve that interface order; the canonical interface flux remains $O(h)$ unless additional cancellation occurs.

[NOT ESTABLISHED] A second-order interface flux for unequal diffusivities.

[NOT ESTABLISHED] A pointwise vanishing truncation error in the cells directly adjacent to a discontinuous interface.

[NOT ESTABLISHED] The global spatial convergence order of the conservative scheme.

\subsubsection*{18A.8 Cell-centred Robin boundary consistency}

Consider the outermost OPyC cell with representative coordinate

\begin{equation}
r_P=R-d,
\tag{TRISO-ACC-220}
\label{eq:triso-acc-220}
\end{equation}

where

\begin{equation}
d=\delta r_R=O(h).
\tag{TRISO-ACC-221}
\label{eq:triso-acc-221}
\end{equation}

Let

\begin{equation}
c_R=c(R)
\tag{TRISO-ACC-222}
\label{eq:triso-acc-222}
\end{equation}

denote the exact physical surface concentration.

The exact Robin condition is

\begin{equation}
J_R=h(c_R-c_\infty).
\tag{TRISO-ACC-223}
\label{eq:triso-acc-223}
\end{equation}

The same exact outward flux also satisfies

\begin{equation}
J_R=-D_5c_r(R).
\tag{TRISO-ACC-224}
\label{eq:triso-acc-224}
\end{equation}

The numerical cell-centred closure is

\begin{equation}
J_R^h
=
\frac{
C_P-c_\infty
}{
\dfrac d{D_5}+\dfrac1h
}.
\tag{TRISO-ACC-225}
\label{eq:triso-acc-225}
\end{equation}

The objective is to compare (TRISO-ACC-225) with the exact $J_R$.

\#\#\#\# 18A.8.1 Exact point-value expansion from the surface to the cell representative point

Taylor-expand the exact OPyC solution from $R$ inward to $r_P=R-d$:

\begin{equation}
c(r_P)
=
c_R
-
dc_r(R)
+
\frac{d^2}{2}c_{rr}(R)
-
\frac{d^3}{6}c_{rrr}(R)
+
O(h^4).
\tag{TRISO-ACC-226}
\label{eq:triso-acc-226}
\end{equation}

Use the exact flux relation

\begin{equation}
c_r(R)
=
-\frac{J_R}{D_5}.
\tag{TRISO-ACC-227}
\label{eq:triso-acc-227}
\end{equation}

Substitute:

\begin{equation}
c(r_P)
=
c_R
+
\frac d{D_5}J_R
+
\frac{d^2}{2}c_{rr}(R)
-
\frac{d^3}{6}c_{rrr}(R)
+
O(h^4).
\tag{TRISO-ACC-228}
\label{eq:triso-acc-228}
\end{equation}

The Robin law gives

\begin{equation}
c_R-c_\infty
=
\frac{J_R}{h}.
\tag{TRISO-ACC-229}
\label{eq:triso-acc-229}
\end{equation}

Subtract $c_\infty$ from (TRISO-ACC-228):

\begin{equation}
c(r_P)-c_\infty
=
(c_R-c_\infty)
+
\frac d{D_5}J_R
+
\frac{d^2}{2}c_{rr}(R)
-
\frac{d^3}{6}c_{rrr}(R)
+
O(h^4).
\tag{TRISO-ACC-230}
\label{eq:triso-acc-230}
\end{equation}

Substitute (TRISO-ACC-229):

\begin{equation}
c(r_P)-c_\infty
=
J_R
\left(
\frac1h+\frac d{D_5}
\right)
+
\frac{d^2}{2}c_{rr}(R)
-
\frac{d^3}{6}c_{rrr}(R)
+
O(h^4).
\tag{TRISO-ACC-231}
\label{eq:triso-acc-231}
\end{equation}

Define the exact cell-centre-to-bulk resistance

\begin{equation}
R_B
=
\frac1h+\frac d{D_5}.
\tag{TRISO-ACC-232}
\label{eq:triso-acc-232}
\end{equation}

For fixed finite $h>0$,

\begin{equation}
R_B=O(1)
\tag{TRISO-ACC-233}
\label{eq:triso-acc-233}
\end{equation}

as $h\to0$ in the mesh-refinement sense $d\to0$; here $h$ is the physical transfer coefficient and is held fixed.

Equation (TRISO-ACC-231) becomes

\begin{equation}
c(r_P)-c_\infty
=
J_RR_B
+
\frac{d^2}{2}c_{rr}(R)
+
O(h^3).
\tag{TRISO-ACC-234}
\label{eq:triso-acc-234}
\end{equation}

Solve for $J_R$:

\begin{equation}
J_R
=
\frac{
c(r_P)-c_\infty
}{
R_B
}
-
\frac{
d^2c_{rr}(R)
}{
2R_B
}
+
O(h^3).
\tag{TRISO-ACC-235}
\label{eq:triso-acc-235}
\end{equation}

Therefore the point-value series-resistance flux

\begin{equation}
J_R^{h,\mathrm{pt}}
=
\frac{
c(r_P)-c_\infty
}{
R_B
}
\tag{TRISO-ACC-236}
\label{eq:triso-acc-236}
\end{equation}

satisfies

\begin{equation}
\boxed{
J_R^{h,\mathrm{pt}}
=
J_R+O(h^2)
}
\tag{TRISO-ACC-237}
\label{eq:triso-acc-237}
\end{equation}

for fixed finite $h>0$.

\#\#\#\# 18A.8.2 Effect of the exact cell average

The numerical closure uses the exact cell average $C_P$, not $c(r_P)$.

Write

\begin{equation}
C_P=c(r_P)+\eta_P.
\tag{TRISO-ACC-238}
\label{eq:triso-acc-238}
\end{equation}

From the centroid analysis,

\begin{equation}
\eta_P=O(h^2).
\tag{TRISO-ACC-239}
\label{eq:triso-acc-239}
\end{equation}

Substitute into the numerical boundary flux:

\begin{equation}
J_R^h
=
\frac{
c(r_P)+\eta_P-c_\infty
}{
R_B
}.
\tag{TRISO-ACC-240}
\label{eq:triso-acc-240}
\end{equation}

Separate the point-value part:

\begin{equation}
J_R^h
=
\frac{
c(r_P)-c_\infty
}{
R_B
}
+
\frac{\eta_P}{R_B}.
\tag{TRISO-ACC-241}
\label{eq:triso-acc-241}
\end{equation}

For fixed finite $h>0$,

\begin{equation}
R_B=O(1).
\tag{TRISO-ACC-242}
\label{eq:triso-acc-242}
\end{equation}

Therefore

\begin{equation}
\frac{\eta_P}{R_B}=O(h^2).
\tag{TRISO-ACC-243}
\label{eq:triso-acc-243}
\end{equation}

Combine with (TRISO-ACC-237):

\begin{equation}
\boxed{
J_R^h
=
J_R+O(h^2)
}
\tag{TRISO-ACC-244}
\label{eq:triso-acc-244}
\end{equation}

for the finite-transfer Robin boundary under smooth OPyC data and fixed physical $h$.

\subsubsection*{18A.8.3 Boundary amount-rate consistency}

The exact outer area is

\begin{equation}
A_R=4\pi R^2.
\tag{TRISO-ACC-245}
\label{eq:triso-acc-245}
\end{equation}

The exact outward amount rate is

\begin{equation}
\dot N_R=A_RJ_R.
\tag{TRISO-ACC-246}
\label{eq:triso-acc-246}
\end{equation}

The numerical amount rate is

\begin{equation}
\dot N_R^h=A_RJ_R^h.
\tag{TRISO-ACC-247}
\label{eq:triso-acc-247}
\end{equation}

Subtract:

\begin{equation}
\dot N_R^h-\dot N_R
=
A_R(J_R^h-J_R).
\tag{TRISO-ACC-248}
\label{eq:triso-acc-248}
\end{equation}

Because $A_R$ is fixed under mesh refinement and (TRISO-ACC-244) gives $J_R^h-J_R=\(O(h^2)$,

\begin{equation}
\boxed{
\dot N_R^h-\dot N_R
=
O(h^2).
}
\tag{TRISO-ACC-249}
\label{eq:triso-acc-249}
\end{equation}

Thus the total Robin release rate is second-order consistent for fixed finite $h$.

\subsubsection*{18A.8.4 Outer-cell local residual scaling}

The outer-cell volume satisfies

\begin{equation}
V_P=O(h).
\tag{TRISO-ACC-250}
\label{eq:triso-acc-250}
\end{equation}

If the boundary amount-rate error is

\begin{equation}
O(h^2),
\tag{TRISO-ACC-251}
\label{eq:triso-acc-251}
\end{equation}

then its contribution to the outer cell-average time-derivative residual can scale as

\begin{equation}
\frac{O(h^2)}{O(h)}
=
O(h).
\tag{TRISO-ACC-252}
\label{eq:triso-acc-252}
\end{equation}

Therefore the outermost cell can have only first-order pointwise local truncation consistency even though the physical boundary release rate itself is second-order accurate.

This is analogous to the distinction already identified at material interfaces: a lower pointwise residual in $O(1)$ special cells does not by itself determine the global solution convergence order.

\subsubsection*{18A.8.5 Neumann and Dirichlet limiting regimes}

For the insulating limit

\begin{equation}
h=0,
\tag{TRISO-ACC-253}
\label{eq:triso-acc-253}
\end{equation}

the exact boundary condition is

\begin{equation}
J_R=0.
\tag{TRISO-ACC-254}
\label{eq:triso-acc-254}
\end{equation}

The resistance formula is interpreted by its limit

\begin{equation}
\frac1h\to\infty,
\tag{TRISO-ACC-255}
\label{eq:triso-acc-255}
\end{equation}

which gives

\begin{equation}
J_R^h\to0.
\tag{TRISO-ACC-256}
\label{eq:triso-acc-256}
\end{equation}

Thus the no-flux boundary is represented exactly as a limiting boundary law.

The absorbing Dirichlet limit is different.

If

\begin{equation}
h\to\infty,
\tag{TRISO-ACC-257}
\label{eq:triso-acc-257}
\end{equation}

then

\begin{equation}
R_B
=
\frac d{D_5}.
\tag{TRISO-ACC-258}
\label{eq:triso-acc-258}
\end{equation}

Now

\begin{equation}
R_B=O(h)
\tag{TRISO-ACC-259}
\label{eq:triso-acc-259}
\end{equation}

where here $h$ denotes the mesh-size order symbol, not the transfer coefficient.

To avoid this notational collision, denote the mesh scale by $\mathfrak h$.

Then

\begin{equation}
d=O(\mathfrak h),
\tag{TRISO-ACC-260}
\label{eq:triso-acc-260}
\end{equation}

and

\begin{equation}
R_B=O(\mathfrak h).
\tag{TRISO-ACC-261}
\label{eq:triso-acc-261}
\end{equation}

The $O(\mathfrak h^2)$ cell-average representation error divided by $R_B=\(O(\mathfrak h)$ can contribute

\begin{equation}
O(\mathfrak h)
\tag{TRISO-ACC-262}
\label{eq:triso-acc-262}
\end{equation}

to the Dirichlet-limit boundary flux.

Therefore the finite-$h$ second-order flux result (TRISO-ACC-244) must \textbf{not} be transferred automatically to the absorbing Dirichlet limit.

For a true absorbing boundary, a separate Dirichlet-boundary consistency analysis is required.

\subsubsection*{18A.8.6 Robin boundary consistency status}

\textbf{[VERIFIED]} For fixed finite physical transfer coefficient $0<h<\infty$, the cell-centred series-resistance Robin flux is $O(\mathfrak h^2)$ consistent under smooth OPyC data and the centroid representation assumptions.

\textbf{[VERIFIED]} The total outer release amount rate is also $O(\mathfrak h^2)$ consistent.

\textbf{[VERIFIED]} The outer-cell pointwise time-derivative residual may be only $O(\mathfrak h)$ because the boundary amount-rate error is divided by a cell volume $O(\mathfrak h)$.

\textbf{[VERIFIED]} The no-flux Neumann limit is recovered.

[NOT ESTABLISHED] Second-order boundary-flux accuracy in the absorbing Dirichlet limit $h\to\infty$.

[NOT ESTABLISHED] Global spatial convergence order.

\subsubsection*{18A.9 Global stability and conditional convergence framework}

The local consistency results are not uniform over the particle:

\noindent\textbullet\ smooth same-material cells have conditional $O(\mathfrak h^2)$ divergence consistency;\\
\noindent\textbullet\ a fixed number of cells adjacent to the four material interfaces can have $O(1)$ pointwise residuals under the present two-point transmission flux;\\
\noindent\textbullet\ the outer Robin cell can have an $O(\mathfrak h)$ pointwise residual for fixed finite physical $h$.\\

A global convergence argument must therefore use a norm that respects cell volumes rather than taking the maximum pointwise residual as the only consistency measure.

\#\#\#\# 18A.9.1 Volume-weighted discrete norm

Define the discrete volume-weighted norm

\begin{equation}
\boxed{
\|\mathbf x\|_V^2
=
\mathbf x^T\mathbf V\mathbf x
=
\sum_{P=0}^{M-1}
V_Px_P^2.
}
\tag{TRISO-ACC-263}
\label{eq:triso-acc-263}
\end{equation}

This is the natural discrete analogue of the spherical $L^2$ norm because

\begin{equation}
V_P
=
\int_{\Omega_P}dV.
\tag{TRISO-ACC-264}
\label{eq:triso-acc-264}
\end{equation}

For a shape-regular radial refinement of a fixed particle,

\begin{equation}
V_P=O(\mathfrak h)
\tag{TRISO-ACC-265}
\label{eq:triso-acc-265}
\end{equation}

for cells away from the origin.

The central cell is smaller:

\begin{equation}
V_0=O(\mathfrak h^3),
\tag{TRISO-ACC-266}
\label{eq:triso-acc-266}
\end{equation}

because its radius is itself $O(\mathfrak h)$.

\subsubsection*{18A.9.2 Semi-discrete error equation}

Let

\begin{equation}
\overline{\mathbf c}(t)
\tag{TRISO-ACC-267}
\label{eq:triso-acc-267}
\end{equation}

denote the vector of exact spherical cell averages of the continuum solution on the numerical mesh.

Define the semi-discrete residual $\boldsymbol\tau_h(t)$ by inserting these exact cell averages into the numerical operator:

\begin{equation}
\boxed{
\mathbf V
\frac{d\overline{\mathbf c}}{dt}
=
\mathbf K\overline{\mathbf c}
+
\mathbf V\mathbf S
+
\mathbf b_\infty
+
\mathbf V\boldsymbol\tau_h.
}
\tag{TRISO-ACC-268}
\label{eq:triso-acc-268}
\end{equation}

The numerical semi-discrete solution satisfies

\begin{equation}
\mathbf V
\frac{d\mathbf C}{dt}
=
\mathbf K\mathbf C
+
\mathbf V\mathbf S
+
\mathbf b_\infty.
\tag{TRISO-ACC-269}
\label{eq:triso-acc-269}
\end{equation}

Define the error

\begin{equation}
\boxed{
\mathbf e
=
\mathbf C-\overline{\mathbf c}.
}
\tag{TRISO-ACC-270}
\label{eq:triso-acc-270}
\end{equation}

Subtract (TRISO-ACC-268) from (TRISO-ACC-269):

\begin{equation}
\mathbf V
\frac{d\mathbf e}{dt}
=
\mathbf K\mathbf e
-
\mathbf V\boldsymbol\tau_h.
\tag{TRISO-ACC-271}
\label{eq:triso-acc-271}
\end{equation}

\subsubsection*{18A.9.3 Energy identity}

Premultiply by $\mathbf e^T$:

\begin{equation}
\mathbf e^T\mathbf V
\frac{d\mathbf e}{dt}
=
\mathbf e^T\mathbf K\mathbf e
-
\mathbf e^T\mathbf V\boldsymbol\tau_h.
\tag{TRISO-ACC-272}
\label{eq:triso-acc-272}
\end{equation}

Because $\mathbf V$ is time independent,

\begin{equation}
\mathbf e^T\mathbf V
\frac{d\mathbf e}{dt}
=
\frac12
\frac{d}{dt}
\left(
\mathbf e^T\mathbf V\mathbf e
\right).
\tag{TRISO-ACC-273}
\label{eq:triso-acc-273}
\end{equation}

Use the norm definition:

\begin{equation}
\boxed{
\frac12
\frac{d}{dt}
\|\mathbf e\|_V^2
=
\mathbf e^T\mathbf K\mathbf e
-
\langle\mathbf e,\boldsymbol\tau_h\rangle_V.
}
\tag{TRISO-ACC-274}
\label{eq:triso-acc-274}
\end{equation}

From the previously proved conductance identity,

\begin{equation}
\mathbf e^T\mathbf K\mathbf e
\le0.
\tag{TRISO-ACC-275}
\label{eq:triso-acc-275}
\end{equation}

Therefore

\begin{equation}
\frac12
\frac{d}{dt}
\|\mathbf e\|_V^2
\le
-
\langle\mathbf e,\boldsymbol\tau_h\rangle_V.
\tag{TRISO-ACC-276}
\label{eq:triso-acc-276}
\end{equation}

Take absolute value of the forcing term:

\begin{equation}
-\langle\mathbf e,\boldsymbol\tau_h\rangle_V
\le
\left|
\langle\mathbf e,\boldsymbol\tau_h\rangle_V
\right|.
\tag{TRISO-ACC-277}
\label{eq:triso-acc-277}
\end{equation}

Apply Cauchy-Schwarz in the $V$-inner product:

\begin{equation}
\left|
\langle\mathbf e,\boldsymbol\tau_h\rangle_V
\right|
\le
\|\mathbf e\|_V
\|\boldsymbol\tau_h\|_V.
\tag{TRISO-ACC-278}
\label{eq:triso-acc-278}
\end{equation}

Hence

\begin{equation}
\frac12
\frac{d}{dt}
\|\mathbf e\|_V^2
\le
\|\mathbf e\|_V
\|\boldsymbol\tau_h\|_V.
\tag{TRISO-ACC-279}
\label{eq:triso-acc-279}
\end{equation}

For $\|\mathbf e\|_V>0$,

\begin{equation}
\frac{d}{dt}
\|\mathbf e\|_V^2
=
2\|\mathbf e\|_V
\frac{d}{dt}\|\mathbf e\|_V.
\tag{TRISO-ACC-280}
\label{eq:triso-acc-280}
\end{equation}

Substitute into (TRISO-ACC-279):

\begin{equation}
\|\mathbf e\|_V
\frac{d}{dt}\|\mathbf e\|_V
\le
\|\mathbf e\|_V
\|\boldsymbol\tau_h\|_V.
\tag{TRISO-ACC-281}
\label{eq:triso-acc-281}
\end{equation}

Cancel $\|\mathbf e\|_V$:

\begin{equation}
\boxed{
\frac{d}{dt}\|\mathbf e\|_V
\le
\|\boldsymbol\tau_h\|_V.
}
\tag{TRISO-ACC-282}
\label{eq:triso-acc-282}
\end{equation}

The same inequality follows by continuity through instants at which the error norm is zero.

Integrate from $0$ to $t$:

\begin{equation}
\|\mathbf e(t)\|_V-\|\mathbf e(0)\|_V
\le
\int_0^t
\|\boldsymbol\tau_h(s)\|_V\,ds.
\tag{TRISO-ACC-283}
\label{eq:triso-acc-283}
\end{equation}

Therefore

\begin{equation}
\boxed{
\|\mathbf e(t)\|_V
\le
\|\mathbf e(0)\|_V
+
\int_0^t
\|\boldsymbol\tau_h(s)\|_V\,ds.
}
\tag{TRISO-ACC-284}
\label{eq:triso-acc-284}
\end{equation}

This proves semi-discrete energy stability and shows that convergence follows if the volume-weighted residual norm tends to zero and the initial discrete representation converges.

\subsubsection*{18A.9.4 Residual scaling from the established local results}

Assume the number of material interfaces remains fixed at four as the mesh is refined.

For $O(\mathfrak h^{-1})$ ordinary smooth cells,

\begin{equation}
\tau_P=O(\mathfrak h^2).
\tag{TRISO-ACC-285}
\label{eq:triso-acc-285}
\end{equation}

Each such cell contributes to the squared $V$-norm

\begin{equation}
V_P\tau_P^2
=
O(\mathfrak h)
O(\mathfrak h^4)
=
O(\mathfrak h^5).
\tag{TRISO-ACC-286}
\label{eq:triso-acc-286}
\end{equation}

Summing $O(\mathfrak h^{-1})$ smooth cells gives

\begin{equation}
O(\mathfrak h^{-1})
O(\mathfrak h^5)
=
O(\mathfrak h^4).
\tag{TRISO-ACC-287}
\label{eq:triso-acc-287}
\end{equation}

Therefore the smooth-region contribution to the residual norm is

\begin{equation}
O(\mathfrak h^2).
\tag{TRISO-ACC-288}
\label{eq:triso-acc-288}
\end{equation}

Now consider the fixed number of interface-adjacent cells.

The local analysis permits

\begin{equation}
\tau_P=O(1).
\tag{TRISO-ACC-289}
\label{eq:triso-acc-289}
\end{equation}

Each such cell has

\begin{equation}
V_P=O(\mathfrak h).
\tag{TRISO-ACC-290}
\label{eq:triso-acc-290}
\end{equation}

Therefore each contributes

\begin{equation}
V_P\tau_P^2
=
O(\mathfrak h).
\tag{TRISO-ACC-291}
\label{eq:triso-acc-291}
\end{equation}

A fixed number of such cells still contributes

\begin{equation}
O(\mathfrak h)
\tag{TRISO-ACC-292}
\label{eq:triso-acc-292}
\end{equation}

to the squared residual norm.

Hence the interface-region contribution can be only

\begin{equation}
\boxed{
O(\mathfrak h^{1/2})
}
\tag{TRISO-ACC-293}
\label{eq:triso-acc-293}
\end{equation}

in the $V$-norm under the currently proved local bounds.

For the single outer Robin cell,

\begin{equation}
\tau_{M-1}=O(\mathfrak h).
\tag{TRISO-ACC-294}
\label{eq:triso-acc-294}
\end{equation}

Its volume is

\begin{equation}
V_{M-1}=O(\mathfrak h).
\tag{TRISO-ACC-295}
\label{eq:triso-acc-295}
\end{equation}

Therefore its squared contribution is

\begin{equation}
V_{M-1}\tau_{M-1}^2
=
O(\mathfrak h^3).
\tag{TRISO-ACC-296}
\label{eq:triso-acc-296}
\end{equation}

and its contribution to the $V$-norm is

\begin{equation}
O(\mathfrak h^{3/2}).
\tag{TRISO-ACC-297}
\label{eq:triso-acc-297}
\end{equation}

The conservative interface region therefore dominates the currently provable residual estimate:

\begin{equation}
\boxed{
\|\boldsymbol\tau_h\|_V
=
O(\mathfrak h^{1/2})
}
\tag{TRISO-ACC-298}
\label{eq:triso-acc-298}
\end{equation}

under the local bounds established so far.

\subsubsection*{18A.9.5 Conditional semi-discrete convergence bound}

If the exact initial condition is represented by exact cell averages, then

\begin{equation}
\mathbf e(0)=\mathbf0.
\tag{TRISO-ACC-299}
\label{eq:triso-acc-299}
\end{equation}

Use (TRISO-ACC-284):

\begin{equation}
\|\mathbf e(t)\|_V
\le
\int_0^t
\|\boldsymbol\tau_h(s)\|_V\,ds.
\tag{TRISO-ACC-300}
\label{eq:triso-acc-300}
\end{equation}

If the residual bound is uniform for $0\le s\le T$,

\begin{equation}
\|\boldsymbol\tau_h(s)\|_V
\le
C_T\mathfrak h^{1/2},
\tag{TRISO-ACC-301}
\label{eq:triso-acc-301}
\end{equation}

then

\begin{equation}
\|\mathbf e(t)\|_V
\le
\int_0^t
C_T\mathfrak h^{1/2}\,ds.
\tag{TRISO-ACC-302}
\label{eq:triso-acc-302}
\end{equation}

Evaluate the integral:

\begin{equation}
\boxed{
\|\mathbf e(t)\|_V
\le
tC_T\mathfrak h^{1/2},
\qquad
0\le t\le T.
}
\tag{TRISO-ACC-303}
\label{eq:triso-acc-303}
\end{equation}

Thus the present energy argument supports at least a \textbf{conditional $O(\mathfrak h^{1/2})$ upper-bound convergence rate in the volume-weighted norm}, given the established local residual estimates and sufficient regularity.

This is a conservative bound, not a prediction of the observed numerical order.

Diffusive smoothing, transmission structure, cancellation, or a sharper negative-norm/interface analysis may yield a higher actual rate.

\subsubsection*{18A.9.6 Fully discrete temporal error}

The explicit Euler method has local temporal truncation error

\begin{equation}
O(\Delta t^2)
\tag{TRISO-ACC-304}
\label{eq:triso-acc-304}
\end{equation}

per step.

Over a fixed time interval, under stability, its global temporal order is

\begin{equation}
O(\Delta t).
\tag{TRISO-ACC-305}
\label{eq:triso-acc-305}
\end{equation}

Therefore the fully discrete error should be separated conceptually as

\begin{equation}
\boxed{
\text{error}
=
\text{spatial error}
+
O(\Delta t),
}
\tag{TRISO-ACC-306}
\label{eq:triso-acc-306}
\end{equation}

rather than inferring spatial order from a refinement study in which $\Delta t$ is not reduced sufficiently.

The previously derived explicit positivity/stability condition requires

\begin{equation}
\Delta t
\le
\Delta t_{\max}(h).
\tag{TRISO-ACC-307}
\label{eq:triso-acc-307}
\end{equation}

For diffusion on a regular mesh,

\begin{equation}
\Delta t_{\max}=O(\mathfrak h^2)
\tag{TRISO-ACC-308}
\label{eq:triso-acc-308}
\end{equation}

away from pathological coefficient scaling.

Choosing

\begin{equation}
\Delta t\propto\mathfrak h^2
\tag{TRISO-ACC-309}
\label{eq:triso-acc-309}
\end{equation}

therefore both respects the expected explicit-diffusion stability scaling and makes the first-order temporal error

\begin{equation}
O(\Delta t)=O(\mathfrak h^2).
\tag{TRISO-ACC-310}
\label{eq:triso-acc-310}
\end{equation}

Under that refinement path, temporal error should not dominate a spatial rate lower than second order.

\subsubsection*{18A.9.7 What the proof does not establish}

The energy estimate proves stability of the semi-discrete error equation in the $V$-norm.

It does not establish that the $O(\mathfrak h^{1/2})$ bound is sharp.

It does not establish an $L^\infty$ convergence rate.

It does not prove second-order global convergence.

It does not replace numerical grid refinement.

The four material interfaces are a measure-shrinking set under refinement, but their stronger local residuals require either a sharper transmission-problem estimate or empirical convergence evidence before a better global rate is claimed.

\subsubsection*{18A.10 Accuracy/convergence study specification}

The subsequent numerical study should distinguish three questions.

\#\#\#\# Spatial refinement

Use a sequence of interface-aligned meshes with refinement ratio

\begin{equation}
q=\frac{h_k}{h_{k+1}}>1.
\tag{TRISO-ACC-311}
\label{eq:triso-acc-311}
\end{equation}

Keep the physical geometry, diffusivities, source, $h$, and final observation time fixed.

Choose time steps satisfying

\begin{equation}
\Delta t_k
=
C h_k^2
\tag{TRISO-ACC-312}
\label{eq:triso-acc-312}
\end{equation}

with $C$ below the explicit stability bound on every mesh.

Compare numerical solutions using at least:

1. the volume-weighted concentration norm;
2. total particle inventory;
3. outer release rate;
4. selected interface-adjacent concentrations.

When an exact analytical benchmark is available, define

\begin{equation}
E_k
=
\|\mathbf C_{h_k}-\overline{\mathbf c}_{\mathrm{exact},h_k}\|_V.
\tag{TRISO-ACC-313}
\label{eq:triso-acc-313}
\end{equation}

The observed order between successive meshes is

\begin{equation}
\boxed{
p_{\mathrm{obs}}
=
\frac{
\log(E_k/E_{k+1})
}{
\log q
}.
}
\tag{TRISO-ACC-314}
\label{eq:triso-acc-314}
\end{equation}

No target value of $p_{\mathrm{obs}}$ is assumed in advance for the discontinuous-$D$ five-layer problem.

\#\#\#\# Temporal refinement

On a sufficiently fine fixed spatial mesh, refine

\begin{equation}
\Delta t_k
\to
\frac{\Delta t_k}{q_t}.
\tag{TRISO-ACC-315}
\label{eq:triso-acc-315}
\end{equation}

The expected explicit-Euler temporal order is one:

\begin{equation}
p_t\approx1
\tag{TRISO-ACC-316}
\label{eq:triso-acc-316}
\end{equation}

once spatial error is subdominant.

\#\#\#\# Conservation refinement

For every run, evaluate the discrete inventory residual

\begin{equation}
\mathcal R_{\mathrm{cons}}
=
\frac{dN_h}{dt}
-
\sum_PS_PV_P
+
G_R(C_{M-1}-c_\infty).
\tag{TRISO-ACC-317}
\label{eq:triso-acc-317}
\end{equation}

For the algebraically conservative semi-discrete formulation this should vanish up to time-discretisation and floating-point effects when evaluated consistently.

The study must report the norm definition, mesh geometry, interface alignment, time-step scaling, and whether the comparison uses exact cell averages or point samples.

\subsubsection*{18A.11 Executed refinement evidence}

The canonical refinement specification has now been executed with the standalone Rust driver \path{verification/fv_convergence.rs} through GitHub Actions run \texttt{36654437525}.

The persisted raw results are \path{verification/fv_convergence_results.txt}.

The benchmark uses five aligned unit-thickness layers,

\begin{equation}
(r_1,r_2,r_3,r_4,R)=(1,2,3,4,5).
\tag{TRISO-ACC-318}
\label{eq:triso-acc-318}
\end{equation}

with

\begin{equation}
(D_1,D_2,D_3,D_4,D_5)=(1,0.5,2,0.25,1.5),
\tag{TRISO-ACC-319}
\label{eq:triso-acc-319}
\end{equation}

and

\begin{equation}
S_0=1,\qquad h=0.8,\qquad c_\infty=0.
\tag{TRISO-ACC-320}
\label{eq:triso-acc-320}
\end{equation}

These are normalized verification parameters, not claimed physical TRISO material data.

\#\#\#\# Spatial refinement result

The volume-weighted steady errors were

\begin{equation}
E_{25}=1.832317223290\times10^{-2},
\tag{TRISO-ACC-321}
\label{eq:triso-acc-321}
\end{equation}

\begin{equation}
E_{50}=4.621922061067\times10^{-3},
\tag{TRISO-ACC-322}
\label{eq:triso-acc-322}
\end{equation}

\begin{equation}
E_{100}=1.158126274643\times10^{-3},
\tag{TRISO-ACC-323}
\label{eq:triso-acc-323}
\end{equation}

\begin{equation}
E_{200}=2.896983800097\times10^{-4},
\tag{TRISO-ACC-324}
\label{eq:triso-acc-324}
\end{equation}

\begin{equation}
E_{400}=7.243504635747\times10^{-5}.
\tag{TRISO-ACC-325}
\label{eq:triso-acc-325}
\end{equation}

The corresponding observed orders were

\begin{equation}
1.987104,\quad1.996700,\quad1.999169,\quad1.999792.
\tag{TRISO-ACC-326}
\label{eq:triso-acc-326}
\end{equation}

Therefore this aligned five-layer steady benchmark exhibits asymptotic behavior consistent with

\begin{equation}
\boxed{E_h=O(\mathfrak h^2).}
\tag{TRISO-ACC-327}
\label{eq:triso-acc-327}
\end{equation}

This observed second-order behavior is substantially sharper than the conservative analytical $O(h^{1/2})$ bound in TRISO-ACC-303. It demonstrates that the local interface-adjacent residual estimate is not predictive of the observed global steady error for this benchmark. It does not prove second-order convergence for every discontinuous-D problem.

\#\#\#\# Temporal refinement result

On the fixed N=200 spatial mesh, successive-step Richardson differences gave observed temporal orders

\begin{equation}
0.999700,\quad0.999780,\quad1.000066,\quad0.999962.
\tag{TRISO-ACC-328}
\label{eq:triso-acc-328}
\end{equation}

Therefore the executed result is consistent with

\begin{equation}
\boxed{p_t=1.}
\tag{TRISO-ACC-329}
\label{eq:triso-acc-329}
\end{equation}

\#\#\#\# Conservation result

The exact benchmark generation rate is

\begin{equation}
\frac{4\pi}{3}=4.188790204786\ldots
\tag{TRISO-ACC-330}
\label{eq:triso-acc-330}
\end{equation}

Across the steady spatial sequence, the computed release rate remained equal to this value to approximately $10^{-12}$ or better.

The steady conservation residual ranged from approximately

\begin{equation}
2.7\times10^{-15}
\tag{TRISO-ACC-331}
\label{eq:triso-acc-331}
\end{equation}

to

\begin{equation}
2.9\times10^{-12}.
\tag{TRISO-ACC-332}
\label{eq:triso-acc-332}
\end{equation}

During the temporal study, the maximum discrete inventory residual remained below

\begin{equation}
\boxed{4\times10^{-10}.}
\tag{TRISO-ACC-333}
\label{eq:triso-acc-333}
\end{equation}

The increase in the reported absolute residual as the time step becomes very small is consistent with floating-point cancellation in the finite-difference evaluation of the inventory time derivative; no systematic conservation drift is observed in this study.

\#\#\#\# Scope of the executed evidence

[EXECUTED / VERIFIED] The canonical aligned five-layer FV benchmark converges approximately second order in the steady volume-weighted spatial norm over the tested meshes.

[EXECUTED / VERIFIED] Forward Euler converges first order in time on the tested fixed spatial mesh.

[EXECUTED / VERIFIED] The conservative FV balance closes to floating-point accuracy in the tested steady and transient runs.

[NOT ESTABLISHED] The same spatial order for non-interface-aligned meshes.

[NOT ESTABLISHED] The same spatial order for arbitrary diffusivity ratios, geometries, or source distributions.

[NOT ESTABLISHED] The absorbing Dirichlet boundary convergence order.

[NOT APPLICABLE] These deterministic FV results do not validate the production WOS algorithm.
\subsubsection*{18A.12 What has and has not been proved}

\textbf{[VERIFIED]} The exact spherical cell average differs from the point value at the spherical volume centroid by $O(h^2)$ for a smooth field.

[CONDITIONALLY VERIFIED] The two-point same-material face gradient is $O(h^2)$ consistent when the refining mesh is shape regular, adjacent cell-average representation errors vary smoothly, and the face is locally centred between representative coordinates to $O(h^2)$.

[CONDITIONALLY VERIFIED] Under those same assumptions, the same-material diffusive face flux is $O(h^2)$ consistent.

[CONDITIONALLY VERIFIED] The complete smooth same-material cell divergence is $O(h^2)$ consistent when the leading $O(h^2)$ face-flux error has a smooth coefficient across neighbouring faces.

\textbf{[VERIFIED]} The canonical resistance-weighted interface flux is conservative and generically $O(h)$ accurate for piecewise-smooth unequal-D transmission data; second-order interface accuracy is not established.

\textbf{[VERIFIED]} For fixed finite physical h, the cell-centred Robin FV boundary flux and total release rate are $O(mesh^2)$ consistent; the outer-cell pointwise residual may remain $O(mesh)$.

[CONDITIONALLY ESTABLISHED] The energy argument gives a conservative $O(mesh^1/2)$ bound from local residual estimates. [EXECUTED] The aligned five-layer steady benchmark instead exhibits asymptotic $O(mesh^2)$ volume-weighted convergence over N=25–400; this observed rate is benchmark-specific rather than a universal theorem.

The prescribed aligned five-layer spatial/temporal/conservation refinement study has now executed successfully. The next verification decision is whether to broaden the parameter/interface-alignment study or submit this deterministic accuracy/convergence milestone for independent audit; no WOS conclusion follows from these FV results.

\subsection*{19. Implementation provenance}

\textbf{[IMPORTANT]} The supervisor production repository remains read-only. The production WOS source provenance used by this research is recorded here for reproducibility; all research verification code and manuscript changes remain in Ray's repository.

\noindent\textbullet\ Production geometry: \path{constructive_solid_geometry/mod.rs} — \texttt{TrisoCell::new}, \texttt{TrisoCell::new\_crp6\_geometry}, \texttt{TrisoCell::get\_triso\_region}, \texttt{TrisoCell::try\_get\_diffusion\_coefficient}.\\
\noindent\textbullet\ Production WOS: \path{first_passage/walk_on_spheres.rs} — \texttt{WoSWalker}, \texttt{step\_multilayer}, \texttt{walk\_until\_released}, \texttt{nearest\_interface\_distance}, \texttt{shell\_bounds}, \texttt{sample\_uniform\_in\_ball}.\\
\noindent\textbullet\ Production interface decision: \path{first_passage/interface.rs} — \texttt{does\_transmit}.\\
\noindent\textbullet\ Production homogeneous sphere FPT: \path{first_passage/sphere_fpt.rs} — first-passage distribution and lookup/interpolation.\\
\noindent\textbullet\ Historical production verification: \path{verification_and_validation/crp6_case1_kernel_release_vs_crank.md}.\\
\noindent\textbullet\ Historical interface-overshoot analysis: \path{docs/buffer_clt_failure_analysis.md}.\\

The research-side equation-to-code function map is maintained in \path{ray to outram park/TRACEABILITY.md}.

\subsection*{20. Foundation gate}

All 33 original displayed equations are now accounted for in the central register.

The mathematical middle is now continuous from physical problem through conservation, Fickian transport, spherical reduction, source/decay convention, conditions, Part-I steady and transient analysis, Part-II steady analytical solution, Part-II transient eigenvalue framework, and Part-I FTCS.

Remaining foundation gaps are explicitly retained:

\noindent\textbullet\ physical source/decay/trapping closure;\\
\noindent\textbullet\ physical partition/interfacial-resistance choice;\\
\noindent\textbullet\ five-layer transient modal coefficient convergence;\\
\noindent\textbullet\ complete FTCS spectral stability proof;\\
\noindent\textbullet\ complete equation-to-code verification;\\
\noindent\textbullet\ WOS-to-continuum transient verification.\\

The independent continuous-mathematics audit has passed the project for discretisation, and the independent deterministic discrete-mathematics audit has passed the canonical finite-volume model for accuracy/convergence study. Review finding R2-D01 remains open specifically for the retained FTCS Robin benchmark until grid-refinement evidence is produced; it does not block the canonical FV accuracy/convergence track.

\section*{21. Method 2 stochastic diffusion and first-passage formulation}

\subsection*{21.1 Diffusion process and backward generator}

\textbf{[EXACT]} In one homogeneous material with constant diffusivity $D$,

\begin{equation}
\frac{\partial c}{\partial t}=D\nabla^2c.
\tag{TRISO-FPT-001}
\label{eq:triso-fpt-001}
\end{equation}

[DEFINITION] Let $X_t$ be the diffusion process generated by the same operator.

Its infinitesimal generator is

\begin{equation}
\mathcal L f=D\nabla^2f.
\tag{TRISO-FPT-002}
\label{eq:triso-fpt-002}
\end{equation}

[DEFINITION] Let $T$ be the first exit time from a prescribed region.

Define

\begin{equation}
u(\mathbf x,s)
=
\mathbb E_{\mathbf x}\left[e^{-sT}g(X_T)\right].
\tag{TRISO-FPT-003}
\label{eq:triso-fpt-003}
\end{equation}

The backward first-exit equation is

\begin{equation}
\mathcal Lu=su.
\tag{TRISO-FPT-004}
\label{eq:triso-fpt-004}
\end{equation}

Substituting Eq.~\eqref{eq:triso-fpt-002} into Eq.~\eqref{eq:triso-fpt-004} gives:

\begin{equation}
D\nabla^2u=su.
\tag{TRISO-FPT-005}
\label{eq:triso-fpt-005}
\end{equation}

\textbf{[ASSUMPTION]} Under spherical symmetry, $u=u(r,s)$.

The radial Laplacian is

\begin{equation}
\nabla^2u
=
\frac{1}{r^2}\frac{d}{dr}\left(r^2\frac{du}{dr}\right).
\tag{TRISO-FPT-006}
\label{eq:triso-fpt-006}
\end{equation}

Substituting the radial Laplacian into Eq.~\eqref{eq:triso-fpt-005} gives:

\begin{equation}
D\frac{1}{r^2}\frac{d}{dr}(r^2u')=su.
\tag{TRISO-FPT-007}
\label{eq:triso-fpt-007}
\end{equation}

Apply the product rule:

\begin{equation}
\frac{d}{dr}(r^2u')=2ru'+r^2u''.
\tag{TRISO-FPT-008}
\label{eq:triso-fpt-008}
\end{equation}

Divide by $r^2$:

\begin{equation}
\frac{1}{r^2}\frac{d}{dr}(r^2u')
=
u''+\frac{2}{r}u'.
\tag{TRISO-FPT-009}
\label{eq:triso-fpt-009}
\end{equation}

Hence

\begin{equation}
\boxed{
D\left(u''+\frac{2}{r}u'\right)=su.
}
\tag{TRISO-FPT-010}
\label{eq:triso-fpt-010}
\end{equation}

This is the backward equation associated with the same diffusion operator as the forward continuum model.

\subsection*{21.2 Uniform-volume initial position}

\textbf{[ASSUMPTION]} Production histories are initially uniform in kernel volume.

Let $U\sim\mathcal U(0,1)$.

The enclosed-volume fraction at radius $r$ is

\begin{equation}
U
=
\frac{(4\pi/3)r^3}{(4\pi/3)R_1^3}.
\tag{TRISO-WOS-001}
\label{eq:triso-wos-001}
\end{equation}

Cancel $4\pi/3$:

\begin{equation}
U=\frac{r^3}{R_1^3}.
\tag{TRISO-WOS-002}
\label{eq:triso-wos-002}
\end{equation}

Multiply by $R_1^3$:

\begin{equation}
r^3=R_1^3U.
\tag{TRISO-WOS-003}
\label{eq:triso-wos-003}
\end{equation}

Take the real cube root:

\begin{equation}
\boxed{r=R_1U^{1/3}.}
\tag{TRISO-WOS-004}
\label{eq:triso-wos-004}
\end{equation}

\subsection*{21.3 Production interface law}

\textbf{[INFERRED FROM CODE]} For the frozen $K=1$ interface rule,

\begin{equation}
\boxed{
p_{i\rightarrow j}
=
\frac{D_j}{D_i+D_j}.
}
\tag{TRISO-WOS-005}
\label{eq:triso-wos-005}
\end{equation}

The reflection probability is

\begin{equation}
p_{i\rightarrow i}
=
1-p_{i\rightarrow j}.
\tag{TRISO-WOS-006}
\label{eq:triso-wos-006}
\end{equation}

Substituting Eq.~\eqref{eq:triso-wos-005} into the reflection identity gives:

\begin{equation}
p_{i\rightarrow i}
=
1-\frac{D_j}{D_i+D_j}.
\tag{TRISO-WOS-007}
\label{eq:triso-wos-007}
\end{equation}

Write the unit term over the same denominator:

\begin{equation}
p_{i\rightarrow i}
=
\frac{D_i+D_j}{D_i+D_j}
-
\frac{D_j}{D_i+D_j}.
\tag{TRISO-WOS-008}
\label{eq:triso-wos-008}
\end{equation}

Subtract the numerators:

\begin{equation}
\boxed{
p_{i\rightarrow i}
=
\frac{D_i}{D_i+D_j}.
}
\tag{TRISO-WOS-009}
\label{eq:triso-wos-009}
\end{equation}

[DEFINITION] The finite capture distance is $\epsilon$.

[DEFINITION] The reinsertion distance is

\begin{equation}
\boxed{\delta=\alpha\epsilon.}
\tag{TRISO-WOS-010}
\label{eq:triso-wos-010}
\end{equation}

\textbf{[INFERRED FROM CODE]} The interface decision contributes zero physical time.

\textbf{[IMPORTANT]} These finite-capture semantics define Process A and remain distinct from the exact-interface process below.

\subsection*{21.4 Centred-ball first-passage transform}

\textbf{[EXACT]} Consider a homogeneous ball of radius $b$ with absorbing boundary at $r=b$.

Define

\begin{equation}
H(r,s)=\mathbb E_r[e^{-sT_b}].
\tag{TRISO-FPT-011}
\label{eq:triso-fpt-011}
\end{equation}

The backward radial equation is

\begin{equation}
D\left(H''+\frac{2}{r}H'\right)=sH.
\tag{TRISO-FPT-012}
\label{eq:triso-fpt-012}
\end{equation}

Introduce

\begin{equation}
v(r)=rH(r,s).
\tag{TRISO-FPT-013}
\label{eq:triso-fpt-013}
\end{equation}

The same derivative cancellation used for the shell gives

\begin{equation}
v''-\lambda^2v=0,
\qquad
\lambda=\sqrt{\frac{s}{D}}.
\tag{TRISO-FPT-014}
\label{eq:triso-fpt-014}
\end{equation}

The general solution is

\begin{equation}
v(r)=A\sinh(\lambda r)+B\cosh(\lambda r).
\tag{TRISO-FPT-015}
\label{eq:triso-fpt-015}
\end{equation}

Regularity of $H=v/r$ at $r=0$ requires

\begin{equation}
v(0)=0.
\tag{TRISO-FPT-016}
\label{eq:triso-fpt-016}
\end{equation}

Substitute $r=0$ into TRISO-FPT-015:

\begin{equation}
0=A\sinh0+B\cosh0.
\tag{TRISO-FPT-017}
\label{eq:triso-fpt-017}
\end{equation}

Therefore

\begin{equation}
B=0.
\tag{TRISO-FPT-018}
\label{eq:triso-fpt-018}
\end{equation}

At the absorbing sphere,

\begin{equation}
H(b,s)=1.
\tag{TRISO-FPT-019}
\label{eq:triso-fpt-019}
\end{equation}

Hence

\begin{equation}
v(b)=b.
\tag{TRISO-FPT-019A}
\label{eq:triso-fpt-019a}
\end{equation}

Using $v(b)=A\sinh(\lambda b)$,

\begin{equation}
A=\frac{b}{\sinh(\lambda b)}.
\tag{TRISO-FPT-019B}
\label{eq:triso-fpt-019b}
\end{equation}

Therefore

\begin{equation}
v(r)=b\frac{\sinh(\lambda r)}{\sinh(\lambda b)}.
\tag{TRISO-FPT-019C}
\label{eq:triso-fpt-019c}
\end{equation}

Divide by $r$:

\begin{equation}
\boxed{
H(r,s)=
\frac{b}{r}
\frac{\sinh(\lambda r)}{\sinh(\lambda b)}.
}
\tag{TRISO-FPT-019D}
\label{eq:triso-fpt-019d}
\end{equation}

At the centre, use $\sinh(\lambda r)\sim\lambda r$:

\begin{equation}
H(0,s)
=
\frac{b\lambda}{\sinh(\lambda b)}.
\tag{TRISO-FPT-019E}
\label{eq:triso-fpt-019e}
\end{equation}

Thus the centred-ball kernel used by the accelerated renewal is derived from the same backward diffusion equation, with regularity at the origin and absorption at the ball surface.

\section*{22. Exact spherical-shell first-passage theory}

\subsection*{22.1 Outer-exit joint transform}

[DEFINITION] Consider a homogeneous shell

\begin{equation}
a<r<b.
\tag{TRISO-FPT-020}
\label{eq:triso-fpt-020}
\end{equation}

Define the first exit time $T$ and the outer-exit joint transform

\begin{equation}
G_b(r,s)
=
\mathbb E_r\left[e^{-sT}\mathbf 1_{\{R_T=b\}}\right].
\tag{TRISO-FPT-021}
\label{eq:triso-fpt-021}
\end{equation}

TRISO-FPT-010 gives

\begin{equation}
D\left(G_b''+\frac{2}{r}G_b'\right)=sG_b.
\tag{TRISO-FPT-022}
\label{eq:triso-fpt-022}
\end{equation}

The outer-exit boundary conditions are

\begin{equation}
G_b(a,s)=0,
\tag{TRISO-FPT-023}
\label{eq:triso-fpt-023}
\end{equation}

\begin{equation}
G_b(b,s)=1.
\tag{TRISO-FPT-024}
\label{eq:triso-fpt-024}
\end{equation}

Introduce

\begin{equation}
v(r)=rG_b(r,s).
\tag{TRISO-FPT-025}
\label{eq:triso-fpt-025}
\end{equation}

Then

\begin{equation}
G_b=\frac{v}{r}.
\tag{TRISO-FPT-026}
\label{eq:triso-fpt-026}
\end{equation}

Differentiate:

\begin{equation}
G_b'
=
\frac{v'}{r}
-
\frac{v}{r^2}.
\tag{TRISO-FPT-027}
\label{eq:triso-fpt-027}
\end{equation}

Differentiate again:

\begin{equation}
G_b''
=
\frac{v''}{r}
-
\frac{2v'}{r^2}
+
\frac{2v}{r^3}.
\tag{TRISO-FPT-028}
\label{eq:triso-fpt-028}
\end{equation}

Also,

\begin{equation}
\frac{2}{r}G_b'
=
\frac{2v'}{r^2}
-
\frac{2v}{r^3}.
\tag{TRISO-FPT-029}
\label{eq:triso-fpt-029}
\end{equation}

Add TRISO-FPT-028 and TRISO-FPT-029:

\begin{equation}
G_b''+\frac{2}{r}G_b'
=
\frac{v''}{r}.
\tag{TRISO-FPT-030}
\label{eq:triso-fpt-030}
\end{equation}

Substituting the transformed derivative into Eq.~\eqref{eq:triso-fpt-022} gives:

\begin{equation}
D\frac{v''}{r}
=
s\frac{v}{r}.
\tag{TRISO-FPT-031}
\label{eq:triso-fpt-031}
\end{equation}

Multiply by $r$:

\begin{equation}
Dv''=sv.
\tag{TRISO-FPT-032}
\label{eq:triso-fpt-032}
\end{equation}

Divide by $D$:

\begin{equation}
v''=\frac{s}{D}v.
\tag{TRISO-FPT-033}
\label{eq:triso-fpt-033}
\end{equation}

Define

\begin{equation}
\lambda=\sqrt{\frac{s}{D}}.
\tag{TRISO-FPT-034}
\label{eq:triso-fpt-034}
\end{equation}

Then

\begin{equation}
v''-\lambda^2v=0.
\tag{TRISO-FPT-035}
\label{eq:triso-fpt-035}
\end{equation}

A convenient general solution measured from $a$ is

\begin{equation}
v(r)
=
A\sinh[\lambda(r-a)]
+
B\cosh[\lambda(r-a)].
\tag{TRISO-FPT-036}
\label{eq:triso-fpt-036}
\end{equation}

At $r=a$,

\begin{equation}
v(a)=aG_b(a,s)=0.
\tag{TRISO-FPT-037}
\label{eq:triso-fpt-037}
\end{equation}

Substitute $r=a$ into TRISO-FPT-036:

\begin{equation}
0=A\sinh0+B\cosh0.
\tag{TRISO-FPT-038}
\label{eq:triso-fpt-038}
\end{equation}

Use $sinh0=0$ and $cosh0=1$:

\begin{equation}
B=0.
\tag{TRISO-FPT-039}
\label{eq:triso-fpt-039}
\end{equation}

Thus

\begin{equation}
v(r)=A\sinh[\lambda(r-a)].
\tag{TRISO-FPT-040}
\label{eq:triso-fpt-040}
\end{equation}

At $r=b$,

\begin{equation}
v(b)=bG_b(b,s)=b.
\tag{TRISO-FPT-041}
\label{eq:triso-fpt-041}
\end{equation}

Substituting Eq.~\eqref{eq:triso-fpt-040} at $r=b$ gives:

\begin{equation}
b=A\sinh[\lambda(b-a)].
\tag{TRISO-FPT-042}
\label{eq:triso-fpt-042}
\end{equation}

Divide:

\begin{equation}
A
=
\frac{b}{\sinh[\lambda(b-a)]}.
\tag{TRISO-FPT-043}
\label{eq:triso-fpt-043}
\end{equation}

Insert this coefficient into TRISO-FPT-040:

\begin{equation}
v(r)
=
b
\frac{\sinh[\lambda(r-a)]}
{\sinh[\lambda(b-a)]}.
\tag{TRISO-FPT-044}
\label{eq:triso-fpt-044}
\end{equation}

Use $G_b=v/r$:

\begin{equation}
\boxed{
G_b(r,s)
=
\frac{b}{r}
\frac{\sinh[\lambda(r-a)]}
{\sinh[\lambda(b-a)]}.
}
\tag{TRISO-FPT-045}
\label{eq:triso-fpt-045}
\end{equation}

\subsection*{22.2 Inner-exit joint transform}

Define

\begin{equation}
G_a(r,s)
=
\mathbb E_r\left[e^{-sT}\mathbf 1_{\{R_T=a\}}\right].
\tag{TRISO-FPT-046}
\label{eq:triso-fpt-046}
\end{equation}

The boundary values are

\begin{equation}
G_a(a,s)=1,
\tag{TRISO-FPT-047}
\label{eq:triso-fpt-047}
\end{equation}

\begin{equation}
G_a(b,s)=0.
\tag{TRISO-FPT-048}
\label{eq:triso-fpt-048}
\end{equation}

The same substitution $v=rG_a$ gives

\begin{equation}
v''-\lambda^2v=0.
\tag{TRISO-FPT-049}
\label{eq:triso-fpt-049}
\end{equation}

Choose a form that satisfies the zero condition at $b$:

\begin{equation}
v(r)=C\sinh[\lambda(b-r)].
\tag{TRISO-FPT-050}
\label{eq:triso-fpt-050}
\end{equation}

At $r=a$,

\begin{equation}
v(a)=aG_a(a,s)=a.
\tag{TRISO-FPT-051}
\label{eq:triso-fpt-051}
\end{equation}

Therefore

\begin{equation}
a=C\sinh[\lambda(b-a)].
\tag{TRISO-FPT-052}
\label{eq:triso-fpt-052}
\end{equation}

Divide:

\begin{equation}
C
=
\frac{a}{\sinh[\lambda(b-a)]}.
\tag{TRISO-FPT-053}
\label{eq:triso-fpt-053}
\end{equation}

Substitute:

\begin{equation}
v(r)
=
a
\frac{\sinh[\lambda(b-r)]}
{\sinh[\lambda(b-a)]}.
\tag{TRISO-FPT-054}
\label{eq:triso-fpt-054}
\end{equation}

Divide by $r$:

\begin{equation}
\boxed{
G_a(r,s)
=
\frac{a}{r}
\frac{\sinh[\lambda(b-r)]}
{\sinh[\lambda(b-a)]}.
}
\tag{TRISO-FPT-055}
\label{eq:triso-fpt-055}
\end{equation}

\subsection*{22.3 Exit probabilities and conditional time}

Use

\begin{equation}
\sinh z\sim z
\qquad(z\rightarrow0).
\tag{TRISO-FPT-056}
\label{eq:triso-fpt-056}
\end{equation}

As $s\rightarrow0$, $lambda\rightarrow0$.

Apply TRISO-FPT-056 to TRISO-FPT-045:

\begin{equation}
G_b(r,0)
=
\frac br
\frac{\lambda(r-a)}
{\lambda(b-a)}.
\tag{TRISO-FPT-057}
\label{eq:triso-fpt-057}
\end{equation}

Cancel $lambda$:

\begin{equation}
\boxed{
P_r(R_T=b)
=
\frac{b(r-a)}{r(b-a)}.
}
\tag{TRISO-FPT-058}
\label{eq:triso-fpt-058}
\end{equation}

Likewise,

\begin{equation}
G_a(r,0)
=
\frac ar
\frac{\lambda(b-r)}
{\lambda(b-a)}.
\tag{TRISO-FPT-059}
\label{eq:triso-fpt-059}
\end{equation}

Cancel $lambda$:

\begin{equation}
\boxed{
P_r(R_T=a)
=
\frac{a(b-r)}{r(b-a)}.
}
\tag{TRISO-FPT-060}
\label{eq:triso-fpt-060}
\end{equation}

The conditional outer-exit transform is

\begin{equation}
\boxed{
\mathbb E_r[e^{-sT}\mid R_T=b]
=
\frac{G_b(r,s)}{G_b(r,0)}.
}
\tag{TRISO-FPT-061}
\label{eq:triso-fpt-061}
\end{equation}

The conditional inner-exit transform is

\begin{equation}
\boxed{
\mathbb E_r[e^{-sT}\mid R_T=a]
=
\frac{G_a(r,s)}{G_a(r,0)}.
}
\tag{TRISO-FPT-062}
\label{eq:triso-fpt-062}
\end{equation}

Differentiate the outer joint transform:

\begin{equation}
\frac{\partial G_b}{\partial s}
=
\mathbb E[-Te^{-sT}\mathbf1_b].
\tag{TRISO-FPT-063}
\label{eq:triso-fpt-063}
\end{equation}

Set $s=0$:

\begin{equation}
\left.\frac{\partial G_b}{\partial s}\right|_{s=0}
=
-\mathbb E[T\mathbf1_b].
\tag{TRISO-FPT-064}
\label{eq:triso-fpt-064}
\end{equation}

Hence

\begin{equation}
\boxed{
\mathbb E[T\mathbf1_b]
=
-
\left.\frac{\partial G_b}{\partial s}\right|_{s=0}.
}
\tag{TRISO-FPT-065}
\label{eq:triso-fpt-065}
\end{equation}

\textbf{[VERIFIED]} The shell exit probabilities, conditional moments and conditional CDF were checked against analytical references and direct WOS before multilayer coupling.


\section*{23. Accelerated exact-interface renewal}

[DEFINITION] Process B replaces repeated homogeneous-region WOS wandering by exact first-passage events while retaining the frozen stochastic interface law.

\textbf{[IMPORTANT]} Process A stops when the distance to an interface is at most $\epsilon$. Process B stops at the physical interface. The two stopping times are not identical.

\subsection*{23.1 Initial finite-capture mass}

For uniform births in an inner sphere of radius $a$, the total sphere volume is

\begin{equation}
V_a=\frac{4\pi}{3}a^3.
\tag{TRISO-WOS-020}
\label{eq:triso-wos-020}
\end{equation}

The volume inside radius $a-\epsilon$ is

\begin{equation}
V_{a-\epsilon}
=
\frac{4\pi}{3}(a-\epsilon)^3.
\tag{TRISO-WOS-021}
\label{eq:triso-wos-021}
\end{equation}

The capture-shell volume is

\begin{equation}
V_{\rm cap}
=
V_a-V_{a-\epsilon}.
\tag{TRISO-WOS-022}
\label{eq:triso-wos-022}
\end{equation}

The capture-shell probability is

\begin{equation}
P_{\rm cap}
=
\frac{V_{\rm cap}}{V_a}.
\tag{TRISO-WOS-023}
\label{eq:triso-wos-023}
\end{equation}

Substituting Eqs. (TRISO-WOS-020)–(TRISO-WOS-022) gives:

\begin{equation}
P_{\rm cap}
=
\frac{a^3-(a-\epsilon)^3}{a^3}.
\tag{TRISO-WOS-024}
\label{eq:triso-wos-024}
\end{equation}

Divide by $a^3$:

\begin{equation}
\boxed{
P_{\rm cap}
=
1-\left(1-\frac{\epsilon}{a}\right)^3.
}
\tag{TRISO-WOS-025}
\label{eq:triso-wos-025}
\end{equation}

For $a=50\,\mu\mathrm m$ and $epsilon=0.1\,\mu\mathrm m$,

\begin{equation}
\frac{\epsilon}{a}=0.002.
\tag{TRISO-WOS-026}
\label{eq:triso-wos-026}
\end{equation}

Therefore

\begin{equation}
P_{\rm cap}=1-(0.998)^3.
\tag{TRISO-WOS-027}
\label{eq:triso-wos-027}
\end{equation}

Numerically,

\begin{equation}
\boxed{P_{\rm cap}=0.005988008.}
\tag{TRISO-WOS-028}
\label{eq:triso-wos-028}
\end{equation}

\textbf{[EMPIRICAL COMPATIBILITY]} The executed Process-A/Process-B controlled discrepancy is small at the declared finite-(epsilon) statistical precision, but it is not an identity.

\subsection*{23.2 Two-layer accelerated verification}

\textbf{[VERIFIED]} The corrected two-layer accelerator uses the exact centred-ball kernel in the inner region and the exact shell kernel in the outer region.

The executed accelerated/FV RMS CDF difference is

\begin{equation}
\boxed{
\mathrm{RMS}(F_B-F_{FV})
=
0.002895.
}
\tag{TRISO-VER-200}
\label{eq:triso-ver-200}
\end{equation}

The maximum absolute difference is

\begin{equation}
\boxed{
\max_t|F_B-F_{FV}|
=
0.005215.
}
\tag{TRISO-VER-201}
\label{eq:triso-ver-201}
\end{equation}

The direct 100-nm WOS benchmark required approximately (1207.2) production steps per history.

The accelerated benchmark required approximately (139.04) renewal events per history.

The event-count ratio is

\begin{equation}
\frac{1207.2}{139.04}
=
8.68.
\tag{TRISO-VER-202}
\label{eq:triso-ver-202}
\end{equation}

Thus the controlled event-count reduction is approximately

\begin{equation}
\boxed{8.68\times.}
\tag{TRISO-VER-203}
\label{eq:triso-ver-203}
\end{equation}

\section*{24. Five-layer computational pathology}

\textbf{[VERIFIED]} Direct production WOS produced 8/8 censored histories at $10^6$ steps per history in the targeted five-layer diagnostic.

\textbf{[VERIFIED]} The exact-shell accelerated five-layer diagnostic produced 0/16 releases and 16/16 capped histories at 100000 renewals per history.

The renewal counts by layer were

\begin{equation}
(N_K,N_B,N_I,N_S,N_O)
=
(24,1599974,2,0,0).
\tag{TRISO-WOS-030}
\label{eq:triso-wos-030}
\end{equation}

The total number of renewals was

\begin{equation}
N_{\rm tot}=1600000.
\tag{TRISO-WOS-031}
\label{eq:triso-wos-031}
\end{equation}

The Buffer renewal fraction is

\begin{equation}
f_B
=
\frac{1599974}{1600000}.
\tag{TRISO-WOS-032}
\label{eq:triso-wos-032}
\end{equation}

Numerically,

\begin{equation}
\boxed{
f_B=0.99998375.
}
\tag{TRISO-WOS-033}
\label{eq:triso-wos-033}
\end{equation}

\textbf{[VERIFIED]} Removing homogeneous-region wandering did not remove the computational pathology.

[INFERRED] The dominant remaining cost is repeated rare interface-state recurrence.

\section*{25. Interface-state Markov-renewal reduction}

\subsection*{25.1 Eight transient states}

There are four physical internal interfaces.

Each interface has two post-interface material sides.

Therefore

\begin{equation}
N_{\rm states}
=
4\times2.
\tag{TRISO-MR-001}
\label{eq:triso-mr-001}
\end{equation}

Hence

\begin{equation}
\boxed{N_{\rm states}=8.}
\tag{TRISO-MR-002}
\label{eq:triso-mr-002}
\end{equation}

Define

\begin{equation}
S_0=\text{Kernel side of Kernel/Buffer}.
\tag{TRISO-MR-003}
\label{eq:triso-mr-003}
\end{equation}

\begin{equation}
S_1=\text{Buffer side of Kernel/Buffer}.
\tag{TRISO-MR-004}
\label{eq:triso-mr-004}
\end{equation}

\begin{equation}
S_2=\text{Buffer side of Buffer/IPyC}.
\tag{TRISO-MR-005}
\label{eq:triso-mr-005}
\end{equation}

\begin{equation}
S_3=\text{IPyC side of Buffer/IPyC}.
\tag{TRISO-MR-006}
\label{eq:triso-mr-006}
\end{equation}

\begin{equation}
S_4=\text{IPyC side of IPyC/SiC}.
\tag{TRISO-MR-007}
\label{eq:triso-mr-007}
\end{equation}

\begin{equation}
S_5=\text{SiC side of IPyC/SiC}.
\tag{TRISO-MR-008}
\label{eq:triso-mr-008}
\end{equation}

\begin{equation}
S_6=\text{SiC side of SiC/OPyC}.
\tag{TRISO-MR-009}
\label{eq:triso-mr-009}
\end{equation}

\begin{equation}
S_7=\text{OPyC side of SiC/OPyC}.
\tag{TRISO-MR-010}
\label{eq:triso-mr-010}
\end{equation}

The exterior release state is absorbing.

\subsection*{25.2 First-step equation from the Kernel-side state}

Define

\begin{equation}
\Phi_i(s)
=
\mathbb E_i[e^{-sT_{\rm rel}}].
\tag{TRISO-MR-011}
\label{eq:triso-mr-011}
\end{equation}

From $S_0$, the exact centred-ball transform to the Kernel/Buffer interface is $H_K(s)$.

At that interface, reflection returns to $S_0$.

Transmission enters $S_1$.

Thus

\begin{equation}
\Phi_0
=
H_Kp_{K\rightarrow K}\Phi_0
+
H_Kp_{K\rightarrow B}\Phi_1.
\tag{TRISO-MR-012}
\label{eq:triso-mr-012}
\end{equation}

Identify

\begin{equation}
K_{00}
=
H_Kp_{K\rightarrow K}.
\tag{TRISO-MR-013}
\label{eq:triso-mr-013}
\end{equation}

Identify

\begin{equation}
K_{01}
=
H_Kp_{K\rightarrow B}.
\tag{TRISO-MR-014}
\label{eq:triso-mr-014}
\end{equation}

All other row-0 entries are zero.

\subsection*{25.3 First-step equation from Buffer side of the inner interface}

From $S_1$, the Buffer shell can first exit inward with transform $G_B^-$.

It can first exit outward with transform $G_B^+$.

If it exits inward, transmission enters Kernel state $S_0$.

If it exits inward and reflects, it returns to Buffer state $S_1$.

If it exits outward and reflects, it moves to Buffer state $S_2$.

If it exits outward and transmits, it enters IPyC state $S_3$.

Therefore

\begin{equation}
\Phi_1
=
G_B^-p_{B\rightarrow K}\Phi_0
+
G_B^-p_{B\rightarrow B}^{(I_0)}\Phi_1
+
G_B^+p_{B\rightarrow B}^{(I_1)}\Phi_2
+
G_B^+p_{B\rightarrow I}\Phi_3.
\tag{TRISO-MR-015}
\label{eq:triso-mr-015}
\end{equation}

The first coefficient is

\begin{equation}
K_{10}
=
G_B^-p_{B\rightarrow K}.
\tag{TRISO-MR-016}
\label{eq:triso-mr-016}
\end{equation}

The second is

\begin{equation}
K_{11}
=
G_B^-p_{B\rightarrow B}^{(I_0)}.
\tag{TRISO-MR-017}
\label{eq:triso-mr-017}
\end{equation}

The third is

\begin{equation}
K_{12}
=
G_B^+p_{B\rightarrow B}^{(I_1)}.
\tag{TRISO-MR-018}
\label{eq:triso-mr-018}
\end{equation}

The fourth is

\begin{equation}
K_{13}
=
G_B^+p_{B\rightarrow I}.
\tag{TRISO-MR-019}
\label{eq:triso-mr-019}
\end{equation}

For $S_2$, the physical destinations are identical to $S_1$, but the shell transforms are evaluated from the outer-side Buffer reinsertion radius. Denote them $G_{B,2}^-$ and $G_{B,2}^+$.

\begin{equation}
\Phi_2=
G_{B,2}^-p_{B\rightarrow K}\Phi_0+
G_{B,2}^-p_{B\rightarrow B}^{(I_0)}\Phi_1+
G_{B,2}^+p_{B\rightarrow B}^{(I_1)}\Phi_2+
G_{B,2}^+p_{B\rightarrow I}\Phi_3.
\tag{TRISO-MR-034}
\label{eq:triso-mr-034}
\end{equation}

Hence

\begin{equation}
K_{20}=G_{B,2}^-p_{B\rightarrow K},
\quad
K_{21}=G_{B,2}^-p_{B\rightarrow B}^{(I_0)}.
\tag{TRISO-MR-035}
\label{eq:triso-mr-035}
\end{equation}

and

\begin{equation}
K_{22}=G_{B,2}^+p_{B\rightarrow B}^{(I_1)},
\quad
K_{23}=G_{B,2}^+p_{B\rightarrow I}.
\tag{TRISO-MR-036}
\label{eq:triso-mr-036}
\end{equation}

\subsection*{25.4 IPyC state rows}

From $S_3$, inward IPyC exit reaches the Buffer/IPyC interface and outward exit reaches IPyC/SiC:

\begin{equation}
\Phi_3=
G_{I,3}^-p_{I\rightarrow B}\Phi_2+
G_{I,3}^-p_{I\rightarrow I}^{(I_1)}\Phi_3+
G_{I,3}^+p_{I\rightarrow I}^{(I_2)}\Phi_4+
G_{I,3}^+p_{I\rightarrow S}\Phi_5.
\tag{TRISO-MR-037}
\label{eq:triso-mr-037}
\end{equation}

Therefore

\begin{equation}
K_{32}=G_{I,3}^-p_{I\rightarrow B},
\quad
K_{33}=G_{I,3}^-p_{I\rightarrow I}^{(I_1)},
\tag{TRISO-MR-038}
\label{eq:triso-mr-038}
\end{equation}

\begin{equation}
K_{34}=G_{I,3}^+p_{I\rightarrow I}^{(I_2)},
\quad
K_{35}=G_{I,3}^+p_{I\rightarrow S}.
\tag{TRISO-MR-039}
\label{eq:triso-mr-039}
\end{equation}

From $S_4$,

\begin{equation}
\Phi_4=
G_{I,4}^-p_{I\rightarrow B}\Phi_2+
G_{I,4}^-p_{I\rightarrow I}^{(I_1)}\Phi_3+
G_{I,4}^+p_{I\rightarrow I}^{(I_2)}\Phi_4+
G_{I,4}^+p_{I\rightarrow S}\Phi_5.
\tag{TRISO-MR-040}
\label{eq:triso-mr-040}
\end{equation}

Thus

\begin{equation}
K_{42}=G_{I,4}^-p_{I\rightarrow B},
\quad
K_{43}=G_{I,4}^-p_{I\rightarrow I}^{(I_1)},
\tag{TRISO-MR-041}
\label{eq:triso-mr-041}
\end{equation}

\begin{equation}
K_{44}=G_{I,4}^+p_{I\rightarrow I}^{(I_2)},
\quad
K_{45}=G_{I,4}^+p_{I\rightarrow S}.
\tag{TRISO-MR-042}
\label{eq:triso-mr-042}
\end{equation}

\subsection*{25.5 SiC state rows}

From $S_5$,

\begin{equation}
\Phi_5=
G_{S,5}^-p_{S\rightarrow I}\Phi_4+
G_{S,5}^-p_{S\rightarrow S}^{(I_2)}\Phi_5+
G_{S,5}^+p_{S\rightarrow S}^{(I_3)}\Phi_6+
G_{S,5}^+p_{S\rightarrow O}\Phi_7.
\tag{TRISO-MR-043}
\label{eq:triso-mr-043}
\end{equation}

Hence

\begin{equation}
K_{54}=G_{S,5}^-p_{S\rightarrow I},
\quad
K_{55}=G_{S,5}^-p_{S\rightarrow S}^{(I_2)},
\tag{TRISO-MR-044}
\label{eq:triso-mr-044}
\end{equation}

\begin{equation}
K_{56}=G_{S,5}^+p_{S\rightarrow S}^{(I_3)},
\quad
K_{57}=G_{S,5}^+p_{S\rightarrow O}.
\tag{TRISO-MR-045}
\label{eq:triso-mr-045}
\end{equation}

From $S_6$,

\begin{equation}
\Phi_6=
G_{S,6}^-p_{S\rightarrow I}\Phi_4+
G_{S,6}^-p_{S\rightarrow S}^{(I_2)}\Phi_5+
G_{S,6}^+p_{S\rightarrow S}^{(I_3)}\Phi_6+
G_{S,6}^+p_{S\rightarrow O}\Phi_7.
\tag{TRISO-MR-046}
\label{eq:triso-mr-046}
\end{equation}

Therefore

\begin{equation}
K_{64}=G_{S,6}^-p_{S\rightarrow I},
\quad
K_{65}=G_{S,6}^-p_{S\rightarrow S}^{(I_2)},
\tag{TRISO-MR-047}
\label{eq:triso-mr-047}
\end{equation}

\begin{equation}
K_{66}=G_{S,6}^+p_{S\rightarrow S}^{(I_3)},
\quad
K_{67}=G_{S,6}^+p_{S\rightarrow O}.
\tag{TRISO-MR-048}
\label{eq:triso-mr-048}
\end{equation}

\subsection*{25.5 OPyC state and direct release}

From $S_7$, inner shell exit reaches the SiC/OPyC interface.

Outer shell exit reaches the absorbing particle exterior.

The inward OPyC exit reaches the SiC/OPyC interface. Transmission to SiC gives $S_6$; reflection in OPyC gives $S_7$. The outward exit releases directly. Therefore

\begin{equation}
\Phi_7=
G_O^-p_{O\rightarrow S}\Phi_6+
G_O^-p_{O\rightarrow O}\Phi_7+
G_O^+.
\tag{TRISO-MR-049}
\label{eq:triso-mr-049}
\end{equation}

The two transient entries are

\begin{equation}
K_{76}=G_O^-p_{O\rightarrow S},
\qquad
K_{77}=G_O^-p_{O\rightarrow O}.
\tag{TRISO-MR-050}
\label{eq:triso-mr-050}
\end{equation}

The direct release transform is

\begin{equation}
\boxed{
B_7(s)=G_O^+(s).
}
\tag{TRISO-MR-020}
\label{eq:triso-mr-020}
\end{equation}

For every other transient state,

\begin{equation}
B_i(s)=0,
\qquad i\ne7.
\tag{TRISO-MR-021}
\label{eq:triso-mr-021}
\end{equation}

\subsection*{25.6 Matrix rearrangement without skipped algebra}

Collect the eight first-step equations:

\begin{equation}
\boldsymbol\Phi
=
\mathbf K\boldsymbol\Phi
+
\mathbf B.
\tag{TRISO-MR-022}
\label{eq:triso-mr-022}
\end{equation}

Subtract $mathbf K\boldsymbol\Phi$ from both sides:

\begin{equation}
\boldsymbol\Phi
-
\mathbf K\boldsymbol\Phi
=
\mathbf B.
\tag{TRISO-MR-023}
\label{eq:triso-mr-023}
\end{equation}

Write $oldsymbol\Phi=\mathbf I\boldsymbol\Phi$:

\begin{equation}
\mathbf I\boldsymbol\Phi
-
\mathbf K\boldsymbol\Phi
=
\mathbf B.
\tag{TRISO-MR-024}
\label{eq:triso-mr-024}
\end{equation}

Factor $oldsymbol\Phi$:

\begin{equation}
(\mathbf I-\mathbf K)\boldsymbol\Phi
=
\mathbf B.
\tag{TRISO-MR-025}
\label{eq:triso-mr-025}
\end{equation}

When $mathbf I-\mathbf K$ is nonsingular, left-multiply by its inverse:

\begin{equation}
(\mathbf I-\mathbf K)^{-1}
(\mathbf I-\mathbf K)
\boldsymbol\Phi
=
(\mathbf I-\mathbf K)^{-1}\mathbf B.
\tag{TRISO-MR-026}
\label{eq:triso-mr-026}
\end{equation}

Use the inverse identity:

\begin{equation}
\boxed{
\boldsymbol\Phi
=
(\mathbf I-\mathbf K)^{-1}\mathbf B.
}
\tag{TRISO-MR-027}
\label{eq:triso-mr-027}
\end{equation}

\subsection*{25.7 Neumann-series path interpretation}

If

\begin{equation}
\rho(\mathbf K)<1,
\tag{TRISO-MR-028}
\label{eq:triso-mr-028}
\end{equation}

then

\begin{equation}
(\mathbf I-\mathbf K)^{-1}
=
\mathbf I
+
\mathbf K
+
\mathbf K^2
+
\mathbf K^3
+\cdots.
\tag{TRISO-MR-029}
\label{eq:triso-mr-029}
\end{equation}

Multiply by (mathbf B):

\begin{equation}
\boldsymbol\Phi
=
\mathbf B
+
\mathbf K\mathbf B
+
\mathbf K^2\mathbf B
+
\mathbf K^3\mathbf B
+\cdots.
\tag{TRISO-MR-030}
\label{eq:triso-mr-030}
\end{equation}

The term (mathbf B) is direct absorption without another transient renewal.

The term $mathbf K\mathbf B$ is absorption after one transient renewal.

The term $mathbf K^2\mathbf B$ is absorption after two transient renewals.

Thus the inverse sums arbitrarily long repeated interface-renewal paths without sampling each path individually.

\subsection*{25.8 Uniform-volume initial source}

The kernel radial probability density is

\begin{equation}
f_R(r)
=
\frac{3r^2}{R_1^3},
\qquad
0\le r\le R_1.
\tag{TRISO-MR-031}
\label{eq:triso-mr-031}
\end{equation}

The exact kernel first-exit transform from radius $r$ is $H_K(r,s)$.

After reaching Kernel/Buffer, reflection gives $S_0$ and transmission gives $S_1$.

Define

\begin{equation}
Q(s)
=
p_{K\rightarrow K}\Phi_0(s)
+
p_{K\rightarrow B}\Phi_1(s).
\tag{TRISO-MR-032}
\label{eq:triso-mr-032}
\end{equation}

Average over the birth distribution:

\begin{equation}
\boxed{
\Phi_{\rm init}(s)
=
\int_0^{R_1}
\frac{3r^2}{R_1^3}
H_K(r,s)
Q(s)\,dr.
}
\tag{TRISO-MR-033}
\label{eq:triso-mr-033}
\end{equation}

\subsection*{25.9 Process A, B and C distinction}

[DEFINITION] Process A is finite-capture production WOS.

[DEFINITION] Process B is explicit accelerated exact-interface renewal.

[DEFINITION] Process C is the deterministic matrix reduction of Process B.

[EXACT TARGET] The mathematical reduction target is $B=C$.

\textbf{[EMPIRICAL COMPATIBILITY]} Process A and Process B/C have controlled finite-(epsilon) compatibility evidence.

\textbf{[IMPORTANT]} Neither statement is a proof of the $epsilon\rightarrow0$ continuum limit.


\section*{26. Verification linked to the matrix derivation}

\subsection*{26.1 Controlled two-layer B/C reconciliation}

\textbf{[VERIFIED]} The explicit Process-B sample size was

\begin{equation}
N_B=20000.
\tag{TRISO-VER-300}
\label{eq:triso-ver-300}
\end{equation}

At every predeclared positive transform point, define the standardized matrix discrepancy

\begin{equation}
z(s)=
\frac{\Phi_B(s)-\Phi_C(s)}
{\mathrm{SE}[\Phi_B(s)]}.
\tag{TRISO-VER-301}
\label{eq:triso-ver-301}
\end{equation}

The executed maximum magnitude was

\begin{equation}
\boxed{
\max_s|z(s)|=1.004.
}
\tag{TRISO-VER-302}
\label{eq:triso-ver-302}
\end{equation}

At zero transform frequency,

\begin{equation}
\boxed{
\Phi_{\rm init}(0)=1
}
\tag{TRISO-VER-303}
\label{eq:triso-ver-303}
\end{equation}

to the predeclared numerical tolerance.

Thus the deterministic matrix is compatible with the explicit exact-interface renewal at the controlled Monte-Carlo precision.

\subsection*{26.2 Genuine three-layer multistate verification}

The controlled radii were

\begin{equation}
(R_1,R_2,R_3)
=
(50,75,100)\,\mu\mathrm m.
\tag{TRISO-VER-304}
\label{eq:triso-ver-304}
\end{equation}

The controlled diffusivities were

\begin{equation}
(D_1,D_2,D_3)
=
(1,2,5)\times10^{-9}\,\mathrm{m^2s^{-1}}.
\tag{TRISO-VER-305}
\label{eq:triso-ver-305}
\end{equation}

The four-state Process-B/Process-C transform comparison gave

\begin{equation}
\boxed{
\max_s|z(s)|=0.729.
}
\tag{TRISO-VER-306}
\label{eq:triso-ver-306}
\end{equation}

Against the independently refined FV reference,

\begin{equation}
\boxed{
\mathrm{RMS}(F_B-F_{FV})
=
1.0395\times10^{-3}.
}
\tag{TRISO-VER-307}
\label{eq:triso-ver-307}
\end{equation}

The maximum absolute CDF discrepancy was

\begin{equation}
\boxed{
\max_t|F_B-F_{FV}|
=
1.7794\times10^{-3}.
}
\tag{TRISO-VER-308}
\label{eq:triso-ver-308}
\end{equation}

\textbf{[VERIFIED]} The FV spatial/time refinement changes were smaller than the stochastic uncertainty, so the continuum-reference discretization error did not dominate this comparison.

\bibliographystyle{plain}
\bibliography{references}
\end{document}
